\documentclass[book, paper=a4paper, 12pt]{jlreq}
\usepackage[dvipdfmx]{graphicx,xcolor,hyperref}%graphics
\usepackage{float}
\usepackage{amsmath}
\usepackage{amsthm}
\usepackage{verbatim}
\usepackage{latexsym}
\usepackage{amssymb}
\usepackage{dynkin-diagrams}
\usepackage{ccicons}
\hypersetup{% hyperrefオプションリスト
   setpagesize=false,
   bookmarksnumbered=true,
   bookmarksopen=true,
   colorlinks=true,
   linkcolor=blue,
   citecolor=purple,
} 
\input xy
\xyoption{all}

\def\emptyline{\vspace{12pt}}
\newcommand*{\term}[2][]{\textcolor{red}{#2}%
	\ifx\relax#1\relax\index{#2}\else\index{#1@#2}\fi}

\usepackage{makeidx}
\makeindex

\newtheorem{PP}{授業の目標}
\newtheorem{Thm}{定理}
\newtheorem{Prop}{命題}
\newtheorem{Q}{考えたい課題}
\newtheorem{Quiz}{問題}
\theoremstyle{definition}
\newtheorem{Def}{定義}

\theoremstyle{remark}
\newtheorem{ex}{例}
\newtheorem{exQ}{例題}

\DeclareMathOperator{\Ker}{Ker}
\DeclareMathOperator{\Image}{Im}
\DeclareMathOperator{\stab}{stab}
\DeclareMathOperator{\sign}{sign}
\DeclareMathOperator{\id}{id}
\DeclareMathOperator{\Is}{Is}
\DeclareMathOperator{\tr}{tr}
\DeclareMathOperator{\vol}{vol}
\DeclareMathOperator{\const}{const}
\DeclareMathOperator{\GL}{GL}
\DeclareMathOperator{\SL}{SL}
\DeclareMathOperator{\PSL}{PSL}
\DeclareMathOperator{\Ot}{O}
\DeclareMathOperator{\SO}{SO}
\DeclareMathOperator{\U}{U}
\DeclareMathOperator{\SU}{SU}
\DeclareMathOperator{\Sp}{Sp}
\DeclareMathOperator{\IM}{Im}
\DeclareMathOperator{\RE}{Re}
\DeclareMathOperator{\Res}{Res}
\DeclareMathOperator{\arcSin}{Sin^{-1}}


\title{ユークリッド幾何学の展開\par \normalsize{---歴史からの現代数学入門---}}
%\title{Beyond the Euclidean Geometry}
\author{snow00two}

\begin{document}
\maketitle

\frontmatter

\chapter{イントロダクション}

現代の大学の数学科の数学教育では
数学の歴史的な発展を学ぶことは重視されておらず，
しかも，遠近法や射影幾何学や一次分数変換や非ユークリッド幾何学などについては説明すらなされないまま，
一般位相や関数論，多様体論，位相幾何学，代数を学び，
その応用の中で突然射影空間，射影変換，楕円曲線，保形形式が現れるような構成になっている．
時間の関係でこのような形になるのは仕方が無いのかもしれないが，
そのために数学に挫折してしまう学生の多いことは悲しいことである．
そこでこの本ではユークリッド幾何学の完成以降の
現代数学に繋がるユークリッド幾何学の発展について解説し，
数学科で学ぶ数学へのイントロダクションになることを期待している．
具体的には，$\SO(3)$や$\Ot(3)$の有限群，
ユークリッド平面の装飾群の(特にココンパクトなもの17個の)分類と
完備平坦リーマン多様体の分類を与える．
これは群作用の観点からまとめてあるので群論の入門になると思うが，
リーマン幾何学の入門ともなっている．
そのあと遠近法や3DCG，射影幾何学，
実射影空間，リーマン球面$\mathbf CP^1$と一次変換の関係(分類)，
トーラスのモジュライと上半平面の双曲幾何，楕円曲線とモジュラー関数
などを紹介しており，関数論との関係も強調しておきたい．
必ずしも歴史的に正確なことにはこだわらないが，
おおよその流れはおえるような形で記述を与えたい．

数学的な予備知識については数学科の大学2年生ぐらいまでで習うことは仮定しているが，
逆にこの本の内容を深く理解しようという視点からそれらの教科書を復習して
もらえれば幸いである．

大学の数学を始めると絵や図を描いて説明してはいけないなどの不思議なルールが
課せられることがあるが，それは初学者への方便であり，
絵や図を使った方が分かり易いところでは沢山の絵や図を入れる方針で本文は書いた．
ホームページにこの本の図のソースファイルや関連した動画やアプリを
\href{https://snow00two.github.io/mySite/}{私のサイト}に用意したので
それらを動かしながら理解を深めてもらいたい．
\tableofcontents

\mainmatter

%%LectureNote01 pre
%\begin{flushright}
%   2024年10月7日
%   
%   九大・数理 大津幸男
%\end{flushright}
%
%\begin{center}
%   \large{微分幾何学大意・数学特論4\small{
%   \href{https://creativecommons.org/licenses/by-nc-nd/4.0/}{\ccbyncnd}}
%   \par ---直交群とその有限部分群 1---}
%\end{center}
%
%\begin{PP}
%   ユークリッド幾何から派生した幾何学・群論・関数論を解説し，
%   現代的な微分幾何への橋渡しを行う．更にプログラミングとアート
%   との関連を探る．
%\end{PP}

\chapter{$2$, $3$次元直交群の有限部分群の分類}
\section{$2$次元直交群の有限部分群}
\subsection{$\mathbf R^2$の線形合同変換の例}
ユークリッド平面$\mathbf R^2$の全ての点$(x,y)$に$\mathbf R^2$の別の点$(x',y')$を対応させる対応を（$\mathbf R^2$の）\term[]{変換}という．
平面$\mathbf R^2$の点$(x,y)$と複素平面$\mathbf C$の点$x+iy$を同一視する．（高校でやったように）原点を中心にした\term[]{回転}(rotation)は$\cos\theta+i\sin\theta$を掛けることに対応するので
\[
     (\cos\theta+i\sin\theta)(x+iy)
     =x\cos\theta-y\sin\theta+i(x\sin\theta+y\cos\theta)
\]
つまり，
\[
     \begin{pmatrix}
          x\cos\theta-y\sin\theta\\
          x\sin\theta+y\cos\theta\\
     \end{pmatrix}
     =
     \begin{pmatrix}
          \cos \theta &-\sin \theta\\
          \sin \theta& \cos \theta\\
     \end{pmatrix}
     \begin{pmatrix}
          x\\
          y\\
     \end{pmatrix}
     =S_\theta\mathbf x
\]
と行列の積で表すことが出来る，ここで，
\[
     S_\theta=
     \begin{pmatrix}
          \cos \theta &-\sin \theta\\
          \sin \theta& \cos \theta\\
     \end{pmatrix},\quad
     \mathbf x
     =
     \begin{pmatrix}
          x\\
          y\\
     \end{pmatrix}.
\]
このように行列の積で表される変換を\term[]{線形変換}という．
また，距離を保つ変換を\term[]{合同変換}という，回転は合同変換の一つである（距離については次の節で詳しく述べる）．
また，
\begin{equation}\label{rel1}
     S_\theta S_\tau=S_{\theta+\tau},\quad S_0=E
\end{equation}
特に
\[
     S_{-\theta}S_\theta=E,\quad S_{\theta}^{-1}
     =S_{-\theta}={}^tS_{\theta}
\]
ただし，$E$は単位行列$(\delta_{ij})$．
\begin{figure}[htb]
   \centering
   \begin{minipage}{0.4\columnwidth}
      \centering
        \includegraphics{~/Documents/GitHub/bEuclideanGeo/ch1/images/figureGroup.1}
        \caption{原点での回転}\label{}
   \end{minipage} 
   \begin{minipage}{0.4\columnwidth}
      \centering
        \includegraphics{~/Documents/GitHub/bEuclideanGeo/ch1/images/figureGroup.2}
        \caption{$x$軸での鏡映}\label{}
   \end{minipage} 
\end{figure}

（複素共役$z\mapsto \bar z$に対応する）$x$軸での\term[]{鏡映（折り返し）}(refrection)は$(x,y)$に$(x,-y)$を対応させることなので，
\[
      \begin{pmatrix}
          x\\
          -y\\
     \end{pmatrix}
     =
     \begin{pmatrix}
          1&0\\
          0&-1\\
     \end{pmatrix}
     \begin{pmatrix}
          x\\
          y\\
     \end{pmatrix}
\]
と線形変換で表すことが出来る，これも合同変換でもある．一般に，
直線$\ell: -x\sin\theta+y\cos\theta=0$についての鏡映の，
行列の積を使った表現を求める：
$\ell$の単位法線ベクトルは
\[
     \mathbf e
     =\begin{pmatrix}
          -\sin\theta\\
          \cos\theta\\
     \end{pmatrix}
\]
であるので$\mathbf x$との
内積は$ (\mathbf e,\mathbf x)=-x\sin\theta+y\cos\theta$．
$\mathbf x$を$\mathbf e$方向の成分とその直交成分にわけると
\[
     \mathbf x= (\mathbf e,\mathbf x)\,\mathbf e
     +\{\mathbf x- (\mathbf e,\mathbf x)\,\mathbf e\}.
\]
\begin{figure}[htb]
   \centering
   \begin{minipage}{0.4\columnwidth}
        \centering
        \includegraphics{~/Documents/GitHub/bEuclideanGeo/ch1/images/figureGroup.3}
        \caption{直線$\ell$での鏡映}\label{}\end{minipage} 
   \begin{minipage}{0.4\columnwidth}
      \centering
      \includegraphics{~/Documents/GitHub/bEuclideanGeo/ch1/images/figureGroup.4}
      \caption{鏡映と鏡映の合成}\label{compRefs}
   \end{minipage} 
\end{figure}
$ (\mathbf e,\mathbf x)\,\mathbf e$ は$\mathbf x$の$\mathbf e$成分，
$\mathbf x- (\mathbf e,\mathbf x)\, \mathbf e$ は直交成分．
$\ell$についての鏡映は（図3）
\[
     \{\mathbf x- (\mathbf e,\mathbf x)\,\mathbf e\}
          - (\mathbf e,\mathbf x)\,\mathbf e
     =\mathbf x-2 (\mathbf e,\mathbf x)\,\mathbf e.
\]
これをベクトルの成分で計算すると
\[
\begin{split}
     \mathbf x-2 (\mathbf e,\mathbf x)\,\mathbf e
     &=
     \begin{pmatrix}
          x\\
          y\\
     \end{pmatrix}
     -2(-x\sin\theta+y\cos\theta)
     \begin{pmatrix}
          -\sin\theta\\
          \cos\theta\\
     \end{pmatrix}\\
     &=
     \begin{pmatrix}
          x(1-2\sin^2\theta)+2y\cos\theta\sin\theta\\
          2x\sin \theta \cos\theta+y(1-2\cos^2\theta)\\
     \end{pmatrix}\\
     &=\begin{pmatrix}
          \cos 2\theta &\sin 2\theta\\
          \sin 2\theta&-\cos 2\theta\\
     \end{pmatrix}
     \begin{pmatrix}
          x\\
          y\\
     \end{pmatrix}
     =R_{\theta}\mathbf x,
\end{split}
\]
ここで
\[
     R_{\theta}
     =
     \begin{pmatrix}
          \cos 2\theta &\sin 2\theta\\
          \sin 2\theta&-\cos 2\theta\\
     \end{pmatrix}.
\]
直線$\ell$を鏡映の軸とも言う．

同じ直線についての鏡映を2回続けると元に戻るので
\begin{equation}\label{rel2}
     R_{\theta}^2=R_{\theta}R_{\theta}=E,
     \quad\text{言い換えると}\quad 
     R_{\theta}^{-1} =R_{\theta}={}^tR_{\theta}.
\end{equation}
$\mathbf 0$を通る直線に沿っての二つの鏡映の合成は$\mathbf 0$についての回転になる：行列の積で表すと
\begin{equation}\label{rel3}
     R_{\theta}R_{\tau}=S_{2(\theta-\tau)}
     \quad\text{よって}\quad
     R_{\theta}=S_{2(\theta-\tau)}R_{\tau}.
\end{equation}
これは初等幾何学を使うことでわかる（図~\ref{compRefs}）が，単に行列の積を計算することでもわかる．
%アプリ./reflections/p5の実行
%アプリ./reflections/p5の実行
%アプリ./reflections/p5の実行

%
%問　回転と回転の合成は回転なので
%\[
%     S_\theta S_\tau=S_{\theta+\tau}
%\]
%が成り立つ．これを行列の積を計算することで確かめよ．
%
%
%問　二つの直線が交わらないとき，つまり，平行なとき，これらの直線に沿っての二つの鏡映の合成は何になるか？
%

\subsection{一般の直交変換群}
$n$を自然数とし$n$次元ユークリッド空間$\mathbf R^n$を考える．
\[
     (\mathbf u,\mathbf v)={}^t\mathbf u\mathbf v
     =u_1v_1+u_2v_2+\cdots+u_nv_n\in \mathbf R
     \quad
     \mathbf u
      =\begin{pmatrix}
              u_1\\u_2\\\vdots\\u_n
     \end{pmatrix},
     \mathbf v
      =\begin{pmatrix}
              v_1\\v_2\\\vdots\\v_n
     \end{pmatrix}
     \in \mathbf R^n
\]
と置き，$\mathbf R^n$の\term[]{標準内積}という．
${}^tA=(b_{ij})$は$A=(a_{ij})$の転置行列，つまり，$b_{ij}=a_{ji}$．
$(\mathbf u,\mathbf u)\ge 0$なので
\[
     \|\mathbf u\| 
     = 
     \sqrt{(\mathbf u,\mathbf u)}
     \ge 0
\]
と置き，\term[]{ノルム}（\term[]{長さ}）という．

$T:\mathbf R^n\to \mathbf R^n$が\term[]{直交変換}とは
\begin{equation}\label{orthogonalTransform}
     (T(\mathbf u),T(\mathbf v))
     =
     (\mathbf u,\mathbf v)\quad
     \mathbf u,\mathbf v\in \mathbf R^n
\end{equation}
が成り立つことである．この定義から直交変換はノルムを保存することが分かる：
\begin{equation}\label{preserveNorm}
     \|T(\mathbf u)\| = \|\mathbf u\| 
\end{equation}
逆にこれが成り立てば直交変換になることも簡単に分かる．したがって
直交変換は合同線形変換である．

$n$次正方行列$A$が\term[]{直交行列}とは
\[
     {}^tAA=E
\]
が成り立つことであり，これは
\[
     A^{-1}={}^tA
\]と同値である．
したがって，回転$S_\theta$や鏡映$R_{\theta}$は直交行列である．

$n$次正方行列$A$に対して$T_A(\mathbf x)=A\mathbf x$で線形変換を定義すると
\[
     T_A\text{ が直交変換}
     \Leftrightarrow 
     A\text{ が直交行列}
\]
$\Leftarrow)$は
\[
\begin{split}
     (T_A(\mathbf u),T_A(\mathbf v))
     &=
     (A\mathbf u,A\mathbf v)
     =
     {}^t(A\mathbf u)A\mathbf v
     =
     {}^t\mathbf u{}^tAA\mathbf v\\
     &
     =
     {}^t\mathbf u({}^tAA)\mathbf v
     =
     {}^t\mathbf uE\mathbf v
     =
     {}^t\mathbf u\mathbf v
     =
     (\mathbf u,\mathbf v)
\end{split}
\]
よりわかる．

$\Rightarrow)$は，
\[
     \mathbf e_1=
     \begin{pmatrix}
              1\\0\\0\\\vdots\\0
     \end{pmatrix},
     \mathbf e_2=
     \begin{pmatrix}
              0\\1\\0\\\vdots\\0
     \end{pmatrix},     
     \cdots,
     \mathbf e_n=
     \begin{pmatrix}
              0\\0\\\vdots\\0\\1
     \end{pmatrix}
\]
を$\mathbf R^n$の標準基とし，
\eqref{orthogonalTransform}の式で
$\mathbf u=\mathbf e_i$, $\mathbf v=\mathbf e_j$ ととると
${}^tAA$と$E$の$(i,j)$成分が一致することが分かるので示せた．
\emptyline

集合$G$が\term[]{群}であるとは任意の$a$, $b\in G$に対して積$ab\in G$が定まり，次の三条件を満たすことである：
\[
     (ab)c = a(bc)\quad \forall a, b, c\in G
\]
を満たす．
単位元$e\in G$が存在し，
\[
     ea=ae=a\quad \forall a\in G
\]
を満たす．また任意の$a\in G$に対して逆元$a^{-1}$が存在して
\[
     aa^{-1}= a^{-1}a = e
\]
を満たす．単位元や逆元の一意性は直ぐに分かる．
\emptyline

$H\subset G$が\term[]{部分群}とは$H$が$G$の積に関して群となることである．これは
\[
     a,b\in H\Rightarrow ab\in H\quad\text{かつ}\quad 
     a\in H\Rightarrow a^{-1}\in H
\]
を満たすことと同値である．
\emptyline

正方行列$A$の行列式を$|A| $あるいは$\det A$と表すと
\[
     |AB|=|A||B|
\]
で，$|A|\neq0$が$A$が正則であることと同値なので
\[
     \GL(n;\mathbf R)=\{A\mid n\text{次正則行列} \}
\]
は
群になる．

直交行列$A$については
$1=|E|=|{}^tA||A|=|A|^2$より$|A|=\pm 1$．
\[
     \Ot(n)=\{A\in \GL(n;\mathbf R)\mid \text{直交行列}\},\quad
     \SO(n)=\{A\in \Ot(n)\mid |A|=1\}
\]
と置くと，
$A$, $B\in \Ot(n)$に対して
\[
     {}^t(AB)AB={}^tB{}^tA AB ={}^tBEB={}^tBB=E
\]
となり，
\[
     A^{-1}={}^tA \quad\text{より}\quad (A^{-1})^{-1}
     =A
     ={}^t({}^tA)={}^t(A^{-1})
\]
なるので$\Ot(n)$は$\GL(n;\mathbf R)$の部分群になる．$A$, $B\in \SO(n)$に対して
\[
     |AB|=|A||B| =1\times 1=1,\quad
     |A^{-1}| =|{}^tA| = |A| =1
\]
となるので$\SO(n)$も群で$\Ot(n)$の部分群である．
\emptyline

$\GL(n;\mathbf R)$を$n$次元\term[]{一般線形群}(General linear group)，
$\Ot(n)$を$n$次元\term[]{直交群}(Orthogonal group)，$\SO(n)$を$n$次元\term[]{特殊直交群}(Special Orthogonal group)という．

%\emptyline
%
%群が多様体で，積を与える写像と逆を与える写像が微分可能になっているとき
%\term[]{Lie 群}(Lie group)というが，これらはLie群の例になっている．
%また，（係数体が$\mathbf R$と限らない）行列の作る群を\term[]{古典群}(Classical group)ということもある．


\subsection{二次元直交群とその有限部分群の分類}
回転と鏡映が$\Ot(2)$の元であることを上で示したが，
$\Ot(2)$の元が回転と鏡映に限ることを示す：
$A\in \Ot(2)$とし
$\mathbf R^2$の標準基
$
     \mathbf e_1
     =
     \begin{pmatrix}
              1\\0
     \end{pmatrix}$, 
$     \mathbf e_2
     =
     \begin{pmatrix}
              0\\1
     \end{pmatrix}
$を作用させる．\eqref{preserveNorm}より
$\|A\mathbf e_1\|=\|\mathbf e_1\|=1$なので，ある$\theta$を使って
\[
     A\mathbf e_1
     =\begin{pmatrix}
              \cos\theta\\\sin\theta
     \end{pmatrix}
\]
と書ける．\eqref{preserveNorm}, \eqref{orthogonalTransform}より
\[
     \|A\mathbf e_2\|=1,\quad 
     (A\mathbf e_1,A\mathbf e_2)=(\mathbf e_1,\mathbf e_2)=0
\]
なので
\[
     A\mathbf e_2
     =
     \begin{pmatrix}
         -\sin\theta\\\cos\theta
     \end{pmatrix},
     \quad\text{あるいは}\quad
     A\mathbf e_2
     =
     \begin{pmatrix}
         \sin\theta\\-\cos\theta
     \end{pmatrix}
\]
最初の場合が回転$S_\theta$で，次の場合が鏡映$R_{\theta/2}$である．
$|S_\theta|=1$, $|R_{\theta/2}|=-1$なので
\[
     \SO(2)=\{S_\theta\mid 0\le \theta<2\pi\},\quad
     \Ot(2)=\SO(2)\cup \{R_{\theta/2}\mid 0\le \theta<2\pi\}
\]
がわかる．
$f:\mathbf R\ni\theta\mapsto S_\theta\in \SO(2)$は\eqref{rel1}より群の準同形で
$\Ker f= 2\pi \mathbf Z$．これは正規部分群なので剰余群が構成でき
\[
     \SO(2)\sim \mathbf R/2\pi \mathbf Z
     \sim S^1=\{z\in \mathbf C\mid z\bar z=1\}
\]
となる．
%（正確には群としての同形だが，位相空間や，多様体，Lie群としての同形も導いている）．
%同様に$\Ot(2)$も多様体であることも分かる．

\begin{Q}
     $\Ot(2)$の有限部分群の分類
\end{Q}

まず$G$を$\SO(2)$の有限部分群とする．$A\neq E\in G$について
$A=S_\theta$のとき $\theta(A)=\theta\in (0,\pi)$と置く．
$G$は有限集合なので$\{\theta(A)\in (0,\pi)\mid A\in G\}$に
最小値$\theta_0 > 0$ が存在する．
$A_0\in G$を$\theta_0=\theta(A_0)$となるものとする．
$A\neq E\in G$について
\[
     \theta(AA_0^{-k})=\theta(A)-k\theta_0\quad k\in \mathbf N
\]
なので，ある$k\in \mathbf N$について
\[
     0\le \theta(A)-k\theta_0 < \theta_0
\]
となり，
$\theta_0 > 0$ の最小性から$ \theta(A)-k\theta_0=0$となり
\[
     AA_0^{-k}=E\quad\text{つまり}\quad A=A_0^{k}=S_{\theta_0}^{k}
\]
$m$を$G$の個数とすると
\[
     G=\{E, S_{\theta_0},S_{\theta_0}^2,\dots,S_{\theta_0}^{m-1}\}
     =\langle S_{\theta_0}\rangle,\quad \theta_0 = 2\pi/m
\]
最後の表示は$S_{\theta_0}$が$G$を生成していることを表す．
このような回転一つから生成される群を$m$\term[]{次巡回群(Cyclic group)}といい，$C_m$と表す．
\emptyline

$G$が$\Ot(2)$の有限部分群で$\SO(2)$部分群でないとする．
するとある鏡映$M\in G$が存在する．
他の鏡映$M'\in G$があると\eqref{rel2}より$M'M^{-1}\in C_m$となり
\[
     M'= MS_{\theta_0}^{k}\in G \quad  k=0,\dots,m-1
\]
と書けることが分かる．逆にこう書ける元は$G$にはいるので，$G$は
\[
\begin{split}
     G
     &=\{E, S_{\theta_0},S_{\theta_0}^2,\dots,S_{\theta_0}^{m-1}\}
     \cup\{M, MS_{\theta_0},MS_{\theta_0}^2, \dots, MS_{\theta_0}^{m-1}\}\\
     &=C_m\cup MC_m     =
     \langle S_{\theta_0}, M\rangle
\end{split}
\]
このような群を
\term[にめんたいぐん]{二面体群}(Dihedral group)といい，$D_{m}$と表す．
ただし，$C_1=\{E\}$, $D_1=\{E, M\}$,  $C_2=\{E, R\}$, $D_2=C_2\cup MC_2$
とする.　また人によっては$D_{2m}$と書く人もいるので文脈で判断してほしい．

\begin{figure}[htb]
   \centering
   \begin{minipage}{0.4\columnwidth}
        \centering
        \includegraphics{~/Documents/GitHub/bEuclideanGeo/ch1/images/figureGroup.7}
        \caption{$m=3$の二面体群$D_3$}\label{}
   \end{minipage} 
   \begin{minipage}{0.4\columnwidth}
        \centering
        \includegraphics{~/Documents/GitHub/bEuclideanGeo/ch1/images/figureGroup.8}
        \caption{$m=4$の二面体群$D_4$}\label{}
   \end{minipage} 
\end{figure}
\begin{figure}[htb]
   \centering
   \begin{minipage}{0.4\columnwidth}
        \centering
        \includegraphics{~/Documents/GitHub/bEuclideanGeo/ch1/images/figureGroup.9}
        \caption{$m=5$の二面体群$D_5$}\label{}
   \end{minipage} 
   \begin{minipage}{0.4\columnwidth}
        \centering
        \includegraphics{~/Documents/GitHub/bEuclideanGeo/ch1/images/figureGroup.10}
        \caption{$m=6$の二面体群$D_6$}\label{}
   \end{minipage} 
\end{figure}
\begin{figure}[htb]
   \centering
   \begin{minipage}{0.4\columnwidth}
        \centering
        \includegraphics{~/Documents/GitHub/bEuclideanGeo/ch1/images/figureGroup.5}
        \caption{$C_2$と$D_2$}\label{}
   \end{minipage} 
   \begin{minipage}{0.4\columnwidth}
        \centering
        \includegraphics{~/Documents/GitHub/bEuclideanGeo/ch1/images/figureGroup.6}
        \caption{$C_1=\{E\}$と$D_1$}\label{}
   \end{minipage} 
\end{figure}


\begin{Quiz}
   $n$次対称群\footnote{$\{1,2,\dots,n\}$の置換全体のなす群}を 
   $S_n$と書く．$D_{3}$ は$S_3$と群として同形であることを示せ．
\end{Quiz}


%この授業の成績は最後のレポートで付けます．
%適宜授業の中で問題を出すので解いてください．
%この授業ではSTEAM （=Science, Technology, Engineering, Art and Mathematics）教育の観点を取り入れていくつもりです．
%そこで\term[]{レポートは \TeX で作成し，pdf ファイルとそのソースファイルを一緒にMoodle に提出}してください．metapost, TikZ, Matlab, mathematica, maple  等で制作した図も適宜入れて貰えると評価が上がります．それ以外（手書きのノート等）は受け付けません．diff を使って直ぐにコピーしたことが明らかになるようなレポートは全て$0$点です．
%また，授業の内容に関連してプログラムしたアプリ，動画も受け付けます．
%STEAM:
% classRoom/General/TeX/introductionTex.tex
% classRoom/General/JavaScriptP5/P5reflection/P5Sierpinski/JavaScriptP5js1.tex
%のデモ
% classRoom/General/TeX/introductionTex.tex
% classRoom/General/JavaScriptP5/P5reflection/P5Sierpinski/JavaScriptP5js1.tex
%のデモ




%\documentclass[a4paper,12pt]{amsart}
%\usepackage[dvipdfmx]{graphicx,xcolor,hyperref}%graphics
%\usepackage{float}
%\usepackage{amsmath}
%\usepackage{verbatim}
%\usepackage{latexsym}
%\usepackage{amssymb}
%\usepackage{ccicons}
%%\usepackage{emath}
%%\usepackage{maplestd2e}
%%\usepackage{fancybox}%enclose by box
%\hypersetup{% hyperrefオプションリスト
%setpagesize=false,
% bookmarksnumbered=true,
% bookmarksopen=true,
% colorlinks=true,
%linkcolor=red,
%citecolor=blue,
%} 
%
%\renewcommand{\figurename}{図}
%\renewcommand{\proofname}{証明}
%\renewcommand{\tablename}{表}
%\newtheorem*{PP}{授業の目標}
%\newtheorem*{Thm}{定理}
%\newtheorem*{Prop}{命題}
%\newtheorem*{Q}{考えたい課題}
%\newtheorem*{Quiz}{問題}
%
%\theoremstyle{definition}
%\newtheorem*{Def}{定義}
%
%\theoremstyle{remark}
%\newtheorem*{ex}{例}
%\newtheorem*{exQ}{例題}
%
%\def\emptyline{\vspace{12pt}}
%
%\DeclareMathOperator{\Ker}{Ker}
%\DeclareMathOperator{\Image}{Im}
%\DeclareMathOperator{\stab}{stab}
%
%
%\begin{document}
%%\pagestyle{empty}
%\begin{flushright}
%
%2024年10月21日
%
%九大・数理 大津幸男
%\end{flushright}
%\begin{center}
%\large{微分幾何学大意・数学特論4\small{
%\href{https://creativecommons.org/licenses/by-nc-nd/4.0/}{\ccbyncnd}}\par ---直交群とその有限部分群 2---}
%\end{center}
%
%
%\subsection{復習}
%$A$を$n$次正方行列とするとき
%$T_A(\mathbf x)=A\mathbf x$が直交変換であること
%\begin{equation}\label{orthogonalTransform}
%     (T_A(\mathbf u),T_A(\mathbf v))
%     =
%     (\mathbf u,\mathbf v)\quad
%     \mathbf u,\mathbf v\in \mathbf R^n
%\end{equation}
%は$A$が直交行列であること
%\[
%     {}^tAA=E
%\]
%と同値であった．また，これは
%\[
%     \|T_A(\mathbf x)\|=\|\mathbf x\|
%\]
%とも同値である．ただし
%$\|\mathbf x\|= \sqrt{(\mathbf x,\mathbf x)}\ge 0$ である．
%このとき$|A|=\pm1$であり，$n=2$のとき
%$|A|=1$の元は回転$S_\theta$，$|A|=-1$の元は鏡映$R_{\theta/2}$になる：
%\[
%     S_\theta
%     =
%     \begin{pmatrix}
%          \cos \theta &-\sin \theta\\
%          \sin \theta& \cos \theta\\
%     \end{pmatrix},\quad
%     R_{\theta/2}
%     =\begin{pmatrix}
%          \cos \theta &\sin \theta\\
%          \sin \theta&-\cos \theta\\
%     \end{pmatrix}.
%\]
%ただし
%$|A|$は$A$の行列式を表す．$n\in \mathbf N$ に対して
%\[
%     \Ot(n)=\{A\in \GL(n;\mathbf R)\mid \text{直交行列}\},\quad
%     \SO(n)=\{A\in \Ot(n)\mid |A|=1\}
%\]
%と置く．
%\begin{Thm}
%$\Ot(2)$の有限部分群は
%巡回群$C_m$と
%二面体群$D_{m}$に限る．ただし
%$m\in \mathbf N$．
%\end{Thm}
%
%
\begin{figure}[htb]
\centering
\begin{minipage}{0.4\columnwidth}
\centering
     \includegraphics{~/Documents/GitHub/bEuclideanGeo/ch1/images/figureGroup2.1}
     \caption{巡回群$C_4$}\label{cyclic}
\end{minipage} 
\begin{minipage}{0.4\columnwidth}
\centering
     \includegraphics{~/Documents/GitHub/bEuclideanGeo/ch1/images/figureGroup2.2}
     \caption{二面体群$D_4$}\label{dihedral}
\end{minipage} 
\end{figure}

\section{3次元直交群とその有限部分群}
\subsection{3次元特殊直交群の有限部分群}
\begin{Thm}[Euler]
$S\neq E\in \SO(3)$ は原点を通るある直線を軸とする回転になる．
\end{Thm}
\begin{proof}
\begin{figure}[htb]
     \centering
     \includegraphics{~/Documents/GitHub/bEuclideanGeo/ch1/images/figureGroup2.3}
     \caption{}\label{}
\end{figure}

$S\neq E\in \SO(3)$ の固有多項式
\[
     g_S(t)
     =
     |tE-S|
     =t^3 + \alpha t^2+ \beta t +\gamma
\]
は$3$次の実係数多項式なので必ず一つは実数解$\lambda_1$を持つ．$\lambda\in \mathbf R$を実数の固有値，$\mathbf x\neq \mathbf 0$をその固有ベクトルとすると，
\[
\begin{split}
     &(T\mathbf x,T\mathbf x)
     =
     (\lambda\mathbf x,\lambda\mathbf x)
     =
     \lambda^2(\mathbf x,\mathbf x)\\
     &
     =(\mathbf x,\mathbf x)
\end{split}
\]
最後で$T$が直交変換であることを使った．$(\mathbf x,\mathbf x)>0$なので
$\lambda^2=1$，つまり，
\[
     \lambda=\pm1
\]
が分かる．


$\lambda_2$, $\lambda_3\in \mathbf C$を他の固有値とする．
$\lambda_2$が実数ならば$\lambda_3$も実数で，
\[
     |S|=-g_S(0)=\lambda_1\lambda_2\lambda_3
\]
$S\in \SO(3)$なので$\lambda_1\lambda_2\lambda_3=1$となり，
適当に順番を入れ換えると
\begin{equation}\label{specialCase}
     \lambda_1=\lambda_2=\lambda_3=1,\quad
     \text{あるいは}\quad
     \lambda_1=1, \lambda_2=\lambda_3=-1\tag{*}
\end{equation}
のいずれか．
%最初の場合は$S=E$となり，仮定に反する．次の場合は
%固有値$1$の固有ベクトルを軸とする$\pi$回転である．
\emptyline

$\lambda_2$が実数でないならば$\lambda_3$はその共役複素数
となり，全ての固有値が相異なるので，複素正則行列で対角化出来て
\[
     1 = |S|=\lambda_1\lambda_2\bar \lambda_2
\]
$\lambda_2\bar \lambda_2>0$より$\lambda_1=1$が分かる．
固有値$1$の固有ベクトルを$\mathbf x\neq\mathbf 0 $と置く．
$\mathbf y\in\mathbf R^3$が$(\mathbf y,\mathbf x)=0$を満たすとすると
\[
     (S\mathbf y,\mathbf x)
     =
     (S\mathbf y,S\mathbf x)
     =
     (\mathbf y,\mathbf x)
     =
     0
\]
となり，$S$は平面$T=\{\mathbf y\mid (\mathbf y,\mathbf x)=0\}$を保つことが分かる．$S|_T$は$2$次元直交群の元だが，$T$をユークリッド平面と同一視して
\[
     \det S|_T=
     \lambda_2\bar \lambda_2=1
\]
\[
     \det S|_T=
     \begin{cases}
          \lambda_2\bar \lambda_2=1& \lambda_2\notin \mathbf R\\
          (\pm1)^2=1  & \lambda_2,\lambda_3  \text{ is }\eqref{specialCase}
     \end{cases}
\]
なので，$T$は原点を通り$\mathbf x$を通る直線を軸とする回転になる．\eqref{specialCase} の最後の場合は$T$内の$\pi$回転の場合であり，最初の場合は$S=E$となり仮定からこの場合は起こらない．
\end{proof}

\begin{Prop}\label{transOfAxis}
$A$が$\mathbf u\neq \mathbf 0$を通る直線を回転の軸とする回転とすると，
任意の回転$B$に対して
$BAB^{-1}$は$B\mathbf u$を通る直線を回転の軸とする回転である．
\end{Prop}
\begin{proof}
$BAB^{-1}$が回転であることは$\SO(3)$が群であることからあきらか．
\[
     BAB^{-1}B\mathbf u=BA\mathbf u=B\mathbf u
\]
なので，$BAB^{-1}$は$B\mathbf u$を通る直線を回転の軸とする回転である．
\end{proof}

\begin{Q}
$\SO(3)$の有限部分群の分類
\end{Q}

上の定理より，$S\neq E\in \SO(3)$ に対してある$\mathbf x\neq \mathbf 0\in \mathbf R^3$が存在して
\[ 
     S\mathbf x=\mathbf x
\]
となる．
このような$\mathbf x$を$S$の\textcolor{red}{極}という．
$G$を$\SO(3)$の有限部分群とする．
$S^2(1)=\{\mathbf x\in \mathbf R^3\mid \|\mathbf x\|=1\}$
と置き，
\[
     P
     =
     \{\mathbf x\in S^2(1)\mid 
     \text{ある}S\neq E\in G\text{についての極, i.e.~}
     S\mathbf x=\mathbf x\}
\]
とすると$P$は有限集合である．

また， $\mathbf x\in P$ を$S\in G$の極とする．
命題~\ref{transOfAxis}より$T\in G$に対して$T\mathbf x$は
$TST^{-1}\in G$の極になる．
よって$P\ni\mathbf x\mapsto T\mathbf x\in  P$が定まる．明らかに$T, U,E\in G$に対して
\[
     (UT)\mathbf x=U(T\mathbf x),\quad 
     E\mathbf x=\mathbf x \quad (\mathbf x\in  P)
\]
を満たす．このようなとき$G$は$P$に\textcolor{red}{作用}するといい
$G\curvearrowright P$ と表す．

$\mathbf x\in P$ に対して
\[
     G\mathbf x
     =
     \{S\mathbf x\in P\mid S\in G\}
\]
を$G$\term[]{軌道}という．また
\[
     \stab(\mathbf x)
     =
     \{ S\in G\mid
          S\mathbf x=\mathbf x
     \}
\]
と置き\term[]{安定部分群}(stabilizer, or  isotropy group)という．確かに
$T$, $U\in \stab(\mathbf x)$に対して
\[
     TU\mathbf x=T\mathbf x=\mathbf x,\quad
     \mathbf x=T^{-1}T\mathbf x=T^{-1}\mathbf x
\]
なので$G$の部分群になっている．
もちろん，$\stab(\mathbf x)$は正規部分群とは限らないないので商群は定まらないが，以下のように剰余集合$G/H$が定まる．
\emptyline

$H$が有限群$G$の部分群であるとき，
$g,h\in G$に対して
\[
      g\sim h \quad 
     \Leftrightarrow \quad
     g^{-1}h\in H
\]
と関係を定めると同値関係になり，剰余集合$G/H$が定まる：
これは
\[
     G = g_1H\cup g_2H\cup \dots \cup g_m H,\quad
     g_iH\cap g_jH = \emptyset 
     \quad i\neq j
\]
という$G$の類別を与える．特に，
\begin{equation}\label{num}
     |G|=|G/H||H|
\end{equation}
ただし，$|X|$は有限集合$X$の元の個数を表す，行列式と同じ記号だが混乱は生じないと思う．また，$|G:H|=|G/H|$と表すこともある．
\emptyline

$G$が$\SO(3)$の有限部分群で$H=\stab(\mathbf x)$のとき
\[
     f:G/\stab(\mathbf x)\ni g H\mapsto g\mathbf x\in G \mathbf x
\]
が全単射となる．
なぜなら，$gH=g'H$ ならば $g^{-1}g'\in H=\stab(\mathbf x) $
$g^{-1}g'\mathbf x=\mathbf x$となり，$g'\mathbf x=g\mathbf x$
となる．よってwell-defined. 
逆に$g'\mathbf x=g\mathbf x$ならば$g^{-1}g'\mathbf x=\mathbf x$なので $g^{-1}g'\in H=\stab(\mathbf x) $．$H$は群なので$gH=g'H$ となり，単射となる．全射は定義から明か．
\emptyline

$P$を$G$軌道に分割する：
\[
     P = G\mathbf x_1\cup G\mathbf x_2\cup \dots \cup G\mathbf x_k,\quad
     G\mathbf x_i\cap G\mathbf x_j = \emptyset 
     \quad i\neq j
\]
\begin{Prop}
$G\neq \{E\}$
のとき，$k=2,3$
\end{Prop}

どのような状況になっているのかを理解するため例を挙げる：
まず図~\ref{cyclic}の正方形を固定する回転群では正方形の中心に直交する二つの単位ベクトルが2つの異なる軌道となっている．
図~\ref{dihedral}の正方形を固定する二面体群では
正方形の中心に直交する二つの単位ベクトルが一つの軌道で
正方形の頂点全体が一つの軌道，辺の中点全体ももう一つの軌道で3っつの軌道があることが分かる．
ここでは，極の長さは$1$ではなく，二面体の表面の方が見やすいのでそうする．

立方体の像を固定する回転の作る群$G$を考える(図 \ref{cube1},  \ref{cube2},  \ref{cube3})．
ここでも極の長さは$1$ではなく，立方体の表面の方が見やすいのでそうする．$\stab$の位数が三つそれぞれ異なっているので，これらの$G$軌道は3つの異なる軌道である．これらの軌道中の極が回転で移り合うことも確かめてほしい．
\begin{figure}[htb]
   \centering
   \begin{minipage}{0.4\columnwidth}
      \centering
      \includegraphics{~/Documents/GitHub/bEuclideanGeo/ch1/images/figureGroup2.4}
        \caption{平行な二つの面の中心を通る直線に沿っての$\pi/2$回転の極}\label{cube1}
   \end{minipage} 
   \begin{minipage}{0.4\columnwidth}
      \centering
     \includegraphics{~/Documents/GitHub/bEuclideanGeo/ch1/images/figureGroup2.5}
     \caption{最も遠い二つの頂点を通る直線に沿っての$2\pi/3$回転の極}\label{cube2}
   \end{minipage} 
\end{figure}
\begin{figure}[htb]
   \centering
   \includegraphics{~/Documents/GitHub/bEuclideanGeo/ch1/images/figureGroup2.6}
   \caption{平行な最も遠い辺の中心を通る直線に沿っての$\pi$回転の極}\label{cube3}
\end{figure}


\begin{proof}
$n=|G|\ge 2$と置く．
\[
     \tilde P 
     = 
     \{
          (S,\mathbf x)\mid S\neq E\in G, 
          \mathbf x\in S^2(1)\text{は} S \text{の極}
     \}
\]
とすると，各$S\neq E\in G$毎に極は二つあるので
\begin{equation}\label{numTP}
     |\tilde P| = 2(n-1)
\end{equation}
この式の中の$-1$は単位元を除くことに対応している．
$i=1,\dots,k$について
\[
     \alpha_i = |\stab(\mathbf x_i)|,\quad
     \beta_i = |G\mathbf x_i|
\]
と置くと，
\eqref{num} より
\begin{equation}\label{numTP1}
     n=\alpha_i \beta_i
\end{equation}
各$\mathbf x\in G\mathbf x_i$について$\stab(\mathbf x)$は
平面の回転が作る群で，（上で考えた作用で）すべて同形なので
\[
     \stab(\mathbf x)
     \sim
     \{
          E, T, \dots, T^{\alpha_i-1}
     \}
\]
となり
\begin{equation}\label{numObt}
     |\tilde P| 
     =
     \sum_{i=1}^k (\alpha_i-1)\beta_i
     =
     \sum_{i=1}^k (n - \beta_i)
\end{equation}
ここでも単位元を除いて数えていて，\eqref{numTP1}も使った．

 \eqref{numTP}を使うと
\begin{equation}\label{numObt1}
     2(n-1) = (n - \beta_1)+\dots+(n - \beta_k)
\end{equation}
なので, $n$で割り\eqref{numTP1}を使うと
\begin{equation}\label{numDenomi}
     2-\frac2n = (1-\frac1{\alpha_1})+\dots+ (1-\frac1{\alpha_k})
\end{equation}
$n\ge 2$より，左辺は$\ge 1$．もし$k=1$ならば右辺は$<1$となり矛盾が生じる．
よって$k\ge 2$が分かる，特に，これは$\alpha_i\ge 2$であることを意味する．

$k\ge 4$とすると$\alpha_i\ge 2$と\eqref{numDenomi}から
\[
\begin{split}
     2-\frac2n 
     &\ge 
          (1-\frac1{\alpha_1})+(1-\frac1{\alpha_2})
          +(1-\frac1{\alpha_3})+ (1-\frac1{\alpha_4})\\
     & \ge 
          \frac12+\frac12+\frac12+\frac12
     =2
\end{split}
\]
となり矛盾．よって$k=2$, $3$が分かる．
\end{proof}

以下ではこれまでの議論を進めて分類を与える：
$k=2$とすると， \eqref{numObt1}より
\[
     2n-2 = (n- \beta_1)+(n-\beta_2)
\]
つまり
\[
     2=\beta_1+\beta_2
\]
となり，$\beta_1=\beta_2=1$が分かる．\eqref{num} より
\[
     \alpha_1=\alpha_2=n
\]
$T$を$\mathbf x$を極とする
$\frac{2\pi}{n}$の回転として
\[
     G=
     \{
          E, T, \dots, T^{n-1}
     \}
     =\langle T\rangle
\]
でこれが巡回群の場合で，
確かに$\mathbf x$と$-\mathbf x$が異なる$G$軌道である．

次に$k=3$のとき\eqref{numDenomi}より
\[
     2-\frac2n 
     =
     (1-\frac1{\alpha_1})+(1-\frac1{\alpha_2})+(1-\frac1{\alpha_3})
\]
纏めると
\[
     1+\frac2n 
     =
     \frac1{\alpha_1}+\frac1{\alpha_2}+\frac1{\alpha_3}
\]
$\alpha_1\le\alpha_2\le \alpha_3$と仮定して良いのでそうする．
もし，$\alpha_1\ge 3$とすると右辺は$\le 1$だが，
左辺は$\ge1$なので矛盾する.
したがって，$\alpha_1=2$がわかる．よって
\[
     \frac1{2}+\frac2n 
     =
     \frac1{\alpha_2}+\frac1{\alpha_3}
\]

$\alpha_2\ge 4$とすると右辺は$\le \frac14+\frac14=\frac1{2}$となり左辺は$>\frac1{2}$なので矛盾する.つまり，$\alpha_2=2$, $3$のいずれか．

$\alpha_2=2$のとき
\[
     \frac2n 
     =
     \frac1{\alpha_3}
\]
となり，$2\alpha_3=n$がわかる．これが正多面体群$D_{\alpha_3}$の場合である．

$\alpha_2=3$のとき
\[
     \frac1{6}+\frac2n 
     =
     \frac1{\alpha_3}
\]
$\alpha_3\ge 6$とすると右辺は$\le \frac16$となり左辺は$>\frac1{6}$なので矛盾する.つまり，$\alpha_3=3$, $4$, $5$のいずれか．
$\alpha_3=3$のとき
\[
     \frac2n 
     =
     \frac13-\frac1{6}
     =
     \frac1{6}\quad\Rightarrow n=12
\]
$\alpha_3=4$のとき
\[
     \frac2n 
     =
     \frac1{4}-\frac1{6}
     =
     \frac1{12}\quad\Rightarrow n=24
\]
$\alpha_3=5$のとき
\[
     \frac2n 
     =
     \frac1{5}-\frac1{6}
     =
     \frac1{30}\quad\Rightarrow n=60
\]
\emptyline

\begin{figure}[htb]
\centering
\begin{minipage}{0.4\columnwidth}
\centering
     \includegraphics{~/Documents/GitHub/bEuclideanGeo/ch1/images/figureGroup2.7}
     \caption{$\pm2\pi/3$回転}
\label{}
\end{minipage} 
\begin{minipage}{0.4\columnwidth}
\centering
     \includegraphics{~/Documents/GitHub/bEuclideanGeo/ch1/images/figureGroup2.8}
     \caption{$\pi$回転}\end{minipage} 
\end{figure}
%\begin{figure}[htb]
%\centering
%\begin{minipage}{0.4\columnwidth}
%\centering
%     \includegraphics[height=5cm]{dodecahedron.jpg}
%     \caption{正十二面体}\label{}
%\end{minipage} 
%\begin{minipage}{0.4\columnwidth}
%\centering
%     \includegraphics{~/Documents/GitHub/bEuclideanGeo/ch1/images/figureGroup2.9}
%     \caption{正十二面体の模式図}\label{}
%\end{minipage} 
%\end{figure}

\noindent
纏めると以下の表のようになる：
\begin{table}[h]
  \centering
  \begin{tabular}{l|llllllll} 
       \noalign{\hrule height 1pt}
      $G$ &$k$& $\alpha_1$ &$\alpha_2$&$\alpha_3$& $\beta_1$&$\beta_2$&$\beta_3$& $n$ \\ \hline \hline
     $C_n$&2&$n$&$n$&$ $  & $1$&$1$ & $ $& $n$\\ \hline
     $D_m$&3&$2$&$2$&$m$  & $m$&$m$ & $2$& $2m$\\ \hline
     $ T$&3&$2$& $3$ & $3$ &$6$ &$4$ &$4$ & $12$ \\ \hline
     $W$&3& $2$&$3$  & $4$ & $12$&$8$&$6$ &$24$ \\ \hline
     $I$& 3&$2$& $3$ & $5$&$30$ &$20$ &$12$ & $60$\\ \hline
  \noalign{\hrule height 1pt}
\end{tabular}
\caption{$\SO(3)$の有限部分群の分類}
\end{table}

ユークリッドの原論の第13巻にあるように，プラトンの正多面体は

\begin{quote}
正四面体，立方体，正八面体，正一二面体，正二十面体
\end{quote}

\noindent
の$5$種類である．
立方体と正八面体は面と頂点を入れ換える変換で双対であり，正一二面体と正二十面体も双対である．正四面体は自己双対であり，上の表はそれらの像を保つ回転のなす有限部分群の分類に対応していて，
$T$を正四面体群，$Q$を正八面体%（立方体）
群，$I$を正二十面体%（正一二面体）
群という．



%以下の参考文献が%（数理内からアクセスすれば）
%\href{https://link.springer.com/}{ここから}
%ダウンロード可能．
%\begin{thebibliography}{9}
%\item
%L.C.~Grove, C.T~Benson, \textit{Finite reflection groups 2nd Ed.}, GTM99, 1996, Springer
%\end{thebibliography}

%\begin{Quiz}$3$
%$Q$が$S_4$と同形になることを示せ．
%\end{Quiz}
%\begin{Quiz}$4$
%$I$が$A_5$と同形になることを示せ．ただし，$A_5$は5次の交代群である．
%\end{Quiz}


まず，$D_m$が$\Ot(2)$の有限部分群では${}^t(\cos\theta,\sin\theta)$の向きの軸に関する鏡映
\[
     R
     =
     \begin{pmatrix}
          \cos 2\theta &\sin 2\theta\\
          \sin 2\theta&-\cos 2\theta\\
     \end{pmatrix}
\]
から構成されていて，${}^t(\cos\theta,\sin\theta)$が固有値$1$, ${}^t(\sin\theta,-\cos\theta)$が固有値$-1$の固有ベクトルである．
行列$R$の行列式は$-1$である．
$\SO(3)$では
\[
     \tilde R
     =
     \begin{pmatrix}
          \cos 2\theta &\sin 2\theta&0\\
          \sin 2\theta&-\cos 2\theta&0\\
          0&0&-1\\
     \end{pmatrix}
\]
と拡張されていて，
${}^t(\cos\theta,\sin\theta,0)$が固有値$1$, 
${}^t(\sin\theta,-\cos\theta,0)$と ${}^t(0,0,1)$が固有値$-1$の固有ベクトルである．
この行列$\tilde R$の行列式は$1$である．
これは${}^t(\cos\theta,\sin\theta,0)$の向きの軸に関する$\pi$回転である．つまり$2$次元の鏡映は$3$次元の$\pi$回転になることを注意しておく．

表の
$T$を\term[]{正四面体群}(tetrahedron group), $W$を\term[]{正八面体群}(octahedron group)，$I$を\term[]{正二十面体群}(icosahedron group)という．
これらが名前を付けた正多面体の（等長）群に一致することを示すことが出来るが，
時間の節約のため，ここでは正四面体群$T$の場合のみ確かめる：
$\alpha_2=3$のとき$\beta_2=4$なので $\mathbf x_2\in P$が
\[
     G\mathbf x_2
     =
     \{\mathbf y_1, \mathbf y_2,\mathbf y_3,\mathbf y_4\}
     \ni \mathbf x_2
\]
$S\neq E\in \stab(\mathbf x_2)$は$\alpha_2=3$より，$2\pi/3$回転である．回転の極は$\pm \mathbf x_2$なので，全ての
$\mathbf y_i\in G\mathbf x_2$を止めることはない．
$S$の極が$\mathbf y_1$だとすると適当に順番を付け替えて
\[
     S\mathbf y_1=\mathbf y_1,\quad
     S\mathbf y_2=\mathbf y_3,\quad
     S\mathbf y_3=\mathbf y_4,\quad
     S\mathbf y_4=\mathbf y_2
\]
$S$は直交行列なので
\[
\begin{split}
     \|\mathbf y_1-\mathbf y_2\|
     &=
     \|S\mathbf y_1-S\mathbf y_2\|\\
     &=
     \|\mathbf y_1-\mathbf y_3\|
     =
     \|\mathbf y_1-\mathbf y_4\|
\end{split}
\]
$S\in \stab(\mathbf y_2)$, $\stab(\mathbf y_3)$, $\stab(\mathbf y_3)$について同じ議論をすると，$\{\mathbf y_1, \mathbf y_2,\mathbf y_3,\mathbf y_4\}$が正四面体をなすこと，
すべての二点が$G$の元で移りあうこと分かる．
面については，
\[
     \mathbf z_{ijk}=\tfrac13(\mathbf y_i+\mathbf y_j+\mathbf y_k)\quad (i< j<k)
\]
と置けば，最初の$S$については
\[
     S\mathbf z_{234}=\mathbf z_{234}
\]
となるので，$S$は正四面体の頂点と面の重心を通る回転と一致することが分かる．

\begin{figure}[htb]
   \centering
   \begin{minipage}{0.4\columnwidth}
      \centering
       \includegraphics{~/Documents/GitHub/bEuclideanGeo/ch1/images/figureGroup3.1}
       \caption{$2\pi/3$回転}
      \label{}
   \end{minipage} 
   \begin{minipage}{0.4\columnwidth}
      \centering
       \includegraphics{~/Documents/GitHub/bEuclideanGeo/ch1/images/figureGroup3.11}
       \caption{$2\pi/3$回転}
      \label{}
   \end{minipage} 
\end{figure}
\begin{figure}[htb]
   \centering
   \begin{minipage}{0.4\columnwidth}
      \centering
       \includegraphics{~/Documents/GitHub/bEuclideanGeo/ch1/images/figureGroup3.2}
       \caption{$\pi$回転}
   \end{minipage} 
   \begin{minipage}{0.4\columnwidth}
      \centering
       \includegraphics{~/Documents/GitHub/bEuclideanGeo/ch1/images/figureGroup3.12}
       \caption{$\pi$回転}
   \end{minipage} 
\end{figure}

次に$S\neq E\in \stab(\mathbf x_1)$は$\alpha_1=2$より，$\pi$回転である．$G\mathbf x_2$に$S$を作用させて$S\mathbf y_i=\mathbf y_i$とすると，$S\in \stab(\mathbf x_2)$となり，$2\pi/3$回転となり矛盾．よって
\[
     S\mathbf y_i\neq \mathbf y_i
\]
$SS=E$なので，適当に名前を付け替えると
\[
     S\mathbf y_1=\mathbf y_2,\quad
     S\mathbf y_2=\mathbf y_1,\quad
     S\mathbf y_3=\mathbf y_4,\quad
     S\mathbf y_4=\mathbf y_3
\]
したがって
\[
     \mathbf z_{ij}=\tfrac12(\mathbf y_i+\mathbf y_j)\quad (i< j)
\]
と置くと
\begin{gather*}
     S\mathbf z_{12}=\mathbf z_{12},\quad 
     S\mathbf z_{34}=\mathbf z_{34},\\
     S\mathbf z_{13}=\mathbf z_{24},\quad 
     S\mathbf z_{14}=\mathbf z_{23},\quad
     S\mathbf z_{23}=\mathbf z_{14},\quad 
     S\mathbf z_{24}=\mathbf z_{13}
\end{gather*}
となり，正四面体の辺の変換と一致することが分かり，$G=T$が分かる．ただし，ここで等しいとは四頂点の配置分の自由度を除いて等しいということである．

%\begin{Quiz}$3$
%表~\ref{k=3}
%の$D_m$場合も正二面体群$D_m$に一致することを示せ．
%\end{Quiz}
\begin{Quiz}
   表~\ref{classification}
   の$W$の場合も正八面体群$W$に一致することを示せ．
\end{Quiz}
\begin{Quiz}
   表~\ref{classification}
   の$I$の場合も正二十面体群$I$に一致することを示せ．
\end{Quiz}
\begin{Quiz}
   $n$次交代群\footnote{$S_n$の偶置換全体のなす群，補足に詳しい説明を与えてある}を $A_n$と書く．
   正四面体群$T$が$4$次交代群$A_4$と同形になることを示せ．
\end{Quiz}
\begin{Quiz}
   正八面体群
   $W$が$4$次対称群$S_4$と同形になることを示せ．
\end{Quiz}

\begin{Quiz}
   正二十面体群
   $I$が$5$次交代群$A_5$と同形になることを示せ．
\end{Quiz}


\subsection{$3$次直交群の有限部分群}
\begin{Thm}
   $S\in \Ot(3)$が$|S|=-1$をみたすとき，$S$ はある平面についての鏡映
   と，それに直交する軸についての回転になる．
\end{Thm}
\begin{proof}
   Euler の定理の証明とほぼ同じ．
   \begin{figure}[htb]
      \centering
      \begin{minipage}{0.4\columnwidth}
         \centering
         \includegraphics{~/Documents/GitHub/bEuclideanGeo/ch1/images/figureGroup3.3}
         \caption{鏡映と回転}
      \end{minipage} 
      \begin{minipage}{0.4\columnwidth}
         \centering
         \includegraphics{~/Documents/GitHub/bEuclideanGeo/ch1/images/figureGroup3.5}
         \caption{$Z$}\label{Z}
   \end{minipage} 
   \end{figure}
   
   $S\in \Ot(3)$ の固有多項式
   \[
        g_S(t)
        =
        |tE-S|
        =t^3 + \alpha t^2+ \beta t +\gamma
   \]
   は$3$次の実係数多項式なので必ず一つは実数解$\lambda_1$を持つ．$\lambda\in \mathbf R$を実数の固有値，$\mathbf x\neq \mathbf 0$をその固有ベクトルとすると，
   \[
        (\mathbf x,\mathbf x)
        =(S\mathbf x,S\mathbf x)
        =
        (\lambda\mathbf x,\lambda\mathbf x)
        =
        \lambda^2(\mathbf x,\mathbf x)
   \]
   $(\mathbf x,\mathbf x)>0$なので
   $\lambda^2=1$，つまり，
   \[
        \lambda=\pm1
   \]
   が分かる．
   
   $\lambda_2$, $\lambda_3\in \mathbf C$を他の固有値とする．
   $\lambda_2$が実数ならば$\lambda_3$も実数で，
   \[
        |S|=-g_S(0)=\lambda_1\lambda_2\lambda_3
   \]
   $|S|=-1$なので$\lambda_1\lambda_2\lambda_3=-1$となり，
   適当に順番を入れ換えると
   \begin{equation}\label{specialCase}
        \lambda_1=-1, \lambda_2=\lambda_3=1,\quad
        \text{あるいは}\quad
        \lambda_1=\lambda_2=\lambda_1=-1 \tag{*}
   \end{equation}
   のいずれか．
   
   $\lambda_2$が実数でないならば$\lambda_3$はその共役複素数
   となり，全ての固有値が相異なるので，複素正則行列で対角化出来て
   \[
        -1 = |S|=\lambda_1\lambda_2\bar \lambda_2
   \]
   $\lambda_2\bar \lambda_2>0$より$\lambda_1=-1$, $\lambda_2\bar \lambda_2=1$が分かる．
   
   \emptyline
   固有値$-1$の固有ベクトルを$\mathbf x\neq\mathbf 0 $と置く．
   $\mathbf y\in\mathbf R^3$が$(\mathbf y,\mathbf x)=0$を満たすとすると
   \[
        (S\mathbf y,\mathbf x)
        =
        (S\mathbf y,-S\mathbf x)
        =
        -(\mathbf y,\mathbf x)
        =
        0
   \]
   となり，$S$は平面$T=\{\mathbf y\mid (\mathbf y,\mathbf x)=0\}$を保つことが分かる．$S|_T$は$2$次元直交群の元だが，$T$をユークリッド平面と同一視して
   \[
        \det S|_T=
        \begin{cases}
             \lambda_2\bar \lambda_2=1& \lambda_2\notin \mathbf R\\
             (\pm1)^2=1  & \lambda_2,\lambda_3  \text{ is }
                \eqref{specialCase}
        \end{cases}
   \]
   なので，$T$は原点を通り$\mathbf x$を通る直線を軸とする回転になる．つまり
   \[
        S(\xi \mathbf x+ \eta \mathbf y)
        = -\xi \mathbf x+ \eta S_{\theta}\mathbf y
   \]
   ただし，$S_{\theta}$は$ \mathbf x$を中心にした$\theta$回転．
\end{proof}
\eqref{specialCase}の最初のケースは, $\mathbf x$についての単なる($0$回転の)鏡映で，
最後の場合に対応する行列は
\[
     Z=
     \begin{pmatrix}
          -1&0&0\\
          0&-1&0\\
          0&0&-1\\
     \end{pmatrix}
     =-E\in \Ot(3)
\]
である(図~\ref{Z})．つまり，固有空間が$\mathbf R^3$で全ての$\mathbf x\neq \mathbf 0$が固有値$-1$の固有ベクトルである．
また，任意の$R\in \Ot(3)$の元に対して
\[
     ZR = -R = RZ
\]
つまり，$Z$は全ての元と可換である．

\begin{Q}
   $\Ot(3)$の有限部分群の分類？
\end{Q}
$G$を$\Ot(3)$の有限部分群とする．$G\subset \SO(3)$の場合は既に求まったので，$R\in G\setminus  \SO(3)$が存在するとし，$H=G\cap \SO(3)$と置く．これは
$G$の部分群である．
\[
     H =\{S_0=E, S_1, \dots , S_{m-1}\}
\]
と置くと，
\[
     G\supset \{E, S_1, \dots , S_{m-1},R,RS_1, \dots , RS_{m-1}\}
\]
であるが，$R'\in G\setminus  \SO(3)$とすると
\[
     |R^{-1}R'|=|R^{-1}||R'|=(-1)(-1)=1
\]
なので$R^{-1}R'\in H$となり，ある$S_i\in H$が存在して
\[
     R^{-1}R'=S_i\quad\Rightarrow\quad R'=RS_i
\]
よって等号
\begin{equation}\label{G}
     G= \{E, S_1, \dots , S_{m-1},R,RS_1, \dots , RS_{m-1}\}
\end{equation}
が示せた．
\emptyline

case 1) : $Z\in G$のとき$R=Z$として
\[
     G= \{E, S_1, \dots , S_{m-1},Z,ZS_1, \dots , ZS_{m-1}\}
     =:\bar H
\]
と表す．つまり，$G$は$H$に$Z$を加えた群$\bar H$になる．
\emptyline

case 2) : $Z\notin G$のとき
\eqref{G}より
\[
     G^{+}:= \{E, S_1, \dots , S_{m-1},ZRS_0, \dots , ZRS_{m-1}\}
\]
と表すと，$|ZRS_i|=|Z||R||S_1|=-1\times(-1)\times 1=1$より，
$\SO(3)$の有限部分群で$H$を部分群として含むものである．
\[
     (ZRS_i)^{-1}H(ZRS_i)=H,\quad S_i^{-1}HS_i=H
\]
より，指数$2$の正規部分群であり，$H\lhd G^{+}$と表す．

逆に$K$が$\SO(3)$の有限部分群で，その指数$2$の正規部分群$H$を持つとする．
\[
     H=\{E, S_1, \dots , S_{m-1}\},\quad
     K=\{E, S_1, \dots , S_{m-1},T_{0},\dots, T_{m-1}\}
\]
と表すとき
\[     
     K]H:= \{E, S_1, \dots , S_{m-1},ZT_{0},\dots, ZT_{m-1}\}
\]
と置くと$\SO(3)$に含まれない$\Ot(3)$の有限部分群が得られる．

\emptyline

$|K|=2|H|$であることなどに注意して，$\SO(3)$の有限部分群の分類
（表 \ref{classification}）
を見ると，
指数$2$の正規部分群をもつ$\SO(3)$の有限部分群の分類が分かる：
\[
     C_m \lhd C_{2m},\quad C_m\lhd  D_m,\quad 
     D_m \lhd D_{2m},\quad T\lhd W
\]


以上を纏めると：

\begin{Thm}
   $\Ot(3)$の有限部分群は以下のいずれか：
   \begin{enumerate}
      \item
      \[
           C_m,\quad D_m,\quad  T,\quad W,\quad  I
      \]
      
      \item
      \[
           \bar C_m,\quad \bar D_m,\quad, \bar T \quad \bar W, \quad  \bar I
      \]

      \item
      \[
           C_{2m}]C_m,\quad D_m]C_m,\quad D_{2m}]D_m,\quad W]T
      \]
   \end{enumerate}
\end{Thm}
\begin{figure}[htb]
   \centering
   \begin{minipage}{0.4\columnwidth}
      \centering
      \includegraphics{~/Documents/GitHub/bEuclideanGeo/ch1/images/figureGroup3.6}
      \caption{$D_3\lhd D_6$}\label{$D_3D_6$}
   \end{minipage} 
   \begin{minipage}{0.4\columnwidth}
      \centering
      \includegraphics{~/Documents/GitHub/bEuclideanGeo/ch1/images/figureGroup3.4}
      \caption{$T\lhd W$}\label{$TW$}
   \end{minipage} 
\end{figure}
\begin{Quiz}
$\SO(3)$の群$K$とその指数$2$の正規部分$H$群の組が上で上げたものに以外にないことを示し，上で上げたものが確かに指数$2$の正規部分群であることを示せ（図 \ref{$D_3D_6$}, \ref{$TW$} も参考にせよ）．
\end{Quiz}


%以下の\cite{GB}が%（数理内からアクセスすれば）
%\href{https://link.springer.com/}{ここから}
%ダウンロード可能．
%
%\emptyline
%次回から平面結晶群の分類に入る．
%
%\begin{thebibliography}{9}
%   \bibitem{GB}
%   L.C.~Grove, C.T~Benson, \textit{Finite reflection groups 2nd Ed.}, GTM99, 1996, Springer
%\end{thebibliography}



\section{$\SO(3)$の幾つかの有限部分群の構造}
$\SO(3)$の有限部分群と置換群・交代群の関係を調べる．
置換群・交代群については

正二面体群$D_3$が$3$次対称群$S_3$と同形になること：
正三角形の頂点を$1$, $2$, $3$と呼ぶことにする．
$\sigma\in D_3$は頂点を頂点に移すので，頂点$1$ が頂点$i_1$に移り，
頂点$2$ が頂点$i_2$に移り，頂点$3$ が頂点$i_3$に移るとすると
\[
     D_3\ni \sigma\mapsto 
     \begin{pmatrix}1&2&3\\i_1&i_2&i_3\end{pmatrix}\in S_3
\]
で写像が定まる．\eqref{S3}の任意の$S_3$の置換に対応する正三角形の等長変換が存在する．つまり，
二つの$2\pi/3$回転$\begin{pmatrix}1&2&3\end{pmatrix}$, 
$\begin{pmatrix}3&2&1\end{pmatrix}$と単位元が偶置換で，
それら全体は$C_3$と同形となる．
奇置換$\begin{pmatrix}1&2\end{pmatrix}$, $\begin{pmatrix}2&3\end{pmatrix}$, $\begin{pmatrix}3&1\end{pmatrix}$ が$\pi$回転となっていて，単位元とそれぞれの元が$C_2$と同形な$3$つの部分群となっている．
$|S_3|=|D_3|=6$より全単射になる．等長変換の合成は頂点の変換の合成となるので準同形であることも明か．
表~\ref{again}の$D_3$の場合$\pi$回転の$3$個の軸が表の（異なる二つの$G$軌道の）$3+3=6$個の極に対応し，$2\pi/3$回転の$1$個の軸が表の（図では表現できていないが中心の点の上下にある，一つの$G$軌道の）$2$個の極に対応することが確認が出来る．
\begin{table}[h]
     \centering
     \begin{tabular}{l|lllllll} 
          \noalign{\hrule height 1pt}
          $G$ & $\alpha_1$ &$\alpha_2$&$\alpha_3$& $\beta_1$&$\beta_2$&$\beta_3$& $n$ \\ \hline \hline
          $C_m$&$m$ &$m$& &1 &1  & & $m$\\ \hline
          $D_m$&$2$&$2$&$m$  & $m$&$m$ & $2$& $2m$\\ \hline
          $ T$&$2$& $3$ & $3$ &$6$ &$4$ &$4$ & $12$ \\ \hline
          $W$& $2$&$3$  & $4$ & $12$&$8$&$6$ &$24$ \\ \hline
          $I$& $2$& $3$ & $5$&$30$ &$20$ &$12$ & $60$\\ \hline
      \noalign{\hrule height 1pt}
   \end{tabular}
   \caption{$\SO(3)$の有限部分群の分類(再掲)}\label{again}
\end{table}

\begin{figure}[hb]
   \centering
   \begin{minipage}{0.4\columnwidth}
      \centering
        \includegraphics{~/Documents/GitHub/bEuclideanGeo/ch1/images/figureGroup3.131}
        \caption{$3$本の$\pi$回転の軸}\label{quadPi}
   \end{minipage} 
   \begin{minipage}{0.4\columnwidth}
      \centering
        \includegraphics{~/Documents/GitHub/bEuclideanGeo/ch1/images/figureGroup3.13}
        \caption{$4$本の$2\pi/3$回転の軸}\label{quad2Pi/3}
   \end{minipage} 
\end{figure}

正四面体群$T$が$4$次交代群$A_4$と同形であること：
正四面体の頂点を$1$, $2$, $3$, $4$と呼ぶことにする．
表 \ref{again} によれば$\alpha_1=2$のときは三本の軸の回りの$\pi$回転で図 \ref{quadPi}のように
\[
     \begin{pmatrix}1&2\end{pmatrix}
     \begin{pmatrix}3&4\end{pmatrix}, \quad
     \begin{pmatrix}1&3\end{pmatrix}
     \begin{pmatrix}2&4\end{pmatrix}, \quad
     \begin{pmatrix}1&4\end{pmatrix}
     \begin{pmatrix}2&3\end{pmatrix}
\]
の三つの偶置換からなる．これに単位元を付け加えたのがクライン群である．
図 \ref{quad2Pi/3}のように
4本の軸の回りの$\pm 2\pi/3$回転なので$2\times 4=8$個の偶置換
\[
\begin{split}
     &\begin{pmatrix} 2& 3&  4\end{pmatrix},\quad
     \begin{pmatrix} 1& 3&  4\end{pmatrix},\quad
     \begin{pmatrix} 1& 2&  4\end{pmatrix},\quad
     \begin{pmatrix} 1& 2&  3\end{pmatrix},\\
     &\begin{pmatrix} 4& 3&  2\end{pmatrix},\quad
     \begin{pmatrix} 4& 3&  1\end{pmatrix},\quad
     \begin{pmatrix} 4& 2&  1\end{pmatrix},\quad
     \begin{pmatrix} 3& 2&  1\end{pmatrix}
\end{split}
\]
からなるので，$T$の群の位数は$ 1 + 3 +8 = 12$となる．$A_4$は$12$個からなるので$T$から$A_4$への同型が得られる．

\begin{figure}[hb]
   \centering
   \begin{minipage}{0.4\columnwidth}
      \centering
        \includegraphics{~/Documents/GitHub/bEuclideanGeo/ch1/images/figureGroup3.21}
     \caption{正八面体と立方体}\label{octaQube}
   \end{minipage} 
   \begin{minipage}{0.4\columnwidth}
      \centering
        \includegraphics{~/Documents/GitHub/bEuclideanGeo/ch1/images/figureGroup3.22}
        \caption{立方体と正八面体}\label{qubeOcta}
   \end{minipage} 
\end{figure}


正八面体群$W$が$4$次対称群$S_4$と同形になること：
立方体と正八面体はそれぞれの面を点に取換えることで得られるので互いに双対である（図~\ref{octaQube}, \ref{qubeOcta}）ので立方体を使って考える．
図 \ref{cubePi/2} のように正四面体を立方体に埋め込み，
正四面体の頂点を上と同様に$1$, $2$, $3$, $4$と名付ける．
更に残りの頂点には$1$と重心を通る直線が通るもう一つの頂点を同じ$1$と名付けようにして$1$, $2$, $3$, $4$と名付ける．
そして，$1$と命名された$2$つの頂点の集合を$\bar 1$のように表し，
集合$\{\bar 1, \bar 2 ,\bar 3 ,\bar 4\}$の置換と$W$の元が同一視されることを以下で確かめる．
図 \ref{cube2Pi/3}を見ると
例えば$1$と$1$を結ぶ直線を軸にする$\pm 2\pi/3$回転は$1$を止めるので
\[
\begin{split}
     &\begin{pmatrix}\bar 2&\bar 3&  \bar 4\end{pmatrix},\quad
     \begin{pmatrix} \bar 1& \bar 3&  \bar 4\end{pmatrix},\quad
     \begin{pmatrix} \bar 1& \bar 2&  \bar 4\end{pmatrix},\quad
     \begin{pmatrix} \bar 1& \bar 2&  \bar 3\end{pmatrix},\\
     &\begin{pmatrix} \bar 4& \bar 3&  \bar 2\end{pmatrix},\quad
     \begin{pmatrix} \bar 4& \bar 3&  \bar 1\end{pmatrix},\quad
     \begin{pmatrix} \bar 4& \bar 2&  \bar 1\end{pmatrix},\quad
     \begin{pmatrix} \bar 3& \bar 2&  \bar 1\end{pmatrix}
\end{split}
\]
の$8$個の三換（偶置換）からなる．
図 \ref{cubePi/2} を見ると$\pi/2$回転から生成され
その回転軸が三本あり，そのうちの半数が正四面体を固定する$\pi$回転$3$個であり，置換としては
\[
     \begin{pmatrix}\bar 1&\bar 2\end{pmatrix}
     \begin{pmatrix}\bar 3&\bar 4\end{pmatrix},\quad
     \begin{pmatrix}\bar 1&\bar 3\end{pmatrix}
     \begin{pmatrix}\bar 2&\bar 4\end{pmatrix},\quad
     \begin{pmatrix}\bar 1&\bar 4\end{pmatrix}
     \begin{pmatrix}\bar 2&\bar 3\end{pmatrix}
\]
である．これらに恒等写像を加えた$8+3+1=12$個の偶置換全体が正四面体を固定する等長群をなし，$A_4$と同型である．
さらに正四面体を固定しない$\pi/2$回転は
\[
\begin{split}
     &\begin{pmatrix}\bar 1&\bar 2&\bar 3&\bar 4\end{pmatrix},\quad
     \begin{pmatrix}\bar 4&\bar 3&\bar 2&\bar 1\end{pmatrix},\quad
     \begin{pmatrix}\bar 1&\bar 2&\bar 4&\bar 3\end{pmatrix},\\
     &\begin{pmatrix}\bar 3&\bar 4&\bar 2&\bar 1\end{pmatrix},\quad
     \begin{pmatrix}\bar 1&\bar 3&\bar 2&\bar 4\end{pmatrix},\quad
     \begin{pmatrix}\bar 4&\bar 2&\bar 3&\bar 1\end{pmatrix}
\end{split}
\]
という$6$個の長さ$4$の循環置換からなる奇置換である．
\begin{figure}[th]
   \centering
   \begin{minipage}{0.4\columnwidth}
      \centering
        \includegraphics{~/Documents/GitHub/bEuclideanGeo/ch1/images/figureGroup3.15}
        \caption{$4$本の$2\pi/3$回転の軸}\label{cube2Pi/3}
   \end{minipage} 
   \begin{minipage}{0.4\columnwidth}
      \centering
        \includegraphics{~/Documents/GitHub/bEuclideanGeo/ch1/images/figureGroup3.14}
        \caption{$3$本の$\frac{\pi}{2}$回転の軸} \label{cubePi/2}
      \label{}
   \end{minipage} 
\end{figure}

図 \ref{cubePi} の$\pi$回転は
$\begin{pmatrix}
      4\\2
\end{pmatrix}
= 6$ 個の奇置換(互換)
$\begin{pmatrix}\bar i&\bar j\end{pmatrix}$($i\neq j$)を成している．
言い換えると$S_4$の奇置換全体$12$個を含んでいて，
$W$から$S_4$への同型写像が得られた．

\begin{figure}[hb]
   \centering
   \begin{minipage}{0.4\columnwidth}
   \centering
        \includegraphics{~/Documents/GitHub/bEuclideanGeo/ch1/images/figureGroup3.16}
        \caption{$6$本の$\pi$回転の軸}\label{cubePi}
   \end{minipage} 
   \begin{minipage}{0.4\columnwidth}
   \centering
        \includegraphics{~/Documents/GitHub/bEuclideanGeo/ch1/images/figureGroup3.18}
        \caption{正二十面体と正十二面体}\label{IcosaDual1}
   \end{minipage} 
\end{figure}
\emptyline


\begin{figure}[hb]
   \centering
   \begin{minipage}{0.4\columnwidth}
      \centering
        \includegraphics{~/Documents/GitHub/bEuclideanGeo/ch1/images/figureGroup3.19}
        \caption{正十二面体と正二十面体}\label{IcosaDual2}
   \end{minipage} 
   \begin{minipage}{0.4\columnwidth}
      \centering
        \includegraphics{~/Documents/GitHub/bEuclideanGeo/ch1/images/figureGroup3.20}
     \caption{正十二面体と立方体}\label{dodecaQube}
   \end{minipage} 
\end{figure}

正二十面体群
$I$が$5$次交代群$A_5$と同形になること：
正十二面体と正二十面体はそれぞれの面を点に取換えることで得られるので互いに双対である（図~\ref{IcosaDual1}, \ref{IcosaDual2}）．
そこで
$I$を正十二面体\footnote{作図へのコメント：$\tau$を黄金比とし$(\pm \tau,\pm \tau,\pm \tau)$を立方体の$8$つの頂点の座標とする．この立方体を含む正十二面体の立方体の頂点でない$20-8=12$個の頂点の座標を$(\pm1, 0, \pm(1+\tau))$, $(\pm(1+\tau), \pm1, 0)$, $(0, \pm(1+\tau), \pm1)$とすると，一辺$2$の正十二面体が描ける．正二十面体はそれらの頂点の重心を取れば良い，}の等長群と見て考察する．まず，正十二面体に立方体
を埋め込むことを考える（図~\ref{dodecaQube}）．立方体の辺は$12$個あり，その全てが正十二面体の$12$個の面の中に一本ずつ入っていることを確認してほしい．
次に，正二十面体の$20$個の頂点に$1, 2, 3, 4, 5$ を$4$個ずつ置く．
ただし，各正五角形は全ての異なる$1, 2, 3, 4, 5$が配置されているものを考える．そのため，まず，図~\ref{nameDodeca}の赤の正五角形に$1, 2, 3, 4, 5$を反時計回りに配置し，それらの頂点と未だ名前のない頂点を含む辺の頂点を$ \mod 5$で$2$を足した数で名付けると，新しい頂点には反時計回りに$1, 2, 3, 4, 5$を配置されていることがわかる．未だ名前が一個だけ足りない各正五角形の頂点には足りない$1, 2, 3, 4, 5$を配置することが出来て，これらの新しい頂点にも反時計回りに$1, 2, 3, 4, 5$が配置されていることもわかる．これらの新しい頂点と未だ名前のない頂点含む辺の頂点に$ \mod 5$で$2$を足した数で名付けると，新しい頂点には反時計回りに$1, 2, 3, 4, 5$を配置している
ことがわかる．これで欲しい頂点の配置が求まった．

\begin{figure}[hb]
   \centering
   \begin{minipage}{0.4\columnwidth}
      \centering
        \includegraphics{~/Documents/GitHub/bEuclideanGeo/ch1/images/figureGroup3.26}
        \caption{正十二面体の頂点名}\label{nameDodeca}
   \end{minipage} 
   \begin{minipage}{0.4\columnwidth}
      \centering
        \includegraphics{~/Documents/GitHub/bEuclideanGeo/ch1/images/figureGroup3.27}
     \caption{$\bar 5$を止める$\pi$回転}\label{piRote}
   \end{minipage} 
\end{figure}

\begin{figure}[hb]
     \centering
        \begin{minipage}{0.4\columnwidth}
        \centering
        \includegraphics{~/Documents/GitHub/bEuclideanGeo/ch1/images/figureGroup3.28}
        \caption{$\bar 5$を止める$\pm2\pi/3$回転}\label{2pi/3Rote}
   \end{minipage} 
     \begin{minipage}{0.4\columnwidth}
           \centering
           %\includegraphics{~/Documents/GitHub/bEuclideanGeo/ch1/images/figureGroup3.23}
          \includegraphics{~/Documents/GitHub/bEuclideanGeo/ch1/images/figureGroup3.25}
           \caption{ $\frac{2\pi}{5}$ 回転と $(1\,2\,3\,4\,5)$}\label{qubes5}
   \end{minipage} 
\end{figure}


$1$と命名された頂点の集合を$\bar 1$のように表し，
集合$\{\bar 1, \bar 2 ,\bar 3 ,\bar 4, \bar 5\}$の偶置換と$I$の元が同一視されることを以下で確かめる．
図~\ref{dodecaQube}の立方体の$8$個の頂点には$5$の頂点が$4$個，
残りに$1, 2, 3, 4$が一個づつ配置されている（図 \ref{piRote}, \ref{2pi/3Rote}）．
$I$の元でこの立方体を動かさないのは図 \ref{piRote}, \ref{2pi/3Rote}に示した$\pi$回転と$\pm2\pi/3$回転であり，$\pi$回転は$\{\bar 1, \bar 2 ,\bar 3 ,\bar 4\}$の独立な互換の積
\[
     \begin{pmatrix}\bar 1&\bar 2\end{pmatrix}
     \begin{pmatrix}\bar 3&\bar 4\end{pmatrix},\quad
     \begin{pmatrix}\bar 1&\bar 4\end{pmatrix}
     \begin{pmatrix}\bar 2&\bar 3\end{pmatrix},\quad
     \begin{pmatrix}\bar 1&\bar 3\end{pmatrix}
     \begin{pmatrix}\bar 2&\bar 4\end{pmatrix}
\]
であり，$\pm 2\pi/3$回転も$\{\bar 1, \bar 2 ,\bar 3 ,\bar 4\}$の（偶置換の）三環
\[
\begin{split}
     &\begin{pmatrix}\bar 2&\bar 3& \bar 4\end{pmatrix},\quad
     \begin{pmatrix}\bar 1&\bar 3& \bar 4\end{pmatrix},\quad
     \begin{pmatrix}\bar 1&\bar 2& \bar 4\end{pmatrix},\quad
     \begin{pmatrix}\bar 1&\bar 2& \bar 3\end{pmatrix},\\
     &\begin{pmatrix}\bar 4&\bar 3& \bar 2\end{pmatrix},\quad
     \begin{pmatrix}\bar 4&\bar 3& \bar 1\end{pmatrix},\quad
     \begin{pmatrix}\bar 4&\bar 2& \bar 1\end{pmatrix},\quad
     \begin{pmatrix}\bar 3&\bar 2& \bar 1\end{pmatrix}
\end{split}
\]
を生成する．これらは全ての頂点の組に矛楯なく作用していることを確かめてほしい．上で述べたようにそれらは正四面体群$T$を成し，
$\bar 5$を動かさない$4$次交代群$A_4$と同形となっている．
%更に正十二面体の中の異なる5つの立方体が取れ（図 \ref{qubes1}, \ref{qubes2}），異なる立方体毎に5つの正四面体群を作っている．


%\begin{figure}[hb]
%   \centering
%   \begin{minipage}{0.4\columnwidth}
%      \centering
%        \includegraphics{~/Documents/GitHub/bEuclideanGeo/ch1/images/figureGroup3.27}
%        \caption{正十二面体と立方体}
%   \end{minipage} 
%   \begin{minipage}{0.4\columnwidth}
%        \centering
%        \includegraphics{~/Documents/GitHub/bEuclideanGeo/ch1/images/figureGroup3.21}
%        \caption{$3$本の$\pi$回転の軸}\label{piRote}
%   \end{minipage} 
%\end{figure}

表 \ref{again}（とEuler の定理）によれば, $I$の合同変換には$\pi$回転と$2\pi/3$回転の他に
正5角形の二つの面を保つ$2\pi/5$回転がある．
特に図~\ref{qubes1}の正5角形を回転する回転の生成元として
\[
     \sigma
     =\begin{pmatrix}\bar 1&\bar 2&\bar 3& \bar 4&\bar 5\end{pmatrix}
\]
が取れる．これも全ての頂点の組に矛楯なく作用していて，偶置換であることもアミダクジの時と同じように示せる．
$\sigma^n$ $(n =1, 2, 3, 4)$によりが図~\ref{dodecaQube}の立方体を別の立方体に置換するかを示したのが図~\ref{qubes1}--\ref{qubes4} である．$n=1, 2, 3, 4, 5$として
$\bar n$を止める立方体を保存する$I$の部分群を$A_4^n$と置くと
$g=\sigma^n$の内部自己同形の作用で
\[
     A_4^5\ni \tau\mapsto  g\tau g^{-1}\in A_4^n\subset A_5
\]
となる．$A_4^5$ は$\bar 5$を止める全ての三環を含んでいるので，$n = 1, 2, 3, 4, 5$を動かすと$A_5$は全ての三環を含んでいることがわかる．
$A_5$が可解でないことを示したところでコメントしたように
$A_5$は三環で生成されるので，
$I$は$A_5$への全射となり，同型になる．

\begin{figure}[thb]
   \begin{minipage}{0.4\columnwidth}
           \centering
           \includegraphics{~/Documents/GitHub/bEuclideanGeo/ch1/images/figureGroup3.29}
           \caption{$\bar 1$を止める立方体}\label{qubes1}
   \end{minipage} 
   \begin{minipage}{0.4\columnwidth}
     \centering
      \includegraphics{~/Documents/GitHub/bEuclideanGeo/ch1/images/figureGroup3.30}
      \caption{$\bar 2$を止める立方体}\label{qubes2}
   \end{minipage} 
\end{figure}

\begin{figure}[t]
   \begin{minipage}{0.4\columnwidth}
     \centering
      \includegraphics{~/Documents/GitHub/bEuclideanGeo/ch1/images/figureGroup3.31}
      \caption{$\bar 3$を止める立方体}\label{qubes3}
   \end{minipage} 
   \begin{minipage}{0.4\columnwidth}
           \centering
           \includegraphics{~/Documents/GitHub/bEuclideanGeo/ch1/images/figureGroup3.32}
           \caption{$\bar 4$を止める立方体}\label{qubes4}
   \end{minipage} 
\end{figure}

他の方法でも$I$は$A_5$と同型になることを示すことができる．
まず，$60$個の$I$の元のうち
三環は$5$個から$3$個の数字を選ぶことに対応するので
$\begin{pmatrix}
     5\\
     3
\end{pmatrix}
=10$種類あり，
それらはそれぞれ$3$次循環群をなすので，三環は合計$10\times (3-1)=20$個ある．
独立な互換の積は$5$個の数字から$1$個除き，
残りの$4$個の数字を$3$種のペアに分けることに対応するので合計
$5\times 3=15$個ある．
ここから更に単位元（偶置換）$1$個をそぞいた$60-36=24$個が残っている．
そこで$\bar 1, \bar 2 ,\bar 3 ,\bar 4, \bar 5$の長さ$5$の循環置換（これは偶置換である）を考える．
これらは$5!$個ありそうだが$\begin{pmatrix}\bar 1&\bar 2&\bar 3& \bar 4&\bar 5\end{pmatrix}$と$\begin{pmatrix}\bar 3& \bar 4&\bar 5&\bar 1&\bar 2\end{pmatrix}$
は同じ置換を表しているので$5!/5=4!=24$個ある．
そこで図 \ref{nameDodeca}の$12$ある面のうち
表の$6$個の面の頂点の配置を反時計向きに$1$から始まるように並べると，異なる
長さ$5$の循環置換
\[
\begin{split}
     &\begin{pmatrix}\bar 1&\bar 2&\bar 3&\bar 4&\bar 5\end{pmatrix},\quad
     \begin{pmatrix}\bar 1&\bar 3&\bar 5&\bar 4&\bar 2\end{pmatrix},\quad
     \begin{pmatrix}\bar 1&\bar 5&\bar 3&\bar 2&\bar 4\end{pmatrix},\\
     &\begin{pmatrix}\bar 1&\bar 4&\bar 3&\bar 5&\bar 2\end{pmatrix},\quad
     \begin{pmatrix}\bar 1&\bar 3&\bar 2&\bar 5&\bar 4\end{pmatrix},\quad
     \begin{pmatrix}\bar 1&\bar 5&\bar 2&\bar 4&\bar 3\end{pmatrix}
\end{split}
\]
が見つかる，これらはそれぞれ$I$の$2\pi/5$の回転であり，それらは$5$次の巡回群を生成している．
各面の（一番遠い）反対側の面が逆の循環置換になっていることもチェック出来る．
これらの生成する長さ$5$の循環置換の個数は$6\times(5-1)=24$個（表 \ref{5cyclic}）であり，
これらが残りの$24$個の偶置換となっている．
このことから$I$は$A_5$と同型になることが分かる．
\begin{table}[h]
     \centering
     \begin{tabular}{l|lll} 
          \noalign{\hrule height 1pt}
          $\tau$ 
          & $\begin{pmatrix}\bar 1&\bar 2&\bar 3&\bar 4&\bar 5\end{pmatrix}$ 
          &$\begin{pmatrix}\bar 1&\bar 3&\bar 5&\bar 4&\bar 2\end{pmatrix}$
          &$\begin{pmatrix}\bar 1&\bar 5&\bar 3&\bar 2&\bar 4\end{pmatrix}$ \\ 
          $\tau^2$ 
          & $\begin{pmatrix}\bar 1&\bar 3&\bar 5&\bar 2&\bar 4\end{pmatrix}$ 
          &$\begin{pmatrix}\bar 1&\bar 5&\bar 2&\bar 3&\bar 4\end{pmatrix}$
          &$\begin{pmatrix}\bar 1&\bar 3&\bar 4&\bar 5&\bar 2\end{pmatrix}$ \\ 
          $\tau^3$ 
          & $\begin{pmatrix}\bar 1&\bar 4&\bar 2&\bar 5&\bar 3\end{pmatrix}$ 
          &$\begin{pmatrix}\bar 1&\bar 4&\bar 3&\bar 2&\bar 5\end{pmatrix}$
          &$\begin{pmatrix}\bar 1&\bar 2&\bar 5&\bar 4&\bar 3\end{pmatrix}$ \\   
          $\tau^4$ 
          & $\begin{pmatrix}\bar 1&\bar 5&\bar 4&\bar 3&\bar 2\end{pmatrix}$ 
          &$\begin{pmatrix}\bar 1&\bar 2&\bar 4&\bar 5&\bar 3\end{pmatrix}$
          &$\begin{pmatrix}\bar 1&\bar 4&\bar 2&\bar 3&\bar 5\end{pmatrix}$ \\  
          $\tau^5$ 
          & $\qquad\quad\varepsilon$  &  $\qquad\quad\varepsilon$ & 
               $\qquad\quad\varepsilon$  \\
          \noalign{\hrule height 1pt}
   \end{tabular}
\end{table}

\begin{table}[h]
     \centering
   \begin{tabular}{l|lll} 
          \noalign{\hrule height 1pt}
          $\tau$ 
          &$\begin{pmatrix}\bar 1&\bar 4&\bar 3&\bar 5&\bar 2\end{pmatrix}$ 
          &$\begin{pmatrix}\bar 1&\bar 3&\bar 2&\bar 5&\bar 4\end{pmatrix}$
          &$\begin{pmatrix}\bar 1&\bar 5&\bar 2&\bar 4&\bar 3\end{pmatrix}$ \\ 
          $\tau^2$ 
          & $\begin{pmatrix}\bar 1&\bar 3&\bar 2&\bar 4&\bar 5\end{pmatrix}$ 
          &$\begin{pmatrix}\bar 1&\bar 2&\bar 4&\bar 3&\bar 5\end{pmatrix}$
          &$\begin{pmatrix}\bar 1&\bar 2&\bar 3&\bar 5&\bar 4\end{pmatrix}$ \\ 
          $\tau^3$ 
          &$\begin{pmatrix}\bar 1&\bar 5&\bar 4&\bar 2&\bar 3\end{pmatrix}$ 
          &$\begin{pmatrix}\bar 1&\bar 5&\bar 3&\bar 4&\bar 2\end{pmatrix}$
          &$\begin{pmatrix}\bar 1&\bar 4&\bar 5&\bar 3&\bar 2\end{pmatrix}$ \\   
          $\tau^4$ 
          &$\begin{pmatrix}\bar 1&\bar 2&\bar 5&\bar 3&\bar 4\end{pmatrix}$ 
          &$\begin{pmatrix}\bar 1&\bar 4&\bar 5&\bar 2&\bar 3\end{pmatrix}$
          &$\begin{pmatrix}\bar 1&\bar 3&\bar 4&\bar 2&\bar 5\end{pmatrix}$ \\  
          $\tau^5$ 
          &$\qquad\quad\varepsilon$  &  $\qquad\quad\varepsilon$ & 
               $\qquad\quad\varepsilon$  \\
          \noalign{\hrule height 1pt}
   \end{tabular}
   \caption{24個の長さ5の巡回置換}\label{5cyclic}
\end{table}


%\newpage

%\end{document}




%\documentclass[a4paper,12pt]{amsart}
%\usepackage[dvipdfmx]{graphicx,xcolor,hyperref}%graphics
%\usepackage{float}
%\usepackage{amsmath}
%\usepackage{verbatim}
%\usepackage{latexsym}
%\usepackage{amssymb}
%\usepackage{ccicons}
%%\usepackage{emath}
%%\usepackage{maplestd2e}
%%\usepackage{fancybox}%enclose by box
%\hypersetup{% hyperrefオプションリスト
%setpagesize=false,
% bookmarksnumbered=true,
% bookmarksopen=true,
% colorlinks=true,
%linkcolor=red,
%citecolor=blue,
%} 
%
%\renewcommand{\figurename}{図}
%\renewcommand{\proofname}{証明}
%\renewcommand{\tablename}{表}
%
%\newtheorem*{PP}{授業の目標}
%\newtheorem*{Thm}{定理}
%\newtheorem*{Prop}{命題}
%\newtheorem*{Q}{考えたい課題}
%\newtheorem*{Quiz}{問題}
%
%\theoremstyle{definition}
%\newtheorem*{Def}{定義}
%
%\theoremstyle{remark}
%\newtheorem*{ex}{例}
%\newtheorem*{exQ}{例題}
%
%\def\emptyline{\vspace{12pt}}
%
%\DeclareMathOperator{\Ker}{Ker}
%\DeclareMathOperator{\Image}{Im}
%\DeclareMathOperator{\id}{id}
%\DeclareMathOperator{\Is}{Is}
%\begin{document}
%%\pagestyle{empty}
%\begin{flushright}
%
%2024年11月6日
%
%九大・数理 大津幸男
%\end{flushright}
%
%\begin{center}
%\large{微分幾何学大意・数学特論4\small{
%\href{https://creativecommons.org/licenses/by-nc-nd/4.0/}{\ccbyncnd}}\par ---ユークリッド空間の合同変換群とその離散群 1---}
%\end{center}

\chapter{ユークリッド平面の結晶群と平坦曲面の分類}
\section{ユークリッド空間の離散合同変換群}
\subsection{ユークリッド空間の合同変換の特徴付け}
$\mathbf R^n=(\mathbf R^n,(,))$をユークリッド空間とする．ただし，$\mathbf x=(x_1,\dots,x_n),
	\mathbf y=(y_1,\dots,y_n)\in \mathbf R^n$ に対して
\[
	(\mathbf x,\mathbf y)
	= 
	x_1y_1+\dots +x_ny_n,\quad
	\|\mathbf x\|
	=
	\sqrt{(\mathbf x,\mathbf x)}
\]
であった．簡単な計算で
\begin{gather}
	\|\mathbf x+\mathbf y\|^2 + \|\mathbf x-\mathbf y\|^2
	=
	2(\|\mathbf x\|^2+ \|\mathbf y\|^2)\label{normNorm}\\
	\|\mathbf x+\mathbf y\|^2 -\|\mathbf x-\mathbf y\|^2
	=
	4(\mathbf x,\mathbf y)\label{normInnerP}\
\end{gather}
が分かる．
写像$F:\mathbf R^n\to \mathbf R^n$が\term[]{合同変換}とは
\begin{equation}\label{congruentT}
	  \|F(\mathbf x)-F(\mathbf y)\|
	  =
	  \|\mathbf x-\mathbf y\|
\end{equation}
が全ての$\mathbf x,\mathbf y\in \mathbf R^n$ に対して成り立つことである．

\begin{Prop}
$F:\mathbf R^n\to \mathbf R^n$が合同変換 $ \Leftrightarrow $ 
$A\in O(n)$, $\mathbf a\in \mathbf R^n$が存在して
\[
	F(\mathbf x)
	=
	A\mathbf x+\mathbf a
\]
（このような変換を\term[]{アファイン変換}と言った．）
\end{Prop}
\begin{proof}
$\mathbf a=F(\mathbf 0)$とし
\[
	T(\mathbf x):=F(\mathbf x)-\mathbf a
\]
と置くと$T(\mathbf 0)=\mathbf 0$．
\eqref{congruentT} より
\begin{equation}
	  \|T(\mathbf x)-T(\mathbf y)\|
	  =
	  \|\mathbf x-\mathbf y\|\quad\text{特に}\quad
	\|T(\mathbf x)\|
	  =
	  \|\mathbf x\|
\end{equation}
\eqref{normNorm}より
\[
\begin{split}
	\|T(\mathbf x)+T(\mathbf y)\|^2 
	&=
	2(\|T(\mathbf x)\|^2+ \|T(\mathbf y)\|^2)
	-\|T(\mathbf x)-T(\mathbf y)\|^2\\
	&=
	2(\|\mathbf x\|^2+ \|\mathbf y\|^2)
	-\|\mathbf x-\mathbf y\|^2\\
	&=
	\|\mathbf x+\mathbf y\|^2
\end{split}
\]
\eqref{normInnerP}より
\begin{equation}\label{orthogonal}
\begin{split}
	(T(\mathbf x),T(\mathbf y))
	&=\frac14
	(\|T(\mathbf x)+T(\mathbf y)\|^2 
	-\|T(\mathbf x)-T(\mathbf y)\|^2)\\
	&=\frac14
	(\|\mathbf x+\mathbf y\|^2 -\|\mathbf x-\mathbf y\|^2)
	=
	(\mathbf x,\mathbf y)
\end{split}
\end{equation}

次に，$T$が線形写像であることを示す：
$\{\mathbf e_i\mid i=1,\dots,n\}$を$\mathbf R^n$の標準的な正規直交基底とする．
$\mathbf a_i=T(\mathbf e_i)$ とおくと，\eqref{orthogonal}より
$\{\mathbf a_i\mid i=1,\dots,n\}$も正規直交基底．
$\mathbf x=\sum_{i=1}^n x_i\mathbf e_i$に対して\eqref{orthogonal}より
\[
	(\mathbf a_i,T(\mathbf x))
	=
	(T(\mathbf e_i),T(\mathbf x))
	=
	(\mathbf e_i,\mathbf x)
	=
	x_i
\]
よって
\[
	T(\mathbf x)=\sum_{i=1}^n x_i\mathbf a_i
\]
と表せるので$T$は線形写像．\eqref{orthogonal}より$T$は直交変換で，ある$A\in O(n)$が存在して$T(\mathbf x)=A\mathbf x$．つまり，
\[
	F(\mathbf x)
	=
	A\mathbf x+\mathbf a
\]
\end{proof}
次に
\[
	\Is(\mathbf R^n)
	=\{F:\mathbf R^n\to \mathbf R^n\mid \text{合同変換}\}
\]
と置く． $F\in \Is(\mathbf R^n)$について，上の証明の表記を使い
\[
	t(F)=\mathbf a,\quad r(F)=A\in O(n)
\]
と表す．
\begin{Prop}
$\Is(\mathbf R^n)$は写像の合成で群になる．
\end{Prop}
\begin{proof}
\[
	F(\mathbf x)
	=
	A\mathbf x+\mathbf a, \quad
	G(\mathbf x)
	=
	B\mathbf x+\mathbf b \in \Is(\mathbf R^n)
\]
とすると
\[
	G\circ F(\mathbf x)
	=
	B(A\mathbf x+\mathbf a)+\mathbf b
	=
	B A\mathbf x +B\mathbf a+\mathbf b
\]
で$BA\in O(n)$なのでこれは合同変換．
\[
	\id(\mathbf x)=\mathbf x = E \mathbf x+ \mathbf 0
\]
が合同変換の単位元で，$A^{-1}\in O(n)$なので
\[
	F^{-1}(\mathbf x)
	=
	A^{-1}\mathbf x - A^{-1}\mathbf a
\]
と置くと$F^{-1}\circ F=F\circ F^{-1}=\id$となり逆元となり, 
$\Is(\mathbf R^n)$は群になる（結合律は自らか）．
\end{proof}


\subsection{$n=2$の合同変換の特徴付け}
以下$F\in \Is(\mathbf R^2)$の特徴付けを与える：

Case 1) $r(F)\in SO(2)$ のとき, 第1回の授業で示したように，回転行列
\[
	r(F) = 
	\begin{pmatrix}
		\cos \theta &-\sin \theta\\
		\sin \theta& \cos \theta\\
	\end{pmatrix}
	=:
	S_\theta
\]
となる．$\theta=0$のとき$S_\theta=E$なので
\[
	F(\mathbf x)
	=
	\mathbf x+\mathbf a
\]
これは\term[]{平行移動}である．

$0<\theta<2\pi $のとき$F$はある点$\mathbf x_0$を中心とした\term[]{回転}となる．実際
\[
	S_\theta\mathbf x_0+\mathbf a=\mathbf x_0
	\Rightarrow
	(E-S_\theta)\mathbf x_0=\mathbf a
\]
とすると
\[
\begin{split}
	|E-S_\theta|
	&=
	\begin{vmatrix}
		1-\cos \theta &\sin \theta\\
		-\sin \theta& 1-\cos \theta\\
	\end{vmatrix}
	=
	(1-\cos \theta)^2+\sin^2 \theta\\
	&=
	2(1-\cos\theta)
	\neq
	0\quad (0<\theta<2\pi)
\end{split}
\]
なので，
\[
	\mathbf x_0=(E-S_\theta)^{-1}\mathbf a
\]
となり，一意的に$\mathbf x_0$が定まる．このとき
\[
\begin{split}
	F(\mathbf x)-\mathbf x_0
	&=
	S_\theta\mathbf x+\mathbf a-\mathbf x_0\\
	&=
	S_\theta\mathbf x+(E-S_\theta)\mathbf x_0-\mathbf x_0
	=
	S_\theta(\mathbf x-\mathbf x_0)
\end{split}
\]
なので，$F$は$\mathbf x_0$を中心とした角度$\theta$の回転になる．
\begin{figure}[htb]
   \centering
   \begin{minipage}{0.4\columnwidth}
      \centering
      \includegraphics{~/Documents/GitHub/bEuclideanGeo/ch2/images/figureGroup4.1}
      \caption{平行移動}
   \end{minipage} 
   \begin{minipage}{0.4\columnwidth}
      \centering
      \includegraphics{~/Documents/GitHub/bEuclideanGeo/ch2/images/figureGroup4.2}
      \caption{回転}
   \end{minipage} 
\end{figure}


\begin{Thm}[回転角の加法定理]
角度$\alpha$の回転と角度$\beta$の回転の合成は
角度$\alpha+\beta$の回転である．ただし，回転の中心は異なっていていい．また，角度$0$の回転とは無限遠点を中心にした回転，つまり，平行移動のことである．
\end{Thm}
\begin{proof}
\[
	F(\mathbf x)
	=
	S_\alpha\mathbf x+\mathbf a, \quad
	G(\mathbf x)
	=
	S_\beta\mathbf x+\mathbf b 
\]
とすると
\[
	G\circ F(\mathbf x)
	=
	S_\beta S_\alpha\mathbf x +S_\beta\mathbf a+\mathbf b
\]
$S_\beta S_\alpha=S_{\alpha+\beta}$なので角度$\alpha+\beta$の回転である．
\end{proof}
\begin{itemize}
\item
平行移動を回転の一つと見ないと，命題の場合分けが複雑になるだけである．
また群論の立場からも平行移動を回転の一つとみることは自然である．

\item
平行移動を無限遠点を中心にした角度$0$の回転と見なす意味を確認する：
\[
	F(\mathbf x)=
	S_\theta\mathbf x+\mathbf a\to  \mathbf x+\mathbf a \quad
	 (\theta \to 0)
\]
であり
\[
\begin{split}
	\mathbf x_0&=(E-S_\theta)^{-1}\mathbf a\\
	&=
	\frac1{2(1-\cos\theta)}
	\begin{pmatrix}
		1-\cos \theta &-\sin \theta\\
		\sin \theta& 1-\cos \theta\\
	\end{pmatrix}\mathbf a\\
	&=
	\begin{pmatrix}
		\frac12 &-\frac{\sin \theta}{2(1-\cos\theta)}\\
		\frac{\sin \theta}{2(1-\cos\theta)}& \frac12\\
	\end{pmatrix}\mathbf a
\end{split}
\]
なので，$\mathbf a={}^t(a,b)\neq \mathbf 0$と表すと，Maclaurin展開により$\mathbf x_0$は$\theta \to +0$で
$\frac1{\theta}$の早さで$(-b,a)$の方角へ，$\theta \to -0$で
$\frac1{\theta}$の早さで$(b,-a)$の方角へ遠ざかることがわかる（図~\ref{limitRot}）．
\end{itemize}
\begin{figure}
   \begin{minipage}{0.4\columnwidth}
      \centering
      \includegraphics{~/Documents/GitHub/bEuclideanGeo/ch2/images/figureGroup4.3}
      \caption{回転の極限}\label{limitRot}
      \end{minipage} 
   \begin{minipage}{0.4\columnwidth}
      \centering
      \includegraphics{~/Documents/GitHub/bEuclideanGeo/ch2/images/figureGroup4.4}
      \caption{鏡映}\label{ref1}
   \end{minipage} 
\end{figure}

Case 2) $r(F)\in O(2)\setminus SO(2)$ とすると, 第1回の授業で示したように，
\[
	r(F)=
	\begin{pmatrix}
		\cos 2\theta &\sin 2\theta\\
		\sin 2\theta&-\cos 2\theta\\
	\end{pmatrix}
	=:R_{\theta}
\]
つまり，原点を通る直線$\ell$についての\term[]{鏡映}となる．$\ell$の向きの単位ベクトルを$\mathbf e$, $\ell$の単位法線ベクトルを$\mathbf f$とすると
\[
	\mathbf e
	=\begin{pmatrix}
		\cos\theta\\
		\sin\theta\\
	\end{pmatrix},\quad
	\mathbf f
	=\begin{pmatrix}
		-\sin\theta\\
		\cos\theta\\
	\end{pmatrix}
\]
となり(図~\ref{ref1}, \ref{ref2}) $\mathbf e\perp \mathbf f$で
\[
	R_{\theta}\mathbf e = \mathbf e,\quad
	R_{\theta}\mathbf f = -\mathbf f
\]
\begin{figure}[htb]
   \centering
   \begin{minipage}{0.4\columnwidth}
      \centering
      \includegraphics{~/Documents/GitHub/bEuclideanGeo/ch2/images/figureGroup4.5}
      \caption{鏡映}\label{ref2}
   \end{minipage} 
   \begin{minipage}{0.4\columnwidth}
      \centering
      \includegraphics{~/Documents/GitHub/bEuclideanGeo/ch2/images/figureGroup4.6}
      \caption{並進鏡映}\label{paraRef}
   \end{minipage} 
\end{figure}
この$F$について
\[
	F(\mathbf x_0)-\mathbf x_0 \perp \mathbf f
\]
が成り立つ$\mathbf x_0$が存在することを示す：このような$\mathbf x_0$は
\[
	0=(R_{\theta}\mathbf x_0+\mathbf a-\mathbf x_0,\mathbf f)
	=-(\mathbf x_0,\mathbf f) + (\mathbf a,\mathbf f) 
	-(\mathbf x_0,\mathbf f)
\]
をみたすので，
\[
	(\mathbf x_0,\mathbf f) = \frac12 (\mathbf a,\mathbf f) 
\]
をみたす必要がある．逆に
\[
	\mathbf x_0=\frac12 (\mathbf a,\mathbf f) \mathbf f
\]
と置くと
\[
\begin{split}
	F(\mathbf x_0)
	&=
	R_{\theta}\mathbf x_0+ \mathbf a
	=\frac12 (\mathbf a,\mathbf f) R_{\theta}\mathbf f+ \mathbf a\\
	&=
	-\frac12 (\mathbf a,\mathbf f) \mathbf f
	+ (\mathbf a,\mathbf f) \mathbf f+ (\mathbf a,\mathbf e) \mathbf e\\
	&=
	\frac12 (\mathbf a,\mathbf f) \mathbf f
	+(\mathbf a,\mathbf e) \mathbf e
	=
	\mathbf x_0+\alpha_0\mathbf e
\end{split}
\]
より，欲しかった性質をもつことが分かる．
ここで$\alpha_0=(\mathbf a,\mathbf e)$と置いた．

$\ell$を$\mathbf x_0$を通り $\mathbf e$の向きに平行な直線とする．
$\mathbf x=\mathbf x_0+\alpha\mathbf e+\beta \mathbf f$と表すと
\[
\begin{split}
	F(\mathbf x)
	&=
	F(\mathbf x_0 )+\alpha R_{\theta}\mathbf e + \beta R_{\theta}\mathbf f
	=\mathbf x_0+\alpha_0\mathbf e+\alpha \mathbf e -\beta \mathbf f\\
	&=\mathbf x_0+\alpha \mathbf e +\color{red}{\alpha_0\mathbf e-\beta \mathbf f}
\end{split}
\]
つまり，$F$は直線$\ell$ に関する\term[]{並進鏡映}である(図~\ref{paraRef})．
特に，$\alpha_0=0 $ のときが直線$\ell$ に関する鏡映である．

\[
\begin{split}
	F\circ F(\mathbf x)
	&=F(\mathbf x_0+(\alpha +\alpha_0)\mathbf e -\beta \mathbf f)\\
	&=
	\mathbf x_0+(\alpha +\alpha_0+\alpha_0)\mathbf e +\beta \mathbf f
	=\mathbf x +2\alpha_0\mathbf e
\end{split}
\]
(図~\ref{ref3})
なので
\begin{Prop}
並進鏡映を$2$回合成すると平行移動になる．
\end{Prop}
以上を纏めると：
\begin{Thm}[シャール]
   $\mathbf R^2$の合同変換は
   平行移動，ある点を中心とした回転，並進鏡映（鏡映）のいずれか．
\end{Thm}
\begin{figure}[htb]
   \centering
   \begin{minipage}{0.4\columnwidth}
      \centering
      \includegraphics{~/Documents/GitHub/bEuclideanGeo/ch2/images/figureGroup4.16}
      \caption{並進鏡映の合成}\label{ref3}
   \end{minipage} 
   \begin{minipage}{0.4\columnwidth}
      \centering
      \includegraphics{~/Documents/GitHub/bEuclideanGeo/ch2/images/figureGroup4.7}
      \caption{$\mathbf Z^n$}\label{paraRef}
   \end{minipage} 
\end{figure}

\subsection{離散部分群}
\begin{gather*}
	\Is^+(\mathbf R^n)
	=\{F\in\Is(\mathbf R^n)\mid r(F)\in SO(n)\},\\
	T(\mathbf R^n)
	=\{F\in\Is^+(\mathbf R^n)\mid \text{平行移動}\}
\end{gather*}
と置く．次は明らかだろう
\begin{Prop}
   $\Is^+(\mathbf R^n)$ は$\Is(\mathbf R^n)$の部分群，
   $T(\mathbf R^n)$ は$\Is^+(\mathbf R^n)$, $\Is(\mathbf R^n)$の部分群である．
\end{Prop}

\begin{Q}
   $O(n)$ときは有限部分群を考察しだが，
   $T(\mathbf R^n)$, $\Is(\mathbf R^n)$はコンパクトでないので，
   どのような部分群が有限部分群の代わりになるだろうか？
\end{Q}

$G\subset \Is(\mathbf R^n)$を部分群とする．
$G$が\term[]{離散部分群} とは，任意の$\mathbf x\in \mathbf R^n$に対して，$G$-軌道
\[
	G\mathbf x
	=
	\{F(\mathbf x)\in \mathbf R^n\mid F\in G\}
	\subset \mathbf R^n
\]
が離散集合であること，つまり
\[
	\forall \mathbf x\in \mathbf R^n, 
	\exists \varepsilon>0 \text{ such that }
	B(\mathbf x, \varepsilon)\cap G\mathbf x
	=\{\mathbf x\}
\]
ただし，$
	B(\mathbf x, \varepsilon)
	=\{\mathbf y\in \mathbf R^n\mid
	 	\|\mathbf y-\mathbf x\|<\varepsilon
	 \}
$である．

\begin{ex} $\mathbf Z$を整数全体とする．
\[
	\mathbf Z^n
	=
	\{
		(z_1,\dots,z_n)\mid z_1,\dots, z_n \in \mathbf Z
	\}
	\subset T(\mathbf R^n)
\]
つまり，$\mathbf x\mapsto \mathbf x+ \mathbf z$ ($ \mathbf z \in \mathbf Z^n$)により$\mathbf R^n$に作用する．
このとき$\mathbf x\in \mathbf R^n$ に対して
$0<\varepsilon<1 $ と取れば
\[
	B(\mathbf x, \varepsilon)\cap \mathbf Z^n\mathbf x
	=\{\mathbf x\}
\]
なので，$\mathbf Z^n$は離散部分群である．
\end{ex}

%
%Case 0)
%\[
%	G\cap T(\mathbf R^2)=\{\id\}
%\]
%シャールの定理より$G$は回転と鏡映のみからなる，言い換えると，($\id$以外の)平行移動や．(鏡映以外の)並進鏡映を含まない．このとき

%\emptyline
%
%これについては来週解説
%\emptyline
%
%
%以下の参考文献が（数理内からアクセスすれば）
%\href{https://link.springer.com/}{ここから}
%ダウンロード可能．
%\begin{thebibliography}{9}
%\item
%Viacheslav V. Nikulin, Igor R. Shafarevich,\textit{Geometries and Groups}, UTX, 1996, Springer
%\end{thebibliography}


\section{$\mathbf R^2$の離散合同変換群の分類}\label{class_2dim}
%\subsection{復習}
%$\mathbf R^n=(\mathbf R^n,(,))$をユークリッド空間とする．写像$F:\mathbf R^n\to \mathbf R^n$が\term[]{合同変換}とは
%\begin{equation}\label{congruentT}
%	  \|F(\mathbf x)-F(\mathbf y)\|
%	  =
%	  \|\mathbf x-\mathbf y\|
%\end{equation}
%が全ての$\mathbf x,\mathbf y\in \mathbf R^n$ に対して成り立つこと．
%\[
%	\Is(\mathbf R^n)
%	=\{F:\mathbf R^n\to \mathbf R^n\mid \text{合同変換}\}
%\]
%とは群であり，\term[]{等長変換群}という． 
%\begin{Prop}
%   $F:\mathbf R^n\to \mathbf R^n\in \Is(\mathbf R^n)$とすると，
%   $A\in O(n)$, $\mathbf a\in \mathbf R^n$が存在して
%   \[
%   	F(\mathbf x)
%   	=
%   	A\mathbf x+\mathbf a
%   \]
%\end{Prop}
%
%$F\in \Is(\mathbf R^n)$について，上の表記を使い
%\[
%	t(F)=\mathbf a,\quad r(F)=A\in O(n)
%\]
%と置く．
%\begin{gather*}
%%	\Is^+(\mathbf R^n)
%%	=\{F\in\Is(\mathbf R^n)\mid r(F)\in SO(n)\}, \\
%	T(\mathbf R^n)
%	=\{F\in\Is^+(\mathbf R^n)\mid \text{平行移動}\}
%\end{gather*}
%と置くと
%%$\Is^+(\mathbf R^n)$ は$\Is(\mathbf R^n)$の部分群，
%$T(\mathbf R^n)$ は$\Is(\mathbf R^n)$の部分群である．
%
%$r(F)\in SO(2)$のとき$F$は回転：
%\begin{Thm}[回転角の加法定理]
%   角度$\alpha$の回転と角度$\beta$の回転の合成は
%   角度$\alpha+\beta$の回転である．ただし，角度$0$の回転とは
%   無限遠点を中心にした回転，つまり，平行移動のことである．
%\end{Thm}
%
%\begin{Thm}[シャール]
%   $\mathbf R^2$の合同変換は
%   平行移動，ある点を中心とした回転，並進鏡映（鏡映）のいずれか．
%\end{Thm}
%
%\begin{figure}[htb]
%   \centering
%   \begin{minipage}{0.4\columnwidth}
%      \centering
%      \includegraphics{~/Documents/GitHub/bEuclideanGeo/ch2/images/figureGroup5.1}
%      \caption{平行移動}
%   \end{minipage} 
%   \begin{minipage}{0.4\columnwidth}
%      \centering
%      \includegraphics{~/Documents/GitHub/bEuclideanGeo/ch2/images/figureGroup5.2}
%      \caption{回転}
%\end{minipage} 
%\end{figure}
%
%\begin{figure}[htb]
%   \centering
%   \begin{minipage}{0.4\columnwidth}
%      \centering
%      \includegraphics{~/Documents/GitHub/bEuclideanGeo/ch2/images/figureGroup5.3}
%      \caption{並進鏡映}\label{ref3}
%   \end{minipage} 
%   \begin{minipage}{0.4\columnwidth}
%      \centering
%      \includegraphics{~/Documents/GitHub/bEuclideanGeo/ch2/images/figureGroup5.4}
%      \caption{鏡映}
%   \end{minipage} 
%\end{figure}
%
%
%$G\subset \Is(\mathbf R^n)$を部分群とする．
%$G$が\term[]{離散部分群}とは，任意の$\mathbf x\in \mathbf R^n$に対して，$G$-軌道
%\[
%	G\mathbf x
%	=
%	\{F(\mathbf x)\in \mathbf R^n\mid F\in G\}
%	\subset \mathbf R^n
%\]
%が離散集合であること，つまり
%\[
%	\forall \mathbf x\in \mathbf R^n, 
%	\exists \varepsilon>0 \text{ such that }
%	B(\mathbf x, \varepsilon)\cap G\mathbf x
%	=\{\mathbf x\}
%\]
%
%\begin{Q}
%   $\Is(\mathbf R^2)$の離散部分群の分類
%\end{Q}
\subsection{平行移動が$\id$以外にない場合}
\begin{Q}
	$\Is(\mathbf R^2)$の離散部分群の分類？
\end{Q}

方針：$G\subset \Is(\mathbf R^2)$を離散部分群として，$G\cap T(\mathbf R^2)$の大きさに従って場合分け：

Case 0)
\[
	G\cap T(\mathbf R^2)=\{\id\}
\]
つまり，平行移動が$\id$以外にない場合．
シャールの定理より$G$は回転と鏡映のみからなる，言い換えると，($\id$以外の)平行移動や．(鏡映以外の)並進鏡映を含まない．このとき
\begin{Prop}
ある$\mathbf x_0\in \mathbf R^n$が存在して，
\[
	G\mathbf x_0=\{\mathbf x_0\}
\]
\end{Prop}
\begin{proof}
・回転の中心は一意的：

$ \mathbf x_0$を中心とする角度$\theta\in(0,2\pi)$の回転 $S\neq \id \in G$と$ \mathbf y_0\neq  \mathbf x_0$を中心とする角度$\theta'\in(0,2\pi)$の回転 $T\neq \id \in G$とが存在したとする．
$S'=TST^{-1}$と置くと
\[
	S'(T \mathbf x_0)=TST^{-1}(T \mathbf x_0)=TS \mathbf x_0=T \mathbf x_0
\]
で，回転角の加法定理より$S'=TST^{-1}$の回転角は
\[
	\theta' + \theta - \theta' = \theta 
\]
なので$S'$は$T \mathbf x_0$を中心とする角度$\theta$の回転．
\begin{figure}[htb]
   \centering
   \begin{minipage}{0.4\columnwidth}
      \centering
      \includegraphics{~/Documents/GitHub/bEuclideanGeo/ch2/images/figureGroup5.5}
      \caption{}\label{rotationAnotherCenter}
      \end{minipage} 
   \begin{minipage}{0.4\columnwidth}
      \centering
      \includegraphics{~/Documents/GitHub/bEuclideanGeo/ch2/images/figureGroup5.6}
      \caption{}\label{refRotRef}
   \end{minipage} 
\end{figure}


$ \mathbf y_0\neq \mathbf x_0$より$Tx_0\neq \mathbf x_0$で，
$S'\mathbf x_0\neq \mathbf x_0$である(図~\ref{rotationAnotherCenter})．$S'S^{-1}$は
\[
	S'S^{-1} \mathbf x_0 = S' \mathbf x_0\neq  \mathbf x_0
\]
をみたし，$S'S^{-1}$の回転角は回転角の加法定理より
\[
	 \theta - \theta =0
\]
となる．よって，上の結果から$\id $と異なる平行移動となり，仮定に矛盾．

\emptyline

・$G$が
$\mathbf x_0$を中心とする角度$\theta\in(0,2\pi)$の回転 $S\neq \id$
と直線$\ell$に関する鏡映$R$を含めば$\mathbf x_0\in \ell$：

もし，$\mathbf x_0\notin \ell$ とすると$RSR^{-1}$
\[
	RSR^{-1}R\mathbf x_0=RS\mathbf x_0=R\mathbf x_0\neq \mathbf x_0
\]
をみたし，$RSR^{-1}$は$R\mathbf x_0$を中心とした回転(図~\ref{refRotRef})となり
最初述べたことに矛盾する．

・$G$が回転を含まないときは高々一つの鏡映しか含まない：

回転を含まないとし，
直線$\ell$に関する鏡映$R$と
直線$m$に関する鏡映$R'\neq R$を含むとする．
$\ell$と$m$が点$\mathbf b$で交わると$R'R$は回転(図~\ref{refRef1})になり矛盾．
$\ell$と$m$が平行とすると，$R'R$はは平行移動(図~\ref{refRef2})になり矛盾．よって，このようなことは起こらない．

$G=\{\id\}$のときは$x_0$は任意，$G=\{\id, R\}$で$R$が直線$\ell$に関する鏡映のときは$x_0\in \ell$とすれば良い．
\end{proof}
\begin{figure}[htb]
   \centering
   \begin{minipage}{0.4\columnwidth}
      \centering
      \includegraphics{~/Documents/GitHub/bEuclideanGeo/ch2/images/figureGroup5.7}
      \caption{鏡映の合成は回転}\label{refRef1}
   \end{minipage} 
   \begin{minipage}{0.4\columnwidth}
      \centering
      \includegraphics{~/Documents/GitHub/bEuclideanGeo/ch2/images/figureGroup5.8}
      \caption{鏡映の合成が平行移動}\label{refRef2}
   \end{minipage} 
\end{figure}

\begin{figure}[htb]
   \centering
   \begin{minipage}{0.4\columnwidth}
      \centering
      \includegraphics{~/Documents/GitHub/bEuclideanGeo/ch2/images/figureGroup5.9}
      \caption{$C_4$の基本領域}\label{C4}
   \end{minipage} 
   \begin{minipage}{0.4\columnwidth}
      \centering
      \includegraphics{~/Documents/GitHub/bEuclideanGeo/ch2/images/figureGroup5.10}
      \caption{$D_4$の基本領域}\label{D4}
   \end{minipage} 
\end{figure}

・$G$は$O(2)$の有限部分群：

命題の$\mathbf x_0$を固定する．
$G$の元は合同変換であるので$G$は
\[
	S^1(\mathbf x_0, 1)
\]
に作用してる．ただし
$
	S^1(\mathbf x_0, R)
$ は$\mathbf x_0$を中心とした半径$R>0$の円．
点$y_0\in S^1(\mathbf x_0, 1)$を固定すると
$G$は離散なので，$Gy_0$は$ S^1(\mathbf x_0, 1)$の有限集合になる．
$G$の元は回転か鏡映であり，回転は回転角のみで定まるので
$G\cap SO(2)$は回転のなす有限群である．
$G$が鏡映を含めば$G\cap SO(2)$にその鏡映を加えたものが$G$であるので
$G$は$O(2)$の有限部分群である．
\emptyline

第1回の結果から，$G$は
\[
	C_m (m\ge 1),\quad D_m  (m\ge 1)
\]
に同形となる．
\emptyline

$G$を$\Is(\mathbf R^n)$の部分群とするとき，ある領域（つまり，連結開集合）$D\subset \mathbf R^n$が存在し
\[
	\bigcup_{F\in G} \overline{F(D)}=\mathbf R^n,\quad
	F(D)\cap F'(D)=\emptyset\quad (F\neq F'\in G)
\]
のとき$D$を$G$の\term[]{基本領域}という．ただし，$\overline{F(D)}$は$F(D)$の閉包を表す．
図~\ref{C4}, \ref{D4}に$C_4$, $D_4$の基本領域の例を図示している．


\subsection{平行移動が一次元に限る場合}
$G$に平行移動が$\id$以外にある場合．その平行移動のみを取り出した
\[
	\Gamma= G^* = G\cap T(\mathbf R^2)
\]
は群になるので，\term[]{格子群}という．
また，後で示すように$r(G)\subset O(2)$は有限部分群になることが分かるので，
これを\term[]{点群}と言う．

まず$\Gamma$が一次元の場合を考える．
また，$D_m$は$C_m$を正規部分群として含むので，最初は
\[
	r(G)\subset SO(2)
\]
と仮定する．つまり回転のみからなる場合を考える．
\emptyline

case 1 )
$\Gamma=G\cap T(\mathbf R^2) $がある直線に平行な平行移動のみからなる

とする．

\emptyline
case 1-1 ) $G\subset T(\mathbf R^2) $のとき：

$T\in G$に対して$\mathbf a=T(\mathbf 0)$と置くと，$T=T_{\mathbf a}$である．仮定からある直線$\ell$が存在して
\[
	G\mathbf 0=\{\mathbf a\in \mathbf R^2\mid T_{\mathbf a}\in G\}
	\subset \ell
\]
$G$は離散なので，ある$\mathbf a\in G\mathbf 0$が存在して
\[
	\|\mathbf a\|
	=\min_{\mathbf b\neq \mathbf 0\in G\mathbf 0}\|\mathbf b\|
\]
このとき，任意の$\mathbf b\in G\mathbf 0$に対して，
ある$m\in \mathbf Z$が存在して $\mathbf b=m\mathbf a$となる：

もし，そうで無い$\mathbf b\in G\mathbf 0$が存在するとすると，
ある$m\in \mathbf Z$が存在して
\[
	\mathbf b \text{ は } m\mathbf a \text{ と }(m+1)\mathbf a\text{ を結ぶ線分の内部にある}
\]
つまり，
\[
	\|\mathbf b- m\mathbf a\|<\| \mathbf a\|
\]
$\mathbf b- m\mathbf a\in G\mathbf 0$より，
$\mathbf b=m\mathbf a$となり矛盾．
\begin{figure}[htb]
   \centering
   \begin{minipage}{0.4\columnwidth}
      \centering
      \includegraphics{~/Documents/GitHub/bEuclideanGeo/ch2/images/figureGroup5.11}
      \caption{}\label{}
   \end{minipage} 
   \begin{minipage}{0.4\columnwidth}
      \centering
      \includegraphics{~/Documents/GitHub/bEuclideanGeo/ch2/images/figureGroup5.12}
      \caption{}\label{}
   \end{minipage} 
\end{figure}

\emptyline 

このような$\mathbf a$を\term[]{基本ベクトル}とよぶ．よって
\[
	G=\Gamma
	=\{T_{n\mathbf a}\mid n \in \mathbf Z\}
	=\{T_{\mathbf a}^n \mid n \in \mathbf Z\}
	=\langle T_{\mathbf a}\rangle
\]
このような群を$F_1$と表す．
図 \ref{F1} が$F_1$の基本領域の例である．
\emptyline

case 1-2 )  $G\subset \Is^{+}(\mathbf R^2) $が回転を含むとき：

・上の結果から$\mathbf a\neq \mathbf 0$が存在して
\[
	G^* =\Gamma=\{T_{\mathbf a}^n \mid n \in \mathbf Z\}
\]

・$S$を点$\mathbf b$を中心にした回転角$\theta$の回転, $T_{\mathbf a}\in G$ を平行移動とする．
$ST_{\mathbf a}S^{-1}$の回転角は$\theta+0-\theta=0$なので平行移動．$S(\mathbf x )= S_{\theta}(\mathbf x -\mathbf b)+\mathbf b$なので
\[
	ST_{\mathbf a}S^{-1}\mathbf b
	=ST_{\mathbf a}\mathbf b
	=S(\mathbf b+\mathbf a)
	=\mathbf b+S_{\theta}\mathbf a
\]
つまり，
$S_{\theta}\mathbf a$の平行移動である．
\begin{figure}[htb]
   \centering
   \begin{minipage}{0.4\columnwidth}
      \centering
      \includegraphics{~/Documents/GitHub/bEuclideanGeo/ch2/images/figureGroup5.13}
      \caption{$F_1$の基本領域}\label{F1}
   \end{minipage} 
   \begin{minipage}{0.4\columnwidth}
      \centering
      \includegraphics{~/Documents/GitHub/bEuclideanGeo/ch2/images/figureGroup5.14}
      \caption{$ST_{\mathbf a}S^{-1}$}\label{}
   \end{minipage} 
\end{figure}

・$S\neq \id\in G$を回転とすると$\pi$回転

∵）$S\neq \id\in G$を$\theta$-回転とすると，$G\cap T(\mathbf R^2)=\langle T_{\mathbf a}\rangle$なので，上の性質から
\[
	 S_{\theta}\mathbf a=-\mathbf a,\quad\text{つまり}\quad
	 S(\mathbf b +\mathbf x) = \mathbf b -\mathbf x
\]
つまり$\theta=\pi$(図~\ref{halfRot})．
\emptyline

$S\neq \id\in G$を$\pi$-回転とし，その回転の中心を$\mathbf x_0$と置き
\[
	\mathbf x_n = T^n\mathbf x_0 \quad n\in \mathbf Z
\]
と置く．

・$T^nS$は$\mathbf x_0\mathbf x_n$の中点を中心とする$\pi$回転

∵）$T^nS$は回転角は$n\cdot0+\pi=\pi$より$\pi$回転で
\[
	T^nS\mathbf x_0 = T^n\mathbf x_0=\mathbf x_n,\quad
	T^nS\mathbf x_n = T^n\mathbf x_{-n} = \mathbf x_0
\]
なので$\mathbf x_0\mathbf x_n$の中点を中心とする$\pi$回転(図~\ref{$T^nS$})．

$\mathbf y_n$ を$\mathbf x_n\mathbf x_{n+1}$の中点とすると，
$\mathbf y_n$は$\pi$回転の中心．
\begin{figure}[htb]
   \centering
   \begin{minipage}{0.4\columnwidth}
      \centering
      \includegraphics{~/Documents/GitHub/bEuclideanGeo/ch2/images/figureGroup5.15}
      \caption{}\label{halfRot}
   \end{minipage} 
   \begin{minipage}{0.4\columnwidth}
      \centering
      \includegraphics{~/Documents/GitHub/bEuclideanGeo/ch2/images/figureGroup5.16}
      \caption{$T^nS$}\label{$T^nS$}
   \end{minipage} 
\end{figure}

・$R\in G$を任意の$\pi$回転とすると
$RS^{-1}$の回転角は$\pi-\pi=0$なので平行移動．
$\Gamma=G\cap T(\mathbf R^2)=\langle T_{\mathbf a}\rangle$なので
\[
	\exists\,  T_{\mathbf a}^m \text{ such that } 
	RS^{-1}=T_{\mathbf a}^m\quad
	\Rightarrow R=T_{\mathbf a}^mS
\]

・$x_0$を$y_0$に写す$\pi$回転$S$が存在すると，その中心は$x_0$と$y_0$の中点なので，$x_0$中心の$\pi$回転と$\pi$回転$S$を合成すると$x_0$を$y_0$に写す平行移動が存在することになり，$\mathbf a$の最小性に反する．$x_n$同士，$y_n$同士が$G$で移りあうことは明らかなので, 
$Gx_0$, $Gy_0$は二つの異なる$G$軌道である．

・$T_{\mathbf a}^mST_{\mathbf a}^mS=\id$なので
\[
	T_{\mathbf a}^mS=ST_{\mathbf a}^{-m}
\]
したがって
\[
	T_{\mathbf a}^3ST_{\mathbf a}^2ST_{\mathbf a}S
	=T_{\mathbf a}^3T_{\mathbf a}^{-2}SST_{\mathbf a}S
	=T_{\mathbf a}^{2}S
\]
のように計算して，積を$T_{\mathbf a}^m$, $T_{\mathbf a}^nS$の形に表すこことが出来る．

\begin{figure}[htb]
   \centering
   \includegraphics{~/Documents/GitHub/bEuclideanGeo/ch2/images/figureGroup5.17}
   \caption{$F_2$の基本領域二つ}\label{F2}
\end{figure}

以上から
\[
	G = \{T_{\mathbf a}^m, T_{\mathbf a}^nS\mid m,n \in \mathbf Z\}
	= \langle T_{\mathbf a}, S\rangle
\]
このような群を$F_2$と呼ぶ．図 \ref{F2}に$F_2$の基本領域の例を二つ与える．
\emptyline

次に
\subsection{$2$次元の平行移動からなる場合}
%case 2 )
%$\Gamma=G\cap T(\mathbf R^2) $ は
%一次独立な二つ以上のベクトルからなる平行移動からなる．
%
%暫く，$r(G)\subset SO(2)$の場合のみ考える．
%$\Gamma=G\cap T(\mathbf R^2) $を\term[]{格子群}といい，
%集合としての$\Gamma$を\term[]{格子}という．
%\emptyline
%
%case 1 )
%$\Gamma=G\cap T(\mathbf R^2) $ がある直線に平行な平行移動からなる場合
%
%$G$は$F_1$, $F_2$のいずれかに同形．
%
%\begin{figure}[htb]
%   \centering
%   \begin{minipage}{0.4\columnwidth}
%      \centering
%      \includegraphics{~/Documents/GitHub/bEuclideanGeo/ch2/images/figureGroup6.5}
%      \caption{$F_1$}
%   \end{minipage} 
%   \begin{minipage}{0.4\columnwidth}
%      \centering
%      \includegraphics{~/Documents/GitHub/bEuclideanGeo/ch2/images/figureGroup6.6}
%      \caption{$F_2$}
%   \end{minipage} 
%\end{figure}
%
%\emptyline
%
case 2)
$\Gamma=G\cap T(\mathbf R^2) $ は
一次独立な二つ以上のベクトルからなる平行移動からなる場合

%\[
%	L=G\cap T(\mathbf R^2)
%\]
\emptyline

$T\in \Gamma$に対して$\mathbf a=T(\mathbf 0)$と置くと，$T=T_{\mathbf a}$である．$G$は離散なので，ある$\mathbf a\in \Gamma \mathbf 0\setminus\{\mathbf 0\}$が存在して
\[
	\|\mathbf a\|
	=\min_{\mathbf b\in \Gamma \mathbf 0\setminus\{\mathbf 0\} }\|\mathbf b\|
\]
このとき，仮定から
\[
	\Gamma \mathbf 0\supsetneqq \{n\mathbf a\mid n\in \mathbf Z\}
\]
なので，$\mathbf b\in G\mathbf 0\setminus  \{n\mathbf a\mid n\in \mathbf Z\}
$ で
\[
	\|\mathbf b\|
	=
	\min_{
		\mathbf c\in 
		\Gamma \mathbf 0
		\setminus  \{n\mathbf a\mid n\in \mathbf Z\}
	}
	\|\mathbf c\|
\]
成るものが存在する．このとき
\[
	\|\mathbf a\|\le \|\mathbf b\|,
	\quad
	 -\frac12\|\mathbf a\|
	 \le
	 (\mathbf b, \mathbf a)
	 \le 
	 \frac12\|\mathbf a\|
\]
が成り立ち(図~\ref{elementaryVectors},~\ref{elementaryVectors2})

\begin{figure}[htb]
   \centering
   \begin{minipage}{0.4\columnwidth}
      \centering
      \includegraphics{~/Documents/GitHub/bEuclideanGeo/ch2/images/figureGroup6.7}
      \caption{基本ベクトル}\label{elementaryVectors}
   \end{minipage} 
   \begin{minipage}{0.4\columnwidth}
      \centering
      \includegraphics{~/Documents/GitHub/bEuclideanGeo/ch2/images/figureGroup6.8}
      \caption{基本ベクトル}\label{elementaryVectors2}
   \end{minipage} 
\end{figure}

\[
	\Gamma =
	G\cap T(\mathbf R^2) 
	= 
	\{T_{m\mathbf a+n\mathbf b}\mid m,n \in \mathbf Z\}
\]
となる．%このような$\Gamma $を\term[]{格子}, 
このような$\mathbf a$, $\mathbf b$を$\Gamma $の\term[]{基本ベクトル}とよぶ．
\emptyline

case 2-1) $G$が回転を含まないとき，つまり，$G\subset T(\mathbf R^2)$のとき 
\[
	G
	= \Gamma =
	\{T_{m\mathbf a+n\mathbf b}\mid m,n \in \mathbf Z\}
	=\langle T_{\mathbf a}, T_{\mathbf b}\rangle
\]
となる．
このような$G$を$M_1$と呼ぶ(図~\ref{M1})．
\begin{figure}[htb]
   \centering
   \begin{minipage}{0.4\columnwidth}
      \centering
      \includegraphics{~/Documents/GitHub/bEuclideanGeo/ch2/images/figureGroup6.9}
      \caption{$M_1$}\label{M1}
      \end{minipage} 
   \begin{minipage}{0.4\columnwidth}
      \centering
      \includegraphics{~/Documents/GitHub/bEuclideanGeo/ch2/images/figureGroup6.10}
      \caption{$M_2$}\label{M2}
   \end{minipage} 
\end{figure}

\emptyline

case 2-2) $G$が回転としては$\pi$回転しか含まないとき.

$S\neq \id\in G$を$\pi$-回転とし，$S$の回転の中心を$\mathbf x_{0,0}$と置く．
\[
	\mathbf x_{m,n} 
	=T_{m\mathbf a+n\mathbf b}\mathbf x_{0,0} \quad 
	T_{m\mathbf a+n\mathbf b}\in \Gamma 
\]
と置くと，
先週示したのと同じ議論で：
\emptyline

・$T_{m\mathbf a+n\mathbf b}S$は$\mathbf x_{0,0}\mathbf x_{m,n} $の中点を中心とする$\pi$回転
\emptyline


・$R\in G$を任意の$\pi$回転とすると
$RS^{-1}$の回転角は$\pi-\pi=0$なので平行移動，つまり
\[
	\exists\,  T_{m\mathbf a+n\mathbf b}\in \Gamma  
	\text{ such that } 
	RS^{-1}=T_{m\mathbf a+n\mathbf b}\quad
	\Rightarrow R=T_{m\mathbf a+n\mathbf b}S
\]

・$T_{m\mathbf a+n\mathbf b}ST_{m\mathbf a+n\mathbf b}S=\id$なので
\[
	T_{m\mathbf a+n\mathbf b}S=ST_{-m\mathbf a-n\mathbf b}
\]

以上から
\[
	G = \{T_{m\mathbf a+n\mathbf b}, T_{m'\mathbf a+n'\mathbf b}S\mid m,n, m', n'\in \mathbf Z\}
	=\langle T_{\mathbf a}, T_{\mathbf b}, S \rangle
\]
と表せる．
このような群を$M_2$と呼ぶ．この$\pi$回転の中心は$4$個の$G$軌道に分かれていることを，前回と同じように示すことが出来る(図~\ref{M2}の$4$色の丸がそれを示している)．
\emptyline


case 2-3) $G$が回転角が$\pi$でない回転を含むとき：

回転$S$の回転の中心$\mathbf x_0$を$\mathbf 0$とみなし，$\mathbf 0$を起点とした格子$\Gamma $を考える．
\begin{Prop}[結晶条件]
回転
\[
	S=\begin{pmatrix}
		\cos \theta &-\sin \theta\\
		\sin \theta& \cos \theta\\
	\end{pmatrix}\neq E
\]
が$\Gamma $を保つ（$\Gamma $の$S$による像が$\Gamma $に移る）とすると，$S$の回転角は
\[
	\theta=\pm 2\pi/p\quad p=2, 3,4,6
\]
の何れかである．
\end{Prop}
\begin{proof}
$\mathbf a$, $\mathbf b$を$\Gamma $の基本ベクトルとし，
$U=\begin{pmatrix}\mathbf a&\mathbf b\end{pmatrix}$とおくと
\[
	SU
	=
	\begin{pmatrix}S\mathbf a& S\mathbf b\end{pmatrix}
	=
	(k\mathbf a+l\mathbf b, m\mathbf a+n\mathbf b)
	=
	U\begin{pmatrix}
		k &m\\
		l& n\\
	\end{pmatrix}
\]
となる$k,l,m,n\in \mathbf Z$が存在する．
\[
	\tr S
	=
	\tr (U\begin{pmatrix}
		k &m\\
		l& n\\
	\end{pmatrix}U^{-1})
	=
	\tr (\begin{pmatrix}
		k &m\\
		l& n\\
	\end{pmatrix}U^{-1}U)
	=
	k+n \in \mathbf Z
\]
なので
\[
	2\cos \theta \in \mathbf Z
\]
$-1\le \cos \theta\le 1$より$\theta$は
\[
	\frac{\pi}{3},\quad\frac{\pi}{2},\quad\frac{2\pi}{3},\quad
	\pi,\quad-\frac{\pi}{3},\quad -\frac{\pi}{2},\quad-\frac{2\pi}{3}
\]
のいずれかである．
\end{proof}


%$\pi$回転でない$\frac{2q\pi}{p}$回転$S$が見つかった場合，$p$と$q$は互いに素と取れるので，ユークリッドの互除法より$mp+kq=1$なる$m, k\in \mathbf Z$が見つかる．
%$S^m$の回転角は
%\[
%	\frac{2\pi kq}{p}=\frac{2\pi(1-mp)}{p}=\frac{2\pi}{p} + 2\pi m
%\]
%なので$\frac{2\pi}{p}$回転が得られる．

$\pi$回転でない点$\mathbf p$を中心とした回転は巡回群をなすので，
$P$を$\mathbf p$を中心とした($P$の回転が時計回りの時は逆元に取り替えれば良いので，反時計回りの)$\frac{2\pi}{p}$回転で，
仮定から$p\ge 3$として良い．ここで
$T\neq \id \in \Gamma =G\cap T(\mathbf R^2) $を平行移動とすると，
$TP$の回転角は$0+\frac{2\pi}{p}=\frac{2\pi}{p}$なので，$\pi$回転でない回転であり，$T\neq \id$より
\[
	TP\mathbf p=T\mathbf p \neq \mathbf p
\]
より，$G$は$\mathbf p$以外の点を中心とした$\pi$回転でない回転を含むことが分かる．まず，$\frac{2\pi}{p}$回転$P$の$p$はこのようなものの最小なものとする．
そこで$Q$を$\mathbf q$を中心とした$\frac{2\pi}{q}$回転とし, $q\ge 3$で，更に
\[
	\|\mathbf p-\mathbf q\|
	=\min_{%\mathbf q'\text{:$q\ge 3$の $\frac{2\pi}{q}$-回転中心}
	\mathbf q': {2\pi}/{q'}\text{ 回転の中心}, q'\ge p\ge 3}
		\|\mathbf p-\mathbf q'\|
\]
をみたすとして良い．ここでも反時計回りの$Q$を取る．

$R=(QP)^{-1}$と置くと，
\[
	0<\frac{2\pi}{p}+\frac{2\pi}{q} 
	\le \frac{4\pi}{3}<2\pi
\]
なので$R$の回転角は($2\pi$を加算して計算することで)
\[
	2\pi>2\pi-\frac{2\pi}{p}-\frac{2\pi}{q} 
	\ge 2\pi- \frac{4\pi}{3}
	=\frac{2\pi}{3}
\]
よって（$\pi$回転も含む）回転になる．その中心を$\mathbf r$と置く．
上の命題より回転角は$\frac{2\pi}{r} $で，$r=2,3$としてよい．

$RQP=\id$に注意して$\mathbf r$を探す．
図~\ref{comp1}, ~\ref{comp2}の向きを考慮し，
$\angle \mathbf r\mathbf p\mathbf q$ が回転$P$の回転角の半分で，
$\angle \mathbf p\mathbf q\mathbf r$が回転$Q$の回転角の半分となるように$\mathbf r$を選ぶことが出来る．ただし，$\angle \mathbf p \mathbf q \mathbf r$は線分
$\mathbf q\mathbf p $と$\mathbf q \mathbf r$の頂点$\mathbf q$での（反時計回りを正の方向とする）角度を表す．
すると
\[
\begin{split}
	&2\angle \mathbf r \mathbf p \mathbf q
	=\angle \mathbf x \mathbf p P\mathbf x
	= \frac{2\pi}{p},\\
	& 2\angle \mathbf p \mathbf q \mathbf r
	=\angle P\mathbf x \mathbf q QP\mathbf x
	= \frac{2\pi}{q},\\
	&2\angle \mathbf q \mathbf r\mathbf p 
	=\angle QP\mathbf x \mathbf r \mathbf x
	= \frac{2\pi}{r}
\end{split}
\]
となり，三角形$ \triangle\mathbf p \mathbf q \mathbf r$の内角の和は$\pi$なので
\[
	\frac{\pi}{p}+\frac{\pi}{q}+\frac{\pi}{r} = \pi
\]
が分かる．

\begin{figure}[htb]
   \centering
   \begin{minipage}{0.4\columnwidth}
      \centering
      \includegraphics{~/Documents/GitHub/bEuclideanGeo/ch2/images/figureGroup6.11}
      \caption{}\label{comp1}
   \end{minipage} 
   \begin{minipage}{0.4\columnwidth}
      \centering
      \includegraphics{~/Documents/GitHub/bEuclideanGeo/ch2/images/figureGroup6.12}
      \caption{}\label{comp2}
   \end{minipage} 
\end{figure}


なので
\begin{equation}
	\frac{1}{p}+\frac{1}{q}+\frac{1}{r} =1,\quad 
	p, q =3, 4, 6,\quad r= 2,3\tag{*}\label{cristal}
\end{equation}
まず，$q\ge p\ge 3$として良い．
\begin{enumerate}
\item 
$p=q=3$のとき
\[
	\frac{1}{r} =1-\frac{1}{3}-\frac{1}{3}=\frac{1}{3}
\]
より，$r=3$.
\item
$q\ge p\ge 3$で$p=q=3$でないとき
\[
	\frac{1}{r} >1-\frac{1}{3}-\frac{1}{3}=\frac{1}{3}
\]
で$r=2, 3$なので$r=2$．すると
\[
	\frac{1}{p}+\frac{1}{q} =1-\frac{1}{2}=\frac{1}{2}
\]
$q\ge p\ge 5$とすると
\[
	\frac{1}{2}=\frac{1}{p}+\frac{1}{q}
	\le \frac{1}{5}+\frac{1}{5}\le \frac{2}{5}
\]
となり矛盾するので，$p=3,4$の何れか．
\begin{enumerate}
\item 
$p=3$ のとき
\[
	\frac{1}{q}=\frac{1}{2}-\frac{1}{3}=\frac{1}{6}
\]
より$q=6$．

\item
$p=4$のとき
\[
	\frac{1}{q}=\frac{1}{2}-\frac{1}{4}=\frac{1}{4}
\]
より$q=4$．
\end{enumerate}

\end{enumerate}
まとめると
\begin{table}[h]
  \centering
  \begin{tabular}{lll|l} 
  	\noalign{\hrule height 1pt}
	 $p$&$q$&$r$& 格子 \\ \hline \hline
	 3&3&3&正六角形格子\\ \hline
	 3&6&2&正六角形格子 \\ \hline
	  4&4&2& 正方格子\\
     \noalign{\hrule height 1pt}
   \end{tabular}
  \caption{可能な$p, q, r$}
\end{table}

注）第2回で説明した$SO(3)$の有限群の極の位数から導かれる関係式
\[
	1+\frac2n 
	=
	\frac1{\alpha_1}+\frac1{\alpha_2}+\frac1{\alpha_3}
\]
や類等式は条件\eqref{cristal}との類似になっている．ただし，$n=|G|$, $\alpha_i = |\stab(\mathbf x_i)|$.
\emptyline

これらの場合は回転で格子が保たれているので，常に
\[
	\|\mathbf a \|=\|\mathbf b \|
\]
が成り立つことも指摘しておく．
%\begin{figure}[htb]
%\centering
%\includegraphics{~/Documents/GitHub/bEuclideanGeo/ch2/images/figureGroup6.15}
%\caption{$M_3$}\label{comp1}
%\end{figure}
%\begin{figure}[htb]
%\centering
%\includegraphics{~/Documents/GitHub/bEuclideanGeo/ch2/images/figureGroup6.14}
%\caption{正方格子}\label{comp2}
%\end{figure}
%\emptyline

%\begin{figure}[htb]
%\centering
%\includegraphics{~/Documents/GitHub/bEuclideanGeo/ch2/images/figureGroup6.18}
%\caption{格子}\label{}
%\end{figure}
%\emptyline
%
%\cite{C}の3, 4章，
%\cite{HC}の2章，\cite{W} の3章に数式をあまり用いない解説が与えられている．いずれの本も（非専門家に向けた）近代幾何学の古典と言って良い本である．
%
%\begin{thebibliography}{9}
%\bibitem{C}
%H.S.M. コクセター, 幾何学入門〈上, 下〉 (ちくま学芸文庫) 2009 
%
%\bibitem{HC}
%D・ヒルベルト, S・コーン=フォッセン, 直観幾何学【新装版】, 2019 , 
%みすず書房 %(D. Hilbert, S. Cohn-Vossen, Geometry and the Imagination, 2020, Ams Chelsea Publishing) 
%
%\bibitem{W}
%ヘルマン・ワイル, シンメトリー (ちくま学芸文庫) 2022
%(Hermann Weyl, Symmetry, Princeton Science Library, 2016)
%
%%\bibitem{E}
%%C.H.マックギラフィ, 有馬 朗人 , エッシャー—シンメトリーの世界
%\end{thebibliography}
%

\subsection{鏡映や並進鏡映を含まない場合}
暫く鏡映や並進鏡映を含まない場合を考える．

case 2 )
$\Gamma=G\cap T(\mathbf R^2) $ は
一次独立な二つ以上のベクトルからなる平行移動からなる

%\[
%	\Gamma=
%	G\cap T(\mathbf R^2) 
%	= 
%	\{T_{m\mathbf a+n\mathbf b}\mid m,n \in \mathbf Z\}
%\]
%となる．このような$\Gamma$を群としては\term[]{格子群}, 
%集合としては\term[]{格子}, 
%このような$\mathbf a$, $\mathbf b$を\term[]{基本ベクトル}とよぶ．
%\emptyline

case 2-1) $G$が回転を含まないとき，つまり，$G\subset T(\mathbf R^2)$のとき 
\[
	G
	= \Gamma=
	\{T_{m\mathbf a+n\mathbf b}\mid m,n \in \mathbf Z\}
\]
このような$G$を$M_1$ (P1\footnote{以下国際記号(International nomenclature)による群の表示を括弧内に追記する. 他にも})と呼ぶ(図~\ref{M1})．
\begin{figure}[htb]
   \centering
   \begin{minipage}{0.4\columnwidth}
      \centering
      \includegraphics{~/Documents/GitHub/bEuclideanGeo/ch2/images/figureGroup7.1}
      \caption{ $M_1$ ($P1$), Oblique, $C_1$}\label{M1}
   \end{minipage} 
   \begin{minipage}{0.4\columnwidth}
      \centering
      \includegraphics[trim=20 0  0 35, clip]{~/Documents/GitHub/bEuclideanGeo/ch2/images/figureGroup7.2}
      \caption{ $M_2$ ($P2$), Oblique, $C_2$}\label{M2}
   \end{minipage} 
\end{figure}

\emptyline

case 2-2) $G$が回転としては$\pi$回転のみを含むとき.
このような群を$M_2$ (P2)と呼ぶ(図~\ref{M2}, 格子群$\Gamma$ の基本領域を灰色でしめした)．
\emptyline


case 2-3)  $G$が$\pi$回転以外の回転を含むとき：
\emptyline

\begin{Prop}[結晶条件]
$\mathbf 0$を中心とした回転
\[
	S=\begin{pmatrix}
		\cos \theta &-\sin \theta\\
		\sin \theta& \cos \theta\\
	\end{pmatrix}
\]
が$\Gamma$を保つとすると，$S$の回転角は
\[
	\theta=\pm 2\pi/p\quad p=2, 3,4,6
\]
の何れか．
\end{Prop}

\emptyline

$G$は異なる点を中心とした$\pi$回転でない回転を含む．まず，$\frac{2\pi}{p}$回転$P$の$p$はこのようなものの最小なもので，$p\ge 3$として良い．
$Q$を$\mathbf q$を中心とした$\frac{2\pi}{q}$回転とし, $q\ge p\ge 3$で，更に
\[
	\|\mathbf p-\mathbf q\|
	=\min_{\mathbf q': {2\pi}/{q'}\text{回転の中心}, q'\ge p\ge 3}\|\mathbf p-\mathbf q'\|
\]
をみたすとして良い．

$R=(QP)^{-1}$と置くと，
（$\pi$回転も含む）回転になる．その中心を$\mathbf r$と置く．
上の命題より回転角は$\frac{2\pi}{r} $で，$r=2,3$としてよい．
このとき，
\begin{equation}\label{cristal}
	\frac{1}{p}+\frac{1}{q}+\frac{1}{r} =1,\quad 
	p\le q =3, 4, 6,\quad r= 2,3
\end{equation}
がわかり，組み合わせを調べることで次の場合しかないことが分かった．
\begin{table}[h]
   \centering
   \begin{tabular}{lll|l} 
  	\noalign{\hrule height 1pt}
	 $p$&$q$&$r$& 格子 \\ \hline \hline
	 3&3&3&正六角形格子\\ \hline
	 3&6&2&正六角形格子 \\ \hline
	  4&4&2& 正方格子\\
  \noalign{\hrule height 1pt}
\end{tabular}
  \caption{可能な$p, q, r$}
\end{table}

\subsection{case 2-3の分類}
\begin{Prop} 
$P\in G$を$\mathbf p$を中心とする回転角$\theta$の回転とし，
$Q\in G$を回転角$\theta'$の回転とすると
$QPQ^{-1}$は$Q\mathbf p$を中心とする回転角$\theta$の回転．
特に，$G$は回転の中心を回転の中心に写す．
\end{Prop}
\begin{proof}
$QPQ^{-1}$の回転角は$\theta'+\theta -\theta'=\theta$なので，回転角$\theta$の回転．
\[
	QPQ^{-1}Q\mathbf p=QP\mathbf p=Q\mathbf p
\]
なので，$Q\mathbf p$が中心．
\end{proof}

$Q\in G$の群$G$への作用を
\[
	\iota_Q(P)= QPQ^{-1}\in G \quad (P\in G)
\]
で定めると
\[
\begin{split}
	\iota_R(PQ) = R(PQ)R^{-1}&=RPR^{-1}RQR^{-1}=\iota_R(P)\iota_R(Q),\\
	\iota_R\id& = R\id R^{-1}=\id
\end{split}
\]
から$\iota_R:G\to G$は群の準同形になることが分かる．これを\term[]{内部自己同形}という．内部自己同形全体は
\[
	\iota_{RQ}=\iota_R\iota_Q,\quad \iota_{\id} = \id
\]
が成り立つので群を成すので，\term[]{内部自己同形群}という．
上の命題は群$G$の内部自己同形群が
$\mathbf R^2$のなかの$G$の回転の中心全体に作用していることを表している．

\emptyline

Case 2-3) $p=q=r=3$の場合：

$P$を$\mathbf p$を中心とした$\frac{2\pi}{3}$回転とし，$Q$を$\mathbf q$を中心とした$\frac{2\pi}{3}$回転,  $R$を$\mathbf r$を中心とした$\frac{2\pi}{3}$回転とする．
$S=QPQ^{-1}$と置くと命題より
$S$は$\mathbf s=Q\mathbf p$を中心とする角度$\frac{2\pi}{3}$の回転．
$S^{-1}Q$は
\[
	S^{-1}Q\mathbf p = S^{-1}\mathbf s = \mathbf s 
\]
をみたし，
$S^{-1}Q$の回転角は$-\frac{2\pi}{3}+\frac{2\pi}{3}=0$なので，
$\mathbf s-\mathbf p$の平行移動．これを$\mathbf a$ と置く．

\begin{figure}[htb]
   \centering
   \begin{minipage}{0.4\columnwidth}
      \centering
      \includegraphics{~/Documents/GitHub/bEuclideanGeo/ch2/images/figureGroup7.3}
      \caption{}\label{}
   \end{minipage} 
   \begin{minipage}{0.4\columnwidth}
      \centering
      \includegraphics{~/Documents/GitHub/bEuclideanGeo/ch2/images/figureGroup7.4}
      \caption{}\label{}
   \end{minipage} 
\end{figure}
\emptyline

$S'=Q^{-1}PQ$と置くと$\mathbf s'=Q^{-1}\mathbf p$を中心とする角度$\frac{2\pi}{3}$の回転．
$S'Q^{-1}$は
\[
	S'Q^{-1}\mathbf p = S'\mathbf s' = \mathbf s' 
\]
をみたす
$\mathbf s'-\mathbf p=:\mathbf b$の平行移動．

$\mathbf a$と$\mathbf b$は一次独立で平行移動を生成する．そこで
格子群
\[
	 \Gamma=
	\{T_{m\mathbf a+n\mathbf b}\mid m,n \in \mathbf Z\}
	\subset
	G \cap T(\mathbf R^2)
\]
を得る．
\emptyline 

・$P$, $Q$, $R$以外から生成される回転を$G$は含まない．

∵）もし$\pi$-回転を含めば，$-\frac{2\pi}{3}$回転と合成することで
$\frac{\pi}{3}$回転が得られる．そこで$P$, $Q$, $R$以外の回転は$\frac{2\pi}{m}$ $(m\ge 3)$回転
と仮定していい．その回転の中心は
$T(\mathbf R^2)\cap G$で平行移動して，三角形$\mathbf p\mathbf q\mathbf r$から作った$\mathbf p$を中心とした正六角形の内部と境界を含めた部分から頂点を除いた集合に移すことが出来る(図~\ref{fundDom2})が，
これは$\mathbf q$を$\mathbf p$にもっとも近い$\frac{2\pi}{m}$ $(m\ge 3)$回転の中心とした仮定に反する．

・$T_{\mathbf a}$, $T_{\mathbf b}$より短い平行移動は存在しない．

∵）平行移動で$\frac{2\pi}{3}$の回転は$\frac{2\pi}{3}$の回転に移るので，
そのような平行移動があるとすれば$\mathbf r -\mathbf p$(または$\mathbf q -\mathbf p$)の形の平行移動$T$である．このとき$RT$は
\[
	RT\mathbf p=R\mathbf r=\mathbf r,\quad
	RT\mathbf r=\mathbf q,\quad
	RT\mathbf q=\mathbf p
\]
となる(図~\ref{rotationAroundc})ので回転角$\frac{2\pi}{3}$の回転．したがって，三角形$\mathbf p\mathbf q\mathbf r$
の重心$\mathbf c$を中心とした$\frac{2\pi}{3}$の回転となり，$\mathbf q$を$\mathbf p$にもっとも近い回転の中心とした仮定に反する．

\emptyline
したがって，
\[
	 \Gamma=
	\{T_{m\mathbf a+n\mathbf b}\mid m,n \in \mathbf Z\}
	=
	G \cap T(\mathbf R^2)
\]
が格子群．
\emptyline


\begin{figure}[htb]
   \centering
   \begin{minipage}{0.4\columnwidth}
      \centering
      \includegraphics{~/Documents/GitHub/bEuclideanGeo/ch2/images/figureGroup7.5}
      \caption{}\label{rotationAroundc}
   \end{minipage} 
   \begin{minipage}{0.4\columnwidth}
      \centering
      \includegraphics{~/Documents/GitHub/bEuclideanGeo/ch2/images/figureGroup7.6}
      \caption{}\label{fundDom2}
   \end{minipage} 
\end{figure}



\begin{figure}[htb]
   \centering
   \includegraphics[trim=60 60 60 100, clip]{~/Documents/GitHub/bEuclideanGeo/ch2/images/figureGroup7.7}
   \caption{ $M_3$ ($P3$), 正六角形格子, $C_3$ }\label{comp1}
\end{figure}

以上より
\[
	G = \langle P, Q, R\rangle
\]
このような群$G$を$M_3$ (P3)という(図~\ref{comp1}回転の頂点の色は異なる$3$つの$G$軌道を表す)．
\emptyline

Case 2-4) $p=q=4, r=2$の場合：
$P$を$\mathbf p$を中心とした$\frac{2\pi}{4}=\frac{\pi}{2}$回転とし，$Q_1$を$\mathbf q_1$を中心とした$\frac{\pi}{2}$回転,  $R$を$\mathbf r$を中心とした$\pi$回転とする．
命題より
$Q_2=PQ_1P^{-1}$ は$\mathbf q_2=P\mathbf q_1$を中心とした$\frac{\pi}{2}$回転,
$Q_3=P^{2}Q_1P^{-2}$は$\mathbf q_3=P^{2}\mathbf q_1$を中心とした$\frac{\pi}{2}$回転，
$Q_4=P^{3}Q_1P^{-3}$は$\mathbf q_4=P^{3}\mathbf q_1$を中心とした$\frac{\pi}{2}$回転

$PQ_1^{-1}$, $P^{-1}Q_1$は 平行移動で
\[
	PQ_1^{-1}\mathbf q_1= P\mathbf q_1=\mathbf q_2,\quad
	P^{-1}Q_1\mathbf q_1= P^{-1}\mathbf q_1=\mathbf q_4
\]
より，基本ベクトルを$\mathbf a = \mathbf q_2-\mathbf q_1$, $\mathbf b = \mathbf q_4-\mathbf q_1$と取ることが出来る．


\begin{figure}[htb]
   \centering
   \begin{minipage}{0.4\columnwidth}
      \centering
      \includegraphics{~/Documents/GitHub/bEuclideanGeo/ch2/images/figureGroup7.8}
      \caption{}\label{}
   \end{minipage} 
   \begin{minipage}{0.4\columnwidth}
      \centering
      \includegraphics{~/Documents/GitHub/bEuclideanGeo/ch2/images/figureGroup7.9}
      \caption{}\label{rotM4}
   \end{minipage} 
\end{figure}

・$P$, $Q$, $R$以外から生成されない回転を$G$は含まない．

∵）$M_3$と同じ理由
\emptyline


・$T_{\mathbf a}$, $T_{\mathbf b}$より短い平行移動は存在しない．

∵）平行移動は$\frac{\pi}{2}$回転を$\frac{\pi}{2}$回転に写すので，そのような平行移動があれば $\mathbf q - \mathbf p$の形の平行移動 $T$のはずである($\mathbf r' - \mathbf r$は$\mathbf a$か$\mathbf b$の整数倍なので無視していい)．
$QT$は図~\ref{rotM4}より，$\mathbf r$ を中心とした$\frac{\pi}{2}$回転となり，$\mathbf q$を$\mathbf p$にもっとも近い$\frac{2\pi}{m}$ $(m\ge 3)$回転の中心とした仮定に反する．
\emptyline
したがって，
\[
	 \Gamma=
	\{T_{m\mathbf a+n\mathbf b}\mid m,n \in \mathbf Z\}
	=
	G \cap T(\mathbf R^2)
\]
が格子群．
\emptyline

\begin{figure}[htb]
   \centering
   \includegraphics[trim=10 10 10 10, clip]{~/Documents/GitHub/bEuclideanGeo/ch2/images/figureGroup7.10}
   \caption{$M_4$ (P4), 正方格子, $C_4$ }\label{comp2}
\end{figure}



このような群を$M_4$ (P4)という(図~\ref{comp2}回転の頂点の色は異なる$3$つの$G$軌道を表す)．．
\emptyline 

Case 2-3) $p=3, q=6, r=2$の場合：
$P$を$\mathbf p$を中心とした$\frac{2\pi}{3}$回転とし，$Q_1$を$\mathbf q_1$を中心とした$\frac{2\pi}{6}=\frac{\pi}{3}$回転,  $R$を$\mathbf r$を中心とした$\pi$回転とする．
命題より
$Q_2=PQ_1P^{-1}$ は$\mathbf q_2=P\mathbf q_1$を中心とした$\frac{\pi}{3}$回転,
$Q_3=P^{-1}Q_1P$は$\mathbf q_3=P^{-1}\mathbf q_1$を中心とした$\frac{\pi}{3}$回転．

$\mathbf a = \mathbf q_2-\mathbf q_1$ ,$\mathbf b = \mathbf q_3-\mathbf q_1$を基本ベクトルに取ることができ，
\emptyline


・$P$, $Q$, $R$以外から生成されない回転を$G$は含まない．

∵）$M_3$と同じ理由
\emptyline


・$T_{\mathbf a}$, $T_{\mathbf b}$より短い平行移動は存在しない．

∵）平行移動で$\frac{2\pi}{m}$の回転は$\frac{2\pi}{m}$の回転に移るので，
そのような平行移動があるとすれば$\mathbf p'-\mathbf p$, 
$\mathbf r' -\mathbf r$の形の短い平行移動$T$であるが，どちらの場合も矛盾が生じる．

\emptyline

\begin{figure}[htb]
   \centering
   \begin{minipage}{0.4\columnwidth}
      \centering
      \includegraphics{~/Documents/GitHub/bEuclideanGeo/ch2/images/figureGroup7.11}
      \caption{}\label{}
   \end{minipage} 
   \begin{minipage}{0.4\columnwidth}
      \centering
      \includegraphics[width = 5cm]{~/Documents/GitHub/bEuclideanGeo/ch2/images/figureGroup7.12}
      \caption{}\label{}
   \end{minipage} 
\end{figure}

\emptyline


このような群を$M_5$ (P6)という(図~\ref{comp3}回転の頂点の色は異なる$3$つの$G$軌道を表す)．
\emptyline 

このようにして，case 2-3の離散群の分類が得られた．
\emptyline 

\begin{Quiz}
$8$ $M_4$の場合に$($授業中に$)$残った部分に詳しい証明を与えよ．
\end{Quiz}

\begin{Quiz}$9$ $M_5$の場合に$($授業中に$)$残った部分に詳しい証明を与えよ．
\end{Quiz}


\emptyline 
\begin{figure}[htb]
   \centering
   \includegraphics[trim=67 65  40 150, clip]{~/Documents/GitHub/bEuclideanGeo/ch2/images/figureGroup7.13}
   \caption{$M_5$ (P6), 正三角形格子, $C_6$}\label{comp3}
\end{figure}
\emptyline


\subsection{ラフな離散群の分類}
\[
	\Is(\mathbf R^n)
	=\{F:\mathbf R^n\to \mathbf R^n\mid \text{合同変換}\}
	\ni F, F'
\]
とすると
\[
	F(\mathbf x) =A\mathbf x+\mathbf b,\quad
	F'(\mathbf x) =A'\mathbf x+\mathbf b'
\]
なので
\begin{equation}\label{semiproduct}
	F'\circ F(\mathbf x)
	=F'(A\mathbf x+\mathbf b)
	=A'(A\mathbf x+\mathbf b)+\mathbf b'
	=A'A\mathbf x+A'\mathbf b+\mathbf b'
\end{equation}
特に
\[
	F^{-1}(\mathbf x)=A^{-1}\mathbf x - A^{-1}\mathbf b
\]
$r(F)=A\in O(n)$と置くと
\[
	r(F'\circ F)=A'A=r(F')r(F)\in O(n)
\] 
となり，$r:\Is(\mathbf R^n)\to O(n)$は準同形写像．

$G$を$\Is(\mathbf R^n)$の離散部分群とする．
\begin{gather*}
	T(\mathbf R^n)
	=\{F\in\Is(\mathbf R^n)\mid \text{平行移動}\}
\end{gather*}
とすると
\[
	G^* = G\cap T(\mathbf R^n)
\]
は$G$の部分群であり，$G$の\term[]{格子群}と呼ばれる．
$G^*$は(離散な)可換群なので，一次独立な$\mathbf a_1,\cdots,\mathbf a_k$が存在して
\[
	G^*
	= 
	\{T_{m_1\mathbf a_1+\cdots +m_k\mathbf a_k}\mid m_1,\cdots ,m_k \in \mathbf Z\}
\]
と書ける．
この$k$を次元という．

更に，$G^*$は$G$の正規部分群である．なぜなら，
$T_{\mathbf a}\in G^*$と $F\in G$に対して
\[
\begin{split}
	F\circ T_{\mathbf a}\circ F^{-1}(\mathbf x)
	&=
	F\circ T_{\mathbf a}(A^{-1}\mathbf x - A^{-1}\mathbf b)\\
	&=
	F(A^{-1}\mathbf x - A^{-1}\mathbf b+\mathbf a)\\
	&=
	A(A^{-1}\mathbf x - A^{-1}\mathbf b+\mathbf a)+\mathbf b
	=
	\mathbf x  +A\mathbf a
\end{split}
\]
なので$F\circ T_{\mathbf a}\circ F^{-1}= T_{A\mathbf a}\in G^*$, 
つまり$F G^* F^{-1}\subset G^*$
だからである．
\emptyline

$r(F)=A\in r(G)$なので，この作用で$r(G)\subset O(n)$は$G^*$に作用している．

\emptyline

$H$, $G$, $K$を群として，$f:H\to G$, $g:G\to K$が群の準同形で
\[
	H\overset{f}{\to} G\overset{g}{\to} H
\]
が\term[]{完全}とは
\[
	\Image f =\{f(a)\in G\mid a\in H\},\quad
 	\Ker g=\{b\in G\mid g(b)=0(=e)\}
 \]と置くとき
 $\Image f=\Ker g$が成り立つことである．
 例えば
 \[
	0\overset{f}{\to} G\overset{g}{\to} H
\]
が完全なときは$g$は単射であり，
 \[
	H\overset{f}{\to} G\overset{g}{\to} 0
\]
が完全なときは$f$は全射である．
 幾つかの群の系列にがあるとき，
 連続する全ての準同形のペアが完全なとき完全という．
\emptyline

・次の群の系列
\[
	0\to G^*\overset{i}{\to} G \overset{r}{\to} r(G)\to 0
\]
は完全系列である．ただし$i$は自然な埋め込み，最後の$0$は正確には積の単位元 $E$のこと．

∵）$G^*\subset G$なので$i$は単射．
$A\in r(G)\subset O(n)$ に対して$F\in G$が存在して$A=r(F)$は定義なので$r$は全射．$r(F)=E$と$F$が平行移動は同値なので完全．

注意）$\Is(\mathbf R^n)$は\eqref{semiproduct}で積が定まる
$T(\mathbf R^n)$と$O(n)$の\term[]{半直積}であり，このことを
\[
	\Is(\mathbf R^n)=T(\mathbf R^n)\rtimes O(n)
\]
と表す．一般の半直積の定義は一見すると分かりづらいが，この状況を一般化するように定式化されたものである．
\emptyline

%$n=2$のとき以下を示した：
%\begin{Prop}[結晶条件]
%$\mathbf 0$を中心とした回転$S\neq E$
%が$2$次元の格子$G^*$を保つとすると，$S$の回転角$\theta$は
%\[
%	\theta=\pm 2\pi/p\quad p=2, 3,4,6
%\]
%の何れか．
%\end{Prop}
%\noindent
%なので，
$r(G)$は$O(2)$の有限部分群である．このような群を\term[]{点群}という．

点群の分類は既に終わっているので，
格子の次元毎に格子の形とその作用を調べれば離散部分群の分類が終わる．

\begin{figure}[htb]
   \centering
   \begin{minipage}{0.4\columnwidth}
      \centering
      \includegraphics{~/Documents/GitHub/bEuclideanGeo/ch2/images/figureGroup8.1}
      \caption{並進鏡映}\label{ref3}
   \end{minipage} 
   \begin{minipage}{0.4\columnwidth}
      \centering
      \includegraphics{~/Documents/GitHub/bEuclideanGeo/ch2/images/figureGroup8.2}
      \caption{鏡映}
   \end{minipage} 
\end{figure}


$r(G)$が$\subset SO(2)$のときは既に終わったので，$
 r(G)\setminus SO(2)\neq\emptyset$のときを考える．
 \begin{Thm}[シャール]
$\mathbf R^2$の合同変換は
平行移動，ある点を中心とした回転，並進鏡映（鏡映）のいずれか．
\end{Thm}
\noindent
なので，これは$G$が鏡映か並進鏡映（あるいはその両方）を含むことを意味している．


case 0)
格子$G^*$が$0$次元の時は以前示したように一点を固定するので，
$C_m$, $D_m$($m\ge 1$)を点群とする群で，この場合は
\[
	0\to 0\to G \overset{\simeq }{\to} r(G)\to 0
\]
より$G$と一致する．
\emptyline

case 1)
格子$G^*$が$1$次元の時，$SO(2)$の有限部分群で点群になるのは$C_1$, $C_2$で，$C_1$が点群の時が$F_1$, $C_2$が点群の時が$F_2$であった．$D_n$は$C_n$を正規部分群として含むので，
%$O(2)$の有限部分群の分類と，鏡映が$1$次元格子を保つことより，
この場合の点群は
\[
	C_1,\quad C_2,\quad
	D_1,\quad D_2
\]
のいずれかとなる．

\begin{figure}[htb]
   \centering
   \begin{minipage}{0.4\columnwidth}
   	\centering
   	\includegraphics{~/Documents/GitHub/bEuclideanGeo/ch2/images/figureGroup8.16}
   	\caption{$C_1=\{E\}$と$D_1$}\label{}
   \end{minipage} 
   \begin{minipage}{0.4\columnwidth}
   	\centering
   	\includegraphics{~/Documents/GitHub/bEuclideanGeo/ch2/images/figureGroup8.15}
   	\caption{$C_2$と$D_2$}\label{}
   \end{minipage} 
\end{figure}


\begin{figure}[htb]
   \centering
   \begin{minipage}{0.4\columnwidth}
      \centering
      \includegraphics{~/Documents/GitHub/bEuclideanGeo/ch2/images/figureGroup8.3}
      \caption{$F_1$}
   \end{minipage} 
   \begin{minipage}{0.4\columnwidth}
      \centering
      \includegraphics{~/Documents/GitHub/bEuclideanGeo/ch2/images/figureGroup8.4}
      \caption{$F_2$}
   \end{minipage} 
\end{figure}

点群が$D_1$または$D_2$のとき
一個の鏡映の格子$G^*=\{T_{n\mathbf a}\}$への作用で可能なものは，
鏡映の軸を$\ell$とするとき，$\ell$が$\mathbf a$に平行なものと，
$\ell$が$\mathbf a$に直交したものの二つがある．
また，鏡映でない並進鏡映の合成は平行移動なので，鏡映でない並進鏡映の軸は$\mathbf a$に平行でなければならない．また，鏡映や並進鏡映と平行移動を合成すると鏡映や並進鏡映になる．これらを使い，群の可能な作用をしらみつぶしで調べると次のことがわかる．

\begin{Thm}
格子$G^*$が$1$次元のとき，離散部分群$G$は表 \ref{7}の
7通りある．
\end{Thm}

\begin{table}[h]
   \centering
   \begin{tabular}{lll} 
  	\noalign{\hrule height 1pt}
	 点群& 格子&7個の離散群（以下国際記号） \\ \hline \hline
	$C_1$&一次元格子&$P1$\\ \hline
	$C_2$ &一次元格子&$P2$\\ \hline
	$D_1$ &一次元格子&$P1m1$, $P11m$, $P11g$\\	 \hline
	$D_2$ &一次元格子&$P2mm$, $P2mg$\\	
	\noalign{\hrule height 1pt}
   \end{tabular}
      \caption{格子が$1$次元のときの離散部分群}\label{7}
\end{table}


\begin{Quiz}$10$群の表示の意味を調べ，（図~\ref{F1}--\ref{F5}の図も参考にして）この分類を示せ．もちろん一部分でも構わない．
\end{Quiz}

%\cite{NS} の p179〜183に図形による分類表がある．
%\cite{NS}は（数理内からアクセスすれば）
%\href{https://link.springer.com/}{ここから}
%ダウンロード可能．

\emptyline


case 2)
格子$G^*$が$2$次元の時，$O(2)$の有限部分群の分類と結晶条件より，点群$r(G)$は
\[
	C_1,\quad C_2,\quad C_3,\quad C_4,\quad C_6,\quad
	D_1,\quad D_2,\quad D_3,\quad D_4,\quad D_6
\]
のいずれかである．


一次元の場合と異なり，二次元の格子自体が幾何学的情報を持っているので，格子自体の構造を求める必要がある．
点群が$C_1=\{E\}$の群は$M_1$, 点群が$C_2=\{E,P\}$($P$は$\pi$回転 )の群は$M_2$であり，それらの格子は\term[]{一般格子(oblique lattice)}と呼ばれるもので，基本ベクトルは大きさが異なり，なす角度が
$\pi/2$, $\pi/3$, $2\pi/3$でもないものである．

\begin{figure}[htb]
   \centering
   \begin{minipage}{0.4\columnwidth}
      \centering
      \includegraphics{~/Documents/GitHub/bEuclideanGeo/ch2/images/figureGroup8.5}
      \caption{$M_1$}\label{M1}
      \end{minipage} 
   \begin{minipage}{0.4\columnwidth}
      \centering
      \includegraphics{~/Documents/GitHub/bEuclideanGeo/ch2/images/figureGroup8.6}
      \caption{$M_2$}\label{M2}
   \end{minipage} 
\end{figure}


$D_m$($m=3, 4, 6$)は必ず$C_m$を正規部分群として含むので，既に分類済みの$M_3$, $M_4$, $M_6$に（並進）鏡映を付け加えたものである．したがって，この場合の格子$G^*$は\term[]{正六角形格子}と\term[]{正方格子}の何れかである．

残りは点群が$D_1$, $D_2$の場合である．任意の2次元格子$G^*$は$\pi$回転で不変であり，$D_1$は鏡映一個から成るが$\pi$回転と合成すると新たな鏡映が生まれる．これは$D_1$の元ではないが格子$G^*$はこの対称性を持っていないといけない．従って，
点群が$D_2$の場合に不変な格子を探せば良い．

$D_2$の二つの異なる2つ鏡映の合成は$\pi$回転なので，
二つの異なる2つ鏡映の軸は直交していて，
鏡映と$\pi$回転の合成はもう一つの鏡映を生み出すので，この対称性のある格子を見つける必要がある．

\begin{figure}[htb]
   \centering
   \includegraphics{~/Documents/GitHub/bEuclideanGeo/ch2/images/figureGroup8.7}
   \caption{鏡映と$\pi$回転の合成}\label{}
\end{figure}

以前示したように基本ベクトルは以下のように見つけることが出来る：
$T\in G^*$に対して$\mathbf a=T(\mathbf 0)$と置くと，$T=T_{\mathbf a}$である．$G$は離散なので，ある$\mathbf a\neq\mathbf 0\in G^*\mathbf 0$が存在して
\[
	\|\mathbf a\|
	=\min_{\mathbf b\neq\mathbf 0 \in G^*\mathbf 0}\|\mathbf b\|
\]
このとき，仮定から
\[
	G^*\mathbf 0\supsetneqq \{n\mathbf a\mid n\in \mathbf Z\}
\]
なので，$\mathbf b\in G\mathbf 0\setminus  \{n\mathbf a\mid n\in \mathbf Z\}
$ で
\[
	\|\mathbf b\|
	=
	\min_{
		\mathbf c\in 
		G^*\mathbf 0
		\setminus  \{n\mathbf a\mid n\in \mathbf Z\}
	}
	\|\mathbf c\|
\]
と選ぶことが出来る．このとき
\[
	\|\mathbf a\|\le \|\mathbf b\|,
	\quad
	 -\frac12\|\mathbf a\|
	 \le
	 (\mathbf b, \mathbf a)
	 \le 
	 \frac12\|\mathbf a\|
\]

\begin{figure}[htb]
   \centering
   \begin{minipage}{0.4\columnwidth}
      \centering
      \includegraphics{~/Documents/GitHub/bEuclideanGeo/ch2/images/figureGroup8.8}
      \caption{基本ベクトル}\label{elementaryVectors}
   \end{minipage} 
   \begin{minipage}{0.4\columnwidth}
      \centering
      \includegraphics{~/Documents/GitHub/bEuclideanGeo/ch2/images/figureGroup8.9}
      \caption{基本ベクトルの特別な配置}\label{elementaryVectors2}
   \end{minipage} 
\end{figure}


\begin{figure}[htb]
   \centering
   \begin{minipage}{0.4\columnwidth}
      \centering
      \includegraphics{~/Documents/GitHub/bEuclideanGeo/ch2/images/figureGroup8.10}
      \caption{長方形格子}\label{}
   \end{minipage} 
   \begin{minipage}{0.4\columnwidth}
      \centering
      \includegraphics{~/Documents/GitHub/bEuclideanGeo/ch2/images/figureGroup8.11}
      \caption{菱形格子}\label{}
   \end{minipage} 
\end{figure}

$\|\mathbf a\|< \|\mathbf b\|$とする．
$D_2$の2つの鏡映を$R$, $S$とする．仮定より$\mathbf a$は最小なので
\[
	R\mathbf a  = -\mathbf a,\quad\text{あるいは}\quad
	R\mathbf a  = \mathbf a,
\]
が成り立つが，
後者の場合，更に，$S\mathbf a=\mathbf a$とすると二次元の鏡映なので，$R=S$となり矛盾．よって$S\mathbf a=-\mathbf a$でなければならない．そこで名前を入れ替えて常に$R\mathbf a  = -\mathbf a$と仮定してよい．そのとき, 格子$G^*$を保存することから
$R\mathbf b$の可能性を検討すると，図~\ref{elementaryVectors}, \ref{elementaryVectors2} より
\[
	R\mathbf b=\mathbf b,\quad\text{あるいは}\quad
	R\mathbf b=\mathbf b-\mathbf a,\quad\text{あるいは}\quad
	R\mathbf b=\mathbf b+\mathbf a
\]
の何れかが成り立つ．このとき
\[
	S\mathbf b=-\mathbf b,\quad\text{あるいは}\quad
	S\mathbf b=-\mathbf b+\mathbf a, 
	\quad\text{あるいは}\quad
	S\mathbf b=-\mathbf b-\mathbf a
\]
最初の格子が\term[]{長方形格子(rectangular lattice)}，最後の格子が(第2番目と第3番目は実は同じ格子である)\term[]{菱形格子(rhombic  lattice)}である．なお，菱形格子の基本ベクトルとして$\mathbf a$の代わりに$\mathbf b'=\mathbf a-\mathbf b$を取ることがあることを注意しておく．


\noindent
ただし，$\|\mathbf a\|= \|\mathbf b\|$の場合の
正六角形格子と正方格子はその特殊な場合と解釈できる．

\begin{figure}[htb]
   \centering
   \includegraphics{~/Documents/GitHub/bEuclideanGeo/ch2/images/figureGroup8.14}
   \caption{正六角形格子と$M_3$}\label{comp1}
\end{figure}
\begin{figure}[htb]
   \centering
   \includegraphics{~/Documents/GitHub/bEuclideanGeo/ch2/images/figureGroup8.13}
   \caption{正方格子と$M_4$}\label{comp2}
\end{figure}


これと点群$r(G)$の格子群への作用を見ることで
\begin{Thm}
   格子$G^*$が$2$次元のとき，離散部分群$G$は表 \ref{17}の
   $17$通りある．
\end{Thm}

\begin{table}[bht]
   \centering
   \begin{tabular}{lll} 
  	\noalign{\hrule height 1pt}
	 点群& 格子&17個の離散群（以下国際記号） \\ \hline \hline
	$C_1$ &一般格子&$P1$\\ \hline
	$C_2$ &一般格子& $P2$\\ \hline
	$D_1$ &長方形格子&$Pm$, $Pg$, $P2gg$  \\ \hline
	$D_2$ &長方形格子&$P2mm$, $P2gm$ \\ \hline
	 $D_1$ &菱形格子& $Cm$\\ \hline
	 $D_2$ &菱形格子& $C2mm$\\ \hline	 
	 $C_4$ &正方格子&$P4$\\  \hline
	 $D_4$ &正方格子&$P4mm$, $P4gm$\\  \hline
	 $C_3$&正六角形格子&$P3$\\   \hline
	 $D_3$&正六角形格子&$P3m1$, $P31m$\\  \hline  
	 $C_6$ &正六角形格子&$P6$\\  \hline  
	 $D_6$&正六角形格子&$P6mm$\\  
	\noalign{\hrule height 1pt}
   \end{tabular}
   \caption{格子が$2$次元のときの離散部分群}\label{17}
\end{table}


\begin{Quiz}$11$
   群の表示の意味を調べ（図~\ref{P1}--\ref{P6mm}も参考にし），この分類を示せ．もちろん一部分でも構わない．
\end{Quiz}



%\begin{thebibliography}{9}
%\bibitem{NS}
%Viacheslav V. Nikulin, Igor R. Shafarevich, \textit{Geometries and Groups}, UTX, 1996, Springer, \href{https://link.springer.com/}{ここからダウンロード可能}
%
%\bibitem{E}
%C.H.マックギラフィ, 有馬 朗人 , エッシャー—シンメトリーの世界
%
%\bibitem{I}
%International Tables for Crystallography: vol. A, Space-Group Symmetry
%\end{thebibliography}
%
%\newpage

\subsection{$2$次元結晶群の分類の図示}\label{2dim}
ユークリッド平面の等長群の離散部分群を結晶群・装飾群といい，
その中でも万華鏡とは生成元が鏡映のみのものを言う．
鏡映を鏡に変え，上から見ることで実際の万華鏡を作ることが出来る．
但し，物理的な万華鏡は光を使うので数学的な理想化とは微妙に異なる性質を持つことがある．

合わせ鏡を誰が最初に行ったか定かではないが，
1816年イギリスの物理学者ブルースター David Brewster(1781‐1868)が二枚の一方の端を合わせた合わせ鏡を上から見ることで万華鏡(kaleidoscope,  図 \ref{D5},  \ref{D6} 参照)を作ったのが万華鏡の最初であるといわれている．

\emptyline

国際記号(International nomenclature\cite{I})と図の作成方法について：

\begin{enumerate}
\item
格子, $P$ :プリミティブ，
$C$ :(centerd)面心 

\item
$1$, $2$, $3$, $4$, $6=:n$は $C_n$($ \lhd\, D_n$)を表す．
回転の中心を色のついた円で表し，同じ色は同じ$G$-軌道を表す．

\item
$m$:鏡映 その軸を太い直線，
$g$:並進鏡映 その軸を破線

\item
基本領域がわかりにくいときは点線で表示．
また，三角形を基本領域の中に描き，群作用でそれがどのように移るかを図示した．幾つかには格子の基本領域を灰色で示している．

\item
キャプションに：群の国際記号, 格子, 点群（ 万華鏡として実現出来るときは万華鏡）と記す.
\end{enumerate}


\emptyline


$2$次元格子を格子とする装飾群を図~\ref{P1}--\ref{P6mm}に，
$1$次元格子を格子とする装飾群を図~\ref{F1}--\ref{F5}に挙げる．
また，
$0$次元格子を格子（つまり自明な平行移動$\mathbf 0$の作る群）とする$C_m$, $D_m$系列の例を図~\ref{C5}--\ref{D6}に挙げる．

\begin{figure}[h]
   \centering
   \begin{minipage}{0.4\columnwidth}
      \centering
      \includegraphics{~/Documents/GitHub/bEuclideanGeo/ch2/images/figureGroup81.1}
      \caption{ $P1$ ( $M_1$ ), 一般, $C_1$}\label{P1}
   \end{minipage} 
   \begin{minipage}{0.4\columnwidth}
      \centering
      \includegraphics[trim=20 0  0 35, clip]{~/Documents/GitHub/bEuclideanGeo/ch2/images/figureGroup81.2}
      \caption{ $P2$ ( $M_2$ ), 一般, $C_2$}\label{}
   \end{minipage} 
\end{figure}
\begin{figure}[h]
   \centering
   \begin{minipage}{0.4\columnwidth}
      \centering
      \includegraphics{~/Documents/GitHub/bEuclideanGeo/ch2/images/figureGroup81.3}
      \caption{ $Pm$, 長方形, $D_1$}\label{}
      \end{minipage} 
   \begin{minipage}{0.4\columnwidth}
      \centering
      \includegraphics{~/Documents/GitHub/bEuclideanGeo/ch2/images/figureGroup81.4}
      \caption{ $Pg$, 長方形, $D_1$}\label{}
   \end{minipage} 
\end{figure}

\begin{figure}[h]
   \centering
   \begin{minipage}{0.4\columnwidth}
      \centering
      \includegraphics{~/Documents/GitHub/bEuclideanGeo/ch2/images/figureGroup81.5}
      \caption{ $Cm$, 菱形, $D_1$}\label{}
      \end{minipage} 
   \begin{minipage}{0.4\columnwidth}
      \centering
      \includegraphics{~/Documents/GitHub/bEuclideanGeo/ch2/images/figureGroup81.6}
      \caption{ $P2mm$, 長方形, $D_2$, 万華鏡}\label{}
   \end{minipage} 
\end{figure}

\begin{figure}[h]
   \centering
   \begin{minipage}{0.4\columnwidth}
      \centering
      \includegraphics[trim=10 0  10 0, clip]{~/Documents/GitHub/bEuclideanGeo/ch2/images/figureGroup81.7}
      \caption{ $P2gm$, 長方形, $D_2$ }\label{}
   \end{minipage} 
   \begin{minipage}{0.4\columnwidth}
      \centering
      \includegraphics[trim=20 30  10 45, clip]{~/Documents/GitHub/bEuclideanGeo/ch2/images/figureGroup81.8}
      \caption{  $P2gg$, 長方形, $D_2$ }\label{}
   \end{minipage} 
\end{figure}

\begin{figure}[h]
   \centering
   \includegraphics{~/Documents/GitHub/bEuclideanGeo/ch2/images/figureGroup81.9}
   \caption{ $C2mm$, 菱形, $D_2$ }\label{}
\end{figure}

\begin{figure}[htb]
	\centering
	\includegraphics[trim=0 10  0 10, clip]{~/Documents/GitHub/bEuclideanGeo/ch2/images/figureGroup81.10}
	\caption{ $P4$ ($M_4$), 正方形, $C_4$ }\label{}
\end{figure}


\begin{figure}[htb]
   \centering
   \includegraphics[trim=10 10  10 10, clip]{~/Documents/GitHub/bEuclideanGeo/ch2/images/figureGroup81.11}
   \caption{ $P4mm$, 正方形, $D_4$, 万華鏡 }\label{}
\end{figure}

\begin{figure}[htb]
   \centering
   \includegraphics[trim=10 20  20 20, clip]{~/Documents/GitHub/bEuclideanGeo/ch2/images/figureGroup81.12}
   \caption{ $P4gm$, 正方形, $D_4$ }\label{}
\end{figure}

\begin{figure}[h]
   \centering
   \includegraphics[trim=30 40  30 145, clip]{~/Documents/GitHub/bEuclideanGeo/ch2/images/figureGroup81.13}
   \caption{ $P3$ ($M_3$), 正六角形, $C_3$ }\label{}
\end{figure}


\begin{figure}[h]
   \centering
   \includegraphics[trim=40 65  40 100, clip]{~/Documents/GitHub/bEuclideanGeo/ch2/images/figureGroup81.14}
   \caption{$P3m1$($P3mm$,), 正六角形, $D_3$, 万華鏡 }\label{}
\end{figure}

\begin{figure}[h]
   \centering
   \includegraphics[trim=29 100  4 100, clip]{~/Documents/GitHub/bEuclideanGeo/ch2/images/figureGroup81.15}%
   \caption{$P31m$, 正六角形, $D_3$ }\label{}
\end{figure}



\begin{figure}[h]
   \centering
   \includegraphics[trim=20 80 30 100, clip]{~/Documents/GitHub/bEuclideanGeo/ch2/images/figureGroup81.16}
   \caption{$P6$（$M5$）, 正六角形, $C_6$ }\label{}
\end{figure}


\begin{figure}[h]
   \centering
   \includegraphics[trim=30 80  30 90, clip]{~/Documents/GitHub/bEuclideanGeo/ch2/images/figureGroup81.17}
   \caption{$P6mm$, 正六角形, $D_6$, 万華鏡 }\label{P6mm}
\end{figure}



\begin{figure}[h]
   \centering
   \begin{minipage}{0.4\columnwidth}
      \centering
      \includegraphics[trim=0 0  0 10, clip]{~/Documents/GitHub/bEuclideanGeo/ch2/images/figureGroup81.21}
      \caption{ $P1$, $\mathbf Z$, $C_1$}\label{F1}
      \end{minipage} 
   \begin{minipage}{0.4\columnwidth}
      \centering
      \includegraphics{~/Documents/GitHub/bEuclideanGeo/ch2/images/figureGroup81.22}
      \caption{ $P2$, $\mathbf Z$, $C_2$}\label{}
      \end{minipage} 
\end{figure}

\begin{figure}[h]
   \centering
   \begin{minipage}{0.4\columnwidth}
      \centering
      \includegraphics[trim=0 20 25 0, clip]{~/Documents/GitHub/bEuclideanGeo/ch2/images/figureGroup81.24}
      \caption{ $P1m1$, $\mathbf Z$, $D_1$}\label{}
   \end{minipage} 
   \begin{minipage}{0.4\columnwidth}
      \centering
      \includegraphics[trim=0 20 0 -40, clip]{~/Documents/GitHub/bEuclideanGeo/ch2/images/figureGroup81.27}
      \caption{$P11g$, $\mathbf Z$, $D_1$ }\label{}
   \end{minipage} 
\end{figure}


\begin{figure}[h]
   \centering
   \includegraphics[trim=0 30 20 20, clip]{~/Documents/GitHub/bEuclideanGeo/ch2/images/figureGroup81.23}
   \caption{ $P11m$, $\mathbf Z$, $D_1$, 万華鏡}\label{}
\end{figure}


\begin{figure}[h]
   \centering
   \begin{minipage}{0.4\columnwidth}
      \centering
      \includegraphics[trim=0 5 25 -10, clip]{~/Documents/GitHub/bEuclideanGeo/ch2/images/figureGroup81.26}
      \caption{ $P2mm$, $\mathbf Z$, $D_2$, 万華鏡}\label{}
   \end{minipage} 
   \begin{minipage}{0.4\columnwidth}
      \centering
      \includegraphics[trim=5 20 0 -10, clip]{~/Documents/GitHub/bEuclideanGeo/ch2/images/figureGroup81.25}
      \caption{ $P2mg$, $\mathbf Z$, $D_2$}\label{F5}
   \end{minipage} 
\end{figure}

\begin{figure}[t]
   \centering
   \begin{minipage}{0.4\columnwidth}
      \centering
      \includegraphics{~/Documents/GitHub/bEuclideanGeo/ch2/images/figureGroup81.31}
      \caption{ $C_{5}$, $\mathbf 0$, $C_{5}$}\label{C5}
   \end{minipage} 
   \begin{minipage}{0.4\columnwidth}
      \centering
      \includegraphics{~/Documents/GitHub/bEuclideanGeo/ch2/images/figureGroup81.32}
      \caption{ $C_{6}$, $\mathbf 0$, $C_{6}$}\label{C6}
   \end{minipage} 
\end{figure}

\begin{figure}[t]
   \centering
   \begin{minipage}{0.4\columnwidth}
      \centering
      \includegraphics{~/Documents/GitHub/bEuclideanGeo/ch2/images/figureGroup81.29}
      \caption{ $D_{5}$, $\mathbf 0$, $D_{5}$, 万華鏡}\label{D5}
   \end{minipage} 
   \begin{minipage}{0.4\columnwidth}
      \centering
      \includegraphics{~/Documents/GitHub/bEuclideanGeo/ch2/images/figureGroup81.30}
      \caption{ $D_{6}$, $\mathbf 0$, $D_{6}$, 万華鏡}\label{D6}
   \end{minipage} 
\end{figure}
\clearpage

\newpage


\subsection{2次平坦多様体の分類}
$G$を$\Is(\mathbf R^2)$の離散部分群とする．
$G$が\term[]{自由に作用}するとは
\[
	\exists \mathbf x_0\in \mathbf R^2, F\in G\quad \text{ such that }
	F(\mathbf x_0)=\mathbf x_0 \Rightarrow F=\id
\]
がなる立つこと．最初の条件$F(\mathbf x_0)=\mathbf x_0$は固定点を持つという条件で，
%\term[]{作用が自由}とは固定点を持たないことである．
$G$が自由に作用するとは，固定点を持つのは自明な作用$\id$に限るということである(英語では fixed point free )．
\emptyline

$\Is(\mathbf R^2)$の離散部分群$G$が自由に作用するとする．
定義から，任意の$\mathbf x\in \mathbf R^2$ と $S\neq \id \in G$に対して
\[
	S\mathbf x\neq \mathbf x
\]
$G$は離散，
つまり
\[
	\forall \mathbf x\in \mathbf R^2, 
	\exists \varepsilon>0 \text{ such that }
	B(\mathbf x, \varepsilon)\cap G\mathbf x
	=\{\mathbf x\}
\]
である．
\begin{figure}[htb]
   \centering
   \begin{minipage}{0.4\columnwidth}
      \centering
      \includegraphics{~/Documents/GitHub/bEuclideanGeo/ch2/images/figureGroup9.1}
      \caption{平行移動}
   \end{minipage} 
   \begin{minipage}{0.4\columnwidth}
      \centering
      \includegraphics{~/Documents/GitHub/bEuclideanGeo/ch2/images/figureGroup9.2}
      \caption{回転}
   \end{minipage} 
\end{figure}

\begin{figure}[htb]
   \centering
   \begin{minipage}{0.4\columnwidth}
      \centering
      \includegraphics{~/Documents/GitHub/bEuclideanGeo/ch2/images/figureGroup9.3}
      \caption{並進鏡映}\label{ref3}
   \end{minipage} 
   \begin{minipage}{0.4\columnwidth}
      \centering
      \includegraphics{~/Documents/GitHub/bEuclideanGeo/ch2/images/figureGroup9.4}
      \caption{鏡映}
   \end{minipage} 
\end{figure}

$\Is(\mathbf R^2)$の離散部分群$G$が自由に作用するとする．
平行移動でない回転と鏡映は明らかに固定点を持つので$G$に含まれない．
$\id$でない平行移動と鏡映でない並進鏡映は固定点がないので$G$に含まれる可能性がある．
よってある$r>0$が存在して
\[
	r=
	\min_{S\neq \id \in G, \mathbf x\in  \mathbf R^2}
		\|S\mathbf x- \mathbf x\|>0
\]

関係を
\[
	\mathbf x\sim \mathbf y
	\Leftrightarrow
	\exists S \in G\text{ such that } \mathbf y=S \mathbf x
\]
で定めると，これは同値関係なので商空間$\mathbf R^2/G$が定まる．
$\mathbf x$の同値類を$ [\mathbf x]$と表し，自然な射影
\[
	\pi: \mathbf R^2\ni \mathbf x \mapsto [\mathbf x]
	\in\mathbf R^2/G
\]
を考える．$\pi$を$B(\mathbf x,r)$に制限すると同相となり，多様体\footnote{多様体とはユークリッド空間の開集合と同相な開集合で被覆されるハウスドルフ空間のことである．これらの同相を座標近傍といい，任意の二つの同相の作るユークリッド空間の開集合の間の変換が全て$C^\infty$写像のとき，それらを$C^\infty$級微分構造と喚び，そのような空間を$C^\infty$級多様体という．}の構造が入る．
さらに内積をリーマン計量と考えると，$\mathbf R^2/G$は完備
な平坦リーマン多様体になる．ただし，\term[]{平坦}とはガウス曲率$K$が$= 0$のことで，\footnote{これは測地線がいくらでも伸ばせるという意味での測地線完備と，リーマン多様体が距離空間として完備, つまり，任意のコーシー列が収束することと同値である}{\term[]{完備}}とは全ての測地線をどこまでも伸ばせることである．
ただし，測地線は局所的に最短線であるのでこの場合は局所的には
直線となっている．
$\mathbf R^2$の直線を$\pi$で射影したものが
$\mathbf R^2/G$の測地線であり，完備であることが分かる．
角度も普通に内積を使ったものである．
$B(\mathbf x,r)$内の最短線からなる測地三角形$\triangle ABC$を考えると，
よく知っているように
\[
	\angle A+\angle B+\angle C=\pi
\]
がなりたつが，これを境界が測地線の場合のガウスボンネの公式
\[
\begin{split}
	\int_{\triangle ABC}Kd\vol 
	&= 2\pi - (\pi-\angle A)-(\pi-\angle B)-(\pi-\angle C)\\
	&=\angle A+\angle B+\angle C-\pi%=0
\end{split}
\]
と較べるとガウス曲率$K$が$=0$であることが分かる．ただし，
$d\vol$は面積要素．

$\mathbf R^2$は単連結で，$\pi$は普遍被覆空間，
$G$は空間$\mathbf R^2/G$の基本群（または，デッキ変換群）となっているので，逆も言えて
\begin{Q}
$2$次元完備平坦リーマン多様体の分類

$\Leftrightarrow$

$\mathbf R^2$に自由に作用する$\Is(\mathbf R^2)$の離散部分群の分類．
\end{Q}
以下この方針で$2$次元完備平坦リーマン多様体の分類を与える．

$\Is(\mathbf R^2)$の離散部分群$G$が自由に作用するとする．
上で指摘したように回転と鏡映は明らかに固定点を持つので$G$に含まれない．つまり$G$の点群は回転を含まない．
$\id$でない平行移動と鏡映でない並進鏡映は固定点がないので$G$に含まれる可能性があるが，$r(G)$が鏡映を含んでもそれは並進鏡映から得られたものであり，
鏡映と鏡映の合成は回転なので，並進鏡映は含んでも同じ直線に沿ったものだけである．
以上から

\begin{Prop}
$\Is(\mathbf R^2)$の自由な作用の離散部分群の点群は
$C_1$, $D_1$のみ．
\end{Prop}

$G$を自由な作用のみからなる$\Is(\mathbf R^2)$の離散部分群とする．
\emptyline

case 0)
$G^*$が$0$次元．もし，点群が$D_1$で並進鏡映を含めば，
その合成が平行移動となるのでそのようなことはおこらない．
鏡映もふくまないので点群は$C_1=\{E\}$のみ．
\[
	0\to 0\to G \overset{\simeq }{\to}0\to 0
\]
つまり，$G$は自明な群である．商空間は$\mathbf R^2$全体である．
\emptyline

case 1)
$G^*$が$1$次元のとき，$\mathbf a\neq\mathbf 0$が存在して
$G^*=\{n\mathbf a\mid n\in \mathbf Z\}$．

点群が$C_1=\{E\}$だと$G^*=\{T_{n\mathbf a}\mid n\in \mathbf Z\}$となり，
\[
	0\to G^* \overset{\simeq }{\to} G \to 0\to 0
\]
なので，$G\simeq G^*$は$F_1$となる．商空間は円柱になる．

\begin{figure}[htb]
\centering
   \begin{minipage}{0.4\columnwidth}
      \centering
      \includegraphics{~/Documents/GitHub/bEuclideanGeo/ch2/images/figureGroup9.5}
      \caption{$F_1$}\label{F_1}
   \end{minipage} 
   \begin{minipage}{0.4\columnwidth}
      \centering
      \includegraphics{~/Documents/GitHub/bEuclideanGeo/ch2/images/figureGroup9.6}
      \caption{$\mathbf R^3$の円柱}\label{cylinder}
   \end{minipage} 
\end{figure}

\begin{figure}[htb]
   \centering
   \begin{minipage}{0.4\columnwidth}
      \centering
      \includegraphics{~/Documents/GitHub/bEuclideanGeo/ch2/images/figureGroup9.7}
      \caption{$F_2$}\label{F_2}
   \end{minipage} 
   \begin{minipage}{0.4\columnwidth}
      \centering
      \includegraphics{~/Documents/GitHub/bEuclideanGeo/ch2/images/figureGroup9.8}
      \caption{$\mathbf R^3$のメビウスの帯}\label{Mobius}
   \end{minipage} 
\end{figure}


点群が$D_1$で並進鏡映を含めば，並進鏡映の合成は平行移動なので
並進鏡映$S$の軸は$\mathbf a$に平行な直線$\ell$で
$\mathbf a/2 $ 進んで$\ell$で反転するものから生成されるとして良い．
点群は$D_1$なので他の並進鏡映は含めない．
商空間はメビウスの帯になる(図~\ref{Mobius} ただし，$\mathbf R^3$に等長に全体を埋め込ことはできない)．
\emptyline

case 2)
$G^*$が$2$次元の場合：
点群が$C_1$のときは$G^*$は一般格子で
$G$はそれに同形で$M_1$となり，商空間はトーラスになる
(図~\ref{Tours} ただし，埋め込みはリーマン多様体としての等長なものではない)．

\emptyline

\begin{figure}[htb]
   \centering
   \begin{minipage}{0.4\columnwidth}
      \centering
      \includegraphics{~/Documents/GitHub/bEuclideanGeo/ch2/images/figureGroup9.9}
      \caption{$M_1$}\label{M_1}
      \end{minipage} 
   \begin{minipage}{0.4\columnwidth}
      \centering
      \includegraphics{~/Documents/GitHub/bEuclideanGeo/ch2/images/figureGroup9.10}
      \caption{$\mathbf R^3$のトーラス}\label{Tours}
   \end{minipage} 
\end{figure}

\begin{figure}[htb]
   \centering
   \begin{minipage}{0.4\columnwidth}
      \centering
      \includegraphics{~/Documents/GitHub/bEuclideanGeo/ch2/images/figureGroup9.11}
      \caption{長方形格子}\label{}
   \end{minipage} 
   \begin{minipage}{0.4\columnwidth}
      \centering
      \includegraphics{~/Documents/GitHub/bEuclideanGeo/ch2/images/figureGroup9.12}
      \caption{菱形格子}\label{rombic}
   \end{minipage} 
\end{figure}



点群が$D_1$のときは点群と格子の分類(表 \ref{classification})より，
$G^*$は長方形格子か，菱形格子しかない．
以下で菱形格子は表れないことを示す：
$\mathbf a$, $\mathbf b$を菱形格子の基本ベクトルとする．
並進鏡映を一種類しか含めず，同じ並進鏡映の合成は
平行移動なので，その並進鏡映を$S$置くと
\[
	SS=T_{\mathbf a}, \quad SS=T_{\mathbf b}
\]
の何れかと仮定して良い．


\begin{figure}[htb]
   \begin{minipage}{0.4\columnwidth}
      \centering
      \includegraphics{~/Documents/GitHub/bEuclideanGeo/ch2/images/figureGroup9.13}
      \caption{$T_{-\mathbf b}S$は鏡映}\label{ref}
      \end{minipage} 
   \begin{minipage}{0.4\columnwidth}
      \centering
      \includegraphics{~/Documents/GitHub/bEuclideanGeo/ch2/images/figureGroup9.14}
      \caption{}\label{proof}
   \end{minipage} 
\end{figure}

最初の場合$S$は$\mathbf a$に平行な直線$\ell$ に沿った
$\mathbf a/2$の並進鏡映なので
$T_{-\mathbf b}$と合成すると
$\ell$を$-\mathbf b/2$だけ平行移動した
直線$\ell'$についての鏡映(図~\ref{ref})になり矛盾する．

次の場合$S$は$\mathbf b$に平行な直線$\ell$ に沿った
$\mathbf b/2$の並進鏡映として良いいので
$S\mathbf x=R\mathbf x +\mathbf b/2$と置く．
$r(G)$は$G^*$に線形に作用しているので
$R\mathbf a$は格子点でないといけないが，
$\|\mathbf a\|=\|\mathbf b\|$, つまり, 
$\mathbf 0$と$\mathbf a$と$\mathbf b$が正三角形の時を除いて一致しない(図~\ref{proof})．しかし，この時は$\mathbf a$が$\mathbf b/2$の真上にあるので，
最初と同じ配置になっていてやはり起こりえない．

点群が$D_1$のときは$G^*$は長方形格子で$G^*$は
直交する$\mathbf a$, $\mathbf b$で生成され（大小関係はここでは無視し）$\mathbf a$に平行な直線に沿った
$\mathbf a/2$の並進鏡映$S$を使って
\[
	G =\langle S, T_{\mathbf b}\rangle
\]
となり，商空間はクラインの壺になる（ 図~\ref{klein1}, \ref{klein2} ただし，この場合も多様体として$\mathbf R^3$に埋め込むことは出来ない）．


\begin{figure}[htb]
   \begin{minipage}{0.4\columnwidth}
      \centering
      \includegraphics{~/Documents/GitHub/bEuclideanGeo/ch2/images/figureGroup9.15}
      \caption{平坦クラインの壺}\label{klein1}
   \end{minipage} 
   \begin{minipage}{0.4\columnwidth}
      \centering
      \includegraphics{~/Documents/GitHub/bEuclideanGeo/ch2/images/figureGroup9.16}
      \caption{$\mathbf R^3$のクラインの壺}\label{klein2}
   \end{minipage} 
\end{figure}

纏めると
\begin{Thm}
$2$次元の完備平坦なリーマン多様体は次のいずれか：
\[
	\text{平面}\mathbf R^2, \quad
	\text{円柱}, \quad
	\text{メビウスの帯}, \quad
	\text{トーラス}, \quad
	\text{クラインの壺}
\]
ただし，メビウスの帯とクラインの壺は向き付け不可能%\footnote{任意の座標変換のヤコビアンが正になるような微分構造をもつ多様体が向き付け可能多様体で，そうできない多様体を向き付け不可能多様体という．}
である．
\end{Thm}

\section{結晶群の分類の歴史と発展}
この\S~\ref{class_2dim}で述べた$n=2$の装飾群の分類については
結果自体は群論が始まる19世紀よりはるか前から知られており，
14世紀に建てられたAlhambra宮殿の壁面の紋様に全て現れているそうである(\cite{NS} p154, ただし，
一般格子のP1の対称性しか持たない模様は発見されていないようなので
「ほとんど全て」が正しい表現であろう)．

$n= 3$の場合の分類：
19世紀の終わりにフョードロフ，シェンフリース，
バーローの3人によって互いに独立に，
格子群が14個，結晶群が
219個(向きを区別すると230個)あることが示された．
これらは現実の結晶の構造と関係していて，
20世紀になり結晶のX線回折(図~\ref{Xray})や場の量子論に応用された(\cite{I}の$125$.結晶群の部分に数学的にもう少し詳しい解説がある,
\cite{IT}も参照)．
$n=4$の時は結晶群の個数が$4783$個であることも示されている．


\begin{figure}[htb]
	\centering
	\includegraphics[width=8cm]{~/Documents/GitHub/bEuclideanGeo/ch2/images/Poederxrd.JPG}
	\caption{\href{https://commons.wikimedia.org/wiki/File:Poederxrd.JPG} Sokpopje at Dutch Wikipedia, CC BY-SA 3.0
	\href{ http://creativecommons.org/licenses/by-sa/3.0/}{\ccbysa}}\label{Xray}
\end{figure}

一般$n\ge 5$次元の場合：
$\mathbf R^n$の結晶群は有限個かというのが
ヒルベルトの第18問題と呼ばれる問題(の一つ)であり，
1910年ビーベルバッハにより肯定的に解決された．
この場合も完全系列
\[
	0\to G^*\overset{i}{\to} G \overset{r}{\to} r(G)\to 0
\]
の考察が重要であった．ただし$r(G)$は$O(n)$の有限部分群になり，$G^*=G\cap T(\mathbf R^n)$は格子群で$r(G)$が作用している\thetag{\cite{H}, \cite{Ch}}．

\begin{Quiz}
シャールの定理により $\mathbf R^2$ の等長変換は分類された．
そこで，$\mathbf R^3$ の等長変換はどのように分類されるかを示し，証明を与えよ\thetag{\cite{NS}, \cite{Y}}．
\end{Quiz}

\begin{Quiz}
$\mathbf R^3$ を張る格子の
点群（つまり，その格子を不変にする$O(3)$ の有限群、$32$個）の分類を与えよ\thetag{\cite{B}, \cite{Y}}．

ヒント：$O(3)$ の有限群の分類のなかで， 結晶条件を満たすものを全て見つけよ．
\end{Quiz}

\begin{Quiz}
それらの点群で不変となる（$14$個の）格子を分類せよ\thetag{\cite{B}, \cite{Y}}.
\end{Quiz}

\begin{Quiz}
$3$ 次元コンパクト平坦多様体（$10$個）の分類を与えよ\thetag{\cite{NS}}.
\end{Quiz}

%\begin{thebibliography}{99}
%\bibitem{B}
%Buerger,
%\newblock {\em Intorduction to crystal geometry}.
%\newblock  McGraw-Hill, 1974.
%
%
%\bibitem{C}
%L.~Charlap, 
%\newblock {\em Bieberbach groups and flat manifolds}.
%\newblock  Springer,
%\newblock 1986, 
%\href{https://link.springer.com/}{ここからダウンロード可能}
%
%\bibitem{H} 
%AMS,
%\newblock  {\em Mathematical developments arising from Hilbert Problems, part 2}, 
%\newblock 1974
%\bibitem{I} 
%日本数学会監修, 
%\newblock 岩波数学辞典 第四版, 
%\newblock 岩波書店, 
%\newblock 2007
%
%\bibitem{IT}
%\newblock {\em International Tables for Crystallography: vol. A, Space-Group Symmetry, new edition},
%\newblock Springer, 
%\newblock 2017.
%
%\bibitem{NS}
%Viacheslav V. Nikulin, Igor R. Shafarevich, \textit{Geometries and Groups}, UTX, 1996, Springer, 
%\href{https://link.springer.com/}{ここからダウンロード可能}
%
%\bibitem{Y}
%P.~Yale,
%\newblock \textit{Geometry and symmetry}.
%\newblock  Dover.
%
%\bibitem{K}
%スピーロ ; 青木 薫 訳,
%\newblock {ケプラー予想 〜四百年の難問が解けるまで〜}
%\newblock  新潮社.
%\end{thebibliography}

 





%\documentclass[a4paper,12pt]{amsart}
%\usepackage[dvipdfmx]{graphicx,xcolor,hyperref}%graphics
%\usepackage{float}
%\usepackage{amsmath}
%\usepackage{verbatim}
%\usepackage{latexsym}
%\usepackage{amssymb}
%\usepackage{ccicons}
%%\usepackage{emath}
%%\usepackage{maplestd2e}
%%\usepackage{fancybox}%enclose by box
%\hypersetup{% hyperrefオプションリスト
%setpagesize=false,
% bookmarksnumbered=true,
% bookmarksopen=true,
% colorlinks=true,
%linkcolor=red,
%citecolor=blue,
%} 
%
%\renewcommand{\figurename}{図}
%\renewcommand{\proofname}{証明}
%\renewcommand{\tablename}{表}
%\newtheorem*{PP}{授業の目標}
%\newtheorem*{Thm}{定理}
%\newtheorem*{Prop}{命題}
%\newtheorem*{Q}{考えたい課題}
%\newtheorem*{Quiz}{問題}
%
%\theoremstyle{definition}
%\newtheorem*{Def}{定義}
%
%\theoremstyle{remark}
%\newtheorem*{ex}{例}
%\newtheorem*{exQ}{例題}
%
%\def\emptyline{\vspace{12pt}}
%\DeclareMathOperator{\const}{const}
%\DeclareMathOperator{\Ker}{Ker}
%\DeclareMathOperator{\Image}{Im}
%\DeclareMathOperator{\id}{id}
%\DeclareMathOperator{\Is}{Is}
%\DeclareMathOperator{\stab}{stab}
%\DeclareMathOperator{\tr}{tr}
%\DeclareMathOperator{\vol}{vol}
%
%\begin{document}
%%\pagestyle{empty}
%\begin{flushright}
%
%2024年12月16日
%
%九大・数理 大津幸男
%\end{flushright}
%
%\begin{center}
%\large{微分幾何学大意・数学特論4\small{
%\href{https://creativecommons.org/licenses/by-nc-nd/4.0/}{\ccbyncnd}}\par ---遠近法と3DCGと射影幾何学1---}
%\end{center}
%

\chapter{ユークリッド幾何学から生まれた様々な幾何学}
\section{透視投影と3DCGと線形代数}
3次元コンピューターグラフィックス(以下3DCG)の数学的側面とプログラミング的側面が線型代数や行列の計算に関連していること，更にそれが射影幾何学や射影空間に関係していることを説明したい．

\subsection{3DCGと遠近法}
3DCGは現実の見え方をシュミレートしたものである．
ここではピラミッドの写真(図\ref{pyramid})を使って3DCGの原理を図示する(図\ref{model})．
\begin{figure}[htb]
   \centering
   \includegraphics[width=7.5cm]{~/Documents/GitHub/bEuclideanGeo/ch3/images/2428333990_d3c913b2b0_c.jpg}
   \caption{\href{https://www.flickr.com/photos/girolame/2428333990/in/photolist-4GvJhV-4GzR6q-68Y3JE-4GzTym-5DUvzA-4Gvzcg-5UXT1K-tAvViv-dmRS35-6j7Tdf-65175S-7hkV8W-7ZP8Wc-f7stxJ-6kADQr-92FxE-9MpvoF-7DWyqL-7DWzkj-7trRb8-7kV2Kh-74NeWm-6fnJRK-gBqRa-A4k2s-5DLk1m-A4jWP-6JPcdB-erLBS-7KLGJq-fWSgy-7bxjK9-5DRPd1-dduMZ-5DUGgw-6iWpt7-oueimu-oeLdfu-6GS2bv-6JCrjr-5Phx2c-7KGHVB-29jwn5-6JQBeo-oeDc4c-5DUBKb-dvw1V-5UXWtP-ouBgy3-gBn4K}{Giza Pyramids} by Jerome BonFollow 
   \href{https://creativecommons.org/licenses/by/2.0/?ref=openverse}{\ccby}}\label{pyramid}
\end{figure}
カメラでは像はレンズの後の壁にあり，しかも180°回転しているためイラストを作図しづらいので（他にもレンズがあると被写界深度が問題になることや，人間の目の構造上網膜が平坦でないことによるカメラや遠近法とのずれなど微妙な違いが有る．詳しくは図~\ref{pinCamera}, \ref{camera}, \ref{eye}を参考にしてほしい )，以下では遠近法を用いた作図をつかい，画家の視点から(透明な)キャンバス面への被写体を投影したものを考える．
\begin{figure}[htb]
   \centering
   \begin{minipage}{0.4\columnwidth}
      \centering
      \includegraphics{~/Documents/GitHub/bEuclideanGeo/ch3/images/figure.14}
      \caption{カメラ・オブスキュラ%フェスメール？
      }\label{pinCamera}
   \end{minipage} 
   \begin{minipage}{0.4\columnwidth}
      \centering
      \includegraphics{~/Documents/GitHub/bEuclideanGeo/ch3/images/figure.15}
      \caption{カメラ}\label{camera}
   \end{minipage} 
\end{figure}
\begin{figure}[htb]
   \centering
   \includegraphics{~/Documents/GitHub/bEuclideanGeo/ch3/images/figure.18}
   \caption{目---カメラとの微妙な違い}\label{eye}
\end{figure}

光は直進するので，視点と被写体上の点を結ぶ直線(視線)とキャンバスの平面との交わりを求めることが遠近法になる．
%これは連立一次方程式を解くことにあたる．
%特に，空間内の直線と視点を結ぶ直線全体は空間内の平面であるので，キャンバスとの交わりはキャンバス上の直線に投影される．

\begin{figure}[htb]
   \centering
   \includegraphics{~/Documents/GitHub/bEuclideanGeo/ch3/images/figure.2}
\caption{遠近法(透視投影)}\label{model}
\end{figure}

\subsection{(透視)投影}
\term[]{遠近法(perspective)}でどのように直線が描かれるかを説明したい：
原点$\mathbf 0$に視点があるとする．
$(x,y,z) \in \mathbf R^3$  ($x\neq 0$ )とし，
ここでは，$x$軸の正の方向に視線を向け，
$x=1$ にキャンバスがある場合を考える．
$U=\{(x,y,z)\in \mathbf R^3\mid x\neq 0\}$として各点$(x,y,z)$と原点$\mathbf 0$を結ぶ直線と$x=1$の面の交点を求める，つまり，連立一次方程式を解くと$(1,\frac yx,\frac zx)$がわかる(図\ref{perspective})．その結果を与える写像
\[
   P(x,y,z)=\biggl(\frac yx,\frac zx\biggr)\qquad (x,y,z)\in U
   \tag{*}\label{projection}
\]
$P:U\to \mathbf R^2$が(透視)\term[]{投影}(あるいは，\term[]{射影})である．
この$ \mathbf R^2$, あるいは，$ \mathbf R^2$を $\mathbf R^3$の$x=1$の面とみたものを\term[]{投影面}と言う(図\ref{model})．

\begin{figure}[htb]
   \centering
   \begin{minipage}{0.4\columnwidth}
      \centering
      \includegraphics{~/Documents/GitHub/bEuclideanGeo/ch3/images/figure.5}
      \caption{透視投影}\label{perspective}
   \end{minipage} 
   \begin{minipage}{0.4\columnwidth}
      \centering
      \includegraphics{~/Documents/GitHub/bEuclideanGeo/ch3/images/figure.17}
      \caption{直線の射影}\label{perspective2}
   \end{minipage} 
\end{figure}

定義域が$U\neq \mathbf R^3$なので線形写像にならないことは明らかだが
\[
   P(ax,ay,az)=P(x,y,z) \quad a\neq0
\]
という性質がある．
\eqref{projection}は$y$, $z$については線形写像で，線形写像を少しひねっただけの簡単なものである．

$ \mathbf R^3$で点$(x,y,z)$から向き$(u,v,w)$に向かう直線は
\[
   (x,y,z)+t(u,v,w)
\]
と表せるので
\[
   P(x+ut,y+vt ,z+wt)=\biggl(\frac{y+vt}{x+ut},\frac{z+wt}{x+ut}\biggr)\tag{!}
\]
これは$t=0$で$(\frac yx,\frac zx)$を通り，$u\neq 0$のとき$t\to \infty$で$(\frac{v}{u},\frac{w}{u})$に収束する．この$(\frac{v}{u},\frac{w}{u})$を\term[]{消失点(vanishing point)}という．
差をみると
\[
   \biggl(\frac{y+vt}{x+ut},\frac{z+wt}{x+ut}\biggr)
   =
   \biggl(\frac{v}{u},\frac{w}{u}\biggr)
   +
   \frac{1}{x+ut}\biggl(y-x\frac{v}{u},z-x\frac{w}{u}\biggr)
\]
となり，最後の項は定数ベクトルを$\tfrac{1}{x+ut}$倍したものなので，投影面内で$(\frac yx,\frac zx)$を通る直線で消失点へ至るものであることが分かる(図\ref{perspective2}に初等幾何学的な証明をあげた)．
線形写像は直線を直線に写すが，投影は直線を直線の一部(半直線)に写すことがわかった．
また，平行な直線は同じ消失点に収束することも分かる．
さらに，$u=0$のとき
$(\frac{y+vt}{x},\frac{z+wt}{x})$となり，等速の直線運動になる．


遠近法の数学と直接関係ないが，地面や海面を表す平面が絵を描く（写真をとる）場合には大切になる．
地面で画家やカメラマンがたっている点を立点，その真上に視点をとり立点と視点を通る軸を$z$軸とする．地面は通常$z=\const$に取る．（これは必ずしも必要でないが）地面に平行などちらかへの視線向きを$x$軸，それに直交する（右手系になる）軸を$y$軸ととると上で考えた座標系と一致する．
地面$z=\const$の消失点全体が地平線をなす．
遠近法の基本的な手段は地平線上の
いくつかの消失点とそこに収束する平行直線群を見いだし，それを元に描画することである(図~\ref{1VanishingPoint}, \ref{2VanishingPoint})．もちろん，地面に垂直な線は$z$軸に平行に描かれる．

\begin{figure}[htb]
   \centering
   \includegraphics{~/Documents/GitHub/bEuclideanGeo/ch3/images/figure.7}
   \caption{1消失点の平行直線群}\label{1VanishingPoint}
\end{figure}

\begin{figure}[htb]
   \centering
   \includegraphics{~/Documents/GitHub/bEuclideanGeo/ch3/images/figure.16}
   \caption{平行直線群の作成}\label{parallelLines}
\end{figure}

\begin{figure}[htb]
   \centering
   \includegraphics[trim=13  0 305 0 , clip]{~/Documents/GitHub/bEuclideanGeo/ch3/images/figure.8}%
   \caption{2消失点の平行直線群}\label{2VanishingPoint}
\end{figure}

よく見られる一消失点と平行直線群($u\neq0$, $v=w=0$ に対応)と直交する交わらない平行直線群($v\neq0$, $u=w=0$ に対応)をどう作図するかを説明する：
平面を$z=-1$に固定する．この平面とカンバスの最底辺$\ell_1$を固定する．その平面上の$\mathbf R^3$の中での直交する平行直線は
\[
    (0,y_0,-1)+s(1,0,0),\quad (x_1,y_1,-1)+t(0,1,0),
\]
なので射影面では
\[
   \biggl(\frac{y_0}{s},\frac{-1}{s}\biggr),\quad 
   \biggl(\frac{y_1+t}{x_1},\frac{-1}{x_1}\biggr)
\]
と表せる．
まず，直線群は幅$a$, $b$に並んでいるとし，
\[
   y_0 = am,\quad
   x_1 = bn\quad  m\in \mathbf Z, n\in \mathbf N
\]
と取り
\[
   {\ell'}_m:\biggl(\frac{am}{s},\frac{-1}{s}\biggr)\quad
    (s\ge b),\quad
   \ell_n: \biggl(\frac{y_1+tv}{bn},\frac{-1}{bn}\biggr)
\]
と置く．${\ell'}_n$の作図はカンバスの最底辺を表す直線$\ell_1$の$0$を通る等分点と消失点を結べば良い．カンバスの外については後にコメントする．

${\ell'}_m$と$\ell_n$の交点は
\[
   s=bn
\]
であり，
$n$, $m$はインデックスに過ぎないので，${\ell'}_n$では
\[
   s=b n\quad 
   \Rightarrow \quad
   \frac{an}{s} = \frac{a}{b}
\]
なので射影面上で$y= \frac{a}{b} $と${\ell'}_n$, $y=- \frac{a}{b} $と${\ell'}_{-n}$の交点を求め，それを通る直線として${\ell}_n$を作図できる．$\ell_1$を$\ell_2$に替え，
その上で得られた等分点を拡張することで，画面内でどこまでも続けられる．

1消失点を使った構図(図~\ref{1VanishingPoint}), 
2消失点を使った構図(図~\ref{2VanishingPoint})，3消失点使った構図(図~\ref{3VanishingPoint})とその写真例を挙げる(図~\ref{ex1}, \ref{ex2}, \ref{ex3})

\begin{figure}[htb]
   \centering
   
\includegraphics[width=7.5cm]{~/Documents/GitHub/bEuclideanGeo/ch3/images/8367600279_ac56e7454d_b.jpg}
   \caption{写真例 \href{https://www.flickr.com/photos/31583394@N04/8367600279}{Vanishing Storm} by David Kingham\href{https://creativecommons.org/licenses/by-nd-nc/2.0/jp/?ref=openverse.}{\ccbyncnd}}\label{ex1}
\end{figure}
\begin{figure}[h]
   \centering
   \includegraphics[width=8cm]{~/Documents/GitHub/bEuclideanGeo/ch3/images/9110846773_dad72ed427_k.jpg}
   \caption{写真例 \href{https://www.flickr.com/photos/30996501@N07/9110846773}{MEZQUITA CATEDRAL DE CORDOBA 6378} by Jose Javier Martin Espartosa \href{https://creativecommons.org/licenses/by-nd-nc/2.0/jp/?ref=openverse.}{\ccbyncsa}}\label{ex2}
\end{figure}

\begin{figure}[htb]
   \begin{minipage}{0.5\columnwidth}
      \centering
      \includegraphics{~/Documents/GitHub/bEuclideanGeo/ch3/images/figure.9}
      \caption{3消失点を使った直方体の作図}\label{3VanishingPoint}
   \end{minipage} 
   \begin{minipage}{0.4\columnwidth}
      \centering
      \includegraphics[width=4cm]{~/Documents/GitHub/bEuclideanGeo/ch3/images/145347226_76283ffb36_c.jpg}
      \caption{写真例 \href{https://www.flickr.com/photos/46965656@N00/145347226}{Barcelona - Gaudi building} by Erick \href{https://creativecommons.org/licenses/by-nd-nc/2.0/jp/?ref=openverse.}{\ccbyncnd}}\label{ex3}
   \end{minipage} 
\end{figure}

\begin{figure}[htb]
   \centering
   \begin{minipage}{0.4\columnwidth}
      \centering
      \includegraphics[width=5.8cm]{~/Documents/GitHub/bEuclideanGeo/ch3/images/figure.10}
      \caption{透視投影と平行投影}\label{paraPerspective1}
   \end{minipage} 
   \begin{minipage}{0.4\columnwidth}
      \centering
      \includegraphics[width=4.7cm]{~/Documents/GitHub/bEuclideanGeo/ch3/images/figure.11}
      \caption{平行投影}\label{paraPerspective2}
   \end{minipage} 
\end{figure}


投影面を$x=0$に止めて，あらためて視点を$(-a,0,0)$にとり，$a\to \infty$として無限遠まで動かしたとき得られる投影を\term[]{平行投影}という．このとき透視投影の極限として
\[
   \biggl(\frac{ a}{x+a}y,\frac{ a}{x+a}z\biggr)
   \to (y,z)\qquad (a\to \infty)
\]
となる(図\ref{paraPerspective1}, \ref{paraPerspective2})ので，これは線形写像である．ただし，平行投影は挿絵(この資料の図もそうである)図学等ではよく使われるが，3DCGでは現実感を出すため透視投影が主に使われる．

\subsection{反射，拡散と影}
ここから3DCGに関連する話題に移る．
まず大事なのは光とその反射面の取り扱いである．ここでは屋外なので光源として太陽を考えている．
スタジオ撮影では多数のライトを用い，光の強度は距離の二乗に反比例することも考慮する必要がある．逆に太陽の場合は無限の彼方にあると考え，距離に依らない強度を持つとする．
光源から出た光は直進し
被写体の面で反射して，また直進してキャンバスと視点に到達するので，
光の平面での\term[]{鏡面反射}を計算する必要がある．
これは与えられた平面での光源の鏡映を求め，その点と視点を結ぶ直線と平面の交点を求めることで求まる(図12)．与えられた平面が$z=0$とすると，点$(x,y,z)$のこの平面での鏡映は行列を用いて
\[
   \begin{pmatrix}
      1&0&0\\
      0&1&0\\
      0&0&-1\\
   \end{pmatrix}
   \begin{pmatrix}
      x\\
      y\\
      z\\
   \end{pmatrix}
   =
   \begin{pmatrix}
      x\\
      y\\
      -z\\
   \end{pmatrix}
\]


\begin{figure}[htb]
   \centering
   \includegraphics{~/Documents/GitHub/bEuclideanGeo/ch3/images/figure.3}
   \caption{鏡映と鏡面反射, 拡散}\label{ref}
\end{figure}

\noindent
と表せる．
光源の位置を$(x,y,z)$とし，視点の位置を$(x_0,y_0,z_0)$
とすると途中の点は
\[
   \begin{split}
      &t(x, y, -z)+(1-t)(x_0, y_0, z_0)\\
      =&
      (tx+(1-t)x_0, ty+(1-t)y_0,-tz+(1-t)z_0)
   \end{split}
\]
と表される．反射する点では$-tz+(1-t)z_0=0$なので，連立一次方程式を解くと
\[
   \biggl(\frac{zx_0+z_0x}{z+z_0},\frac{zy_0+z_0y}{z+z_0},0 \biggr)
\]
のように反射点が求まる(図 \ref{ref})．
さらには，物理特性である面の色・材質による拡散も計算する必要がある．

\begin{figure}[htb]
   \centering
   \includegraphics{~/Documents/GitHub/bEuclideanGeo/ch3/images/figure.20}
   \caption{隠面消去・シェーディング・影付け}\label{shadow}
\end{figure}

また，一つの物体の見えない部分や，二つの物体の重なりを計算し遠くの物体の隠れた部分を消す\term[]{隠面消去}や，
光源によって照らされた部分の明るさや色を求める
\term[]{シェーディング}, 
影をシュミレートする\term[]{影付け}も必要になる(図 \ref{shadow})．

なお，遠近法には\term[]{空気遠近法}という，遠くのものを
薄く描いて遠近感を出す方法も知られている．今回は詳しい説明は与えないが，図 \ref{1VanishingPoint}, \ref{2VanishingPoint} はその原理に
基づいて作図されている．

\subsection{アファイン変換}
$3$次元の\term[]{アファイン変換}とは
\[
   \begin{split}
      \begin{pmatrix}
         c_1\\
         c_2\\
         c_3\\
      \end{pmatrix}
      +
      \begin{pmatrix}
         a_{11}&a_{12}&a_{13}\\
         a_{21}&a_{22}&a_{23}\\
         a_{31}&a_{32}&a_{33}\\
      \end{pmatrix}
      \begin{pmatrix}
         x_1\\
         x_2\\
         x_3\\
      \end{pmatrix}
      &=
      \begin{pmatrix}
         c_1 + a_{11}x_1 + a_{12}x_2+ a_{13}x_3\\
         c_2 + a_{21}x_1 + a_{22}x_2+ a_{23}x_3\\
         c_3 + a_{31}x_1 + a_{32}x_2+ a_{33}x_3\\
      \end{pmatrix}\\
      &=
      A\mathbf x + \mathbf c
   \end{split}
\]
となる写像で，平行移動と行列の積の合成である．
通常$A=(a_{ij})
%\begin{pmatrix}
%        a_{11}&a_{12}&a_{13}\\
%        a_{21}&a_{22}&a_{23}\\
%        a_{31}&a_{32}&a_{33}\\
% \end{pmatrix}
$
は正則行列と仮定する．特に，合同変換はアファイン変換である．
以後アファイン変換はCGの作成で何度も現れる．

これを直線
\begin{equation*}
   \mathbf x = \mathbf a + t \mathbf u
 \end{equation*}
に適用すると
\[
   A(\mathbf a + t \mathbf u)+\mathbf c
   =
   \mathbf c+A\mathbf a+tA\mathbf u\quad A\mathbf u\neq \mathbf 0
\]
つまり，アファイン変換は直線全体を直線全体に写し，平行な直線を平行な直線に写す．
また，直線上の三点の比を同じ比の線分上の点に写すことも分かる．

行列の積で
\[
   \begin{pmatrix}
      1&0&0&0\\
      c_1&a_{11}&a_{12}&a_{13}\\
      c_2&a_{21}&a_{22}&a_{23}\\
      c_3&a_{31}&a_{32}&a_{33}\\
   \end{pmatrix}
   \begin{pmatrix}
      1\\
      x_1\\
      x_2\\
      x_3\\
   \end{pmatrix}
   =
   \begin{pmatrix}
      1\\
      c_1 + a_{11}x_1 + a_{12}x_2+ a_{13}x_3\\
      c_2 + a_{21}x_1 + a_{22}x_2+ a_{23}x_3\\
      c_3 + a_{31}x_1 + a_{32}x_2+ a_{33}x_3\\
   \end{pmatrix}
\]
となるので，その第2第3第4成分がアファイン変換に対応している．
これはアファイン変換の行列の積を使った表現であり，
その意味でアファイン変換は線形変換に近いと言えるが，第1成分が$1$に固定されているので，これは線形写像ではない．このような座標系を\term[]{非同次座標系}という．これは行列の計算なのでコンピュータでの組織的な計算に適している．

\subsection{モデリング，ビューイングパイプライン}
\begin{figure}[htb]
   \centering
   \includegraphics{~/Documents/GitHub/bEuclideanGeo/ch3/images/figure.4}
   \caption{ビューイングパイプライン}\label{viewingPipeline}
\end{figure}


コンピュータの中で被写体の曲面の三次元データを作ることを\term[]{モデリング}という．これは\term[]{モデリング座標系}での，ピラミッドのようなポリゴンの頂点とそれらの張る辺，三角形の作る面の幾何学的情報と色や質感等の物理情報の組からなる．
シュミレートしたい被写体とカメラの位置・向きの指定された三次元(データの)世界を\term[]{ワールド座標系}という．幾つかのモデリングのデータを適当なアファイン変換%(下の補足Aで示すように行列の積で表現される)
でうつすことで，ワールド座標系にシュミレートしたい被写体のデータを移している(図\ref{viewingPipeline})．
ここでは，2つのピラミッドを一つのデータから構成している．
ワールド座標系の中での視点の座標と視線の方向を定めることで，
その視点から見た座標系が定まるが，
それはワールド座標系をアファイン変換したものであり\term[]{カメラ座標系}という．
視点から投影面に（透視投影した）二次元座標系を\term[]{投影座標系}という．
ただし
実際は投影座標系の中の最終的な画面出力となる
\term[]{ウィンドウ}と呼ばれる長方形部分のみを計算すればよい．

ワールド座標系とカメラ座標系のデータ及び，光源，反射，陰影等からウィンドウへの投影を計算し，ウィンドウ内のデジタルデータを作図するのを\term[]{レンダリング}という．
これらを少しずつ動かしながら繰り返すことでアニメーションが構成できる．
視点を動かし，カメラワークをシュミレートすることも可能である．
このような幾つもの操作の流れを\term[]{ビューイングパイプライン}という．
これらのことから，線型代数が3DCGの中で占める役割の大きさが分かると思う．

なお，\href{https://snow00two.github.io/??}{Star**へのオマージュ}にJavaScriptとp5.jsを使ったデモと\href{https://snow00two.github.io/??}{解説}を公開しているので，プログラミングの参考にしてほしい．

\emptyline
CGに興味を持った人には\cite{CG}, 遠近法の歴史については\cite{D},
具体的な遠近法の絵の描き方については\cite{3DT}, \cite{RS}, \cite{RSII} などをお奨めする．

%   \begin{thebibliography}{9}
%      \bibitem{CG}
%         画像情報教育振興協会, コンピューターグラフィックス[改訂新版], 2015.
%      
%      \bibitem{D}
%         下村 耕史訳編, アルブレヒト・デューラー「測定法教則」注解, 中央公論美術出版, 2008.
%      
%      \bibitem{3DT}
%         3DTotal.com, 
%         デジタルアーティストが知っておくべきアートの原則 改訂版 -色、光、構図、解剖学、遠近法、奥行き 改訂版, 
%         2021.
%      
%      \bibitem{RS}
%         スコット・ロバートソン, 
%         スコット・ロバートソンのHow to Draw -オブジェクトに構造を与え、実現可能なモデルとして描く方法, 
%         2014.
%         
%      \bibitem{RSII}
%         スコット・ロバートソン, 
%         スコット・ロバートソンのHow to Render：アイデアを明確に伝える 光と影、反射の描き方, 2015.
%   \end{thebibliography}

%\newpage

\section{射影幾何学と射影空間}
\subsection{投影を使った定理}
デザルグは1636年に出版した本で遠近法の幾何学的研究を発表し，射影幾何学の礎を作った．3つの直線が\term[]{共点}とは3つの直線が同じ点を通ることである．3つの点が\term[]{共線}とは3つの点が同じ直線を通ることである．

\begin{Thm}[デザルグの定理]
平面内の
二つの三角形$\triangle ABC$, $\triangle A'B'C'$が与えられ，
三つの直線$AA'$, $BB'$, $CC'$が共点とする．
直線$AB$と$A'B'$の交点を$P$,
直線$BC$と$B'C'$の交点を$Q$,
直線$CA$と$C'A'$の交点を$R$ とすると，
$P$, $Q$, $R$は共線になる．
\begin{figure}[hb]
	\centering
	\includegraphics{~/Documents/GitHub/bEuclideanGeo/ch3/images/figureProje2.1}
	\caption{デザルグの定理}\label{}
\end{figure}
\end{Thm}
\begin{proof}
以下に投影を使った証明を述べる．
空間の中に三角形$\triangle ABC$, $\triangle A'B'C'$が与えられ，
三つの直線$AA'$, $BB'$, $CC'$が同じ点$O$を通るとする．
$\triangle ABC$を含む平面$H$を取ると
直線$AB$, $BC$, $CA$はこの平面$H$に含まれる．
同様に$\triangle A'B'C'$を含む平面$H'$を取ると
直線$A'B'$, $B'C'$, $C'A'$はこの平面$H'$に含まれる．
平面$H$と$H'$ 交わりの直線$\ell$は$P$, $Q$, $R$を含んでいる．

この状況を平面$H$と$H'$と異なる平面をとり，
そこに投影するとデザルグの定理が得られる．
\end{proof}

\begin{figure}[htb]
   \centering
   \begin{minipage}{0.4\columnwidth}
      \centering
      \includegraphics{~/Documents/GitHub/bEuclideanGeo/ch3/images/figureProje2.8}
      \caption{ジュルゴンヌ点}\label{gerg}
      \end{minipage} 
   \begin{minipage}{0.4\columnwidth}
      \centering
      \includegraphics{~/Documents/GitHub/bEuclideanGeo/ch3/images/figureProje2.9}
      \caption{内接円}\label{cont}
   \end{minipage} 
\end{figure}

\begin{Thm}[ジュルゴンヌ]
   三角形$\triangle ABC$に楕円が内接し，それらの内接点を
   図~\ref{gerg}のように$P$, $Q$, $R$と置くと直線$AP$, $BQ$, $CR$は共点である．
\end{Thm}
\begin{proof}
   図~\ref{gerg}のような平面上の三角形と楕円が与えられたとする．
   視点を適当に動かして，視点が楕円が円に見える位置に置かれたとする．
   視点から出る視線がその円の中心を通るとし，
   視線に直交する平面を与えられた図形より遠くに取る．
   その平面上に楕円を射影した円は三角形$\triangle ABC$を射影した三角形$\triangle A'B'C'$に内接している．
   それらの内接点を
   図~\ref{cont}のように$P'$, $Q'$, $R'$と置くと
   $A'P'$, $B'Q'$, $C'R'$は内接している円の中心を通るので共点である．
   この円の中心と視点を結ぶ直線と元の平面の交わりは
   直線$AP$, $BQ$, $CR$の全てを通るので共点であることが分かる．
   %チェバの定理より
   %\[
   %	\frac{A'R'}{R'B'}\frac{B'P'}{P'C'}\frac{C'Q'}{Q'A'} =1
   %\]
   %と$A'P'$, $B'Q'$, $C'R'$は共点であることは同値であるが，
   %この関係は円が三角形に内接することが
   %\[
   %	A'R'=A'Q',\quad B'R'=B'P', \quad C'P'=C'Q'
   %\]
   %を意味することから分かる．
\end{proof}

\subsection{遠近法の無限遠点}
直線の投影を振り返る：
\begin{figure}[htb]
   \centering
   \begin{minipage}{0.4\columnwidth}
      \centering
      \includegraphics{~/Documents/GitHub/bEuclideanGeo/ch3/images/figureProje2.15}
      \caption{透視投影}\label{perspective}
   \end{minipage} 
   \begin{minipage}{0.4\columnwidth}
      \centering
      \includegraphics{~/Documents/GitHub/bEuclideanGeo/ch3/images/figureProje2.16}
      \caption{}\label{perspective2}
   \end{minipage} 
\end{figure}
原点$\mathbf 0$に視点があるとする．
$(x,y,z) \in \mathbf R^3$  ($x\neq 0$ )とし，
ここでは，$x$軸の正の方向に視線を向け，
$x=1$ にキャンバスがある場合を考える．
$U=\{(x,y,z)\in \mathbf R^3\mid x\neq 0\}$として各点$(x,y,z)$と原点$\mathbf 0$を結ぶ直線と$x=1$の面の交点を求めると$(1,\frac yx,\frac zx)$がわかる．その結果を与える写像
\[
	P(x,y,z)=\biggl(\frac yx,\frac zx\biggr)\qquad (x,y,z)\in U
	\tag{*}\label{projection}
\]
$P:U\to \mathbf R^2$が\term[]{投影}(あるいは，\term[]{射影})である．
\[
	P(ax,ay,az)=P(x,y,z) \qquad a\neq0
\]
という性質がある．$ \mathbf R^3$で$z=-1$上の点から向き$(u,v,0)$  に向かう直線$\ell$は
\[
	(x,y,-1)+t(u,v,0)
\]
と表せる．まず $u > 0$と仮定する．
\[
	P(x+tu,y+tv ,-1)=\biggl(\frac{y+tv}{x+tu},\frac{-1}{x+tu}\biggr)\tag{!}
\]
これは$t=0$で$(\frac yx,\frac {-1}x)$を通り，$t\to \infty$で\term[]{消失点}$(\frac{v}{u},0)$に収束する．消失点は$t\to \infty$で表れるので\term[]{無限遠点}ということもある．

\begin{figure}[thb]
   \centering
   \includegraphics{~/Documents/GitHub/bEuclideanGeo/ch3/images/figureProje2.3}
   \caption{後ろからの透視投影}\label{backProj}
\end{figure}

$t\to -\infty$の場合の無限遠点を検討しよう．
まず$t> \frac{-x}{u}$で消失点からの投影面上の半直線に沿って進み，
$t= \frac{-x}{u}$で投影面$x=1$上に対応する点がなくなるが，
$z=-1$上では$x=0$という平面と直線$\ell$の交点であることに注意する．
空間の直線$\ell$自身は平面$x=0$を越えると，その上の点から視点を通る直線が$x=1$に至る点が表れる．その点は直線$\ell$の
$t<\frac{-x}{u}$の部分に対応する$x=1$上の逆側の無限遠から消失点に伸びる半直線に沿って進み，
$t\to -\infty$で消失点に（同じ角度で）交わることが図 \ref{backProj} からも数式からも分かる．
\emptyline

\begin{figure}[htb]
   \centering
   \includegraphics{~/Documents/GitHub/bEuclideanGeo/ch3/images/figureProje2.17}
   \caption{カメラ・オブスキュラと遠近法の融合}\label{pinCamera}
\end{figure}

\emptyline
\noindent
遠近法で視点の後ろの点を射影面に投影するのは不思議に思えるかもしれないが，図 \ref{pinCamera}のようなカメラ・オブスキュラ（ピンホールカメラ）と遠近法を融合したような道具を考えれば実現できるであろう．
この立場からは空間の直線は射影面の直線（全体）に写っていて，
平行な直線は無限遠点で二つの直線として交叉する（図 \ref{paraAtVanishingPoint}）ことも分かる．
これは原点と直線$\ell$から定まる平面と$x=1$の交わりの上にあることからも確認できる．
つまり，$z=-1$に対応する点がない（無限遠）点は
$x=1$では$z=0$という直線上の点になる．
\begin{figure}[htb]
   \centering
   \begin{minipage}{0.4\columnwidth}
      \centering
      \includegraphics{~/Documents/GitHub/bEuclideanGeo/ch3/images/figureProje2.12}
      \caption{平行線の配置}\label{paraAtVanishingPoint}
      \end{minipage} 
   \begin{minipage}{0.4\columnwidth}
      \centering
      \includegraphics{~/Documents/GitHub/bEuclideanGeo/ch3/images/figureProje2.13}
      \caption{放物線と水平線}\label{}
   \end{minipage} 
\end{figure}

このような考察により， $u < 0$ ( $x<0$ )の場合も同様に扱えることが分かる．
$u=0$の場合も$y=1$への投影を考察すれば良く，$(u,v)\neq\mathbf 0$の条件で十分であることが分かる．さらに，$(u,v,w)\neq\mathbf 0$とすれば任意の直線を考察できることも分かる．


まとめると，平面$z=-1$の無限遠点は平面$x=1$での通常の点とみることが出来き，平行な直線はこの点で交わっている．
一つの平面で無限遠点でも別の平行でない平面では通常点になることが分かる．
我々に必要なのは最低 $x=\const\neq 0$, $y=\const\neq 0$, $z=\const\neq 0$の三っつの平面であることも理解できると思う．
\emptyline

また，$z=-1$上の放物線
$
	x = y^2+1
$
の$x=1$への投影
\[
	s=\frac yx=\frac{y}{y^2+1},\quad
	t=\frac zx=\frac{-1}{y^2+1}
\]
を考えると
\[
	(s,t)
	=
	\begin{cases}
		(0,-1)& y=0\\
		(0,0)& y\to \pm \infty
	\end{cases}
\]
で，
\[
\begin{split}
	s^2 + \biggl(t+\frac12\biggr)^2
	&=\frac{y^2}{(y^2+1)^2} + \frac{(-2 +y^2+1)^2}{4(y^2+1)^2}\\
	&= \frac{1}{4(y^2+1)^2}(4y^2 + (y^2-1)^2)
	=\frac{1}{4}
\end{split}
\]
となり，水平線$t=0$に接する円に見えることが分かる．
3つの円錐曲線が射影幾何学で同一視されることの一端がここに表れている．


\begin{Quiz}$16$
$z=-1$上の放物線
$
	x = y^2
$
の$x=1$への投影を求めよ．このときそれぞれの無限遠点の対応を調べよ．
\end{Quiz}

\begin{Quiz}$17$
$z=-1$上の双曲線
$
	x^2- y^2=1
$
の$x=1$への投影を求めよ．このとき双曲線二つの成分の投影と漸近線の消失点の関係を調べよ．
\end{Quiz}


\subsection{同次座標，実射影平面}
ポンスレーは1812年ナポレオンのロシアへの遠征に従軍し捕虜となり，その捕らわれの間に射影幾何学を発見した．射影幾何学とは射影変換(\S \ref{projTransf}を参照)で不変な性質を研究する幾何学である．その研究の進展に伴い
メビウス，プリュッカーにより同次座標や射影空間が導入された．


$(x_0,x_1,x_2), (x'_0,x'_1,x'_2)\neq(0,0,0)$にたいして
\[
\begin{split}
	&(x_0,x_1,x_2)\sim (x'_0,x'_1,x'_2)\\
	\Leftrightarrow \quad
	&\text{ある }\lambda\neq 0\text{ が存在して }
	 (x'_0,x'_1,x'_2) = \lambda (x_0,x_1,x_2)
\end{split}
\]
で関係を入れると，同値関係になる．$(x_0,x_1,x_2)$の同値類を
$[x_0:x_1:x_2]$と表し，\term[]{同次座標}という．
$\mathbf P^2=\{[x_0:x_1:x_2]\mid (x_0,x_1,x_2)\neq \mathbf 0\}$
を\term[]{実射影平面}
という．
位相は$\mathbf R^3\setminus \{\mathbf 0\}$の商位相，つまり，
\[
	\mu:\mathbf R^3\setminus \{\mathbf 0\} 
	\ni(x_0,x_1,x_2)\mapsto  [x_0:x_1:x_2]\in\mathbf P^2
\]
として，$O\subset \mathbf P^2$が開集合であるとは
$\mu^{-1}(O)\subset\mathbf R^3\setminus \{\mathbf 0\} $が開集合であることとする．
この位相は
$[x_0:x_1:x_2]$, $[x_0':x_1':x_2']\in \mathbf P^2$から定まる
二つの直線$\{t(x_0,x_1,x_2)\in\mathbf R^3\mid t\in \mathbf R\}$と
$\{t(x_0',x_1',x_2')\in\mathbf R^3\mid t\in \mathbf R\}$のなす角（の小さいもの）
を距離とする距離空間の位相と見なせるので，
$\mathbf P^2$ はハウスドルフ空間である．
球面
$S^2(1)=\{(x_0,x_1,x_2)\in\mathbf R^3\mid x^2_0 +x^2_1 +x^2_2 =1\}$
はコンパクトなので$ S^2(1)/\sim$もコンパクトである．
写像
\[
	\pi:S^2(1)/\sim \ni 
	[x_0,x_1,x_2]\mapsto [x_0:x_1:x_2]\in \mathbf P^2
\]
は連続な全単射であり，$S^2(1)/\sim$はコンパクトで$\mathbf P^2$ はハウスドルフなので$\pi$は閉写像でもある．
\[
   \xymatrix{
      S^2(1) \ar[d]_{\mu'} \ar[r]& \mathbf R^3\setminus \{\mathbf 0\}  \ar[d]^{\mu} \\
     S^2(1)/ \sim  \ar[r]_{\pi} & \mathbf P^2 
   }
\]
このことから実射影平面は$S^2(1)/ \sim$ ( 図~\ref{projSpace} ) と同相であることが分かる(図~\ref{projSpace})．
また，半球の境界を図~\ref{projSpaceMebius}のように同一視したとき表れるメビウスの帯の境界に半球を貼り付けた空間でもある．
このことから$\mathbf P^2$は向き付け不可能であることもわかる．
\begin{figure}[htb]
\centering
   \begin{minipage}{0.4\columnwidth}
      \centering
      \includegraphics{~/Documents/GitHub/bEuclideanGeo/ch3/images/figureProje2.4}
      \caption{実射影平面$\mathbf P^2$}\label{projSpace}
   \end{minipage} 
   \begin{minipage}{0.4\columnwidth}
      \centering
      \includegraphics{~/Documents/GitHub/bEuclideanGeo/ch3/images/figureProje2.7}
      \caption{$\mathbf P^2$とメビウスの帯}\label{projSpaceMebius}\label{projSpace2}
   \end{minipage} 
\end{figure}

$U_0=\{[x_0:x_1:x_2]\mid x_0\neq 0\}$に対して
\[
	\varphi_0:U_0\ni [x_0:x_1:x_2]
	\mapsto \biggl(\frac{x_1}{x_0},\frac{x_2}{x_0}\biggr)
	\in \mathbf R^2
\]
と置くとwell-defined な同相写像で，透視投影$P$%\eqref{projection}
と同じ表示になっている．このような座標を（第1成分に$1$を付け加え，同次座標に対して）非同次座標ということがある．

同様に，$(x_1,x_2)$, $ (x'_1,x'_2)\neq (0,0)$に対して
\[
\begin{split}
	&(x_1,x_2)\sim (x'_1,x'_2)\\
	\Leftrightarrow \quad
	&\text{ある }\lambda\neq 0\text{ が存在して }
	 (x'_1,x'_2) = \lambda (x_1,x_2)
\end{split}
\]
は同値関係で，
$(x_1,x_2)$の同値類を
$[x_1:x_2]$と表し，商空間
$\mathbf P^1=\{[x_1:x_2]\mid (x_1,x_2)\neq \mathbf 0\}$
を\term[]{実射影直線}という．このとき
\[
\begin{split}
	\mathbf P^2
	&=\{[x_0:x_1:x_2]\in \mathbf P^2\mid x_0\neq 0\}
	\cup 
	\{[0:x_1:x_2]\in \mathbf P^2\}\\
	&\simeq U_0 \cup \mathbf P^1
	\simeq  \mathbf R^2 \cup \mathbf P^1
\end{split}
\]
$V_i=\{[x_1:x_2]\mid x_i\neq 0\}$と置くと
\[
	\psi_1:V_1\ni [x_1:x_2]
	\mapsto \frac{x_2}{x_1}
	\in \mathbf R,\quad
	\psi_2:V_2\ni [x_1:x_2]
	\mapsto \frac{x_1}{x_2}
	\in \mathbf R
\]
は同相で，上と同様に(図~\ref{proj1})
\[
\begin{split}
	\mathbf P^1
	&=\{[x_1:x_2]\in \mathbf P^1\mid x_1\neq 0\}
	\cup 
	\{[0:x_2]\in \mathbf P^1\}\\
	&= V_1 \cup \{[0:1]\}
	\simeq\mathbf R\cup \{[0:1]\}
	\simeq S^1
\end{split}
\]
\begin{figure}[htb]
\centering
   \begin{minipage}{0.4\columnwidth}
      \centering
      \includegraphics{~/Documents/GitHub/bEuclideanGeo/ch3/images/figureProje2.11}
      \caption{$\mathbf P^1$と$S^1$}\label{proj1}
      \end{minipage} 
   \begin{minipage}{0.4\columnwidth}
      \centering
      \includegraphics{~/Documents/GitHub/bEuclideanGeo/ch3/images/figureProje2.5}
      \caption{$\mathbf P^2$と$\mathbf P^1_0$, $\mathbf P^1_1$, $\mathbf P^1_0$}\label{proj2}
   \end{minipage} 
\end{figure}


ここまでは$x_0$方向を中心に考えてきたが，
$U_1=\{[x_0:x_1:x_2]\mid x_1\neq 0\}$, $U_2=\{[x_0:x_1:x_2]\mid x_2\neq 0\}$ に対しても同様の写像
\[
\begin{split}
	&\varphi_1:U_1\ni [x_0:x_1:x_2]
	\mapsto 
	\biggl(\frac{x_0}{x_1},\frac{x_2}{x_1}\biggr)\in \mathbf R^2,\\
	&\varphi_2:U_2\ni [x_0:x_1:x_2]
	\mapsto 
	\biggl(\frac{x_0}{x_2},\frac{x_1}{x_2}\biggr)\in \mathbf R^2
\end{split}
\]
が定まり，$(x_0,x_1,x_2)\neq \mathbf 0$は $x_i$の少なくとも一つは$0$でないことを意味するので$\{U_0, U_1, U_2\}$は射影平面を覆う．
また，例えば$U_0\cap U_2$では$x_0$, $x_2\neq 0$なので，
 $u=\frac{x_1}{x_0}$, $v=\frac{x_2}{x_0}\neq 0$ と置くと
\[
	\varphi_2\circ\varphi_0^{-1}(u,v)
	=
	\biggl(\frac{1}{v},\frac{u}{v}\biggr)
\]
のように変換され，微分可能であることがわかる．
このようなもの（でハウスドルフなもの）が$C^\infty$多様体であった．
$\mathbf P^1_i=\mathbf P^2\setminus U_i$ $(i=0,1,2)$は射影直線とみなせるので，これらを直線で表し，それぞれ2つが1点で交叉しているので図~\ref{proj2}のように表すことがある．
\emptyline

・二つの直線は射影平面上一点で必ず交わる．

∵）平行でない直線は平面上で必ず交わるが，平行な直線は，上で見たように無限遠点で必ず交わるから．

\subsection{射影変換}\label{projTransf}
正則行列
$A=\begin{pmatrix}
    	    	a&b&c\\
    		d&e&f\\
    		g&h&i\\
 	\end{pmatrix}$を使って
\[
	\begin{pmatrix}
    	    	a&b&c\\
    		d&e&f\\
    		g&h&i\\
 	\end{pmatrix}
	\begin{pmatrix}
    	    	 x_0\\
    		 x_1\\
    		 x_2\\
 	\end{pmatrix}
	=
	\begin{pmatrix}
    	    	 ax_0+bx_1+cx_2\\
    		 dx_0+ex_1+fx_2\\
    		 gx_0+hx_1+ix_2\\
 	\end{pmatrix}
\]
を考えると，$A(\lambda\mathbf x)=\lambda A\mathbf x$ で
$A\mathbf x=\mathbf 0$となるのは$\mathbf x=\mathbf 0$なので
$[x_0:x_1:x_2]\in \mathbf P^2$の代表元の取り方に依らず，写像
\[
\begin{split}
	\rho_A:&\mathbf P^2\ni [x_0:x_1:x_2]\mapsto\\
	&[ax_0+bx_1+cx_2:dx_0+ex_1+fx_2:gx_0+hx_1+ix_2]\in \mathbf P^2
\end{split}
\]
が定まる．この写像$\rho_A$，あるいは元の写像をこのように解釈し\term[]{射影変換}という．
\emptyline

・射影変換の合成は射影変換で，射影変換全体は群を成す．

∵）$\rho_B\circ\rho_A=\rho_{BA}$,  $\rho_E =\id$, 
$(\rho_A)^{-1}=\rho_{A^{-1}}$を見れば良いがこれは簡単．

\emptyline

・射影平面の任意の二点は射影変換で移り合う．

∵）$[a:b:c]\in \mathbf P^2$を$[1:0:0]$に写す射影変換を見つければいいが，$\mathbf u={}^t(a,b,c)\neq \mathbf 0$を含む$\mathbf R^3$の基底$\{\mathbf u,\mathbf v, \mathbf w\}$を作り
$A=	
\begin{pmatrix}
	\mathbf u&\mathbf v&\mathbf w
\end{pmatrix}
$と置くと，これは正則で$\mathbf u=A\mathbf e_1$なので
$\rho_{A^{-1}}$が求めるもの．

\emptyline

・射影平面の任意の二直線は射影変換で移り合う．

∵）射影平面の直線は$(a,b,c)\neq \mathbf 0$を使って$ax_0+bx_1+cx_2=0$と表せる．この平面が$x_0=0$(つまり，$\mathbf P^1_0$)に写せることを示す．
$\mathbf a =(a,b,c)$と置き，$\mathbf y = B \mathbf x$と変換すると
\[
 	0= \mathbf a\mathbf x =  \mathbf a B^{-1}\mathbf y
\]
なので
\[
	\mathbf a B^{-1} = (1,0,0)\quad \Leftrightarrow \quad
	{}^tB\mathbf e_1={}^t\mathbf a=\mathbf u
\]
となり，上と同じ方程式なので解ける．

\emptyline
また， 前回見たアファイン変換
\[
	\begin{pmatrix}
    	    1&0&0\\
    		a&c&d\\
    		b&e&f\\
 	\end{pmatrix}
	\begin{pmatrix}
    	    x_0\\
    		x_1\\
    		x_2\\
 	\end{pmatrix}
	=
	\begin{pmatrix}
    		x_0\\
    	    ax_0 + cx_1 +dx_2\\
    		bx_0+ex_1+ fx_2\\
 	\end{pmatrix}
\]
も射影変換の一つである．
アファイン変換は$\mathbf P^1_0$を$\mathbf P^1_0$に写している．

\subsection{二次曲線}
$
	B=\begin{pmatrix}
    	    	b_{00}&b_{01}&b_{02}\\
    		b_{01}&b_{11}&b_{12}\\
    		b_{02}&b_{12}&b_{22}\\
 	\end{pmatrix}$
は零行列でない実対称行列として
\[
\begin{split}
	&b_{00}x_0^2+b_{11}x_1^2+b_{22}x_2^2 
	+ 2b_{01}x_0x_1+2b_{12}x_1x_2 +2b_{02}x_0x_2
	=
	{}^t\mathbf x B\mathbf x\\
\end{split}
\]
を考える．
\[
	{}^t(\lambda\mathbf x) B(\lambda\mathbf x)
	=
	\lambda^2\, {}^t\mathbf x B\mathbf x
\]
なので${}^t\mathbf x B\mathbf x=0$は$\mathbf P^2$の部分空間を定め，これらを\term[]{二次曲線}(conic)という．
ただし，${}^t\mathbf x B\mathbf x=0$が一点の場合は(次元の関係から)除くこととする．

\begin{Thm}
$\mathbf P^2$の二次曲線は射影変換で次の何れかに写すことが出来る．
\[
	(\mathrm{i})\quad x^2 + y^2= z^2,\quad
	(\mathrm{ii})\quad  xy=0,\quad
	(\mathrm{iii})\quad  x^2=0
\]
\begin{figure}[htb]
	\centering
	\includegraphics{~/Documents/GitHub/bEuclideanGeo/ch3/images/figureProje2.14}
	\caption{$\mathbf P^2$の二次曲線}\label{}
\end{figure}
\end{Thm}
\begin{proof}
実対称行列$B$の固有値を$\lambda_{0}$, $\lambda_{1}$, $\lambda_{2}$とする．これらは（全てが$0$となることはない）実数である．$B$を対角化する直交行列を$S$とすると
\[
	{}^tS B S
	=
	\begin{pmatrix}
    	    	\lambda_{0}&0&0\\
    		0&\lambda_{1}&0\\
    		0&0&\lambda_{2}\\
 	\end{pmatrix}
\]
なので
\[
	c_{i}=
	\begin{cases}
		\frac{1}{\sqrt{|\lambda_{i}|}}& \lambda_{i}\neq 0\\
		1&\lambda_{i}= 0\\
	\end{cases},\quad
	C=(c_{i}\delta_{ij})
\]
と置くと$C$は正則となり，$SC$も正則で
\[
	{}^t(SC) B (SC)
	=
	{}^tC{}^tS B SC
	=
	\begin{pmatrix}
    	    	\varepsilon_{0}&0&0\\
    		0&\varepsilon_{1}&0\\
    		0&0&\varepsilon_{2}\\
 	\end{pmatrix}
\]
ただし，$\varepsilon_{i}$は$1$, $0$, $-1$のいずれか．

順番を入れ換えたり，$-1$倍することで任意の二次曲線は
\begin{align}
	&x_{0}^2 = 0\\
	&x_{0}^2 +x_{1}^2=0\\
	&x_{0}^2 +x_{1}^2+x_{2}^2=0\\
	&x_{0}^2 -x_{1}^2=0\\
	&x_{0}^2 +x_{1}^2-x_{2}^2=0
\end{align}
に何れかに写すことが出来る．(3)は空集合で，(2) は一点になり，この場合は2次曲線でないので除けるので，
(1), (4), (5) の何れか．
(1)は$(\mathrm{iii})$, (4)は$(\mathrm{ii})$, (5)は$(\mathrm{i})$である．
\end{proof}
$(\mathrm{i})$は楕円，放物線，双曲線の円錐曲線が直円錐から得られることに対応している．$(\mathrm{ii})$は二つの直線，$(\mathrm{iii})$は重直線の場合に対応している．

%\begin{Thm}[パスカルの定理]
%
%\end{Thm}
%\begin{Thm}[パップスの定理]
%
%\end{Thm}
%\section{一般化}
%$K$を体として
%上と同様の構成を行い$K\mathbf P^1$, $K\mathbf P^2$を定義することが出来る．
%特に$K$を複素数体$\mathbf C$にすると複素射影直線$\mathbf C\mathbf P^1$が次のように構成される：
%$(z_1,z_2)$, $ (z'_1,z'_2)\neq (0,0)\in \mathbf C^2$に対して
%\[
%\begin{split}
%	&(z_1,z_2)\sim (z'_1,z'_2)\\
%	\Leftrightarrow \quad
%	&\text{ある }\lambda\neq 0\text{ が存在して }
%	 (z'_1,z'_2) = \lambda (z_1,z_2)
%\end{split}
%\]
%は同値関係で，
%$(z_1,z_2)$の同値類を
%$[z_1:z_2]$と表し，商空間
%$\mathbf C\mathbf P^1=\{[z_1:z_2]\mid (z_1,z_2)\neq \mathbf 0\}$
%を\term[]{複素射影直線}という．関数論では\term[]{リーマン球面}といった．
%$\infty=[0:1]$と表し，
%$\hat{\mathbf C}= \mathbf C\cup \infty$と記すことも多い．
%
%$V_i=\{[z_1:z_2]\mid z_i\neq 0\}$ ($i=1,2$)と置くと
%\[
%	\psi_1:V_1\ni [z_1:z_2]
%	\mapsto \frac{z_2}{z_1}=z
%	\in \mathbf C,\quad
%	\psi_2:V_2\ni [z_1:z_2]
%	\mapsto \frac{z_1}{z_2}=w
%	\in \mathbf C
%\]
%は同相で，
%\[
%\begin{split}
%	\psi_1(V_1\cap V_2)\ni z \mapsto 
%	\frac1z=w\in \psi_2(V_1\cap V_2)
%\end{split}
%\]
%は正則写像である．
%これが位相的には球面であることが以下のように分かる：
%$z=x+iy\in \mathbf C$と$\mathbf R^3$の$xy$-平面内の点$(x,y,0)$ と同一視する．$S^2(1)$の元$(0,0,1)$を$\infty$と置き，
%$\infty$と$(x,y,0)$を結ぶ直線と$S^2(1)$の交点を求めると
%連続写像$\varphi:\mathbf C\to S^2(1)\setminus\{\infty\}$が
%\[
%	\varphi(x+iy)=
%	\biggl(\frac{2x}{x^2 + y^2+1},\frac{2y}{x^2 + y^2+1},
%	\frac{x^2 + y^2-1}{x^2 + y^2+1}
%	\biggr)
%\]
%により, $\varphi^{-1}: S^2(1)\setminus\{\infty\}\to \mathbf C$が
%\[
%	\varphi^{-1}(x_1,x_2,x_3)
%	=\frac{x_1+i x_2}{1-x_3}
%\]
%により定まり，
%\[
%	\mathbf C\mathbf P^1\simeq\mathbf C\cup \{\infty\}
%	\simeq S^2
%\]
%\begin{figure}[htb]
%	\centering
%	\includegraphics{../texProjective3/~/Documents/GitHub/bEuclideanGeo/ch3/images/~/Documents/GitHub/bEuclideanGeo/ch3/images/figureSphere.7}
%	\caption{リーマン球面}
%\end{figure}
%同様に複素射影平面$\mathbf C\mathbf P^2$が得られ，
%\[
%	\mathbf C\mathbf P^2
%	\simeq  \mathbf C^2 \cup \mathbf C\mathbf P^1
%\]
%更に$n\in \mathbf Z_{\ge 0}$に対して$n$次射影空間$\mathbf C\mathbf P^n$が定まり，
%その(幾つかの同次多項式の零点全体の)既約射影代数多様体が代数幾何の展開される空間である．

\begin{figure}[htb]
\centering
\begin{minipage}{0.4\columnwidth}
	\centering
	\includegraphics{~/Documents/GitHub/bEuclideanGeo/ch3/images/figureProje2.10}
	\caption{パスカルの定理}\label{Pascal}
\end{minipage} 
\begin{minipage}{0.4\columnwidth}
	\centering
	\includegraphics{~/Documents/GitHub/bEuclideanGeo/ch3/images/figureProje2.2}
	\caption{パップスの定理}\label{Pappus}
\end{minipage} 
\end{figure}

説明は省くが，パスカルの定理(1640, 図 \ref{Pascal})，パップスの定理(図 \ref{Pappus})
も射影幾何学に属すものである．
代数幾何の教科書にもその歴史の片鱗を見つけることが出来る(\cite{F}~p62, \cite{U}~p281, なお，\cite{F}はよく知られた代数幾何の入門書で，
フリーのpdf が
\href{https://dept.math.lsa.umich.edu/~wfulton/CurveBook.pdf}{ここからダウンロード}可能)．射影幾何学に興味を持った人には\cite{N}をお奨めする．また，ここで述べた歴史に興味をもった人は\cite{Kl}をお奨めする．

\begin{Quiz}
パスカルの定理の命題を調べ，射影変換を使って分かり易い形に変更して示せ\thetag{\cite{N}}．
\end{Quiz}

\begin{Quiz}
パップスの定理の命題を調べ，射影変換を使って分かり易い形に変更して示せ\thetag{\cite{N}}．
\end{Quiz}

%
%\begin{thebibliography}{9}
%\bibitem{N}西山 享, 射影幾何学の考え方 (数学のかんどころ 19), 2013 
%
%\bibitem{K}F.~Kline, クライン:19世紀の数学, 彌永 昌吉監修, 1995
%
%\bibitem{F}
%W.~Fulton, 
%ALGEBRAIC CURVES---An Introduction to Algebraic Geometry (3rd edition), 2008
%
%\bibitem{U} 上野健爾, 代数幾何, 岩波書店, 2005
%
%\end{thebibliography}

\section{リーマン球面と一次分数変換}
\subsection{$\mathbf C\mathbf P^1$}
複素数体$\mathbf C$で
射影直線の構成を行う：
$(z_1,z_2)$, $ (z'_1,z'_2)\in \mathbf C^2\setminus\{(0,0)\}$に対して
\[
\begin{split}
	&(z_1,z_2)\sim (z'_1,z'_2)
	\Leftrightarrow \quad
	\text{ある }\lambda\neq 0\in  \mathbf C\text{ が存在して }
	 (z'_1,z'_2) = \lambda (z_1,z_2)
\end{split}
\]
と置くとこれは同値関係で，
$(z_1,z_2)$の同値類を
$[z_1:z_2]$と表す．商空間
$\mathbf C\mathbf P^1=\{[z_1:z_2]\mid (z_1,z_2)\neq \mathbf 0\in \mathbf C^2\}$
を\textcolor{red}{複素射影直線}という．

$p:\mathbf C^2\setminus (0,0)\to \mathbf C\mathbf P^1$を自然な射影とする．
$V_i=\{[z_1:z_2]\mid z_i\neq 0\}$ ($i=1,2$)と置くと
\[
   \hat V_i:=p^{-1}(V_i)
   =\{(z_1,z_2)\in \mathbf C^2\setminus (0,0)\mid z_i\neq 0\}
\]
は開集合なので$V_i$は$\mathbf C\mathbf P^1$の開集合であり，$[z_1:z_2]\in  \mathbf C\mathbf P^1$は$z_1\neq 0$あるいは$z_2\neq 0$を意味するので$\mathbf C\mathbf P^1=V_1\cup V_2$が分かる．
\[
	\psi_1:V_1\ni [z_1:z_2]
	\mapsto \frac{z_2}{z_1}=z
	\in \mathbf C
\]
は同相写像であることが以下のように分かる．
$\hat \psi_1:\hat V_1\to \mathbf C$を
$\hat \psi_1(z_1,z_2)=z_2/z_1$で定義すると明らかに連続で
左の可換図式
\[
   \xymatrix{
      \hat V_1\ar[d]_{p} \ar[dr]^{\hat \psi_1}& \\
      V_1  \ar[r]_{\psi_1} & \mathbf C
   }\qquad\qquad
   \xymatrix{
      \mathbf C\ar[d]_{\id} \ar[r]^{\iota}& \ar[d]^{p}  \hat V_1\\
      \mathbf C & \ar[l]^{\psi_1}   V_1
   }
\]
を見ると$\psi_1$の連続性が分かる．
$\psi_1$が全単射も明らか．
$\iota(z)=(1,z)$と置くと
$\iota:\mathbf C\to \hat V_1$は連続で
右の可換図式より$\psi_1^{-1}=p\circ \iota$が分かり，$\psi_1^{-1}$の連続性も言える．
よって$\psi_1$は同相写像で，同様に
\[
	\psi_2:V_2\ni [z_1:z_2]
	\mapsto \frac{z_1}{z_2}=w
	\in \mathbf C
\]
も同相写像である．
%$\{ (V_1,\psi_1), (V_2,\psi_2)\}$が微分構造であることを示す．
これらの座標変換を(複素変数で)求めると
\[
\begin{split}
	\psi_1(V_1\cap V_2)\ni z \mapsto 
	\frac1z=w\in \psi_2(V_1\cap V_2)
\end{split}
\]
で$\psi_1(V_1\cap V_2)=\mathbf C\setminus \{0\}$で正則写像である．
複素多様体のことは知らないとして実数変数で表してみる．
$z=x + iy$, $w=u + iv$と置き$ \mathbf C= \mathbf R^2$と同一視すると
$w=\frac{1}{z }=\frac{\bar z}{z\bar z}$なので
\[
\begin{split}
	&u = \frac{x}{x^2 + y^2},\quad v = \frac{-y}{x^2 + y^2},\\
	&x = \frac{u}{u^2 + v^2},\quad y = \frac{-v}{u^2 + v^2}
\end{split}
\]
となり$\varphi_i(V_1\cap V_2)=\mathbf R^2\setminus \{(0,0)\}$上
 $C^{\infty}$級であることが分かり，$C^{\infty}$級多様体であることが分かった．
 \[
\begin{split}
	\frac{\partial u}{\partial x}=\frac{y^2-x^2}{(x^2 + y^2)^2}=v^2-u^2,
	\quad &\quad
	\frac{\partial u}{\partial y}=\frac{-2xy}{(x^2 + y^2)^2}=2uv,\\
	\frac{\partial v}{\partial x}=\frac{2xy}{(x^2 + y^2)^2}=-2uv,
	\quad &\quad
	\frac{\partial v}{\partial y}=\frac{y^2-x^2}{(x^2 + y^2)^2}=v^2-u^2
\end{split}	
\]
なので，
そのヤコビアンは
\[
	\frac{\partial (u,v)}{\partial (x,y)}
	%=(v^2-u^2)^2+(2uv)^2
	=(v^2+u^2)^2>0,\quad
	\frac{\partial (x,y)}{\partial (u,v)}=(x^2+y^2)^2>0
\]
となる．このように任意の座標変換のヤコビアンが正になるような微分構造をもつ多様体を向き付けられた多様体という．普通の立体射影では
$v =\frac{y}{x^2 + y^2}$であったが，この場合ヤコビアンの符号は$-$となり向き付けられていない．
$\mathbf C\mathbf P^1$は複素多様体と喚ばれる多様体になっている．
一般に複素多様体は向き付けられた多様体である．

$\mathbf C\mathbf P^1$が位相的には球面であることが以下のようにして分かる：\newline 
$[0:1]\in \mathbf C\mathbf P^1$で
\[
	\mathbf C\mathbf P^1\setminus\{ [0:1]\}=V_1
\]
であり，上で示したように$\psi_1:V_1\to \mathbf C$は同相であった．\linebreak
$[0:1]\in \mathbf C\mathbf P^1$の近傍での位相は
$\psi_2:V_2\to \mathbf C$が同相になるように入っている．
つまり，$w=0$を中心とした任意半径の円板が$[0:1]$の基本近傍系を成している．特に$\mathbf C\mathbf P^1$はハウスドルフであることも分かる．
$z=x+iy\in \mathbf C$と$\mathbf R^3$の$xy$-平面内の点$(x,y,0)$ と同一視する．$S^2(1)$の元$(0,0,1)$を$\infty$と置き，以下$[0:1]\in \mathbf C\mathbf P^1$と同一視されることを示す．
$\infty$と$(x,y,0)$を結ぶ直線と$S^2(1)$の交点を求めると
連続写像$\psi:\mathbf C\to S^2(1)\setminus\{\infty\}$が
\begin{equation}\label{stereoProj}
	\psi(x,y)=
	\biggl(\frac{2x}{x^2 + y^2+1},\frac{2y}{x^2 + y^2+1},
	\frac{x^2 + y^2-1}{x^2 + y^2+1}
	\biggr)
\end{equation}
により, 連続な逆写像$\varphi^{-1}: S^2(1)\setminus\{\infty\}\to \mathbf C$が
\[
	\varphi^{-1}(x_1,x_2,x_3)
	=\frac{x_1+i x_2}{1-x_3}
\]
により定まり同相になる．
$\psi(x,y)$と$\infty$の距離の2乗は
\[
\begin{split}
	&\biggl(\frac{2x}{x^2 + y^2+1}\biggr)^2
	+\biggl(\frac{2y}{x^2 + y^2+1}\biggr)^2
	+\biggl(\frac{x^2 + y^2-1}{x^2 + y^2+1}-1\biggr)^2\\
	&=\frac{4}{x^2 + y^2+1}=\frac{4|w|^2}{1+|w|^2}
\end{split}
\]
となる．最後の変形で$|w|=\frac{1}{|z|}$を使った．上の注意より$\psi$を\newline
$\bar\psi([0:1])=\infty$と拡張すると
$\bar\psi:\mathbf C\mathbf P^1\to S^2(1)$
は同相になる．
つまり，$\mathbf C\mathbf P^1$は$\mathbf C$の一点コンパクト化である．
関数論では$\mathbf C\mathbf P^1$のことを\textcolor{red}{リーマン球面}といった．
$\infty=[0:1]$と表したり$\mathbf C\mathbf P^1$を$\hat{\mathbf C}= \mathbf C\cup \infty$と表したりすることも多い．

\begin{figure}[htb]
   \centering
   \begin{minipage}{0.45\columnwidth}
      \centering
      \includegraphics{~/Documents/GitHub/bEuclideanGeo/ch3/images/figureSphere.7}
      \caption{リーマン球面}
   \end{minipage} 
   \begin{minipage}{0.45\columnwidth}
      \centering
      \includegraphics{~/Documents/GitHub/bEuclideanGeo/ch3/images/figureSphere.8}
      \caption{円周角}\label{argP}
   \end{minipage} 
\end{figure}

$\mathbf C\mathbf P^1$と
同様に複素射影平面$\mathbf C\mathbf P^2$が得られ，$\mathbf P^2$と同様に
\[
	\mathbf C\mathbf P^2
	\simeq  \mathbf C^2 \cup \mathbf C\mathbf P^1
\]
が成り立つ．
更に$n\in \mathbf Z_{\ge 0}$に対して$n$次射影空間$\mathbf C\mathbf P^n$が定まり，
その(幾つかの同次多項式の零点全体のなす)既約射影代数多様体が代数幾何の展開される空間である．

\subsection{一次変換}
前回は実変数の射影変換を考えたが，ここでは
複素正則行列
$A=\begin{pmatrix}
    	    	a&b\\
    		c&d\\
 	\end{pmatrix}$の定める
$ \mathbf C\mathbf P^1$の射影変換
\[
	\rho_A:\mathbf C\mathbf P^1\ni [z_1:z_2]\mapsto
	[az_1+bz_2:c z_1+d z_2]\in \mathbf C\mathbf P^1
\]
を考える．ただし $ad-bc\neq 0, a,b,c,d \in \mathbf C$. 
$\mathbf C^2\setminus\{\mathbf 0\}\ni (z_1,z_2)\to 
(az_1+bz_2, c z_1+d z_2)\in \mathbf C^2\setminus\{\mathbf 0\}$ は連続なので
$\rho_A$も連続である．
$z=z_1/z_2\in \psi_2(V_2)$,
$(az_1+bz_2)/(c z_1+d z_2)\in \psi_2(V_2)$のとき
\begin{equation}\label{linearTrans}
	Sz =\rho_A(z)  = \frac{a z +b}{c z +d}
\end{equation}
となるので，\textcolor{red}{一次(分数)変換}（リーマン球面の双正則写像とみるとき\textcolor{red}{メービウス変換}ともいう）
であることが分かる．
一次変換と一次変換の合成も（前回の議論から直ぐに分かるが，具体的な計算で示すと）
\[
	 \frac{a'  \frac{a z +b}{c z +d} +b'}{c ' \frac{a z +b}{c z +d} +d'} 
	 =
	 \frac{a'  (a z +b)+b'(c z +d)}{c '  (a z +b) +d'(c z +d)} 
	 =
	 \frac{(a' a+ b' c)z +a'b+b' d}{(c ' a+ d'c)z +c 'b +d'd} 
\]
なので一次変換で，合成の法則は行列の積
\[
	\begin{pmatrix}
		a' & b'\\
		c' & d'
	\end{pmatrix}
	\begin{pmatrix}
		a & b\\
		c & d
	\end{pmatrix}
	=
	\begin{pmatrix}
		a'a +b'c & a'b+b'd\\
		c'a+d'c & c'b+d'd
	\end{pmatrix}
\]
と同じである．
このことから( $A^{-1}$ではすべての成分を同じ$ad-bc\neq 0$で割っているので）
\[
	S^{-1} z =\rho_{A^{-1}}(z)  = \frac{d z -b}{-c z +a}
\]
となり，$S$は$\mathbf C\mathbf P^1$から$\mathbf C\mathbf P^1$への全単射であることがわかる．
$S:\mathbf C\mathbf P^1\to\mathbf C\mathbf P^1$は連続写像であり，
$S^{-1}$も一次変換なので$S$は同相写像であることが分かる．
更に一次変換全体が群であることも分かる．これを$\PSL(2,\mathbf C)$と表す．

$z$ に関する複素微分は
\[
	S'(z)=\frac{ad-bc}{(c z +d)^2}
\]
であるので，$c=0$ならば定数となり等角写像である．
$c\neq0$のとき$z\neq \frac{-d}{c}$ならば等角写像であるが$z=\frac{-d}{c}$では微分が定義されていないので等角性については何も言っていない．
$w=\frac{c}{cz+d}$を$\frac{-d}{c}$の近傍での新しい座標系とすると
\[
	S(w)=-(ad-bc)w+\frac{a}{c}
\]
となり，この座標系では微分は定数であるので，任意の点で等角写像であることが分かる．更に
$S:\mathbf C\mathbf P^1\to\mathbf C\mathbf P^1$は双正則写像であることも分かった．

$c\neq0$の時は
\[
	\frac{a z +b}{c z +d} 
	=- \frac{ad-bc}{c}\frac{1}{c z +d}  +  \frac{a}{c}
\]
で，$c=0$のときは$\frac{a}{d} z +\frac{b}{d}$なので一般の一次変換は
\[
	z + c,\quad kz,\quad  \frac{1}{z}
\]
の合成により生成されている．

一次変換$Tz= \frac{a z +b}{c z +d}$が$1$, $0$, $\infty$を固定すると
\[
\begin{split}
	&1=T1= \frac{a +b}{c +d},\quad 0=T0=\frac{b}{d},\quad 
	\infty=T\infty=\frac{a}{c}\\
	&\Rightarrow \quad
	b=0,\quad c=0, \quad a=d
\end{split}
\]
となり$T=\id$がわかる．このことから，
二つの一次変換$Sz= \frac{a z +b}{c z +d}$,
一次変換$Tz= \frac{a' z +b'}{c' z +d'}$同じ変換を定義するとすると，
ある$e\neq 0$が存在して
\[
	\begin{pmatrix}
		a' & b'\\
		c' & d'
	\end{pmatrix}
	=e
	\begin{pmatrix}
		a & b\\
		c & d
	\end{pmatrix}
\]
したがって$ad-bc=1$となるように選ぶことが出来る．
それでも$e=\pm 1$の不確定性は残るので．
一次変換全体は群$\PSL(2,\mathbf C)= \SL(2,\mathbf C)/\{\pm E\}$を成す．
まとめると
\[
   f:
   \SL(2,\mathbf C)
   \ni \begin{pmatrix}
      a & b\\
      c & d
   \end{pmatrix}
   \mapsto
   \frac{az+b}{ca+d}
   \in \PSL(2,\mathbf C)
\]
は全射準同形で$\Ker f=\{\pm E\}$.

以下$\mathbf C\mathbf P^1$の代わりに$\hat{\mathbf C}= \mathbf C\cup \infty$を用いるが，
これは$\psi_2:V_2\to \mathbf C$を中心に考察することになり，複素関数論での一般的記法と一致している．

相異なる$z_2$, $z_3$, $z_4\in 
\hat{\mathbf C}$
を$1$, $0$, $\infty\in 
\hat{\mathbf C} $に写す一次変換は
\[
	Tz=\frac{z-z_3}{z-z_4}/\frac{z_2-z_3}{z_2-z_4}
	\quad z\in  \hat{\mathbf C}
\]
と一意的に表される．ただし，
$z_2=\infty$, $z_3=\infty$, $z_4=\infty$のときは
\[
	Tz=\frac{z-z_3}{z-z_4}, \quad \frac{z_2-z_4}{z-z_4},\quad
	\frac{z-z_3}{z_2-z_3}
\]
と解釈することにする．
実際，$Tz_2=1$, $Tz_3=0$,  $Tz_4=\infty$は直ぐに分かる．
また，$S$を同じ性質をもつ一次変換とすると$S^{-1}T$は
$1$, $0$, $\infty$を固定する一次変換なので上より$S^{-1}T=\id$
となり，一意性も分かる．そこで，$z_1\in 
\hat{\mathbf C}$ に対して
\begin{equation}\label{crossRatio}
	(z_1,z_2,z_3,z_4)
	=
	\frac{(z_1-z_3)(z_2-z_4)}{(z_2-z_3)(z_1-z_4)}
\end{equation}
と置き，
\textcolor{red}{複比（非調和比）}という．
\begin{Thm}
複比は
任意の一次変換$Tz= \frac{az+b}{cz+d}$ で不変，つまり，
\[
	(Tz_1,Tz_2,Tz_3,Tz_4)
	=
	(z_1,z_2,z_3,z_4)
\]
をみたす．
\end{Thm}
\begin{proof}$Sz = (z,z_2,z_3,z_4)$と置くと
$
	S z_2 = 1,  Sz_3 = 0,  S z_4 = \infty
$
である．$ST^{-1}$も一次変換で
\[
	ST^{-1}Tz_2=Sz_2=1,\quad ST^{-1}Tz_3=Sz_3=0,\quad
	ST^{-1}Tz_4=Sz_4 = \infty. 
\]
よって
\[
	(z_1, z_2, z_3, z_4)
	=Sz_1
	=ST^{-1}Tz_1
	=(Tz_1,Tz_2,Tz_3,Tz_4)
\]
%\[
%	Sz = (z,z_2,z_3,z_4),\quad Rz = (z,Tz_2,Tz_3,Tz_4)
%\]
%と置くと，共に一次変換で，
%\[
%\begin{split}
%	&S z_2 = 1,\quad S z_3 = 0,\quad S z_4 = \infty, \\
%	&R T z_2 = 1,\quad R T z_3 = 0,\quad R Tz_4 = \infty.
%\end{split}
%\]
%$ST^{-1}$も一次変換で
%\[
%	ST^{-1}Tz_2=Sz_2=1,\quad ST^{-1}Tz_3=Sz_3=0,\quad
%	ST^{-1}Tz_4=Sz_4 = \infty 
%\]
%一意性より$R=ST^{-1}$なので
%\[
%	(Tz_1,Tz_2,Tz_3,Tz_4)=RTz_1=ST^{-1}Tz_1=Sz_1
%	=(z_1, z_2, z_3, z_4)
%\]
\end{proof}

\begin{Thm}
   $z_1,z_2,z_3,z_4$が直線または円周上にあることと
   $(z_1,z_2,z_3,z_4)$が実数であることは同値である．
\end{Thm}
\begin{proof}
$z_1,z_2,z_3,z_4$が円周上にあれば，偏角を$\arg$と表すと
\[
	\arg(z_1,z_2,z_3,z_4)
	=
	\arg\frac{z_1-z_3}{z_1-z_4}
	-\arg\frac{z_2-z_3}{z_2-z_4}
\]
となり，円周角の定理（とその逆）より明らか(図~\ref{argP})．

$z_1,z_2,z_3,z_4$が直線上にある場合もほぼ自明．
\end{proof}

以下では，直線または円周のことを単に円周ということにする．上の二つの定理から

\begin{Thm}[円円対応]
   任意の一次変換は
   円周を円周に写す．
\end{Thm}

円$C$上の相異なる三つの$z_2$, $z_3$, $z_4$をとる．
$z^*\in \hat{\mathbf C}$が円についての$C$の
$z\in \hat{\mathbf C}$の\textcolor{red}{鏡映}とは
\[
	(z^*, z_2,z_3,z_4)=\overline{(z, z_2,z_3,z_4)}\in \hat{\mathbf C}
\]
が成り立つことである．一次変換の一意性より$z^*$は一意的に定まる．
対称点を与える写像を\textcolor{red}{反転}ということがある．

$z\in C$とすると$(z, z_2,z_3,z_4)$は実数なので$z^*=z$，つまり，
$C$上の点は自分自身と対称である．
$(z^*)^*=z$も簡単に分かる．
円$C$が実軸のとき$(z,1, 0,\infty)=z$なので$z^*=\bar z$ となり，反転は鏡映の一般化であることが分かる．

$C$が 円で$|z-a|^2=R^2$と表されるとき
$
	|z_i-a|^2=R^2
$
($ i =2,3,4$)であり，
複比は任意の一次変換で不変なので
\[
\begin{split}
	(z^*, z_2,z_3,z_4)
	=
	&\overline{(z, z_2,z_3,z_4)}\\
	=
	&\overline{(z-a, z_2-a,z_3-a,z_4-a)}\\
	=
	&(\overline{ z-  a}, \overline{z_2-a} ,\overline{z_3-a},\overline{z_4-a})\\
	=
	&(\bar z- \bar a,\frac{R^2}{z_2-a} ,\frac{R^2}{z_3-a}, \frac{R^2}{z_4-a})\\
	=
	&(\frac{\bar z- \bar a}{R^2}, \frac{1}{z_2-a} , \frac{1}{z_3-a}, \frac{1}{z_4-a})\\
	= 
	&(\frac{R^2}{\bar z- \bar a},z_2-a ,z_3-a, z_4-a)
	= 
	(\frac{R^2}{\bar z- \bar a}+a ,z_2 ,z_3, z_4)
\end{split}
\]
したがって
\[
	z^*=\frac{R^2}{\bar z- \bar a}+a \quad \Leftrightarrow \quad
	z^* -a =\frac{R^2}{| z-  a|^2}(z-a)
\]
よって，$z^*$は$a$と$z$を結ぶ直線の$z$方向の半直線$\ell$にあり，
\[
	|z-a|: R =R:|z^* -a|
\]
なので$z$（あるいは$z^*$）で$\ell$に直交する直線が$C$で交わる点での$C$の接線が$\ell$と交わる点が$z^*$(あるいは$z$)である．

\begin{figure}[htb]
   \centering
   \begin{minipage}{0.45\columnwidth}
      \centering
      \includegraphics{~/Documents/GitHub/bEuclideanGeo/ch3/images/figureSphere.11}
      \caption{反転}
   \end{minipage} 
   \begin{minipage}{0.45\columnwidth}
      \centering
      \includegraphics{~/Documents/GitHub/bEuclideanGeo/ch3/images/figureSphere.9}
      \caption{アポロニウスの円の初等証明}\label{}
   \end{minipage} 
\end{figure}

\begin{Thm}[鏡映原理]
   円$C$を$C'$に写す一次変換は
   円$C$についての対称な点 $z$, $z^*\in \hat{\mathbf C}$を
   円$C'$についての対称な点 $w$, $w^*\in \hat{\mathbf C}$に移す．
\end{Thm}
\begin{proof}
円$C$を実軸に移す一次変換は$C$について対称な$z$と$z^*$を実軸について対称な$w$と$\bar w$に写すので明か．
\end{proof}
\subsection{一次変換の分類}
$\alpha\neq \beta$とする．次の形の一次変換
\begin{equation}\label{st0}
	w =k \frac{z-\alpha}{z-\beta}
\end{equation}
は$\alpha$を$0$に，$\beta$を$\infty$に写す．
$w$平面で原点$0$から出る直線$C'_0$は$\alpha$と$\beta$を通る円$C_0$になる．また
$w$平面での$0$を中心とした同心円$C'_1$上のこの一次変換による像$C_1$は円になる．この円は
\[
	 \frac{|z-\alpha|}{|z-\beta|} = \frac{|w|}{|k|} =\const
\]
を満たすので，$C_1$は$\alpha$と$\beta$を極限点とするアポロニウスの円である．
このような円全体を\textcolor{red}{Steiner circles}%シュタイナーの円族}
という．
\begin{figure}[htb]
   \centering
   \begin{minipage}{0.45\columnwidth}
      \centering
      \includegraphics{~/Documents/GitHub/bEuclideanGeo/ch3/images/figureSphere.13}
      \caption{Steiner circles}
   \end{minipage} 
   \begin{minipage}{0.45\columnwidth}
      \centering
      \includegraphics{~/Documents/GitHub/bEuclideanGeo/ch3/images/figureSphere.12}
      \caption{Steiner circles}\label{}
   \end{minipage} 
\end{figure}

$w$平面で$0$以外で$C'_0$と$C'_1$は直交するので
$\alpha$と$\beta$以外で$C_0$と$C_1$は直交する．
$w$平面で$0$以外で$C'_0$に関する反転は$C'_1$を自分自身に写し$C'_1$に関する反転は$C'_0$を自分自身に写すので，
$\alpha$と$\beta$以外で$C_0$に関する反転は$C_1$を自分自身に写し$C_1$に関する反転は$C_0$を自分自身に写す．
また$w$平面では$C'_1$に関して$0$と$\infty$は対称なので
$\alpha$と$\beta$は任意の$C_1$に関して対称である．

\begin{figure}[htb]
   \centering
   \begin{minipage}{0.45\columnwidth}
      \centering
      \includegraphics{~/Documents/GitHub/bEuclideanGeo/ch3/images/figureSphere.14}
      \caption{楕円型}\label{elliptic1}
   \end{minipage} 
   \begin{minipage}{0.45\columnwidth}
      \centering
      \includegraphics{~/Documents/GitHub/bEuclideanGeo/ch3/images/figureSphere.17}
      \caption{楕円型}\label{elliptic2}
   \end{minipage} 
\end{figure}


一次変換$w=Sz$は$\hat{\mathbf C}$から$\hat{\mathbf C}$への全単射を定義しているが，$S\neq\id$の固定点は
\[
	S(z)= \frac{a z +b}{c z +d}=z \quad
	\Rightarrow \quad
	cz^2 +(d-a)z-b=0
\]
をみたすので，必ず存在し高々2個ある．$\alpha\neq \beta$を異なる二つの固定点とする．\eqref{st0}のような変換を介在することで，$k\neq 0\in \mathbf C$が存在して
\[
	 \frac{w-\alpha}{w-\beta} =k \frac{z-\alpha}{z-\beta}
\]
と書けることが分かる．
$|k|= 1$のとき\textcolor{red}{楕円型}一次変換といい，$k\neq \pm 1,0\in \mathbf R$のとき\textcolor{red}{双曲型}一次変換という．
楕円型は$C_1$を同じ$C_1$に写し，双曲型は$C_0$を同じ$C_0$に写している．楕円型と$\SU(2)$との関係は次の分節で説明する．
\begin{figure}[htb]
   \centering
   \begin{minipage}{0.45\columnwidth}
      \centering
      \includegraphics{~/Documents/GitHub/bEuclideanGeo/ch3/images/figureSphere.15}
      \caption{双曲型}\label{}
   \end{minipage} 
   \begin{minipage}{0.45\columnwidth}
      \centering
      \includegraphics{~/Documents/GitHub/bEuclideanGeo/ch3/images/figureSphere.18}
      \caption{双曲型}\label{}
   \end{minipage} 
\end{figure}

\begin{figure}[htb]
   \centering
   \begin{minipage}{0.45\columnwidth}
      \centering
      \includegraphics{~/Documents/GitHub/bEuclideanGeo/ch3/images/figureSphere.16}
      \caption{放物型}\label{}
   \end{minipage} 
   \begin{minipage}{0.45\columnwidth}
      \centering
      \includegraphics{~/Documents/GitHub/bEuclideanGeo/ch3/images/figureSphere.19}
      \caption{放物型}\label{}
   \end{minipage} 
\end{figure}
$z=\alpha$ を重根とするとき，この一次変換を\textcolor{red}{放物型}一次変換という．
\[
	w =z+c\quad c\neq 0
\]
は$w=z=\infty$を唯一の固定点とする放物型一次変換であり，$w=u+iv$とすると
$u=\const$は無限遠点を通る平行直線群$C_0'$で無限遠点で全てが接している．$v=\const$も同様に平行直線群$C_1'$を成し，$C_0'$直線と全ての点で直交している．
$\omega/(z-\alpha)$で写すと，$\alpha$を固定点とする放物型一次変換は
\[
	\frac{\omega}{w-\alpha}=\frac{\omega}{z-\alpha}+c
\]
をみたしていることが分かる．$C_0'$と$C_1'$をこの変換で写した直線群$C_0$と$C_1$は
それぞれが$\alpha$で接していて，任意の点で直交している．

その他の一次変換は\textcolor{red}{混合型}と言われるが，これも双曲型にいれられることもある．

与えられた一次変換がどの型なのかを係数のみで判定できる：
\begin{Thm}
一次変換
\[
	 Sz=\frac{a z +b}{c z +d}
 \]
 の係数が$ad-bc=1$と正規化されているとする．このとき
 \[
 \begin{split}
    a+d\in \mathbf R\quad \text{and}\quad-2<a+d<2 \quad &\Leftrightarrow\quad S \text{ 楕円型}\\
    a+d\in \mathbf R\quad \text{and}\quad |a+d|>2 \quad &\Leftrightarrow\quad S \text{ 双曲型}\\
    a+d\in \mathbf R\quad \text{and}\quad a+d=\pm 2\quad &\Leftrightarrow\quad S \text{ 放物型}
 \end{split}
 \]
\end{Thm}
%\begin{proof}
%$Sz=z$を満たす$z$は
%\[
%	cz^2 + (d-a)z -b=0
%\]
%を満たすのでその判別式$D$をみると
%\[
%	D=(d-a)^2+4bc=(a+d)^2-4(da-bc)=(a+d)^2-4
%\]
%放物型と$D=0$は同値で，他の場合，根$\alpha$, $\beta$は
%\[
%	\alpha\beta=1\quad \Rightarrow\quad
%	\beta = 1/\alpha\quad \Rightarrow\quad
%\]
%\end{proof}


\section{$S^2(1)$と$S^3(1)$の幾何学}
\subsection{$S^2(1)$の微分幾何}
$\mathbf R^3=(\mathbf R^3, (,))$を$3$次元ユークリッド空間とする．
$(,)$は内積で
$
	\|\mathbf x\|=\sqrt{(\mathbf x,\mathbf x)}
$とおく．
原点を中心とした半径$1>0$の球面を
\[
	S^2(1)=\{\mathbf x\in \mathbf R^3\mid \|\mathbf x\|=1\}
\]
と表す．単に$S^2$と表すこともある．

$S^2(1)$には$\mathbf R^3$の計量（内積）から導かれるリーマン計量から導かれる距離が入るが，それは
二点$A$, $B$と原点を結ぶ線分がなす角度と一致する．また最短測地線は$A$, $B$と原点を通る平面と球面$S^2(1)$の交わり（大円）の$A$, $B$に囲まれた円弧に一致する．測地線のパラメータは長さに取る．
$A$と$B$が対蹠点でなければ最短線は一意的である．$A$と$B$が対蹠点のとき最短測地線は
$A$と$B$を結ぶ（大円の）半円となり，無限の取り方がある．$A$から出る
二つの最短線の角度を$A$での接空間での二つの接ベクトルのなす角とし，誤解が無い場合は$\angle A$などと表す．


\begin{figure}[htb]
   \centering
   \begin{minipage}{0.4\columnwidth}
      \centering
      \includegraphics{~/Documents/GitHub/bEuclideanGeo/ch3/images/figureSphere.1}
      \caption{最短測地線}
   \end{minipage} 
   \begin{minipage}{0.4\columnwidth}
      \centering
      \includegraphics{~/Documents/GitHub/bEuclideanGeo/ch3/images/figureSphere.2}
      \caption{極座標}
   \end{minipage} 
\end{figure}

空間極座標：空間$\mathbf R^3$の点$\mathbf x=(x,y,z)\neq\mathbf 0$と原点との距離を$r>0$，$\mathbf x$と$z$軸とのなす角度を$0\le \theta\le \pi$,
$\mathbf x$の$xy$平面への射影が$x$軸となす角を $\varphi\in (-\pi,\pi]$
とすると，
\[
	x=r\sin \theta\cos \varphi,\quad 
	y=r\sin \theta\sin \varphi, \quad 
	z=r\cos \theta
\]
と表せる．$\theta$は緯度，$\varphi$は経度に対応する．
地理では緯度を赤道を$0$なるように取るが，数学では
$\theta$を北極からの角度にするのが普通である．
ヤコビアンを計算すると
\[
\begin{split}
	\frac{\partial(x,y,z)}{\partial(r,\theta,\varphi)}
	&=\begin{vmatrix}
	 	\frac{\partial x}{\partial r}&\frac{\partial x}{\partial \theta}&\frac{\partial x}{\partial \varphi}\\
		\frac{\partial y}{\partial r}&\frac{\partial y}{\partial \theta}
		&\frac{\partial y}{\partial \varphi}\\
		\frac{\partial z}{\partial r}&\frac{\partial z}{\partial \theta}
		&\frac{\partial z}{\partial \varphi}
	\end{vmatrix}=\begin{vmatrix}
	 	\sin \theta\cos \varphi&r\cos \theta\cos \varphi&-r\sin \theta\sin\varphi\\
		\sin \theta\sin \varphi&r\cos \theta\sin \varphi&r\sin \theta\cos \varphi\\
		\cos \theta&-r\sin \theta&0
	\end{vmatrix}\\
	&=
	(-1)^{3+1}\cos \theta
	\begin{vmatrix}
	 	r\cos \theta\cos \varphi&-r\sin \theta\sin\varphi\\
		r\cos \theta\sin \varphi&r\sin \theta\cos \varphi\\
	\end{vmatrix}\\
	&+(-1)^{3+2}(-r\sin \theta)
	\begin{vmatrix}
	 	\sin \theta\cos \varphi&-r\sin \theta\sin\varphi\\
		\sin \theta\sin \varphi&r\sin \theta\cos \varphi\\
	\end{vmatrix} +0\\
	&=r^2\cos^2 \theta\sin \theta +r^2\sin^3 \theta
	=r^2\sin \theta\\
	\end{split}
\]
ただし，3行目にラプラス（余因子）展開を使った．
よって変数変換の式は
\begin{equation*}
\begin{split}
	&\int_D f(x,y,z)\, dxdydz
	\\
	=&\int_{E} 
		f(r\sin \theta\cos \varphi,
		r\sin \theta\sin \varphi, 
		r\cos \theta)\,
		|r^2\sin \theta|\,  
	drd\theta d\varphi
\end{split}
\end{equation*}
である．
$S^2(1)$上の二つの（大円の）半円に囲まれた二角形の面積(図~\ref{biangle})を計算すると
\begin{equation}\label{areaBiangle}
	\int_0^\pi d\theta \int_0^{\angle A}d\varphi 
	\sin \theta
	=\angle A [-\cos \theta ]_0^\pi=2\angle A
\end{equation}

\begin{figure}[htb]
   \centering
   \begin{minipage}{0.4\columnwidth}
      \centering
      \includegraphics{~/Documents/GitHub/bEuclideanGeo/ch3/images/figureSphere.3}
      \caption{極座標のヤコビアン}
   \end{minipage} 
   \begin{minipage}{0.4\columnwidth}
      \centering
      \includegraphics{~/Documents/GitHub/bEuclideanGeo/ch3/images/figureSphere.4}
      \caption{球面二角形の面積}\label{biangle}
   \end{minipage} 
\end{figure}

球面三角形の面積：$\alpha$を求めたい球面三角形$ABC$の面積とし，
図~\ref{diagram}のように面積に名前を付け, 見えない面の三角形の面積を$\kappa$と置くと，$8$枚の球面三角形で球面全体が覆われるので
\[
	\alpha + \beta+\gamma +\delta+\varepsilon+
	\eta+\zeta+\kappa 
	=4\pi
\]
対頂角は等しいので，\eqref{areaBiangle}より
\begin{gather*}
	\alpha + \beta =2 \angle A,\quad
	\alpha + \delta =2 \angle B,\quad
	\alpha + \zeta =2 \angle C,\\
	\kappa + \varepsilon =2 \angle A,\quad
	\kappa + \eta =2 \angle B,\quad
	\kappa + \gamma =2\angle C,
\end{gather*}
よって
\[
	2\alpha+2\kappa +4\pi= 4(\angle A+\angle B+\angle C)
\]
$\alpha=\kappa$なので
\[
	\alpha = \angle A+\angle B+\angle C-\pi
\]
\begin{figure}[htb]
   \centering
   \begin{minipage}{0.4\columnwidth}
      \centering
      \includegraphics{~/Documents/GitHub/bEuclideanGeo/ch3/images/figureSphere.5}
      \caption{球面三角形}
   \end{minipage} 
   \begin{minipage}{0.4\columnwidth}
      \centering
      \includegraphics{~/Documents/GitHub/bEuclideanGeo/ch3/images/figureSphere.6}
      \caption{模式図}\label{diagram}
   \end{minipage} 
\end{figure}
(このエレガントな証明は Thomas Harriot, 1603 によると言われている, \cite{L}, \cite{S}%\S 17.6 7
).
これを境界が測地線の場合のガウスボンネの公式
\[
\begin{split}
	\int_{\triangle ABC}Kd\vol 
	&= 2\pi - (\pi-\angle A)-(\pi-\angle B)-(\pi-\angle C)\\
	&=\angle A+\angle B+\angle C-\pi=\alpha
\end{split}
\]
と較べるとガウス曲率$K$が$=1$であることも分かる．
\emptyline

2回目の授業で示した$3$種類の極と$\alpha_i = |\stab(\mathbf x_i)|$との関係
\[
	1+\frac2n 
	=
	\frac1{\alpha_1}+\frac1{\alpha_2}+\frac1{\alpha_3}
\]
を書き換えると
\[
	\frac{\pi}{\alpha_1}+\frac{\pi}{\alpha_2}+\frac{\pi}{\alpha_3}-\pi
	=	
	\frac{2\pi}n 
\]
であり，これは異なる軌道に属する三つの極の最も近いものを最短線で結んだ測地三角形の面積の公式となっている(表 \ref{classification}参照)．
\begin{table}[h]
  \centering
  \begin{tabular}{l|llllllll} 
  	\noalign{\hrule height 1pt}
	 $G$ & $\alpha_1$ &$\alpha_2$&$\alpha_3$& $|G|$ \\ \hline \hline
	$C_n$ $(n\in \mathbf N)$&$n$&$n$&$1$  & $n$\\ \hline
	$D_m$ $(m\in \mathbf N)$&$2$&$2$&$m$  &  $2m$\\ \hline
	$ T$&$2$& $3$ & $3$ &$12$ \\ \hline
	$W$& $2$&$3$  & $4$ &$24$ \\ \hline
	$I$&$2$& $3$ & $5$&$60$\\ \hline
     \noalign{\hrule height 1pt}
  \end{tabular}
\caption{$SO(3)$の有限部分群の分類}\label{classification}
\end{table}

%\begin{figure}[htb]
%   \centering
%   \begin{minipage}{0.4\columnwidth}
%      \centering
%      \includegraphics{~/Documents/GitHub/bEuclideanGeo/ch3/images/figureSphere.20}
%      \caption{回転の合成}
%   \end{minipage} 
%   \begin{minipage}{0.4\columnwidth}
%      \centering
%      \includegraphics{~/Documents/GitHub/bEuclideanGeo/ch3/images/figureSphere.23}
%      \caption{オイラー角}\label{EulerAngle}
%   \end{minipage} 
%\end{figure}
回転の合成と軸の変換：
向きを保つ$S^2(1)$の等長変換全体が$\SO(3)$となる．

\subsection{$\SU(2)=\Sp(1)=S^3(1)$}
球面との関係で表れる低次元の古典群の間の関係について述べたい．
$\mathbf a=\begin{pmatrix}a_1\\a_2\\ \vdots\\a_n\end{pmatrix}$, 
$\mathbf b=\begin{pmatrix}b_1\\b_2\\ \vdots\\b_n\end{pmatrix}\in \mathbf C^n$
に対して
\[
	\langle\mathbf a, \mathbf b\rangle
	=\bar a_1b_1+\bar a_2b_2+\cdots +\bar a_nb_n
\]
と定めるとエルミート内積となる．これを標準的エルミート内積という．
複素係数の$ n$次正方行列$A=(a_{ij})$に対して$\bar A=(\bar a_{ij})$とおき ${}^tA$を$A$の転置とするとき
\[
	A^*={}^t\bar A
\]
とおく．
複素係数の$n$次正方行列$U$
が
\[
	U^*U=E
\]
をみたすとき$n$次\textcolor{red}{ユニタリ行列}という．これは$U^{-1}=U^*$と同値である．これはまた標準的エルミート内積$\langle,\rangle $について
\[
	\langle U\mathbf a, U\mathbf b\rangle =\langle\mathbf a, \mathbf b\rangle 
\]
が任意の$\mathbf a, \mathbf b\in \mathbf C^n$について成り立つことと同値である．

$n$次の（複素係数の）正則行列全体$\GL(n,\mathbf C)$ が群をなすことは実の場合と同様に分かる．
$\GL(n,\mathbf C)$は$n$\textcolor{red}{次複素一般線形群}と呼ばれる．
\[
	\SL(n,\mathbf C) = \{A\in \GL(n,\mathbf C)\mid\det A = 1 \}
\]
は$\GL(n,\mathbf C)$の部分群で$n$\textcolor{red}{次複素特殊線形群}という．
このセクションの初めにリーマン球面$\mathbf R\mathbf P^1$の一次変換全体
$\PSL(2,\mathbf C)$が群となることを示した．
これを，$2$\textcolor{red}{次複素射影特殊線形群}という．
\[
	\U(n)= \{A \in \GL(n,\mathbf C)\mid n \text{次ユニタリ行列}\}
\]
も$\GL(n,\mathbf C)$の部分群で， $n$\textcolor{red}{次ユニタリ群}という．特に，
\[
	\U(1)=\{a\in \mathbf C \mid \bar a a =1\}=S^1(1)
\]
である．
\[
	\SU(n)= \U(n)\cap \SL(n,\mathbf C)
\]
も$U(n)$の部分群で，$n$\textcolor{red}{次特殊ユニタリ群}という．

まず$\SL(2,\mathbf C)$, $\SU(2)$,  $\PSL(2,\mathbf C)$, $\SO(3)$の間の関係を調べる．
$U\in \SU(2)$ を 
$U=\begin{pmatrix}
   a &b\\
   c& d\\
\end{pmatrix}$と置くと，$\det U=ad-bc=1$と$U^{-1}=U^{*}$より
\[
   \frac{1}{1}
   \begin{pmatrix}
      d &-b\\
      -c& a\\
   \end{pmatrix}
   =
   \begin{pmatrix}
      \overline{a} & \overline{c}\\
       \overline{b}&  \overline{d}\\
   \end{pmatrix}, 
   \quad
   ad-bc=1
\]
より，
\[
   \overline{a} =d, \quad \overline{b}=-c, \quad a\overline{a}+b\overline{b}=1
\]
$a=x_0 + x_3i$, $b=-x_2 + x_1i$と置くと
\begin{equation}\label{quaternion}
   U
   =
   \begin{pmatrix}
      x_0 + x_3i &-x_2 + x_1i\\
      x_2 + x_1i& x_0 - x_3i\\
   \end{pmatrix},
   \quad
   x_0^2+x_1^2+x_2^2+x_3^2=1
\end{equation}
が分かり
\begin{Prop}
写像
\[
   \SU(2) \ni  
   \begin{pmatrix}
      x_0 + x_3i &-x_2 + x_1i\\
      x_2 + x_1i& x_0 - x_3i\\
   \end{pmatrix}
   \mapsto 
   ( x_0 , x_1, x_2 ,  x_3)\in  S^3(1)
\]
は同相である．
\end{Prop}
\noindent
ことまで分かる．

$\SU(2)$の積構造について調べる．
\eqref{quaternion}より
\[
   \mathbf e_0=E,\quad
   \mathbf e_1=
   \begin{pmatrix}
     0 &i\\
      i& 0\\
   \end{pmatrix},\quad
  \mathbf e_2
   =
   \begin{pmatrix}
     0 &-1\\
      1& 0\\
   \end{pmatrix},\quad
  \mathbf e_3
   =
   \begin{pmatrix}
      i &0\\
      0& -i\\
   \end{pmatrix}
\]
と置くと任意の$A\in \SU(2)$は
\[
   A=
   x_0 \mathbf e_0+ x_1\mathbf e_1
   + x_2\mathbf e_2+ x_3\mathbf e_3,\quad
    x_0^2+x_1^2+x_2^2+x_3^2=1
\]と表せる．
$\mathbf e_1$, 
$\mathbf e_2$, 
$\mathbf e_3$ の積は
\begin{gather*}
	\mathbf e_1\mathbf e_1=\mathbf e_2\mathbf e_2=\mathbf e_3\mathbf e_3=-\mathbf e_0,\\
	\mathbf e_1\mathbf e_2=-\mathbf e_2\mathbf e_1=\mathbf e_3,\quad 
	\mathbf e_2\mathbf e_3=-\mathbf e_3\mathbf e_2=\mathbf e_1,\quad
	\mathbf e_3\mathbf e_1=-\mathbf e_3\mathbf e_1=\mathbf e_2
\end{gather*}
をみたし，
$\mathbf e_0\Rightarrow 1$, 
$\mathbf e_1\Rightarrow i$,
 $\mathbf e_2\Rightarrow j $, 
 $\mathbf e_3\Rightarrow k $
 としたときの$4$元数体$\mathbf H$と同じ積構造を持っている．

$a=x_0 + x_1i+x_2 j+ x_3k\in \mathbf H$に対して
$ \overline{a}=
 x_0 - x_1i-x_2 j-x_3k$と置くと簡単に
 \[
    a\overline{a}=x_0^2+x_1^2+x_2^2+x_3^2
 \]
 がわかる．また
 \[
   \overline{ab}=\overline{b}\overline{a}
 \]
 も分かる．$\mathbf H$係数の$ n$次正方行列$A=(a_{ij})$に対して$\bar A=(\bar a_{ij})$とおき ${}^tA$を$A$の転置とするとき
\[
   A^*={}^t\bar A
\]
とおくと
\[
   (AB)^*=B^*A^*
\]
が分かる．
 $\mathbf H$係数の$n$次正方行列$S$が
\[
	S^*S=E
\]
をみたすとき$n$\textcolor{red}{次シンプレクティック行列}という．
\[
	\Sp(n)
	=\{
      	S\mid n\text{次シンプレクティック行列}
	\}
\]
と置き$n$\textcolor{red}{次シンプレクティック群}という．
上で示したことからこれは確かに群をなしている．
特に
\[
   \Sp(1)=\{a\in \mathbf H\mid  \overline{a}a =1\}=S^3(1)
\]
なので
$\rho: \SU(2)\to\Sp(1)$を
\[
   \rho
   \begin{pmatrix}
      x_0 + x_3i &-x_2 + x_1i\\
      x_2 + x_1i& x_0 - x_3i\\
   \end{pmatrix}
   = 
   x_0 + x_1i+x_2 j+ x_3k
\]
で定めると well-defined で全単射である．
\begin{Thm}
写像$\rho: \SU(2)\to\Sp(1)$
は（積についての群の）同型である．
\end{Thm}
\begin{proof}
$\rho: \SU(2)\to\Sp(1)$は
\[
   \rho(\mathbf e_0)=1, \quad
   \rho(\mathbf e_1)=i, \quad
   \rho(\mathbf e_2)=j, \quad
   \rho(\mathbf e_3)=k
\]
をみたしている．上の計算より
\[
   \rho(\mathbf e_l\mathbf e_m)
   = \rho(\mathbf e_l) \rho(\mathbf e_m)\in \mathbf H
\]
なので
\[
   \rho(U'U)=\rho(U')\rho(U)
\]
が分かる．
\end{proof}

\subsection{$SO(3)=\mathbf R\mathbf P^3$}
$SO(3)$と$S^3(1)$の関係を調べたい．
$A\in \SO(3)$を固定し，標準基を
$\{\mathbf e_i=(\delta_{ij})\in \mathbf R^3\}_{i=1,2,3}$とする，
点$A\mathbf e_3\in S^2(1)$の極座標を
 $\varphi\in (-\pi,\pi]$, $\theta\in[0,\pi]$とし$A\mathbf e_3$を$\varphi$だけ$z$軸について逆回転$R_{-\varphi}$し, 
 $xz$平面の$x\ge 0$の部分に写し，
 $xz$平面で$-\theta$だけ$y$軸について回転$L_{-\theta}$すると
\[
	L_{-\theta}R_{-\varphi}A\mathbf e_3=\mathbf e_3
\]
となる（図~\ref{EulerAngle1}）．
これは$\mathbf e_3$を固定する回転なので，
ある$\psi$が存在して$R_{\psi}$と一致する（図~\ref{EulerAngle2}のように$\mathbf e_1$の動きをみると理解しやすい）．したがって，
\begin{equation}\label{Euler}
	A=R_{\varphi}L_{\theta}R_{\psi}
\end{equation}
と表せる．
上の構成から$0\le \varphi<2\pi$, $0\le \theta\le\pi$, $0\le \psi<2\pi$ である．
これらの角$\varphi$, $\theta$, $\psi$を$A\in \SO(3)$の\textcolor{red}{オイラー角}\footnote{ただし，力学などでは$y$軸での回転の代わりに$x$軸での回転を使っているものも多いが，量子力学では$y$軸回転の方が自然らしい\cite{JJS} \P171--4．どちらを使うかで以下の式は微妙に異なるので注意してほしい．}という．
\begin{figure}[htb]
   \centering
   \begin{minipage}{0.4\columnwidth}
      \centering
      \includegraphics{~/Documents/GitHub/bEuclideanGeo/ch3/images/figureSphere.21}
      \caption{Euler角:$\mathbf e_3$の動き}\label{EulerAngle1}
   \end{minipage} 
   \begin{minipage}{0.4\columnwidth}
      \centering
      \includegraphics{~/Documents/GitHub/bEuclideanGeo/ch3/images/figureSphere.22}
      \caption{Euler角:$\mathbf e_1$の動き}\label{EulerAngle2}
   \end{minipage} 
\end{figure}

\[
   R_{\psi} = 
    \begin{pmatrix}
      \cos \psi& -\sin\psi&0\\
      \sin\psi& \cos\psi&0\\
      0&0&1
   \end{pmatrix},\quad
   L_{\theta} = 
   \begin{pmatrix}
      \cos \theta& 0&\sin \theta\\
      0&1&0\\
      -\sin \theta&0& \cos \theta\\
   \end{pmatrix}
\]
なので，上の行列の積を計算すると$A$は
\[
   \begin{pmatrix}
      \cos \theta\cos \varphi\cos\psi - \sin \varphi\sin \psi&
         -\cos \theta\cos \varphi\sin\psi - \sin \varphi\cos \psi&
         \sin\theta\cos \varphi\\
      \cos \theta\sin \varphi\cos\psi + \cos \varphi\sin \psi&
         -\cos \theta\sin \varphi\sin\psi + \cos \varphi\cos \psi&
         \sin\theta\sin \varphi\\\
      -\sin \theta\cos\psi&\sin \theta\sin \psi&\cos \theta
   \end{pmatrix}
\]
と表される．
$A\mathbf e_3$の表示からその球面極座標は$(\theta, \varphi)$である．
\[
\begin{split}
	&( \mathbf e_1,  \mathbf e_2,  \mathbf e_3)R_{\varphi}
	=( \mathbf e'_1,  \mathbf e'_2,  \mathbf e'_3)\\
	&( \mathbf e'_1,  \mathbf e'_2,  \mathbf e'_3)L_{\theta}
	=( \mathbf e''_1,  \mathbf e''_2,  \mathbf e''_3)\\
	&( \mathbf e''_1,  \mathbf e''_2,  \mathbf e''_3)R_{\psi}
	=( \mathbf e'''_1,  \mathbf e'''_2,  \mathbf e'''_3)
\end{split}
\]
と置くと，
\[
   A
   =( \mathbf e_1,  \mathbf e_2,  \mathbf e_3)
   A
   =( \mathbf e'''_1,  \mathbf e'''_2,  \mathbf e'''_3)
\]
で，$\{\mathbf e_1,  \mathbf e_2,  \mathbf e_3\},\dots, 
\{\mathbf e'''_1,  \mathbf e'''_2,  \mathbf e'''_3\}$はすべて$\mathbf R^3$の正規直交基である．具体的に計算すると
\[
\begin{split}
   &\mathbf e_1'
   =\cos\varphi \,\mathbf e_1+ \sin\varphi \,\mathbf e_2,\quad
   \mathbf e_2'
   =-\sin\varphi \,\mathbf e_1+ \cos\varphi \,\mathbf e_2,\quad
   \mathbf e_3'
   =\mathbf e_3
   \\
   &\mathbf e''_1
   =
   \cos \theta\,\mathbf e_1'-\sin\theta\,\mathbf e'_3, \quad
   \mathbf e''_3
   =
   \sin \theta\,\mathbf e_1'+\cos\theta\,\mathbf e'_3, \quad
   \mathbf e''_2
   =
   \mathbf e'_2
   \\
   &\mathbf e'''_1
   =\cos\psi\, \mathbf e''_1+ \sin\psi\, \mathbf e''_2,\quad
   \mathbf e'''_2
   =-\sin\psi\, \mathbf e''_1+ \cos\psi\, \mathbf e''_2,\quad
   \mathbf e'''_3
   =\mathbf e''_3
\end{split}
\]
よって，まず$\{\mathbf e_1,  \mathbf e_2,  \mathbf e_3\}$を$\mathbf e_3$中心に$\varphi$回転したものが$\{\mathbf e'_1,  \mathbf e'_2,  \mathbf e'_3\}$で，それを$\mathbf e'_2$を中心に$\theta$回転したのが$\{\mathbf e''_1,  \mathbf e''_2,  \mathbf e''_3\}$, それを$\mathbf e''_3$を中心に$\psi$回転したのが$\{\mathbf e'''_1,  \mathbf e'''_2,  \mathbf e'''_3\}$である．これが
オイラー角の幾何学的意味である．
\begin{figure}[htb]
   \centering
   \begin{minipage}{0.4\columnwidth}
      \centering
      \includegraphics{~/Documents/GitHub/bEuclideanGeo/ch3/images/figureSphere.23}
      \caption{オイラー角 I}\label{EulerAngle1}
   \end{minipage} 
   \begin{minipage}{0.4\columnwidth}
      \centering
      \includegraphics{~/Documents/GitHub/bEuclideanGeo/ch3/images/figureSphere.24}
      \caption{オイラー角 II}\label{EulerAngle2}
   \end{minipage} 
\end{figure}


\emptyline
特殊ユニタリ行列
$
   U
   =
   \begin{pmatrix}
      a & -\overline{b}\\
     b&  \overline{a}\\
   \end{pmatrix}
   \in \SU(2)
$ の固有多項式は
\[
   \begin{vmatrix}
      t-a &\overline{b}\\
     -b& t- \overline{a}\\
   \end{vmatrix}
   =( t-a)(t-\overline{a}) + b\overline{b}
   =t^2 - 2\RE a\, t+1
\]
$-1\le \RE a\le 1$より固有値は$\lambda =e^{\theta}$, $e^{-\theta}$.
$\theta \not\equiv 0 \mod 2\pi$のときは異なる2つの固有値になるが，
$\theta \equiv 0 \mod 2\pi$のとき$\lambda =1$は重根で，この場合
$\RE a=1$つまり$a=1$が分かり$b=0$もわかるので$U=E$となる．
よって対角化が出来てユニタリ行列$V\in \U(2)$が存在して
\[
   V^{*}UV
   =
   \begin{pmatrix}
      e^{i\theta} &0\\
     0& e^{-i\theta}\\
   \end{pmatrix}
\]
となる．
$V = 
   \begin{pmatrix}
      \alpha & -\overline{\beta}\\
     \beta&  \overline{\alpha}\\
   \end{pmatrix}
$
とおくと，固有ベクトルとして
\[
   \mathbf e_1
   =
   \begin{pmatrix}
      \alpha \\
      \beta\\
   \end{pmatrix},\quad
   \mathbf e_2
   =\begin{pmatrix}
      -\overline{\beta}\\
      \overline{\alpha}\\
   \end{pmatrix}
\]
が取れて
\[
	U\mathbf e_1=e^{i\theta}\mathbf e_1,\quad
	U\mathbf e_2=e^{-i\theta}\mathbf e_2.
\]
$U$のリーマン球面$\mathbf C\mathbf P^1$への一次分数変換としての作用をみると，
同次座標を考えることで
\[
	U[\alpha : \beta]=[\alpha : \beta],\quad
	U[-\overline{\beta}:\overline{\alpha} ]
	=[-\overline{\beta}:\overline{\alpha} ]
\]
つまり，
$[\alpha : \beta]$
と
$[-\overline{\beta}:\overline{\alpha} ]\in \mathbf C\mathbf P^1$
は一次分数変換$U$の固定点である．
\begin{figure}[htb]
   \centering
   \begin{minipage}{0.4\columnwidth}
      \centering
      \includegraphics[trim=6 6 6 6, clip]{~/Documents/GitHub/bEuclideanGeo/ch3/images/figureSphere.17}
       \caption{図~\ref{elliptic2}再掲}\label{elliptic2'}
   \end{minipage} 
   \begin{minipage}{0.4\columnwidth}
      \centering
      \includegraphics{~/Documents/GitHub/bEuclideanGeo/ch3/images/figureSphere.26}
      \caption{$\SU(2)$の元の$\mathbf C\mathbf P^1$への作用}
   \end{minipage} 
\end{figure}

$\alpha\ne 0$ と仮定して一般性を失わないのでそう仮定する．
もし，$\beta=0$ならば$[1:0]=0$と$[0:1]=\infty$が固定点で
$V=E$と考えてよく，このとき一次変換は
\[
	\frac{e^{i\theta}z+0}{0z+e^{-i\theta}}
	=e^{2i\theta}z
\]
となる．つまり$0$と$\infty$を固定をする楕円型の一次変換であるが，
これは図~\ref{elliptic2'}のように$\mathbf C$の原点中心とした回転で，
立体射影で写すと$S^2$の対蹠点$0$と$\infty$を中心とする
$S^2\subset \mathbf R^3$の回転で，
$\mathbf R^3$の$z$軸を軸とする回転に拡張される．
対応する$\SO(3)$の元は
\[
   \begin{pmatrix}
      \cos 2\theta& -\sin 2\theta&0\\
      \sin 2\theta& \cos 2\theta&0\\
      0&0&1
   \end{pmatrix}
\]

\begin{figure}[htb]
   \centering
   \begin{minipage}{0.4\columnwidth}
      \centering
      \includegraphics[trim=6 6 6 6, clip]{~/Documents/GitHub/bEuclideanGeo/ch3/images/figureSphere.14}
       \caption{図~\ref{elliptic1}再掲}\label{elliptic1'}
   \end{minipage} 
   \begin{minipage}{0.4\columnwidth}
      \centering
      \includegraphics{~/Documents/GitHub/bEuclideanGeo/ch3/images/figureSphere.27}
      \caption{$\SU(2)$の元の$\mathbf C\mathbf P^1$への作用}
   \end{minipage} 
\end{figure}

$\beta\neq0$とすると，$u=\frac{\beta}{\alpha}$と
\[
   \frac{-\overline{\alpha}}{\overline{\beta}}=-\frac{1}{\overline{u}}=-\frac{u}{|u|^2}
\]
が固定点であることが分かるが，\eqref{stereoProj}より
\[
	\varphi(u)=(\frac{2u}{|u|^2+1}, \frac{|u|^2-1}{|u|^2+1}),\quad
	\varphi(-\frac{1}{\overline{u}})=-(\frac{2u}{|u|^2+1}, \frac{|u|^2-1}{|u|^2+1})
\]
となり，やはり$S^2(1)$の対蹠点である．一次変換で
\[
   z'=\frac{z-u}{z+\frac{1}{\overline{u}}},\quad
   w'=\frac{w-u}{w+\frac{1}{\overline{u}}}
\]
変換することで，楕円型の一次変換
\[
   \frac{w-u}{w+\frac{1}{\overline{u}}}
   = e^{2i\theta}\frac{z-u}{z+\frac{1}{\overline{u}}},
\]
になることが分かり，
$\varphi^{-1}(u)$を含む直線を軸とする$\mathbf R^3$の回転になることが分かる．
特に$u=i$と$\frac{-1}{\overline{u}}=-i\in \mathbf C$を固定すると
$
   \frac{1}{\sqrt2}
   \begin{pmatrix}
      1\\
      i\\
   \end{pmatrix}
$, 
$
   \frac{1}{\sqrt2}
   \begin{pmatrix}
      1\\
      -i\\
   \end{pmatrix}
$
が標準エルミート内積について$\mathbf C^2$の正規直交基底となるので
\[
   \frac{1}{\sqrt2}
   \begin{pmatrix}
      1&1\\
      i&-i\\
   \end{pmatrix}
   \begin{pmatrix}
      e^{i\theta} &0\\
     0& e^{-i\theta}\\
   \end{pmatrix}
   \frac{1}{\sqrt2}
   \begin{pmatrix}
      1&-i\\
      1&i\\
   \end{pmatrix}
   =
   \begin{pmatrix}
      \cos{\theta} &-\sin \theta\\
     \sin \theta& \cos{\theta}\\
   \end{pmatrix}
\]
が対応する$\SU(2)$の元で，$\mathbf R^3$の$y$軸を中心とする右手系の回転は$z$軸を$x$軸の方向に回転するので，$\SO(3)$の作用で表すと
\[
   \begin{pmatrix}
      \cos 2\theta& 0&\sin 2\theta\\
      0&1&0\\
      -\sin 2\theta&0& \cos 2\theta\\
   \end{pmatrix}
\]
となる．


$\SU(2)\subset \SL(2,\mathbf C)$ は部分群であることは定義から明らか．上で示したように $f:\SL(2,\mathbf C)\to\PSL(2,\mathbf C)$は全射準同形である．よって
$f|_{\SU(2)}:\SU(2)\to\PSL(2,\mathbf C)$は準同形であり，
上の考察から準同形$f|_{\SU(2)}:\SU(2)\to\SO(3)$が定まる．
オイラー角を使うと$\SO(3)$は$z$軸の回転と$y$軸の回転から生成されることが分かるので$f|_{\SU(2)}$は全射である．
\[
   \xymatrix{
      \SU(2) \ar[d] \ar[r]^{f|_{ \SU(2)}}&\SO(3)\ar[d]&\\
      \SL(2,\mathbf C)\ar[r]^{f}&\PSL(2,\mathbf C) &\curvearrowright 
      \mathbf C\mathbf P^1=S^2(1)\\
   }
\]
$\SU(2)=\Sp(1)= S^3(1)$であり，$\Ker f|_{ \SU(2)}=\{\pm 1\}$なので
\[
	\SO(3) = S^3(1)/\{\pm 1\}=\mathbf R\mathbf P^3
\]
まで分かった．

最後に，同型$\mathbf H\to \mathbf C^2$は$0\in \mathbf H$を$(0,0)$に写すので射影
\[
   \Sp(1)\to \mathbf C\mathbf P^1,\quad
   \SO(3) \to \mathbf C\mathbf P^1
\]
を引き起こす．これらはファイバー束である．
\emptyline

固有ローレンツ群と$\SL(2,\mathbf C) $の関係や,
$\SU(2)$と$\SL(2,\mathbf C) $の表現論も物理との関係からも興味深く，スピンとDirac 作用素との関係も面白いので独習されることをお奨めしたい．


%\section{Coxetor グラフ}
% $\mathbf R^n$の有限個の(線形)鏡映から生成される有限群を
% 有限鏡映群という．
% 有限個の(線形でない合同変換の)鏡映が生成する離散群を
% アファイン鏡映群という．これらはCoxetor群の仲間であり，
% 万華鏡の数学的記述である．
% 
% これまでの議論で幾つか例が出てきたので，簡単に説明したい．
% 
% まずCoxetor グラフについて解説する．まず$\mathbf R^n$で
%$R=R_{\mathbf u}$を
% $\mathbf u\neq \mathbf 0$に直交する平面に付いての鏡映とする．
%$\alpha\neq0$について $R_{\alpha\mathbf u}=R_{\mathbf u}$である．$\{R_{\mathbf u_i}\}$が有限(あるいは離散)群$W$を生成するとする．
%
%$R$をルート系という．
%ポジティブルート，
%シンプルルート
%
% Coxetor グラフの分類：
% まず例を挙げる．$A_n$, $B_n$, $D_n$
% 
% $D_3$は$A_2$ ,  $T$は$A_3$である．
% 
% $D_4$は$B_2$ , $I^*$ は$B_3$, 
% 
%  結晶群の中のアファイン鏡映群
%\[
%   2\dim\quad
%   \dynkin [extended,Coxeter]A2, %\dynkin [extended,Coxeter]F3,
%    %\dynkin [extended,Coxeter]G2
%\]
% 
% $\Ot(2)$ の有限群の中の鏡映群
% $D_n$ は$I_n$ 
% $\Ot(3)$ の有限群の中の鏡映群
%\[
%   3\dim\quad
%   \dynkin [Coxeter]B3, \dynkin [Coxeter]A3, \dynkin [Coxeter]H3
%\]
%
%
%
% \begin{figure}[htb]
%   \centering
%   \begin{minipage}{0.4\columnwidth}
%      \centering
%      \includegraphics{~/Documents/GitHub/bEuclideanGeo/ch3/images/figurecoxetor.1}
%      \caption{ シンプルルート$\dynkin [Coxeter]B3$}\label{}
%   \end{minipage} 
%   \begin{minipage}{0.4\columnwidth}
%      \centering
%      \includegraphics{~/Documents/GitHub/bEuclideanGeo/ch3/images/figurecoxetor.2}
%      \caption{立方体の分割}\label{}
%   \end{minipage} 
%\end{figure}
%
% 
%\begin{figure}[htb]
%   \centering
%   \begin{minipage}{0.4\columnwidth}
%      \centering
%      \includegraphics{~/Documents/GitHub/bEuclideanGeo/ch3/images/figurecoxetor.3}
%      \caption{ シンプルルート$\dynkin [Coxeter]A3$}\label{}
%   \end{minipage} 
%   \begin{minipage}{0.4\columnwidth}
%      \centering
%      \includegraphics{~/Documents/GitHub/bEuclideanGeo/ch3/images/figurecoxetor.4}
%      \caption{立方体の分割}\label{}
%   \end{minipage} 
%\end{figure}
%\begin{figure}[h]
%   \centering
%   \begin{minipage}{0.4\columnwidth}
%      \centering
%      \includegraphics{~/Documents/GitHub/bEuclideanGeo/ch3/images/figurecoxetor.5}
%      \caption{ シンプルルート$\dynkin [Coxeter]H3$}\label{}
%   \end{minipage} 
%   \begin{minipage}{0.4\columnwidth}
%      \centering
%      \includegraphics{~/Documents/GitHub/bEuclideanGeo/ch3/images/figurecoxetor.6}
%      \caption{正12面体の分割の一部}\label{}
%   \end{minipage} 
%\end{figure}
%
%
% リー群の中のルート系
%
%
%\clearpage

\section{楕円曲線と双曲幾何 }
\subsection{一般格子のなす空間の幾何学}
以下の議論では追々明らかになるように平面$\mathbf R^2$を複素平面$\mathbf C$ と見なすのが自然であるのでそう見なすことにする．
$+$に関して$\mathbf C$は可換群である．
$\omega_1$, $\omega_2\in \mathbf C$が $\mathbf R$上一次
独立，つまり，$\IM(\frac{\omega_1}{\omega_2})\neq 0$とする．
\[
	\Gamma= \Gamma_{\omega_1,\omega_2}
	=\{
		m\omega_1+n\omega_2\mid m,n\in \mathbf Z
	\}
\]
が（一般）格子であり，$\mathbf C$の($\mathbf Z$加群として)の可換部分群である．よって剰余群（更に，商空間，商多様体，Lie 群として）
\[
	T=T_{\omega_1,\omega_2} 
	= \mathbf C/ \Gamma_{\omega_1,\omega_2}
\]
が定まり，トーラスになる．とくに$T$はコンパクトな可換群である．
図~\ref{GammaFD} に一つの基本領域を示した．

\begin{figure}[htb]
   \centering
   \begin{minipage}{0.4\columnwidth}
      \centering
      \includegraphics{~/Documents/GitHub/bEuclideanGeo/ch3/images/figureHyperbolic.1}
      \caption{$\Gamma_{\omega_1,\omega_2}$}\label{GammaFD}
   \end{minipage} 
   \begin{minipage}{0.4\columnwidth}
      \centering
      \includegraphics{~/Documents/GitHub/bEuclideanGeo/ch3/images/figureHyperbolic.2}
      \caption{トーラス}\label{Tours}
   \end{minipage} 
\end{figure}

例へば$\omega_1=1$, $\omega_2=i$ とし，
$\omega'_1=2+i$, $\omega'_2=3+2i$ とすると
\[
	\begin{cases}
		\omega'_1=2\omega_1+\omega_2\\
		\omega'_2=3\omega_1+2\omega_2
	\end{cases},
	\quad
	\begin{cases}
		\omega_1=2\omega'_1-\omega'_2\\
		\omega_2=-3\omega'_1+2\omega'_2
	\end{cases}
\]
なので，部分群として$\Gamma_{\omega_1,\omega_2}=\Gamma_{\omega'_1,\omega'_2}$となり，$T_{\omega_1,\omega_2}=T_{\omega'_1,\omega'_2}$
となる．そこで，

\begin{Q}
$\omega_1$, $\omega_2$と$\omega'_1$, $\omega'_2$とする
とき，
$
	T_{\omega_1,\omega_2}=T_{\omega'_1,\omega'_2}
$
となるための条件？
\end{Q}
\[
	T_{\omega_1,\omega_2}=T_{\omega'_1,\omega'_2}\quad
	\Leftrightarrow\quad
	\Gamma_{\omega_1,\omega_2}= 
	\Gamma_{\omega'_1,\omega'_2}
\]
なので，この条件は$\alpha$, $\beta$, $\gamma$, $\delta\in\mathbf Z$と$\alpha'$, $\beta'$, $\gamma'$, $\delta'\in\mathbf Z$が存在して
\begin{equation}\label{moduler}
	\begin{cases}
		\omega'_1=\alpha\omega_1+\beta\omega_2\\
		\omega'_2=\gamma\omega_1+\delta\omega_2
	\end{cases}, 
	\quad
	\begin{cases}
		\omega_1=\alpha'\omega'_1+\beta'\omega'_2\\
		\omega_2=\gamma'\omega'_1+\delta'\omega'_2
	\end{cases}. 
\end{equation}
行列で表すと
\[
	\begin{pmatrix}
		\omega'_1\\
		\omega'_2
	\end{pmatrix}
	=
	\begin{pmatrix}
		\alpha&\beta\\
		\gamma&\delta
	\end{pmatrix}
	\begin{pmatrix}
		\omega_1\\
		\omega_2
	\end{pmatrix}, 
	\quad
	\begin{pmatrix}
		\omega_1\\
		\omega_2
	\end{pmatrix}
	=
	\begin{pmatrix}
		\alpha'&\beta'\\
		\gamma'&\delta'
	\end{pmatrix}
	\begin{pmatrix}
		\omega'_1\\
		\omega'_2
	\end{pmatrix}
\]
で, $\omega_1$, $\omega_2$が $\mathbf R$上一次
独立なので
\[
	\begin{pmatrix}
		\alpha'&\beta'\\
		\gamma'&\delta'
	\end{pmatrix}
	\begin{pmatrix}
		\alpha&\beta\\
		\gamma&\delta
	\end{pmatrix}
	=
	\begin{pmatrix}
		1&0\\
		0&1
	\end{pmatrix}.
\]
行列式を取ると
\[
	(\alpha'\delta' -\beta'\gamma')
	(\alpha\delta -\beta\gamma)
	=1.
\]
$\alpha$, $\beta$, $\gamma$, $\delta$, $\alpha'$, $\beta'$, $\gamma'$, $\delta'\in\mathbf Z$なので
\begin{equation}\label{unimoduler}
	\alpha\delta -\beta\gamma
	=\pm 1
\end{equation}
がわかる．逆に，この条件がみたされるとすると
\[
	\begin{pmatrix}
		\alpha'&\beta'\\
		\gamma'&\delta'
	\end{pmatrix}
	=
	\begin{pmatrix}
		\alpha&\beta\\
		\gamma&\delta
	\end{pmatrix}^{-1}
	=
	\frac1{\pm 1}
	\begin{pmatrix}
		\delta&-\beta\\
		-\gamma&\alpha
	\end{pmatrix}
\]
となり，$\Gamma_{\omega_1,\omega_2}= 
	\Gamma_{\omega'_1,\omega'_2}
$が分かる．
このような変換を\term[]{ユニモジュラー変換}という．
ユニモジュラー変換は群をなすことも簡単に分かる．

\begin{Thm}
$\tau=\frac{\omega_1}{\omega_2}$が次をみたす基底$\omega_1$,$\omega_2$ が存在する．
\[
\begin{split}
	&(1) \quad \IM \tau >0\qquad
	(2) \quad -\frac12<\RE \tau \le \frac12\\
	&(3) \quad |\tau |\ge 1\qquad
	(4) \quad |\tau |=1 \text{ のとき } \RE\tau \ge 0 \qquad
\end{split}
\]
\end{Thm}
\begin{proof}
以前(第6回)行った平面の基本ベクトルの分類とほぼ同じ．
\end{proof}
\begin{figure}[htb]
   \centering
   \begin{minipage}{0.4\columnwidth}
      \centering
      \includegraphics{~/Documents/GitHub/bEuclideanGeo/ch3/images/figureHyperbolic.3}
      \caption{基本領域}\label{fundmentalDomain}
   \end{minipage} 
   \begin{minipage}{0.4\columnwidth}
      \centering
      \includegraphics{~/Documents/GitHub/bEuclideanGeo/ch3/images/figureHyperbolic.4}
      \caption{$\mathbf H^{+}$と曲線}\label{}
   \end{minipage} 
\end{figure}

定理の条件をみたす$\omega_1$, $\omega_2$を\term[]{基本ベクトル}という．

$\omega_1$, $\omega_2$と$\omega'_1$, $\omega'_2$ に
\eqref{moduler}の関係があると$\tau$と$\tau'=\frac{\omega'_1}{\omega'_2}$の間に
\begin{equation}
	\tau'
	=\frac{\alpha\omega_1+\beta\omega_2}
		{\gamma\omega_1+\delta\omega_2}
	=\frac{\alpha\tau+\beta}
		{\gamma\tau+\delta}
\end{equation}
という\term[]{一次(分数)変換}が表れる．
\[
\begin{split}
	\tau'-\bar\tau'
	&=
	\frac{\alpha\tau+\beta}{\gamma\tau+\delta}
	-
	\frac{\alpha\bar\tau+\beta}{\gamma\bar\tau+\delta}\\
	&=
	\frac{(\alpha\tau+\beta)(\gamma\bar\tau+\delta)-
	(\alpha\bar\tau+\beta)(\gamma\tau+\delta)}{|\gamma\bar\tau+\delta|^2}\\
	&=\frac{(\alpha\delta -\beta\gamma)(\tau-\bar\tau)}{|\gamma\bar\tau+\delta|^2}
\end{split}
\]
$\tau-\bar\tau=2i\IM \tau$などから
\[
	\IM \tau' =\frac{\pm \IM\tau}{|\gamma\bar\tau+\delta|^2}
\]
特に，$\alpha\delta -\beta\gamma=1$のときは
\begin{equation}\label{transformation}
	\IM \tau' =\frac{\IM\tau}{|\gamma\bar\tau+\delta|^2}
\end{equation}
となる．
\[
\begin{split}
	SL(2,\mathbf Z)
	&=
	\{
		\begin{pmatrix}
			\alpha&\beta\\
			\gamma&\delta
		\end{pmatrix}
		\mid
		\alpha, \beta,\gamma,\delta \in \mathbf Z,
		\alpha\delta -\beta\gamma=1
	\},\\
	SL(2,\mathbf R)
	&=
	\{
		\begin{pmatrix}
			a&b\\
			c&d
		\end{pmatrix}
		\mid
		a,b,c,d \in \mathbf R,
		ad-bc=1
	\}\\
\end{split}
\]
と置く．
$SL(2,\mathbf R)$を\term[]{二次元実特殊線形群}という．
これは古典群であり，
$SL(2,\mathbf Z)$は$SL(2,\mathbf R)$の(離散)部分群である．

$\mathbf H^{+}=\{z \in \mathbf C\mid \IM z  >0\}$と置き，\term[]{上半平面}という．$A\in SL(2,\mathbf R)$, $z\in \mathbf H^{+}$に対して
\begin{equation}\label{transformation1}
	\rho z  =\rho_{A}z = \frac{a z +b}{c z +d}, \quad
	\rho^{-1}_{A} z  = \frac{d z -b}{-c z +a}
\end{equation}
と置くと
\eqref{transformation}では係数が整数ということは使っていないので
同じ式が成り立つ．
$|c z +d|^2=0$ は$c=d=0$を意味し\eqref{unimoduler}に反するので$|c z +d|^2>0$であり
\[
	\IM \rho z  =\frac{ \IM z }{|c z +d|^2}>0
\]
となり，$\rho_{A} z \in \mathbf H^{+}$．同様に$\rho_{A}^{-1} z \in \mathbf H^{+}$で，これは$\rho_{A}$の逆変換なので，
$SL(2,\mathbf R)$は$\mathbf H^{+}$に作用することが分かる．
全ての$z\in \mathbf H^{+}$ に対して$z= \rho_{A} z$とすると
\[
\begin{split}
	&z= \rho_{A} z =\frac{a z +b}{c z +d} \quad\Rightarrow\quad
	cz^2 + (d-a)z  - b =0\\
	&\Rightarrow\quad
	c=b=0 , d=a
	\quad\Rightarrow\quad
	c=b=0 , d=a=\pm 1
\end{split}
\]
最後に$ad-bc=1$を使った．よって
\[
	\rho_{A} z=\rho_{B} z\quad\Rightarrow\quad
	A =\pm B
\]
であり，
\[
	PSL(2,\mathbf Z)
	=
	SL(2,\mathbf Z)/\pm 1,\quad
	PSL(2,\mathbf R)
	=
	SL(2,\mathbf R)/\pm 1
\]
と置くと，この作用は効果的\footnote{全ての点を止める作用は単位元のもののみと云うこと．因みに，自由な作用とはある点を止める作用は単位元のものに限ると云うことである．}で，
$PSL(2,\mathbf Z)$の$\mathbf H^{+}$への作用の基本領域が図~\ref{fundmentalDomain}である．



\subsection{Poincareモデルと非ユークリッド幾何}
$z = x + i y$と実変数に分解し
\[
	dz = dx + idy,\quad d\bar z = dx - idy
\]
の($d$は上の一次変換の$d$ではなく微分形式を表す$d$)ように計算できるとすると，対称テンソルとして
\[
\begin{split}
	dz d\bar z&=(dx + idy)( dx - idy)\\
	&=dx^2 -idxdy + idydx + dy^2
	=dx^2+dy^2
\end{split}
\]
$w= \frac{az+b}{cz+d}$を微分すると
\[
	dw = \frac{a(cz+d)-c(az+b)}{(cz+d)^2}dz
	=\frac{ad-bc}{(cz+d)^2}dz=\frac{1}{(cz+d)^2}dz
\]
なので
\[
	dw d\bar w
	=\frac{1}{|cz+d|^4}dz d\bar z
\]
\eqref{transformation}より
\[
	\frac{dw d\bar w}{|\IM w|^2}
	=\frac{1}{\frac{ |\IM z|^2}{|cz+d|^4}}
	\frac{1}{|cz+d|^4} dz d\bar z
	=\frac{dz d\bar z}{ |\IM z|^2}
\]
となるので$ds^2=\frac{dx^2+dy^2}{y^2}$は$SL(2,\mathbf R)$で不変な$\mathbf H^{+}$のリーマン計量である．
\emptyline

この計量に関する$\mathbf H^{+}$の測地線を求めたい．普通は測地線の方程式を解けば良いが，ここでは
図~\ref{fundmentalDomain} の境界に表れた半直線と半円が測地線であろう推測して以下のように議論する．
まず，$0$から出る虚軸を半直線$\ell$と見なす．
$0<y_1\le y_2$とおき，$c(t)=(x(t),y(t))$ ($t_1\le t\le t_2$)を$(0,y_1)$と
$(0,y_2)$を結ぶ任意の$\mathbf H^{+}$内の曲線とする．
\[
\begin{split}
	L(c)
	&=
	\int_{t_1}^{t_2} \frac {\sqrt{x'(t)^2 +y'(t)^2}}{|y(t)|}dt
	\ge \int_{t_1}^{t_2} \frac {|y'(t)|}{|y(t)|}dt\\
	&\ge \biggl|\int_{t_1}^{t_2} \frac {y'(t)}{y(t)}dt\biggr|
	=\bigl|\bigl[\log y(t) \bigr]_{t_1}^{t_2} \bigr|
	=\bigl|\log \frac{ y_2}{ y_1}\bigr|
\end{split}
\]
実際，$y(t)$を単調な関数に取り替えた$\ell$上の道をとるとその値が得られる．
つまり
\[
	d(iy_1,iy_2)=\bigl|\log \frac{ y_2}{ y_1}\bigr|
\]
測地線は$\ell$上にあり，
\[
	d(iy_1,iy_2)=d(iy_2,iy_1),\quad
	d(iy_1,iy_2)=0 \Leftrightarrow y_1=y_2
\]
であること，
$y_1\to +0$や$y_2\to \infty$のとき距離が無限に近づくことが分かる．

$\ell$を一次分数変換\eqref{transformation1}で写すと
%$\tilde\ell$を$\ell$を含む直線とすると
\[
\begin{split}
	&z_1=\rho(iy_1)\in \mathbf H^{+},\quad z_2=\rho(iy_2)\in \mathbf H^{+},\\
	&x_1=\rho (0)=\frac{b}{d}\in \mathbf R,\quad
	x_2=\lim_{t\to \infty}\rho (it)=\frac{a}{c}\in \mathbf R
\end{split}
\]
であり，
$
	x_0=\frac12(x_1+x_2)
$と置くと，一次変換で直線は円に写り，等角なので，測地線は実軸に直交する$x_0$を中心とする半円になる．これを計算で確かめてみる：
\[
\begin{split}
	&x_1-x_0
	=\frac{b}{d}-\frac12(\frac ac+\frac bd)
	=\frac12(\frac bd-\frac ac)
	= \frac{-1}{2cd},\\
	&x_2-x_0
	=\frac{a}{c}-\frac12(\frac ac+\frac bd)
	=\frac12(\frac ac-\frac bd)
	= \frac1{2cd}.
\end{split}
\]
$\ell$上の点は$t\mapsto ti $と書けるので
\[
\begin{split}
	&\frac{ati +b}{cti +d} -x_0
	=\frac{ati +b}{cti +d} -\frac12(\frac ac+\frac bd)\\
	&=
	\frac{2cd(ati +b)-(ad+bc)(cti +d)}{2cd(cti +d)}\\
	&=
	\frac{2cd(ati +b)-(1+2bc)(cti +d)}{2cd(cti +d)}\\
	&=
	\frac{2cdati -(1+2bc)cti -d}{2cd(cti +d)}
	=
	\frac{cti - d}{2cd(cti +d)}.
\end{split}
\]
よって
\[
	|g(ti)-x_0| = \tfrac1{2cd}
\]
となり，$x_0$を中心とした半円であることが確かめられた．
\emptyline

距離$d$を$z_1$, $z_2$, $x_1$, $x_2$で表したいが，そのために前回紹介した
\term[]{複比}
\begin{equation}\label{crossRatio}
	(z_1,z_2,z_3,z_4)
	=
	\frac{(z_1-z_3)(z_2-z_4)}{(z_2-z_3)(z_1-z_4)}
\end{equation}
($z_1$, $z_2$, $z_3$, $z_4\in 
\hat{\mathbf C} =\mathbf C\mathbf  P^1$)
を思い出そう．複比の大事な性質は任意の一次分数変換$\rho z= \frac{az+b}{cz+d}$ ($a,b,c,d\in \mathbf C$, $ad-bc\neq0$) で不変なこと，つまり，
\[
	(\rho z_1,\rho z_2,\rho z_3,\rho z_4)
	=
	(z_1,z_2,z_3,z_4)
\]
をみたすことであった．すると，
\[
\begin{split}
	(z_1,z_2,x_1,x_2)
	&=
	(\rho ^{-1}z_1,\rho ^{-1}z_2,\rho ^{-1}x_1,\rho ^{-1}x_2)\\
	&=	
	(y_1,y_2,0,\infty)
	=
	\frac{y_1}{y_2}>0\quad	
\end{split}
\]
したがって，距離$d(z_1,z_2)$ は($SL(2,\mathbf R)$不変であり)
\[
	d(z_1,z_2) = \bigl|\log (z_1,z_2,x_1,x_2)\bigr|
\]
で定まる．


\begin{figure}[htb]
   \centering
   \begin{minipage}{0.4\columnwidth}
      \centering
      \includegraphics{~/Documents/GitHub/bEuclideanGeo/ch3/images/figureHyperbolic.5}
      \caption{ $\mathbf H^{+}$の測地線}\label{}
   \end{minipage} 
   \begin{minipage}{0.4\columnwidth}
      \centering
      \includegraphics{~/Documents/GitHub/bEuclideanGeo/ch3/images/figureHyperbolic.6}
      \caption{$\mathbf H^{+}$での直線の様子}\label{failure5}
   \end{minipage} 
\end{figure}

逆に$z_1\neq z_2\in \mathbf H^{+}$に対し，($z_1$と$z_2$が虚軸と平行な半直線上にないとき)線分$z_1 z_2$の二等分線と実軸との交わりを中心とし，$z_1$と$z_2$を通る半円を書けば
無限遠点$x_1$, $x_2$が見つかる．
\eqref{crossRatio} で$z_1$を$z$に取り直すことで$\rho^{-1}$がもとまり，
任意の二点を結ぶ最短線が一意的に定まることが分かる．
$z_1$と$z_2$が虚軸と平行な半直線上にあるとき，(その実軸との交わりと$\infty$を無限遠点とすれば)最初の公式で距離が定まる．
$\mathbf H^{+}$への作用は推移的であることも分かる．

三角不等式が成り立つことも，円が存在することも分かる．
$(\mathbf H^{+},ds^2)$を(ベルトラミモデルの方が正しいらしい\cite{S}, \cite{S2} が）\term[]{ポアンカレモデル}という．
以上から$(\mathbf H^{+},ds^2)$では
\begin{itemize}
   \item 任意の二点を結ぶ線分が存在して一意的．
   \item 任意の線分はどの方向にもどこまでも伸ばすことが出来る．
   \item 任意の点を中心にして別の任意の点を通る円が存在して一意的．
   \item 任意の直角は全て等しい．
\end{itemize}
という，ユークリッド幾何の四つの公準をみたすことが分かるが，
\emptyline

\emph{平行線の公準}：ある二つの直線が直線と交わり，内角の和が2直角より小さいならば二つの直線は必ず交わる．

\emptyline
\noindent
をみた\textcolor{red}{さない}ことも分かる(図~\ref{failure5})．つまり，非ユークリッド幾何のモデルとなっている．
これはユークリッド幾何に矛盾がなければ非ユークリッド幾何にも矛盾がないことを意味している．
また，$\mathbf H^{+}$は単連結な負の定曲率空間でもあり，\term[]{双曲空間}と呼ばれることもある．

%更に，$PSL(2,\mathbf Z)$は$PSL(2,\mathbf R)$の離散部分群であり，
%トーラスの同形類の空間が商空間
%$\mathbf H^{+}/PSL(2,\mathbf Z)$である(図~\ref{fundmentalDomain})．

\begin{Quiz}
$\mathbf H^{+}$を円板$D=\{z\in \mathbf C\mid |z|<1\}$に写す一次変換を求め，$(\mathbf H^{+},ds^2)$と等長となるような$D$上のリーマン計量を求めよ．
\end{Quiz}

%\begin{thebibliography}{99}
%   \bibitem{A} Ahlfors, Complex analysis 3rd edition, McGraw-Hill, 1979.
%   
%   \bibitem{K} 小林昭七, ユークリッド幾何から現代幾何へ, 日本評論社, 1990.
%   
%   \bibitem{S1}J.  Stillwell, Mathematics and Its History  3rd ed, UTM, Springer, 2010.
%   
%   \bibitem{S2}J. Stillwell, 
%      Sources of Hyperbolic Geometry, History of math.10, AMS. 1996.
%\end{thebibliography}
%

\subsection{楕円関数}
$\omega_1$, $\omega_2\in \mathbf C$が $\mathbf R$上一次
独立，つまり，$\IM(\frac{\omega_1}{\omega_2})\neq 0$とする．
\[
	\Gamma= \Gamma_{\omega_1,\omega_2}
	=\{
		m\omega_1+n\omega_2\mid m,n\in \mathbf Z
	\}
\]
とすると，剰余群%（更に，商空間，商多様体，Lie 群として）
\[
	T=T_{\omega_1,\omega_2} 
	= \mathbf C/ \Gamma_{\omega_1,\omega_2}
\]
が定まり，トーラスになる．$a\in \mathbf C$にたいし
\[
	D=D_a =
	\{
		a+s\,\omega_1+t\,\omega_2\in \mathbf C
		\mid
		0 <s,t<1
	\}
\]
と置き，\term[]{基本周期平行四辺形}という
%（これは$\Gamma$ の Dirichlet 領域とは異なるが，これらは互いに双対になっている）
．これを使うと
$T$が一次元複素多様体であることが直ぐに分かる．

\begin{figure}[htb]
   \centering
   \begin{minipage}{0.4\columnwidth}
      \centering
      \includegraphics{~/Documents/GitHub/bEuclideanGeo/ch3/images/figureHyperbolic2.1}
      \caption{$D_a$}\label{Gamma}
   \end{minipage} 
   \begin{minipage}{0.4\columnwidth}
      \centering
      \includegraphics{~/Documents/GitHub/bEuclideanGeo/ch3/images/figureHyperbolic2.2}
      \caption{トーラス}\label{Tours}
   \end{minipage} 
\end{figure}


簡単に関数論の復習：$\mathbf C$の領域$\Omega$
の有理型関数$f:\Omega\to \mathbf C$
の$a\in \Omega$が$f$の極であるとは
\[
	f(z) = \frac{A_{-m}}{(z-a)^m}+\cdots +\frac{A_{-1}}{z-a}
	+ A_0+A_1(z-a) +\cdots 
\]
とローラン展開出来ることであり，$A_{-1}$をその留数といい$\Res(f,a)$と表した．有理型とは$f$の特異点が極のみでそれらが集積していないことをいう．
$\mathbf C$の領域$\Omega$が有界で，
その境界のなす閉路$C$が区分的滑らかとし，$\Omega$を左に見て進む向きをとる．
$f$の$\Omega$内の極を$a_1, \cdots, a_m$とし，
$C$は極を通らないとするとき，留数定理は
\[
	\frac1{2\pi i}\int_C f(z)dz = \sum_{i=1}^m\Res(f,a_i)
\]
を意味した．

$f:\mathbf C\to\mathbf C $が有理型関数で
2重周期, つまり，
$\mathbf R$上の基$\omega_1$, $\omega_2\in \mathbf C$が存在して
\begin{equation}\label{period}
	f(z+\omega_1)=f(z+\omega_2)=f(z)
\end{equation}
であるのとき\term[]{楕円関数}という．楕円関数は
$T_{\omega_1,\omega_2} $上の有理型関数を定める．
以下，$f$は楕円関数とし，$\bar D$上の関数と見たり，$T$上の関数と見たりするなど問題に応じ便利な見方をとる．

\begin{Prop}\label{const}
楕円関数が極を持たなければ定数関数になる．
\end{Prop}
\begin{proof}
極がなければ$\mathbf C$上の有界な正則関数となり，リュービルの定理より定数関数になる．
\end{proof}


\begin{Prop}\label{sumofRes}
楕円関数の留数の和は$0$である．
\end{Prop}
\begin{proof}
$f$を楕円関数とする．
適当にずらすことで$D$の境界上に極がないとして良い．
図~\ref{path}のように名付けると
\[
	\partial D=\gamma_1+\gamma_2+\gamma_3+\gamma_4,\quad
	\gamma_1=-\gamma_3,\quad \gamma_2=-\gamma_4
\]
これらの関係と留数定理より
\[
\begin{split}
	&2\pi i\sum_{i=1}^m\Res(f,a_i)
	=\int_{\partial D} f(z)dz \\
	&= \int_{\gamma_1} f(z)dz
	-\int_{\gamma_1} f(z)dz
	+\int_{\gamma_2} f(z)dz-\int_{\gamma_2}f(z)dz=0
\end{split}
\]

\end{proof}

\begin{Prop}
楕円関数が一位の極を一つしかもたないと定数関数になる．
\end{Prop}
\begin{proof}
背理法．
$a\in D$を$f$の唯一の一位の極とすると
\[
	f(z)=\frac{A_{-1}}{z-a}
	+ g(z)\quad g(z)\text{：正則}
\]
$\Res(f,a)=A_{-1}$で上の命題\ref{sumofRes}より$A_{-1}=0$.
つまり，$f(z)=g(z)$で正則となり$a$は極でなくなり矛盾．
\end{proof}


\begin{figure}[htb]
   \centering
   \begin{minipage}{0.4\columnwidth}
      \centering
      \includegraphics{~/Documents/GitHub/bEuclideanGeo/ch3/images/figureHyperbolic2.3}
      \caption{積分路}\label{path}
   \end{minipage} 
   \begin{minipage}{0.4\columnwidth}
      \centering
      \includegraphics{~/Documents/GitHub/bEuclideanGeo/ch3/images/figureHyperbolic2.4}
      \caption{$\wp'$の零点}\label{zero}
   \end{minipage} 
\end{figure}

\begin{Prop}\label{arg}
定数でない楕円関数は，重複度込みで同数の極と零点を持つ．
\end{Prop}
\begin{proof}
$f(z)$を楕円関数とすると，$f'(z)$も$f'(z)/f(z)$も同じ周期の楕円関数で
\[
	\Res(f'(z)/f(z),a)
	=
	\begin{cases}
		\text{零点の位数}& >0\\
		-1\times\text{極の位数}& <0\\
	\end{cases}
\]
∵）$f(z)=(z-a)^kg(z)$, $g(a)\neq0$とすると
\[
	f'(z)=k(z-a)^{k-1}g(z)+(z-a)^kg'(z)
\]
なので
\[
	\frac{f'(z)}{f(z)}=\frac{k}{z-a}+\frac{g'(z)}{g(z)}
\]
\emptyline

$N$を重複度を含めた零点の個数，$P$重複度を含めた極の個数とすると，命題~\ref{sumofRes}より
\[
	0 = 
	\frac1{2\pi i}\int_C \frac{f'(z)}{f(z)}dz = N- P\quad\text{(偏極の原理)}
\]
\end{proof}

\begin{Prop}\label{numSol}
任意の$c\in \mathbf C$に対して$f(z)=c$となる$z\in \bar D$の個数は一定．
\end{Prop}
\begin{proof}
$f(z)$を楕円関数とすると，$f(z)-c$も同じ極を持つ楕円関数．
命題~\ref{arg}
$f(z)-c$の零点の個数はその極の個数と一致する．
\end{proof}

$N$個の零点（極）を持つ楕円関数を\term[]{$N$位}の楕円関数という．

$f(z)$を2次の極をただ一つだけ持つ楕円関数とする．適当に平行移動して$z=0$がその極とする．定数倍することで
\[
	f(z) = \frac{1}{z^2}+\frac{A_{-1}}{z}
	+ g(z) \quad g(z)\text{：正則}
\]
上の命題~\ref{sumofRes} より$A_{-1}=0$.

\begin{Prop}\label{even}
この$f(z)$は偶関数．
\end{Prop}
\begin{proof}
$f(z)-f(-z)$も楕円関数だが，
\[
	f(z)-f(-z) = \frac{1}{z^2} +g(z)-  \frac{1}{(-z)^2} -g(-z)
	=g(z)-g(-z)
\]
となり，命題~\ref{const}よりこれは定数関数．つまり，ある$C$が存在して
$f(z)-f(-z)=C$．
ところが
\[
	C = f(\omega_i/2)-f(-\omega_i/2) 
	= f(\omega_i/2)-f(\omega_i/2) =0
\]
なので偶関数．
\end{proof}

偶関数の奇数次の巾は$0$であり，定数を足しても楕円関数であることは変わらないので定数項を引いて
\begin{equation}\label{wp}
	\wp(z)
	=
	 \frac{1}{z^2}
	+ a_1z^2 + a_2z^4 + \cdots 
\end{equation}
を考えればいい(ワイエルシュトラウスの\term[]{ペー関数}と呼ばれる)．
\[
\begin{split}
	\wp'(z)
	&=
	 \frac{-2}{z^3}
	+ 2a_1z + 4a_2z^3 + \cdots \\
	\wp'(z)^2
	&=
	\frac{4}{z^6}-8\frac{a_1}{z^2}-16a_2+\cdots \\
	\wp(z)
	&=
	 \frac{1}{z^2}
	+ a_1z^2 + a_2z^4 + \cdots \\
	\wp(z)^2
	&=
	 \frac{1}{z^4}
	+ 2a_1 + 2a_2z^2 + \cdots \\
	\wp(z)^3
	&=
	 \frac{1}{z^6}
	+ \frac{3a_1}{z^2} + 3a_2 + \cdots \\
	4\wp(z)^3
	&=
	 \frac{4}{z^6}
	+ \frac{12a_1}{z^2} + 12a_2 + \cdots \\
	20\wp(z)
	&=
	 \frac{20}{z^2}
	+ 20a_1z^2 + 20a_2z^4 + \cdots \\
\end{split}
\]
なので
\[
	\wp'(z)^2-4\wp(z)^3+20a_1\wp(z)
	=-28a_2 + h(z)
	=-28a_2
\]
最後に，$h(z)$は定数部分を含まない正則な楕円関数となり，命題~\ref{const}より $0$となることを使った．
$g_2=20a_1$, $g_3=28a_2$と置くと
\begin{equation}\label{dEqWp}
	\wp'(z)^2=4\wp(z)^3-g_2 \wp(z)-g_3
\end{equation}

$e_1$, $e_2$, $e_3$を$4w^3-g_2 w-g_3$の重複度を込めた根とする．
命題~\ref{even}と\eqref{period}より
\[
	\wp(z)=\wp(-z)=\wp(\omega -z)\quad \omega\in\Gamma
\]
なので
\[
	\wp'(z)=-\wp'(\omega -z)
\]
$z=\omega/2$を代入すると
\[
	\wp'(\omega/2)=-\wp'(\omega -\omega/2)=-\wp'(\omega/2)
\]
なので$\wp'(\omega/2)=0$. よって
\begin{equation}\label{zeroWp'}
	\wp'(\omega_1/2)=0,\quad
	\wp'(\omega_2/2)=0,\quad
	\wp'((\omega_1+\omega_2)/2)=0
\end{equation}
$\wp'(z)$は3位の楕円関数で
\[
	\omega_1/2\neq \omega_2/2
	\neq (\omega_1+\omega_2)/2\neq\omega_1/2
\]
は位数$1$の$\wp'$の零点である(図~\ref{zero})．

\begin{Prop}\label{}
\[
	\wp(\omega_1/2)=e_1,\quad
	\wp(\omega_2/2)=e_2,\quad
	\wp((\omega_1+\omega_2)/2)=e_3
\]
と置くと，
$e_1\neq e_2\neq e_3\neq e_1$となるように取れる．特に$4w^3-g_2 w-g_3$は重根を持たない，言い換えると，\eqref{dEqWp}の右辺の$3$次式の
判別式は $g_2^3 -27g_3^2\neq 0$．
\end{Prop}
\begin{proof}\eqref{dEqWp}を書き直すと
\[
	\wp'(z)^2=4(\wp(z)-e_1)(\wp(z)-e_2)(\wp(z)-e_3)
\]
\eqref{zeroWp'}なので$\wp(\omega_1/2)=e_1$として良い．
\[
	\wp'(z)=(z - \omega_1/2)f(z)\quad f( \omega_1/2)\neq 0
\]

\begin{equation}\label{branchWp}
	\wp(z)-e_1=(z - \omega_1/2)^2g(z)\quad g( \omega_1/2)\neq 0
\end{equation}
∵）$\wp(z)-e_1$は$2$位の楕円関数なのでこれが正しくなければ
\[
	\wp(z)-e_1=(z - \omega_1/2)k(z)\quad k( \omega_1/2)\neq 0
\]
だが
\[
\begin{split}
	&\wp'(z)=k(z)+(z - \omega_1/2)k'(z)\\
	\Rightarrow\quad
	&\wp'(\omega_1/2)=k(\omega_1/2)+0\neq 0
	\quad\text{：\eqref{zeroWp'}に矛盾}
\end{split}
\]

$\wp(\omega_1/2)=e_1=e_2=\wp(\omega_2/2)$とすると$\omega_1\neq\omega_2$より
\[
	\wp(z)-e_1=(z - \omega_1/2)^2(z - \omega_2/2)^2h(z)\quad
	 g( \omega_1/2)\neq 0
\]
$\wp(z)-e_1$は位数$4$の零点となり矛盾．$e_2\neq e_3\neq e_1$も同様．
\end{proof}


%\subsection{楕円曲線・双曲幾何に関連した話題}
\emptyline
正則写像$\Phi:T\to\mathbf CP^2$が存在して，$T$は代数多様体になる：
%(その像が$\mathbf CP^2$の射影代数多様体)
\[
	\Phi(z)=[1:\wp(z):\wp'(z)]\in\mathbf CP^2
\]
と置くと(特異でない，つまり，ヤコビ行列が常にフルランクの)正則写像であり，$T$の極$z=0$の近傍では
\[
	[z^3:z^3\wp(z):z^3\wp'(z)]
	=[z^3:z+a_1z^5+\cdots:-2+2a_1z^4+\cdots]
%	\to [0:0:1]=O\in \mathbf CP^2
\]
$z=0$で$O=[0:0:1]$であり，$U_2=\{[x_0:x_1:x_2]\mid x_2\neq 0\}$ で考えると
$z=0$でも正則であることも分かる．$T$は
\[
	\{[x_0:x_1:x_2]\in \mathbf CP^2\mid
		x_0x_2^2-4x_1^3+g_2 x_0^2x_1+g_3x_0^3=0
	\}
\]
と同一視される代数多様体である．$T$は$O$を零元とするアーベル群となる．
これは$z$座標では明らかだが，非同次座標$(\wp(z),\wp'(z))$では図~\ref{addition}, ~\ref{minus}のように定まる演算
\[
	P + Q + R =O, \quad P + P'=O
\]
の群であり\term[]{楕円曲線}と呼ばれる．

\begin{figure}[htb]
   \centering
   \begin{minipage}{0.4\columnwidth}
      \centering
      \includegraphics[width=5cm]{~/Documents/GitHub/bEuclideanGeo/ch3/images/impliciteFun1plot2d4.png}
      \caption{虚軸に平行でない直線との交点$P + Q + R =O$}\label{addition}
   \end{minipage} 
   \begin{minipage}{0.4\columnwidth}
      \centering
      \includegraphics[width=5cm]{~/Documents/GitHub/bEuclideanGeo/ch3/images/impliciteFun1plot2d5.png}
      \caption{虚軸に平行な直線との交点$P + P'=O$}\label{minus}
   \end{minipage} 
\end{figure}
\emptyline

$\wp(z)$には明示的な表示がある：
\[
	\wp(z)=
	\frac1{z^2}+ 
	\sum_{\omega\neq 0\in \Gamma}
	\biggl(
		\frac{1}{(z-\omega)^2}-\frac{1}{\omega^2}
	\biggr)
\]
ただし，$\omega\neq 0 \in \Gamma$についての総和．
\[
	G_k=\sum_{\omega\neq 0\in\Gamma}\frac{1}{\omega^{2k}}
	\quad\text{：アイゼンシュタイン級数}
\]
と置くと，\eqref{dEqWp}の係数は
$g_2=60G_2$, $g_3=140G_3$と表される．また，
\[
	\wp(z)=
	\frac1{z^2}+ 
	\sum_{k=2}^\infty (2k-1)G_kz^{2k-2}
\]


%\emptyline

\begin{figure}[htb]
	\centering
	\includegraphics{~/Documents/GitHub/bEuclideanGeo/ch3/images/figureHyperbolic2.7}
	\caption{$P:T\to \mathbf CP^1\simeq S^2$}\label{}
\end{figure}
$P:T\to \mathbf CP^1\simeq\mathbf C\cup \infty
	\simeq S^2$を
\[
	P(z)=[1:\wp(z)]\in\mathbf CP^1
\]
で定めると，命題~\ref{numSol}より$w\neq e_1$, $e_2$, $e_3\in \mathbf C$とすると($4w^3-g_2 w-g_3\neq 0$のとき\eqref{dEqWp}の$\wp'$の二つの根は異なるので)二つの根は異なる．
\eqref{branchWp}より$e_1$, $e_2$, $e_3\in \mathbf C$ で$\wp'^2=0$の重根になり二重に分岐している．\eqref{dEqWp}は$z=0$の近傍で
\[
	[z^2:z^2\wp(z)]= [z^2:1+ a_1z^4 + a_2z^6 + \cdots]
\]
となり，$z=0$のとき$\wp(z)=[0:1]=\infty$ なので$\infty$で二重に分岐している．つまり，$T$は$S^2$の四点で二重に分岐する被覆となっている．
\emptyline

その代数的な解釈は以下のようである：
$2$変数多項式環$ \mathbf C[X,Y]$の単項イデアル$(4X^3-g_2 X-g_3-Y^2)$は素イデアルで，それによる剰余環
\[
	  \mathbf C[X,Y]/(4X^3-g_2 X-g_3-Y^2)
\]
は整域となるので，その商体$K$を考える．
$1$変数多項式環$ \mathbf C[X]$の商体$\mathbf C(X)$の埋め込み
\[
	  \mathbf C(X)\ni h(X)\mapsto h(X)\in K
\]
は$X^3-g_2 X-g_3-Y^2$による代数拡大(有限次拡大)となり，
体$\mathbf C$の超越次数$1$の拡大体である．これが$P:T\to \mathbf CP^1$の体による解釈である．
\emptyline

上で図~\ref{addition}, ~\ref{minus}のように定まる演算のことをコメントしたが，$\mathbf C$を $\mathbf Q$や有限体$\mathbf F_{q}$($q=p^n$)などに変えても類似のことが成り立つ． 
$\mathbf Q$や有限体の場合は（モーデルの定理やヴェイユ予想など）数論との関係で重要であり，有限体の楕円曲線が楕円曲線暗号の基礎となっている．
\emptyline


\eqref{dEqWp}を
\[
	\frac{dw}{dz}=\sqrt{4w^3-g_2 w-g_3}
\]
と書くと微分方程式であり，変数分離で解くと
\[
	z-z_0=\int^{w}_{w_0} \frac{dw}{\sqrt{4w^3-g_2 w-g_3}}
\]
で分母の平方根の中身$4w^3-g_2 w-g_3$は重根を持たないので，
これは楕円積分である．つまり，楕円積分の逆関数として楕円関数が定まることが分かる．
\begin{figure}[htb]
	\centering
	\includegraphics{~/Documents/GitHub/bEuclideanGeo/ch3/images/figureHyperbolic2.6}
	\caption{lemniscate}\label{lemniscate}
\end{figure}


歴史を振り返ってみる：ガウスは
$\arcSin x=\int_0^x\frac{1}{\sqrt{1-x^{2}}} dx$との類似から，三角関数の拡張を探していた．
$(r,\theta)$を極座標とし，曲線$C$が$r=f(\theta)$で表されるとき，$C$  の$\alpha\le \theta \le \beta$での長さは
\[
	\ell(C)=\int_\alpha^\beta \sqrt{f(\theta)^2+f'(\theta)^2}d\theta
\]
で表されることは知っていると思う．
レムニスケート( 図 \ref{lemniscate} )は%lemniscate
\[
	r^2 = \cos 2\theta
\]
と表せるので，
\[
\begin{split}
	2r\frac{dr}{d\theta}=-2 \sin 2\theta
	\quad&\Rightarrow\quad
	 r^2\biggl(\frac{dr}{d\theta}\biggr)^2 = \sin^2 2\theta
	 \\&\Rightarrow\quad
	 \biggl(\frac{dr}{d\theta}\biggr)^2=\frac{\sin^2 2\theta}{ \cos 2\theta}
\end{split}
\]
である．
 $\theta$から$\pi/4$までの弧の長さを計算し，半径が$r$から$0$になるように変数変換すると
\[
	\int_{\theta}^{\pi/4}\frac{1}{\sqrt{\cos 2\theta}}d\theta
	=
	\int_r^0 \frac{1}{\sqrt{\cos 2\theta}}\frac{-rdr}{\sin 2\theta}
	=
	\int_0^r\frac{dr}{\sin 2\theta}
	=
	\int_0^{r}\frac{dr}{\sqrt{1-r^4}}
\]
となる．$r^4-1$は重根が無いのでこれは楕円積分となる．1797年にガウスはこの積分(関数)の逆関数として楕円関数を発見した．
\emptyline

楕円関数の関数論に興味のわいたひとは\cite{A},  \cite{Si}などを参考にしてほしい．
代数幾何との関連は\cite{F}，数論や保型形式に興味のある人には\cite{Se}, \cite{KKS}, 暗号理論に興味のある人には\cite{Ko}, \cite{ST}, 
数学史に興味のある人には\cite{T}, \cite{Kl}などをお勧めする．


\subsection{Modular 関数}%とKlien群
$\tau=\frac{\omega_1}{\omega_2}$に対して$\lambda$\term[]{不変量}を
\[
	\lambda(\tau)
	=
	\frac{e_3-e_2}{e_1-e_2}
\]
で定める(と$\mod 2$の合同部分群の不変量である．下の問題参照)．
$J$\term[]{不変量}を
\[
	J(\tau)
	=
	\frac{4}{27}
	\frac{(1-\lambda+\lambda^2)^3}{\lambda^2(1-\lambda)^2}%\quad
%	\biggl(=\frac{12^3g_2^2}{g_2^3 -27g_3^2}\biggr)
\]
と置くと\term[]{モジュラー関数}(保型形式)の仲間が表れる．
\begin{Quiz}
$\SL(2,\mathbf Z) ( \PSL(2,\mathbf Z)) $の部分集合
\[
	\begin{pmatrix}
		\alpha&\beta\\
		\gamma&\delta
	\end{pmatrix}
	\equiv
	\begin{pmatrix}
		1&0\\
		0&1
	\end{pmatrix}
	\quad\mod 2
\]
が部分群であることを示せ，ただし，これは各成分毎に
$\alpha\equiv1\mod 2$, $\beta\equiv0\mod 2$などを意味している．これが$\mod 2$の合同部分群である．
この群の元の作る一次変換
$\frac{\alpha\tau+\beta}{\gamma\tau+\delta}$で$\lambda(\tau)$が不変なこと，
つまり，
\[
	\lambda\biggl(\frac{\alpha\tau+\beta}
		{\gamma\tau+\delta}\biggr)
	=
	\lambda(\tau)
\]
となることを示せ．
\end{Quiz}

%\begin{figure}[htb]
%	\centering
%	\includegraphics{~/Documents/GitHub/bEuclideanGeo/ch3/images/figureHyperbolic21.5}
%	\caption{}\label{}
%\end{figure}


%\begin{thebibliography}{99}
%   \bibitem{A} Ahlfors, Complex analysis 3rd edition, McGraw-Hill, 1979
%   
%   \bibitem{S} Siegel, Topics in complex function theory I, II, III, 
%   Wiley-Intersicence, 1969
%   
%   \bibitem{F}
%   W.~Fulton, 
%   ALGEBRAIC CURVES---An Introduction to Algebraic Geometry (3rd edition), 2008, 
%   \href{https://dept.math.lsa.umich.edu/~wfulton/CurveBook.pdf}{ダウンロード}
%   
%   \bibitem{Se}
%   セール, 数論講義, 岩波書店, 1979
%   
%   \bibitem{KKS}
%   加藤和也, 黒川信重, 斎藤毅, 数論 I Fermatの夢と類体論, 岩波書店, 2016
%   
%   黒川 信重, 栗原 将人, 斎藤 毅, 数論II 岩澤理論と保型形式, 岩波書店,  2017
%   
%   \bibitem{Ko} Koblitz, A course in Number theory and Cryptography, GTM,  \href{https://link.springer.com/search?package=11649&facet-content-type=%22Book%22}{ダウンロード}
%   
%   
%   \bibitem{ST} Silverman, Tate,
%   Rational Points on Elliptic Curves, Springer, 2015 
%   \href{https://link.springer.com/search?package=11649&facet-content-type=%22Book%22}{ダウンロード}
%   
%   \bibitem{T}高木貞治, 復刻版 近世数学史談・数学雑談, 共立出版
%   
%   \bibitem{Kl}F.~Kline, クライン:19世紀の数学, 彌永 昌吉監修, 共立出版, 1995
%   
%\end{thebibliography}




\appendix


\chapter{復習と補足}
\section{群論の復習と補足}
\subsection{群論や代数の復習}
授業で取り上げるのは幾何学的な群論であるが，群論の一般論から分かることも多いのでここに簡単に群論や代数の初等的な概念や結果を纏めておくので参考にしてほしい．

%$H\subset G$が\term[]{部分群}とは$H$が$G$の積に関して群となることである．これは
%\[
%     a,b\in H\Rightarrow ab\in H\quad\text{かつ}\quad 
%     a\in H\Rightarrow a^{-1}\in H
%\]
%を満たすことと同値である．

群$G$, $H$ について写像$f:G\to H$が
\[
     f(ab)=f(a)f(b)\quad \forall a,b\in G
\]
をみたすとき，$f$を\term[]{準同形}写像という．
$e_G$を$G$の単位元とすると
\[
     f(e_G)=e_H, \quad f(a^{-1})=f(a)^{-1} 
\]
が成り立つ．
特に準同形$f:G\to H$が全単射である
（自然に逆写像も準同形になる）とき
\term[]{同形}という．
準同形$f:G\to H$に対して
\[
     \Image f =\{f(a)\in H\mid a\in G\},\quad
      \Ker f=\{a\in G\mid f(a)=e\}
 \]
 と置くと，$\Image f$は$H$の部分群，$\Ker f$も$G$の部分群である．

群$G$の部分群$H$が与えられたとき，$g\in G$に対して
\[
     gH=\{gh\mid h\in H\},\quad Hg=\{hg,\mid h\in H\}
     \subset G
\]
と置き
\[
     G/H = \{gH\mid g\in G\}
\]
とする．$G$が有限群の時
\[
     |G : H|= |G/H|
\]
と書き$H$の$G$における\term[]{指数}という．
\begin{Thm}
   有限群$G$の部分群$H$に対して
   \[
        |G| =|G : H| |H|
   \]
\end{Thm}

一般には$gH\neq Hg$であるが，
\[
     \text{任意の}\quad g\in G\quad\text{に対して}\quad gH= Hg
\]
が成り立つとき$H$は$G$の\term[]{正規部分群}という．これは
\[
     \text{任意の}\quad g\in G\quad\text{に対して}\quad gHg^{-1}= H
\]
と同値である．$H\lhd G$と書くことがある．
$N$が$G$の正規部分群とするとき
\[
     G/N = \{gN\mid g\in G\}
\]
と置く．正規部分群の定義より
\[
     gNhN
     =g(hh^{-1})NhN
     =gh(h^{-1}Nh)N=gh NN =gh N
\]
特に
\[
     g^{-1}NgN=g^{-1}gN=N
\]
となり，$G/N$も群になる．これを\term[]{剰余群}という．%商群

$f:G\to H$を準同形とすると，$a\in\Ker f$,  $g\in G$に対して
\[
     f(gag^{-1})=f(g)f(a)f(g^{-1})=f(g)ef(g)^{-1}=f(g)f(g)^{-1}=e
\]
なので$\Ker f$は$G$の正規部分群である．このとき
\[
     G/\Ker f\sim \Image f
\]
が成り立つというのが群の\term[]{準同形定理}であった．
\emptyline

$\mathbf N$ を（$1$以上の）自然数全体のなす集合, $\mathbf Z$を整数全体のなす集合とする．$\mathbf Z$は加法$+$に関して群である．
$\mathbf Z$には掛け算$\cdot$ ( $\times$と書くこともある)もあるが$0$を除いた$\mathbf Z\setminus\{0\}$では割り算ができないこともあるので群ではない．
つまり，$\mathbf Z$は$+$と$\cdot$ に関して環であるが，体\footnote{環$K$の$0$以外の元が$\cdot$についての逆元を持つこと}ではない．

$m+n=n+m$なので可換群（アーベル群）である，加法$+$に関して加群ということがある．
$m\in \mathbf N$を固定し，$x, x'\in \mathbf Z $
\[
     x\equiv x' \mod m \Leftrightarrow 
     \exists q \in  \mathbf Z \text{ such that }x-x'=qm
\]
と置くとこれは同値関係である．
$x\in \mathbf Z $を代表元とする同値類を$\bar x$と置く．
%$\mathbf Z $を加法$+$に関する可換群と見なすと
\[
\begin{split}
     &x\equiv x' \mod m  \text{ and } y\equiv y' \mod m\\
     \Rightarrow \quad
     & x + y \equiv x' + y' \mod m,\quad
     xy\equiv x'y' \mod m
\end{split}
\]
なので
\[
     \bar x + \bar y = \overline{x+y},\quad
     \bar x\cdot \bar y = \overline{xy}
\]
と定め暫く$G$と置く．
$(G, +)$ は$\bar 0$を単位元とする可換群である．
\[
     f:\mathbf Z \ni x \to \bar x \in G
\]
とすると
\[
     \Ker f= m\mathbf Z=\bar 0
\]
見なせる．
$ m\mathbf Z$は$\mathbf Z$の正規部分群であるので，準同形定理より
\[
     G \sim \mathbf Z/m\mathbf Z
\]
となる．
\term[]{巡回群}とは一つの元から生成される群のことであった．$\mathbf Z/m\mathbf Z$は$\bar 1$で生成されているので有限巡回群である．授業では回転で生成される有限部分群を巡回群といったが，普通はこちらの定義の方が一般的である．例えば$D_1$も巡回群で$C_2$と同形である．
また，無限集合になる巡回群を無限巡回群というが
無限巡回群は$\mathbf Z$と同形になる．
\emptyline

以下，
$\mathbf Z/m\mathbf Z$の積$\cdot$についての構造を調べる．
最初にいくつか計算してみる．

\begin{ex}
   $\mathbf Z/4\mathbf Z$の足し算と掛け算の演算表を作ってみる．
   足し算については巡回群$C_4$なのでそのままであるが，
   積に関しては$\bar 2$が零因子となっているので群でなく，
   従って$\mathbf Z/4\mathbf Z$は体でない．
   \begin{table}[htb]
      \centering
      \begin{minipage}{0.4\columnwidth}
           \centering
           \begin{tabular}{l|llll}
                 $+$&$\bar 0$&$\bar1$&$\bar2$&$\bar3$ \\ \hline
                $\bar 0$&$\bar 0$&$\bar1$&$\bar2$&$\bar3$ \\
                $\bar1$& $\bar1$&$\bar2$ &$\bar3$&$\bar0$ \\
                $\bar2$ &$\bar2$&$\bar3$&$\bar 0$&$\bar1$ \\ 
                $\bar3$& $\bar3$ &$\bar0$& $\bar1$&$\bar2$\\ 
           \end{tabular}
           \caption{$+$ $\mod 4$}\label{mod4plus}
      \end{minipage} 
      \begin{minipage}{0.4\columnwidth}
           \centering
           \begin{tabular}{l|llll}
                 $\cdot$&$\bar 0$&$\bar1$&$\bar2$&$\bar3$ \\ \hline
                $\bar 0$&$\bar 0$&$\bar0$&$\bar0$&$\bar0$ \\
                $\bar1$& $\bar0$&$\bar1$ &$\bar2$&$\bar3$ \\
                $\bar2$ &$\bar0$&$\bar2$&$\bar 0$&$\bar2$ \\ 
                $\bar3$& $\bar0$ &$\bar3$& $\bar2$&$\bar1$\\ 
           \end{tabular}
           \caption{$\cdot$ $\mod 4$}\label{mod4times}
      \end{minipage} 
   \end{table}
\end{ex}

\begin{ex}
   $\mathbf Z/5\mathbf Z$の足し算と掛け算の演算表を作ってみる．
   掛け算では零因子がないので$\bar0$を除く$(\mathbf Z/5\mathbf Z)^{\times}=(\mathbf Z/5\mathbf Z)\setminus \{\bar 0\}$は群となる．
   これは巡回群$C_4$と同型である(表~\ref{mod4plus}
   と異なって見えるかもしれないが，$\bar 1_4$, $\bar 2_5$のように表し区別し，$\varphi:(\mathbf Z/5\mathbf Z)^{\times}\to \mathbf Z/4\mathbf Z$を
   \[
        \varphi(\bar 1_5)=\bar 0_4,\quad
        \varphi(\bar 2_5)=\bar 1_4,\quad
        \varphi(\bar 3_5)=\bar 3_4,\quad
        \varphi(\bar 4_5)=\bar 2_4
   \]
   とすると$(\mathbf Z/5\mathbf Z)^{\times}$は掛け算，$\mathbf Z/4\mathbf Z$は足し算として同型になることが分かる)．
   このことから$\mathbf Z/5\mathbf Z$は体でもある．
   
   \begin{table}[htb]
      \centering
      \begin{minipage}{0.4\columnwidth}
           \centering
           \begin{tabular}{l|lllll} 
                 $+$&$\bar 0$&$\bar1$&$\bar2$&$\bar3$ &$\bar4$\\ \hline
                $\bar 0$&$\bar 0$&$\bar1$&$\bar2$&$\bar3$&$\bar4$ \\
                $\bar1$& $\bar1$&$\bar2$ &$\bar3$&$\bar4$& $\bar0$\\
                $\bar2$ &$\bar2$&$\bar3$&$\bar 4$&$\bar0$&$\bar1$ \\ 
                $\bar3$& $\bar3$ &$\bar4$& $\bar0$&$\bar1$&$\bar2$\\ 
                $\bar4$& $\bar4$& $\bar0$&$\bar1$&$\bar2$&$\bar2$ \\
           \end{tabular}
           \caption{$+$ $\mod 5$}
      \end{minipage} 
      \begin{minipage}{0.4\columnwidth}
           \centering
           \begin{tabular}{l|lllll} 
                 $\cdot$&$\bar 0$&$\bar1$&$\bar2$&$\bar3$ &$\bar4$\\ 
                \hline
                $\bar 0$&$\bar 0$&$\bar0$&$\bar0$&$\bar0$&$\bar0$ \\
                $\bar1$& $\bar0$&$\bar1$ &$\bar2$&$\bar3$ &$\bar4$\\
                $\bar2$ &$\bar0$&$\bar2$&$\bar 4$&$\bar1$ &$\bar3$\\ 
                $\bar3$& $\bar0$ &$\bar3$& $\bar1$&$\bar4$&$\bar2$\\ 
                $\bar4$& $\bar0$& $\bar4$&$\bar3$&$\bar2$&$\bar1$ \\
           \end{tabular}
           \caption{$\cdot$ $\mod 5$}
      \end{minipage} 
   \end{table}
\end{ex}

%上で指摘したように$\mathbf Z/m\mathbf Z$には$\cdot$という積があり環となっている．しかし，$\bar 0$を除いても$\cdot$に関して一般には群ではない．


$m, n\in \mathbf Z$に対して$(m,n)$で$m$ と $n$の最大公約数を表す．
\[
     (m,n)=1
     \quad \Leftrightarrow\quad
     \exists a, b\in \mathbf Z \text{ such that } am + bn =1
\]
がユークリッドの互除法から分かる．このことから
\[
     (m,n)=1
     \quad \Leftrightarrow\quad
     \exists  \bar b\in \mathbf Z/m\mathbf Z \text{ such that } 
           \bar b\bar n= \bar 1
\]
$m\in \mathbf N$を固定し$(m,n)=1$なる
$\bar n\in\mathbf Z/m\mathbf Z$を\term[]{既約剰余類}という．
既約剰余類$\bar a\in \mathbf Z/m\mathbf Z$ を固定すると$\bar b\bar a = \bar 1$なる$\bar b\in \mathbf Z/m\mathbf Z$が見つかるので
\[
     f:\mathbf Z/m\mathbf Z\ni \bar n 
      \mapsto \bar{a}\cdot\bar{n}\in  \mathbf Z/m\mathbf Z
\]
を考えると$f$は$+$に関して準同形写像で，$\bar b$を掛けることで同形写像であることが分かる．
\[
     \varphi(m)= |\{n\mid n=1,\dots, m \text{ such that } (m,n)=1\}|
\]
と置くと$\mathbf Z/m\mathbf Z$には$\varphi(m)$個の既約剰余類が存在することが分かる．$\varphi$はオイラー関数と言われる．
したがって$\mathbf Z/m\mathbf Z$には$\varphi(m)$個の生成元が存在し，$\mathbf Z/m\mathbf Z$はどの生成元をとっても加群としては
同型な巡回群となる．

$(m,n)=1$なる$n$を固定する．
$x_1, \cdots, x_{\varphi(m)}$を$\mathbf Z/m\mathbf Z$の既約剰余類全体とする．$\bar n x_1, \cdots , \bar nx_{\varphi(m)}$も既約剰余類全体なので，集合としては一致している．したがって
\[
     x_1\cdots x_{\varphi(m)}=\bar n x_1 \cdots \bar n x_{\varphi(m)}
     =\bar n^{\varphi(m)}x_1\cdots x_{\varphi(m)}
\]
$x_1, \cdots , x_{\varphi(m)}$は逆元をもつので
$
     \bar n^{\varphi(m)}=\bar 1\in\mathbf Z/m\mathbf Z
$,  つまり
\[
     n^{\varphi(m)}\equiv 1 \mod m
\]
$p$を素数とすると，$\varphi(p)=p-1$なので任意の$x\not\equiv 0 \mod p$に対して
\[
     x^{p-1}\equiv 1 \mod p
\]
となり，\term[]{フェルマーの小定理}が得られる．これから任意の$x\in \mathbf Z$に対して
\begin{equation}\label{funPoly}
     x^{p}\equiv x \mod p
\end{equation}
も分かる．
フェルマーの小定理は$(\mathbf Z/p\mathbf Z)^{\times}=(\mathbf Z/p\mathbf Z)\setminus \{\bar 0\}$が積$\cdot$について巡回群であることを意味している．
特に$\mathbf Z/p\mathbf Z$が体であることも意味していている．
このとき$\mathbf F_p$と表すことが多い．これは標数$p$の有限体である．

\begin{Quiz}
   標数$p$の有限体は$\mathbf F_p$の有限次元拡大であること，
   それが$\mathbf F_p$のガロワ拡大であることを示し，
   そのガロワ群を決定せよ（ヒント$:$
    \eqref{funPoly}の形の多項式でガロワ拡大を特徴付けよ）．
\end{Quiz}



\subsection{アミダクジと置換群}
最初に置換と阿弥陀籤が同じものであることを示す．
アミダクジの例を次にあげる．
\begin{figure}[htb]
     \centering
     \includegraphics{~/Documents/GitHub/bEuclideanGeo/appa/images/figure09.1}
     \caption{アミダクジ $\tau$}\label{}
\end{figure}
$4$本の柱に順に$1,\dots,4$と名前を付け，そのアミダクジを上から下にたどるとき何番目の柱から何番目の柱に移るかのリストと同一視することにする．例えば．このアミダクジでは柱$1$から柱$2$に移り，柱$2$から柱$4$に移り，柱$3$から柱$1$に移り，柱$4$から柱$3$に移るので，このアミダクジを
$\begin{pmatrix}1&2&3&4\\2&4&1&3\end{pmatrix}$と表す．
これを置換（表示）ということもある．
上の数字は上の柱の番号で対になる下の番号が移る柱の番号を表している．下の行に出てくる数字に重複がないことを確かめてほしい．
簡単のため，このアミダクジを$\tau$と表す．
\[
     \tau(1)=2, \tau(2)=4,\tau(3)=1,\tau(4)=3
\]
と表すこともある．つまり，アミダクジ$\tau$は$\{1,2,3,4\}$から$\{1,2,3,4\}$への上への一対一対応（写像）を定めている．

アミダクジ$\rho$を
\begin{figure}[htb]
     \centering
     \includegraphics{~/Documents/GitHub/bEuclideanGeo/appa/images/figure09.2}
     \caption{アミダクジ $\rho$}\label{}
\end{figure}

\noindent
とする．$\tau$を行ったあと$\rho$を行う事によって出来るアミダクジ
\vspace{10pt}

\begin{figure}[htb]
     \centering
     \includegraphics{~/Documents/GitHub/bEuclideanGeo/appa/images/figure09.3}
     \caption{$\tau$と$\rho$の合成}\label{}
\end{figure}

\noindent
を$\tau$と$\rho$の\term[]{積}(\term[]{掛け算})といい，$\rho\tau$と表す．つまり，対応（写像）としては
\[
     \rho\tau(i)=\rho(\tau(i))\qquad i=1,2,3,4
\]
（つまり，関数の合成や，(ベクトルを作用させた)行列の積と同じ）である．置換を使い
\[
     \rho\tau
     =
     \begin{pmatrix}1&2&3&4\\ 4&2&1&3\end{pmatrix}
     \begin{pmatrix}1&2&3&4\\2&4&1&3\end{pmatrix}
     =
     \begin{pmatrix}1&2&3&4\\2&3&4&1\end{pmatrix}
\]
のように$\tau$で$1$が$2$にうつり，次に$\rho$で$2$が$2$にうつるので$\rho\tau$では$1$が$2$うつる，$\tau$で$2$が$4$にうつり，次に$\rho$で$4$が$3$にうつるので$\rho\tau$では$2$が$4$うつる
というように計算してもいい．

\begin{Quiz}
\[
     \tau\rho=
     \begin{pmatrix}1&2&3&4\\2&4&1&3\end{pmatrix}
     \begin{pmatrix}1&2&3&4\\ 4&2&1&3\end{pmatrix}
     =
     \begin{pmatrix}1&2&3&4\\ & & & \end{pmatrix}
\]
の空白を埋め，$\tau\rho\neq\rho\tau$を確かめよ．
\end{Quiz}

順番を変えないアミダクジを\term[]{単位}アミダクジといい，$\varepsilon$と表す．つまり，
\begin{figure}[htb]
     \centering
     \includegraphics{~/Documents/GitHub/bEuclideanGeo/appa/images/figure09.4}
     \caption{単位アミダクジ}\label{}
\end{figure}

明らかに任意のアミダクジ$\sigma$に対して
\[
     \varepsilon\sigma=\sigma\varepsilon=\sigma
\]
が成り立つ．

アミダクジ$\sigma$を下から上にたどって出来るアミダクジを\term[]{逆}アミダクジといい，$\sigma^{-1}$と表す．例えば$\tau$の逆は

\begin{figure}[htb]
     \centering
     \includegraphics{~/Documents/GitHub/bEuclideanGeo/appa/images/figure09.5}
     \caption{逆アミダクジ}\label{}
\end{figure}

\noindent
である．左側では置換表示の上の行と下の行を入れ替え，それを並び替えるという操作を行っている．
さらに
\begin{figure}[htb]
     \centering
     \includegraphics{~/Documents/GitHub/bEuclideanGeo/appa/images/figure09.6}
     \caption{アミダクジと逆アミダクジの合成}\label{}
\end{figure}

\noindent
これは行った道を逆にたどれば元いた場所に戻れるということをアミダクジで表しただけなので，任意のアミダクジ$\sigma$に対して
\[
     \sigma^{-1}\sigma=\sigma\sigma^{-1}=\varepsilon
\]
が成り立つ．このようにアミダクジの積はアミダクジの自然な演算を与えている．
\emptyline

$n$本のアミダクジ$\sigma$は
$\{1,2,3,\dots,n\}$から$\{1,2,3,\dots,n\}$への上への一対一対応（写像）
\[
      \sigma=
     \begin{pmatrix}
          1&2&3&\cdots&n\\
          i_1&i_2&i_3&\cdots&i_n
     \end{pmatrix}
\]
で表せる．ただし，$\{i_1,i_2,i_3,\cdots,i_n\}=\{1,2,3,\dots,n\}$ である．
このような写像を$n$\term[]{次置換}という．
同じ置換を表すアミダクジを合同ということにする．
逆に，一般の$n$次置換が$n$本の（ある合同な）アミダクジで表せるかという問題を考えよう．


次のようなアミダクジ（置換）を考える：
\begin{figure}[htb]
     \centering
     \includegraphics{~/Documents/GitHub/bEuclideanGeo/appa/images/figure09.7}
     \caption{循環置換}\label{}
\end{figure}

\noindent
このような置換を（4次の）\term[]{循環置換}といい，
$\begin{pmatrix}1&2&3&4\end{pmatrix}$と表す．
特に$2$次の循環置換を\term[]{互換}という．
例えば

\begin{figure}[htb]
     \centering
     \includegraphics{~/Documents/GitHub/bEuclideanGeo/appa/images/figure09.8}
     \caption{互換}\label{}
\end{figure}

\noindent
である．
上の巡回置換を次のように点線で分けて考えると
\begin{figure}[htb]
     \centering
     \includegraphics{~/Documents/GitHub/bEuclideanGeo/appa/images/figure09.9}
     \caption{アミダクジの互換への分解}\label{}
\end{figure}

\noindent
が成り立つ．
右から左の積の順番はアミダクジの上から下への連結に対応していることをあらためて注意しておく．
また，
別やり方で表すと
\begin{figure}[htb]
     \centering
     \includegraphics{~/Documents/GitHub/bEuclideanGeo/appa/images/figure09.12}
     \caption{アミダクジの互換への別の分解}\label{}
\end{figure}

\noindent
のようになる．
このように置換のアミダクジによる表し方は色々あるので，覚えやすい方法で覚え，わからくなったらアミダクジで表してみることをお奨めする．次に



\begin{figure}[h]
     \centering
     \includegraphics{~/Documents/GitHub/bEuclideanGeo/appa/images/figure09.10}
     \caption{互換のアミダクジでの表示}\label{}
\end{figure}

\noindent
のように任意の柱の互換もアミダクジで表すことが出来る．

\begin{comment}
以上より
循環置換$\begin{pmatrix}i_1&i_2&\dots&i_r\end{pmatrix}$は
\begin{figure}[htb]
     \centering
     \includegraphics{~/Documents/GitHub/bEuclideanGeo/appa/images/figure09.11}
\end{figure}
と互換の積に分解される．
\end{comment}

以上の準備のもと，一般の置換は互換の積で表せることを例を使って示す．
\begin{exQ} 次の置換を循環置換の積に分解し，更に互換の積で表せ．
\[
     \mu=
     \begin{pmatrix}
          1&2&3&4&5&6&7\\
          4&1&6&2&7&5&3
     \end{pmatrix}
\]
\end{exQ}
解答）
$1$が$4$に移り，$4$が$2$に移り，$2$が$1$に移ので，
巡回置換$\begin{pmatrix}1&4&2\end{pmatrix}$になっている．
\[
\begin{split}
     \mu&\Rightarrow 
     \begin{pmatrix}
          1&2&3&4&5&6&7\\\downarrow&&&&&&\\
           4&1&6&2&7&5&3 
      \end{pmatrix}\\
     &\Rightarrow 
     \begin{pmatrix}
          1&2&3&4&5&6&7\\
          \downarrow&&&\downarrow&&&\\
           4&1&6&2&7&5&3 
      \end{pmatrix}\\
     &\Rightarrow 
     \begin{pmatrix}
          1&2&3&4&5&6&7\\
          \downarrow&\downarrow&&\downarrow&&&\\
           4&1&6&2&7&5&3 
      \end{pmatrix}\\
     &\Rightarrow \begin{pmatrix}1&4&2\end{pmatrix}
\end{split}
\]
残りをみると
$3$が$6$に移り，$6$が$5$に移り，$5$が$7$に移り，$7$が$3$に移
ので，巡回置換$\begin{pmatrix}3&6&5&7\end{pmatrix}$になっている．
\[
\begin{split}
     \mu&=\begin{pmatrix}1&2&3&4&5&6&7\\\downarrow&\downarrow&\Downarrow&\downarrow&\Downarrow&\Downarrow&\Downarrow\\
      4&1&6&2&7&5&3 \end{pmatrix}\\
     &=\begin{pmatrix}1&4&2\end{pmatrix}
     \begin{pmatrix}3&6&5&7\end{pmatrix}
\end{split}
\]
したがって，先程のように巡回置換を互換で表すと
\[
\begin{split}
     \mu&=\begin{pmatrix}1&4&2\end{pmatrix}
     \begin{pmatrix}3&6&5&7\end{pmatrix}\\
     &=\begin{pmatrix}1&4\end{pmatrix}
     \begin{pmatrix}4&2\end{pmatrix}
     \begin{pmatrix}3&6\end{pmatrix}
     \begin{pmatrix}6&5\end{pmatrix}
     \begin{pmatrix}5&7\end{pmatrix}.
\end{split}
\]
もちろん
\[
\begin{split}
     \mu&=
     \begin{pmatrix}3&6&5&7\end{pmatrix}
     \begin{pmatrix}1&4&2\end{pmatrix}
     \\
     &=
     \begin{pmatrix}3&6\end{pmatrix}
     \begin{pmatrix}6&5\end{pmatrix}
     \begin{pmatrix}5&7\end{pmatrix}
     \begin{pmatrix}1&4\end{pmatrix}
     \begin{pmatrix}4&2\end{pmatrix}.
\end{split}
\]
でも構わない．

以上のやり方で全ての置換は互換の積で表せることが分かり，
全ての置換はアミダクジで表される．つまり，置換とアミダクジは同じものであることが分かる．

\emptyline
最後に$\sigma$が$m$個の互換の積にかける(つまり，アミダクジの横棒が$m$本の)とき
\[
     \sign \sigma=(-1)^m
\]
で$\sigma$の\term[]{符号}を定める．例えば上の$\mu$については$\sign\mu=(-1)^5=-1$である．
%また，最初の例の$\rho$ではもともとアミダクジで表されていて，互換の数（つまり，横の線の数）は$4$個なので，$\sign\tau=1$である．
\begin{exQ}  次の置換を互換の積で表し，符号を求めよ．
\[
     \mu=\begin{pmatrix}
          1&2&3&4&5&6&7&8&9\\
          7&6&8&2&1&4&9&3&5
     \end{pmatrix}
\]
\end{exQ}
解答）
\[
\begin{split}
     \mu&=
     \begin{pmatrix}
          1&7&9&5
     \end{pmatrix}
     \begin{pmatrix}
          2&6&4
     \end{pmatrix}
     \begin{pmatrix}
          3&8
     \end{pmatrix}\\
     &=\begin{pmatrix}
          1&7
     \end{pmatrix}
     \begin{pmatrix}
          7&9
     \end{pmatrix}
     \begin{pmatrix}
          9&5
     \end{pmatrix}
     \begin{pmatrix}
          2&6
     \end{pmatrix}
     \begin{pmatrix}
          6&4
     \end{pmatrix}
     \begin{pmatrix}
          3&8
     \end{pmatrix}
\end{split}
\]
よって，$\sign \mu=(-1)^6=1$.
\emptyline

単位置換$\varepsilon$は$0$個の互換の積であらわせるので
\[
     \sign\varepsilon=(-1)^0=1
\]
また，定義から
\[
     \sign \sigma\tau=\sign \sigma\sign \tau,\quad \sign \sigma^{-1}= \sign \sigma
\]
が成り立つ．
$\sign \sigma=1$のとき$\sigma$を\term[]{偶置換}，
$\sign \tau=-1$のとき$\tau$を\term[]{奇置換}という．
特に，$\sign:S_n\to \{1, -1\}$は準同形写像であり，
$A_n = \Ker \sign =\{\sigma\in S_n\mid $偶置換$\}$は部分群である．

%\[
%     \mu=\begin{pmatrix}1&2&3&4&5&6&7&8\\4&8&6&3&1&5&2&7\end{pmatrix}
%\]
%とすると，$1$が$4$に移り，$4$が$3$に移り，$3$が$6$に移り，$6$が$5$に移り，$5$が$1$に移ので，巡回置換$\begin{pmatrix}1&4&3&6&5\end{pmatrix}$になっている．残りをみると
%$2$が$8$に移り，$8$が$7$に移り，$7$が$2$に移る
%ので，巡回置換$\begin{pmatrix}2&8&7\end{pmatrix}$になっている．したがって，先程のように巡回置換を互換で表すと
%\[
%\begin{split}
%     \mu&=\begin{pmatrix}1&4&3&6&5\end{pmatrix}\begin{pmatrix}2&8&7\end{pmatrix}\\
%     &=\begin{pmatrix}1&4\end{pmatrix}
%     \begin{pmatrix}4&3\end{pmatrix}
%     \begin{pmatrix}3&6\end{pmatrix}
%     \begin{pmatrix}6&5\end{pmatrix}
%     \begin{pmatrix}2&8\end{pmatrix}
%     \begin{pmatrix}8&7\end{pmatrix}
%\end{split}
%\]
%もちろん数字の大きい方から始めて
%\[
%\begin{split}
%     \mu&=\begin{pmatrix}8&7&2\end{pmatrix}\begin{pmatrix}6&5&1&4&3\end{pmatrix}\\
%     &=
%     \begin{pmatrix}8&7\end{pmatrix}
%     \begin{pmatrix}7&2\end{pmatrix}
%     \begin{pmatrix}6&5\end{pmatrix}
%     \begin{pmatrix}5&1\end{pmatrix}
%     \begin{pmatrix}1&4\end{pmatrix}
%     \begin{pmatrix}4&3\end{pmatrix}
%\end{split}
%\]
%でも構わない．


\begin{Quiz}
　次の置換を互換の積に表せ．
   \[
        \begin{pmatrix}
             1&2&3&4&5\\
             4&5&2&1&3
        \end{pmatrix}
   \]
   それを用いて，この置換を表すアミダクジを一つ作れ．
\end{Quiz}

\subsection{置換群と交代群の構造}
$n$次置換群の構造を調べる．
$n$次の置換全体を$S_n$と置く．$S_n$は群であり，$n$\term[]{次置換群}と呼ばれる．
$S_n$の元は$\begin{pmatrix}1&2&3&\cdots&n\\i_1&i_2&i_3&\cdots&i_n\end{pmatrix}$と表されるので，個数は$n!$ である．これがアミダクジで表されるので，$S_n$は$n-1$個の互換$\begin{pmatrix}1&2\end{pmatrix}$, $\begin{pmatrix}2&3\end{pmatrix}$, \dots, $\begin{pmatrix}n-1&n\end{pmatrix}$
により生成されることが分かる．

$A_n=\{\sigma\in S_n\mid$ 偶置換 $\}$と置くと
 \begin{gather*}
      \sigma,\tau\in A_n
     \Rightarrow
     \sign \sigma\tau=\sign \sigma  \sign\tau=1
     \Rightarrow
      \sigma\tau\in A_n
     \\
       \sigma\in A_n
     \Rightarrow 
     \sign \sigma^{-1}=\sign \sigma=1
     \Rightarrow
     \sigma^{-1}\in A_n
 \end{gather*}
 が成り立つので$A_n$は$S_n$の{部分群}であり，$n$\term[]{次交代群}と呼ばれる．
 \emptyline

 $n=3$のとき$S_3$は
\begin{equation}
\begin{split}
   \begin{pmatrix}1&2&3\\1&2&3\end{pmatrix}=\varepsilon,\quad
   \begin{pmatrix}1&2&3\\1&3&2\end{pmatrix}
   =
   \begin{pmatrix}
      2&3
   \end{pmatrix},\\
   \begin{pmatrix}1&2&3\\2&1&3\end{pmatrix}
   =\begin{pmatrix}
      1&2
   \end{pmatrix},\quad
   \begin{pmatrix}1&2&3\\2&3&1\end{pmatrix}
   =\begin{pmatrix}
      1&2&3
   \end{pmatrix},\\
   \begin{pmatrix}1&2&3\\3&1&2\end{pmatrix}=
   \begin{pmatrix}
      1&3&2
   \end{pmatrix},\quad
   \begin{pmatrix}1&2&3\\3&2&1\end{pmatrix}=
   \begin{pmatrix}
      1&3
   \end{pmatrix}
\end{split}\label{S3}
\end{equation}
の$3!=6$個からなり，$\varepsilon$, $\begin{pmatrix}
          1&2&3
     \end{pmatrix}$, $\begin{pmatrix}
          1&3&2
     \end{pmatrix}$が偶置換で，$\begin{pmatrix}
          1&2
     \end{pmatrix}$, $\begin{pmatrix}
          2&3
     \end{pmatrix}$, $\begin{pmatrix}
          1&3
     \end{pmatrix}$が奇置換である．
よって
\[
     A_3 = \{
     \varepsilon, 
     \begin{pmatrix}
          1&2&3
     \end{pmatrix},
     \begin{pmatrix}
          1&3&2
     \end{pmatrix} \}\subset S_3
\]
で$A_3$は$C_3$と同形．
\emptyline

次に$n=4$のとき
 \[
\begin{split}
     K=\{\varepsilon, &\quad
     \alpha=\begin{pmatrix}1&2\end{pmatrix}
     \begin{pmatrix}3&4\end{pmatrix}, \quad
     \beta=\begin{pmatrix}1&3\end{pmatrix}
     \begin{pmatrix}2&4\end{pmatrix},\\
     &\gamma=\begin{pmatrix}1&4\end{pmatrix}
     \begin{pmatrix}2&3\end{pmatrix}\}
     \subset  A_4
\end{split}
\]
\begin{figure}[thb]
     \centering
     \includegraphics{~/Documents/GitHub/bEuclideanGeo/ch1/figureGroup3.7}
     \caption{$\alpha$, $\beta$, $\gamma$}
\end{figure}


\noindent
とすると，明らかに$\alpha^2=\beta^2=\gamma^2=\varepsilon$で，
図 \ref{betaAlphaGamma}より
\[
     \alpha\beta=\beta\alpha=\gamma
\]

\begin{figure}[htb]
     \centering
     \includegraphics{~/Documents/GitHub/bEuclideanGeo/ch1/figureGroup3.8}
     \caption{$\beta\alpha=\gamma=\alpha\beta$}\label{betaAlphaGamma}
\end{figure}

\noindent
これから
\[
     \gamma\alpha=\beta\alpha\alpha=\beta,\quad
     \alpha\gamma=\alpha\alpha\beta=\beta
\]
のように計算することで次の表を得る．

\begin{table}[htb]
      \centering
      \begin{tabular}{l|llll}\noalign{\hrule height 1pt}
         &$\varepsilon$&$\alpha$&$\beta$&$\gamma$ \\ 
         \hline
         $\varepsilon$&$\varepsilon$&$\alpha$&$\beta$&$\gamma$ \\
         $\alpha$& $\alpha$&$\varepsilon$ &$\gamma$&$\beta$ \\
         $\beta$& $\beta$&$\gamma$&$\varepsilon$&$\alpha$ \\ 
         $\gamma$& $\gamma$ &$\beta$&  $\alpha$&$\varepsilon$\\
        \noalign{\hrule height 1pt}
      \end{tabular}     \caption{クライン群$K$の演算表}
\end{table}

\noindent
$K$は$A_4$の部分群であり，積が順番によらないこともわかる．
$K$はクライン群と呼ばれる可換群である．以後
$\begin{pmatrix}1&2\end{pmatrix}\begin{pmatrix}3&4\end{pmatrix}$
等を独立な互換の積ということがある．

%\emptyline
%$G$を群，$H$を$G$の部分群とする．$g\in G$について
%一般には$gH\neq Hg$であるが，
%\[
%     \text{任意の}\quad g\in G\quad\text{に対して}\quad gH= Hg
%\]
%が成り立つとき$H$は$G$の\term[]{正規部分群}という．これは
%\[
%     \text{任意の}\quad g\in G\quad\text{に対して}\quad gHg^{-1}= H
%\]
%と同値である．$H\lhd G$と書くことがある．
$n$次交代群$A_n$は$S_n$の正規部分群である．なぜなら，任意の
$\sigma \in S_n$と$\tau\in A_n$に対して
\[
     \sign\sigma\tau\sigma^{-1}
     =\sign\sigma\sign\tau\sign\sigma^{-1}
     =\sign\tau=1
\]
なので$\sigma A_n\sigma^{-1}\subset A_n$, つまり，
$\sigma A_n\sigma^{-1}= A_n$となるからである．
\emptyline

%$H\lhd G$のとき正規部分群の定義より
%\[
%     gNhN
%     =g(hh^{-1})NhN
%     =gh(h^{-1}Nh)N=gh NN =gh N
%\]
%特に
%\[
%     g^{-1}NgN=g^{-1}gN=N
%\]
%となり，$G/N$も群になる．これを\term[]{商群}という．

群$G$について正規部分群の列
\[
     G_m=\varepsilon\lhd  G_{m-1}\lhd \cdots  \lhd G_1\lhd  G_0=G
\]
が存在し，その商群$G_{i-1}/G_{i}$が可換群であるとき$G$を
\term[]{可解群}という．ただし，$\varepsilon$ は単位元のみから成る群$\{\varepsilon\}$を意味する．

\emptyline
$S_3$では
\[
     \varepsilon\lhd   A_3 \lhd  S_3
\]
で，上で見たように
$ A_3$は$ S_3$の正規部分群で$S_3/ A_3$は可換群$C_2$であり，$A_3$は$\sigma=\begin{pmatrix}1&2&3\end{pmatrix}$で作られる$3$次循環群$C_3$であり可換群である．よって$S_3$は可解群である．
\emptyline

$S_4$は
\[
     \varepsilon\lhd  K\lhd  A_4 \lhd S_4
\]
という正規部分群の列を持つ．$K$は可換群で$ A_4$と$ S_4$部分は上と同じである．$\sigma\in S_n$はアミダクジ$\tau\in S_n$に$\sigma\tau\sigma^{-1}$で作用するが，これはアミダクジ$\tau$の柱の名前を付け替えることに対応している．$K$は二組の異なる柱の組の互換と単位置換からなるのでこの作用で保存される．したがって$K$は$A_4$の正規部分群であることが分かる．
$\sigma=\begin{pmatrix}1&2&3\end{pmatrix}\notin K$を固定すると，
$A_4/K=\{K, \sigma K, \sigma^{2}K\}$は$C_3$と同形となり，やはり可換群である．よって，$S_4$は可解群である．

\begin{Thm}
   $n\ge 5$のとき$A_n$は可解群でない．
\end{Thm}
\begin{proof}
   まず$A_n$の任意の元が$\begin{pmatrix}i&j&k\end{pmatrix}$の形の循環置換（以下，三環という）の積で表せることを示す．$A_n$は偶置換からなるので
   $\sigma=\begin{pmatrix}i&j\end{pmatrix}\begin{pmatrix}k&l\end{pmatrix}$
   の形($i\neq j$, $k\neq l$)の置換の積で表すことが出来る．
   $\{i,j\}=\{k,l\}$のときは$\sigma = \varepsilon$. 
   $\{i,j\}$と$\{k,l\}$の共通元が一個だけあるとき，それを$j=l$とすると
   \[
        \sigma =
        \begin{pmatrix}i&j\end{pmatrix}\begin{pmatrix}j&k\end{pmatrix}
        =\begin{pmatrix}i&j&k\end{pmatrix}
   \]
   $\{i,j\}$と$\{k,l\}$の共通元がないとき，
   \[
        \sigma =
        \begin{pmatrix}i&j\end{pmatrix}\begin{pmatrix}k&l\end{pmatrix}
        =\begin{pmatrix}i&j\end{pmatrix}\begin{pmatrix}j&k\end{pmatrix}
        \begin{pmatrix}j&k\end{pmatrix}\begin{pmatrix}k&l\end{pmatrix}
        =\begin{pmatrix}i&j&k\end{pmatrix}
        \begin{pmatrix}j&k&l\end{pmatrix}
   \]
   となるので証明できた．
   
   $A_n$の正規部分群$N\neq A_n$で$A_n/N$が可換群であるものが存在したとする．すると，任意の$g, h\in A_n$に対して
   \[
        ghN=gNhN=hNgN=hgN
   \]
   つまり，
   \[
        N= (gh)^{-1}hgN=h^{-1}g^{-1}hgN
   \]
   $\varepsilon\in N$より
   \[
        h^{-1}g^{-1}hg=h^{-1}g^{-1}hg\,\varepsilon\in N
   \]
   となる．
   
   
 \begin{figure}[thb]
      \centering
      \begin{minipage}{0.4\columnwidth}
         \centering
           \includegraphics{~/Documents/GitHub/bEuclideanGeo/ch1/figureGroup3.9}
           \caption{$h^{-1}g^{-1}hg$}
         \label{}
      \end{minipage} 
      \begin{minipage}{0.4\columnwidth}
         \centering
           \includegraphics{~/Documents/GitHub/bEuclideanGeo/ch1/figureGroup3.10}
           \caption{$D_3$}
         \label{}
      \end{minipage} 
   \end{figure}
   
   任意の異なる$1\le i, j, k\le n$に対して，$n\ge 5$よりそれらと異なる
   二つの$1\le l, m \le n$が存在する．
   $g=\begin{pmatrix}i&m&k\end{pmatrix}$, 
   $h=\begin{pmatrix}i&l&j\end{pmatrix}$とすると
   \[
        N \ni
        h^{-1}g^{-1}hg=
        \begin{pmatrix}i&j&l\end{pmatrix}\begin{pmatrix}i&k&m\end{pmatrix}
        \begin{pmatrix}i&l&j\end{pmatrix}\begin{pmatrix}i&m&k\end{pmatrix}
   \]
   
   \[     
        =
        \begin{pmatrix}
             i& m&j&n&k\\
             j&m&k&n&i
        \end{pmatrix}
        =
        \begin{pmatrix}i&j&k\end{pmatrix}
   \]
   となる，最初に調べたことから$N=A_n$となり，仮定に矛盾する．   
\end{proof}

\subsection{類等式}
第一章で使った議論が一般の有限群にも使えるような形に一般化されている：
$G$を群とする．全ての元と可換な元全体
\[
     Z(G) =\{a\in G\mid ax=xa\quad \forall x\in G\}
\]
を$G$の\term[]{中心}という．明らかに$Z(G)$ は$G$の(可換な)部分群である．

$a\in G$は$x\in G$に対して
\[
     x^a = a^{-1}xa
\]
と置くと
\[
     x^e = x,\quad x^{ab}=(x^{a})^b
\]
なので群$G$は集合$G$に作用している．
この作用による$G$の元の関係
\[
     y \sim x \quad \Leftrightarrow \exists a\in G \quad y=x^a
\]
は同値関係である．$a$を含む上の作用の同値類（つまり，$G$軌道）を$a^G$と表し，$a$を含む$G$の共役類という．
$a$の安定部分群は$a$の\term[]{中心化群}とも呼ばれ
$C_G(a)$と書かれる：
\[
     C_G(a)=\{x\in G\mid xa=ax\}
\]
以下$G$を有限群とする．すると
\[
     |a^G| = |G:C_G(a)|=|G|/|C_G(a)|
\]
更に，$|a^G|=1$は$a^G=a$と同値であるので，$ax=xa$が全ての$x\in G$になり立つことを意味し，$a\in Z(G)$と同値である．
$G$を$G$軌道に分解すると，
$a_1, \cdots, a_k\in G$が存在して
\[
     a_1^G \cup \dots \cup a_k^G =G, \quad  
     a_i^G\cap a_j^G=\emptyset \quad( i\neq j)
\]
なる分割を与えるが，
上の考察から，
\[
     |G|=|Z(G) | + \sum_{|a_i^G|>1}|a_i^G|
%     =|Z(G) | + \sum_{|G:C_G(a_i)|>1} |G:C_G(a_i)|
\]
がわかる．
これが\term[]{類等式}であり，
有限群論で非常に重要な事実である．

 
 %\section{Deck 変換と軌道体}

\section{アプリのソースと解説}
\subsection{p5.js}
p5.js は Processing というJava用の描画プログラミング言語を Javascript に移植したものである．以下のウェブ上のエディターで実行可能である．

\href{https://editor.p5js.org}{https://editor.p5js.org} （以下この色はハイパーリンク）

\begin{figure}[htb]
   \centering
   \includegraphics[width=13cm]{~/Documents/GitHub/bEuclideanGeo/appa/images/ScreenCapture3.pdf}
   \caption{ウェブ上の p5.js エディター}\label{}
\end{figure}

\noindent
公開した sketch.js を sketch.js の部分にコピーし，
（ファイル \href{https://www.printables.com/model/411768-octahedron-freecad/files}{octahedron.stl}\footnote{Ray Fox,
This work is licensed under a Creative Commons (International License),
Public Domain} をアップロードし）
左上の三角のボタンを押すと右画面に描画が始まる．
もちろん隣りの四角のボタンで止まる．
%プログラムのソースファイルに是非適宜手を入れながら動かして見て欲しい．
ただし .js はJavaScriptのファイルであることを示す拡張子 ( extension ) である．


また以下のようにすると自分のパソコン上で動かせる．
 p5.js というライブラリーを\href{https://p5js.org/download/}{ダウンロード}するとp5というディレクトリが作られる．
 公開されたsketch.jsをそのディレクトリーの sketch.js と置き換え，
フリーのエディター\href{https://code.visualstudio.com}{Visual Studio Code}でそのディレクトリp5を開き，
 Live Server の拡張をVisual Studio Codeにインストールして
Go Liveのボタンを押すと，chrome で再生を始める．

%\subsection{ p5.js のソース}
以下がソースファイル sketch.js の本質的な内容である．

\begin{verbatim}
/* Cubes landing before the Great Pyramids abbreviated version,
 * snow00two@gmail.com, 
 * \href{https://creativecommons.org/licenses/by-nc-nd/4.0/}{\ccbyncsa}
 */
const widthCanvas = 720 * 3/2  ; //=1080
const heightCanvas = 405 * 3/2 ; //=607.5
const backColor = [150, 200, 250] ;//[200, 220, 200]

let octahedron ; 
function preload(){
octahedron = loadModel('octahedron.stl'); 
//https://www.printables.com/model/411768-octahedron-freecad/files
}

function setup(){//canvas : window of the projection coordinates
createCanvas(widthCanvas, heightCanvas, WEBGL);
background(backColor);
}

let i = 0;//basic parameter for frames

function draw(){
let I = 9425; //change the priod condition to 3x2Pi/0.002=9424.7... 

if (i<I) {
} else {
i = 0;// reset basic parameter
}

  // define the camera coordinates from a camera on the helicopter
  // of the following coordinates in the world coordinates : 
  let cameraX = cos(0.5 + i * 0.002)*250 ;
  let cameraY = sin(0.5 + i * 0.002)*250;
  let cameraZ = 200+sin(0.5 + i * 0.002)*30;
  camera(cameraX,cameraY, cameraZ, 100, 0, 100, 0, 0, -1);
  // the output images are in the projection coordinates.

  //The world coordinates : 
  background(backColor);
  directionalLight(250, 200, 200, 0, -1, -1); //light from the sun 
  fill(250, 200, 200);
  plane(10000,10000);

  stroke(0);
  strokeWeight(2);
  fill(150,150, 150);
  pyramid(3);

  translate(-300,0,0);
  pyramid(3);

  translate(100,-300,0);
  pyramid(2);

  ambientMaterial(125);
  ambientLight(100,100,100);
  shininess(100.0);

  translate(400, 500, max(500 - i * 0.3, 67) );
  sCube(5);
  
  let M = 25;
  for(let m = 0 ; m < M ; m++){
    translate(140,140, max(650 + 150*m - i * 0.3, 0));
    sCube(5);
  } 
  i++;
}

// The modeling coordinates :
function pyramid(scaleN){
  push();
    scale(scaleN);
    model(octahedron);
  pop();
}

function sCube(scaleN){
  push();
    rotateZ(QUARTER_PI);
    scale(scaleN);
    box(25);
  pop();
}
\end{verbatim}

\subsection{JavaScriptの言語仕様}
\begin{itemize}
\item
JavaScript には変数の型はなく，文の終りで ; が必要

\item
let は変数の宣言，以前は var も使われていた．

\item
const は定数の宣言

\item
\verb////はコメント文でこれ以降の文は無視される．

\noindent
/* */はブロックコメント文で途中の文は無視される．
\end{itemize}


よく使う制御構文：
for 文の例：
\begin{verbatim}
   let I = 5;
   for (let i=0; i < I; i++) {
        console.log(i);
   }

\end{verbatim}

if 文の例：
\begin{verbatim}
   let m = 80;
   if ( m<60 ){
        console.log('F');
   } else if ( m<70  ) {
        console.log('C');
   } else if ( m<80  ) {
        console.log('B');
   } else if ( m<90  ) {
        console.log('A');
   } else {
        console.log('S');
   }

\end{verbatim}


function hoge(var1,\dots,var3) は関数hogeの宣言で，続く$\{ \}$の中に関数の定義を書く．var1,\dots,var3が関数の変数で定義の中で使える．
関数の中で定義された変数はローカル変数で，関数外で使われる変数はユニバーサルな変数になる．


%\subsection{p5.js の説明とソースファイルの解説}
関数preload()でディレクトリにある図形 octahedron.stl を読み込み，
関数setup()の中でウィンドウ等を初期化する．

   createCanvas(widthCanvas, heightCanvas) ;

\noindent
で$2$次元の描画を行えるが，ここでは

   createCanvas(widthCanvas, heightCanvas, WEBGL);

\noindent
で$3$次元の描画を行う．ただし，WEBGLとはブラウザ上でCGを描画するためのAPIである．ここでは$3$次元の描画のみ解説するが$2$次元の描画については\cite{P5} を参照してほしい．

関数draw()が呼び出されると変数 i を1フレーム毎に1づつ増やしながら，
その都度フレームを描画し続ける．いじってみるとだんだん理解が深まると思う．

\emptyline

p5.js 特有の組込関数の説明：
\begin{itemize}
   \item
      translate(a, b, c); $(a, b,c)$だけ平行移動し，その点を$(0,0,0)$として描画
   
   \item
      rotateZ(theta); $z$軸を中心に角度 theta だけ直交座標系を回転し，それを新たな直交座標系として描画．rotateX, rotateY も同様．
   \item
      scale(a) 座標系をa 倍する
   
   \item  
      strokeWeight(a) 太さaで描く
     
   \item
      background(a,b,c);  背後の画面をRGB [a,b,c] で埋める．
      ただし$a, b, c$は$0$から$255=2^8-1$ の間の整数．
   
   \item directionalLight(a, b, c, d,e,f) RGB[a, b, c]の光で向き(d,e,f) $-1\le d,e,f\le 1$ で光を照らす
   
   \item
       ambientLight(a, b, c) RGB[a, b, c]の光で背景となる光を照らす
       
    \item 
       ambientMaterial(a, b, c) 次の図形の壁の色をRGB [a,b,c] で設定
   
   \item 
      fill(a,b,c) 次の図形をRGB [a,b,c] で埋める．
   
   \item 
      plane(a,b), box(c) などは3次元の組込立体図形
   
   \item
      camera(a, b, c, a', b', c') カメラの位置$(a, b, c)$ 対象の中心$(a', b', c')$ に向けてウィンドウを描画
   
   \item
      push(), pop() はその間に変えられた設定を元に戻すためのスタックに書き込むための関数
   
   \item
      model(modelName) modelNameを描く

\end{itemize}



\subsection{一般のJavaScriptの動かし方}
JavaScriptの処理の仕方は幾つかある．
JavaScriptはブラウザーに組み込まれているので，まず
\begin{verbatim}

   <!DOCTYPE html>
   <html>
      <head>
         <meta charset="UTF-8" />
         <title>Hello JavaScript</title>
      </head>
      <body>
         <script type="text/javascript" src="test.js"></script>
      </body>
   </html>
   
\end{verbatim}
のようなhtml のファイルを hoge.html として保存し，次に
\begin{verbatim}
   window.alert('hello world!');
\end{verbatim}
という行のみからなるファイルを test.js として hoge.htmlと同じディレクトリに保存する．
ブラウザー Chrome の
メニューバーのFileのメニューOpen Files... からhoge.htmlを開くと

\begin{figure}[htb]
   \centering
   \includegraphics[width=12cm, trim=0 1000  0 0, clip]{~/Documents/GitHub/bEuclideanGeo/appa/images/ScreenCapture1.pdf}
   \caption{}\label{alart}
\end{figure}

\begin{figure}[htb]
   \centering
   \includegraphics[width=13cm, trim=0 800  0 0, clip]{~/Documents/GitHub/bEuclideanGeo/appa/images/ScreenCapture2.pdf}
   \caption{}\label{console}
\end{figure}

\noindent
図 \ref{alart}のような結果が得られる．あるいは
Chrome の
メニューバーのFileのメニュー View から Developer の Developer Tools を開き
\begin{verbatim}
   console.log('hello world!');
\end{verbatim}
をtest.jsとして保存し，またhoge.htmlを開くと図 \ref{console}
のような結果が得られる．
%\newpage

他にも Terminal でtest.jsのあるディレクトリに移動し
\begin{verbatim}
   > osascript test.js
\end{verbatim}
でTerminalに結果が表示される．
ただし，これはMacのみで，逆に言うとMacのスクリプト言語としてJavaScriptを使うことが出来る．
一般的には Node.js を\href{https://nodejs.org/en}{ダウンロード}
しインストールすると
\begin{verbatim}
   > node test.js
\end{verbatim}
でTerminalに結果が表示される．したがってMac を直接コントロールする必要が無いときは osascriptより node.js を使うことをお奨めしたい．
%Node.js + Electron + Chromium で macOS, Windows, スマートホンに対応したデスクトップアプリケーション(One page application)を一度に作ることが出来る (例 VScode 図 \ref{vscode})．
%
%\begin{figure}[htb]
%\centering
%\includegraphics[width=7cm, trim=0 40  0 30, clip]{ScreenCapture4.pdf}
%\caption{}\label{vscode}
%\end{figure}
%
%\begin{thebibliography}{9}
%   \bibitem{P5}
%      松田晃一, p5.jsプログラミングガイド  改訂版, 2021, カットシステム
%   
%   \bibitem{T}
%      山田祥寛, 
%      改訂3版JavaScript本格入門, 2023, 技術評論社
%\end{thebibliography}
%
%
%\end{document}
%
%
%\section{GPUとCPU}
%
%\section{剛体的な図形の作図}





% \chapter{アプリのソースと解説}

\subsection{p5.js}
p5.js は Processing というJava用の描画プログラミング言語を Javascript に移植したものである．以下のウェブ上のエディターで実行可能である．

\href{https://editor.p5js.org}{https://editor.p5js.org} （以下この色はハイパーリンク）

\begin{figure}[htb]
   \centering
   \includegraphics[width=13cm]{~/Documents/GitHub/bEuclideanGeo/appa/images/ScreenCapture3.pdf}
   \caption{ウェブ上の p5.js エディター}\label{}
\end{figure}

\noindent
公開した sketch.js を sketch.js の部分にコピーし，
（ファイル \href{https://www.printables.com/model/411768-octahedron-freecad/files}{octahedron.stl}\footnote{Ray Fox,
This work is licensed under a Creative Commons (International License),
Public Domain} をアップロードし）
左上の三角のボタンを押すと右画面に描画が始まる．
もちろん隣りの四角のボタンで止まる．
%プログラムのソースファイルに是非適宜手を入れながら動かして見て欲しい．
ただし .js はJavaScriptのファイルであることを示す拡張子 ( extension ) である．


また以下のようにすると自分のパソコン上で動かせる．
 p5.js というライブラリーを\href{https://p5js.org/download/}{ダウンロード}するとp5というディレクトリが作られる．
 公開されたsketch.jsをそのディレクトリーの sketch.js と置き換え，
フリーのエディター\href{https://code.visualstudio.com}{Visual Studio Code}でそのディレクトリp5を開き，
 Live Server の拡張をVisual Studio Codeにインストールして
Go Liveのボタンを押すと，chrome で再生を始める．

%\subsection{ p5.js のソース}
以下がソースファイル sketch.js の本質的な内容である．

\begin{verbatim}
/* Cubes landing before the Great Pyramids abbreviated version,
 * snow00two@gmail.com, 
 * \href{https://creativecommons.org/licenses/by-nc-nd/4.0/}{\ccbyncsa}
 */
const widthCanvas = 720 * 3/2  ; //=1080
const heightCanvas = 405 * 3/2 ; //=607.5
const backColor = [150, 200, 250] ;//[200, 220, 200]

let octahedron ; 
function preload(){
octahedron = loadModel('octahedron.stl'); 
//https://www.printables.com/model/411768-octahedron-freecad/files
}

function setup(){//canvas : window of the projection coordinates
createCanvas(widthCanvas, heightCanvas, WEBGL);
background(backColor);
}

let i = 0;//basic parameter for frames

function draw(){
let I = 9425; //change the priod condition to 3x2Pi/0.002=9424.7... 

if (i<I) {
} else {
i = 0;// reset basic parameter
}

  // define the camera coordinates from a camera on the helicopter
  // of the following coordinates in the world coordinates : 
  let cameraX = cos(0.5 + i * 0.002)*250 ;
  let cameraY = sin(0.5 + i * 0.002)*250;
  let cameraZ = 200+sin(0.5 + i * 0.002)*30;
  camera(cameraX,cameraY, cameraZ, 100, 0, 100, 0, 0, -1);
  // the output images are in the projection coordinates.

  //The world coordinates : 
  background(backColor);
  directionalLight(250, 200, 200, 0, -1, -1); //light from the sun 
  fill(250, 200, 200);
  plane(10000,10000);

  stroke(0);
  strokeWeight(2);
  fill(150,150, 150);
  pyramid(3);

  translate(-300,0,0);
  pyramid(3);

  translate(100,-300,0);
  pyramid(2);

  ambientMaterial(125);
  ambientLight(100,100,100);
  shininess(100.0);

  translate(400, 500, max(500 - i * 0.3, 67) );
  sCube(5);
  
  let M = 25;
  for(let m = 0 ; m < M ; m++){
    translate(140,140, max(650 + 150*m - i * 0.3, 0));
    sCube(5);
  } 
  i++;
}

// The modeling coordinates :
function pyramid(scaleN){
  push();
    scale(scaleN);
    model(octahedron);
  pop();
}

function sCube(scaleN){
  push();
    rotateZ(QUARTER_PI);
    scale(scaleN);
    box(25);
  pop();
}
\end{verbatim}

\subsection{JavaScriptの言語仕様}
\begin{itemize}
\item
JavaScript には変数の型はなく，文の終りで ; が必要

\item
let は変数の宣言，以前は var も使われていた．

\item
const は定数の宣言

\item
\verb////はコメント文でこれ以降の文は無視される．

\noindent
/* */はブロックコメント文で途中の文は無視される．
\end{itemize}


よく使う制御構文：
for 文の例：
\begin{verbatim}
   let I = 5;
   for (let i=0; i < I; i++) {
        console.log(i);
   }

\end{verbatim}

if 文の例：
\begin{verbatim}
   let m = 80;
   if ( m<60 ){
        console.log('F');
   } else if ( m<70  ) {
        console.log('C');
   } else if ( m<80  ) {
        console.log('B');
   } else if ( m<90  ) {
        console.log('A');
   } else {
        console.log('S');
   }

\end{verbatim}


function hoge(var1,\dots,var3) は関数hogeの宣言で，続く$\{ \}$の中に関数の定義を書く．var1,\dots,var3が関数の変数で定義の中で使える．
関数の中で定義された変数はローカル変数で，関数外で使われる変数はユニバーサルな変数になる．


%\subsection{p5.js の説明とソースファイルの解説}
関数preload()でディレクトリにある図形 octahedron.stl を読み込み，
関数setup()の中でウィンドウ等を初期化する．

   createCanvas(widthCanvas, heightCanvas) ;

\noindent
で$2$次元の描画を行えるが，ここでは

   createCanvas(widthCanvas, heightCanvas, WEBGL);

\noindent
で$3$次元の描画を行う．ただし，WEBGLとはブラウザ上でCGを描画するためのAPIである．ここでは$3$次元の描画のみ解説するが$2$次元の描画については\cite{P5} を参照してほしい．

関数draw()が呼び出されると変数 i を1フレーム毎に1づつ増やしながら，
その都度フレームを描画し続ける．いじってみるとだんだん理解が深まると思う．

\emptyline

p5.js 特有の組込関数の説明：
\begin{itemize}
   \item
      translate(a, b, c); $(a, b,c)$だけ平行移動し，その点を$(0,0,0)$として描画
   
   \item
      rotateZ(theta); $z$軸を中心に角度 theta だけ直交座標系を回転し，それを新たな直交座標系として描画．rotateX, rotateY も同様．
   \item
      scale(a) 座標系をa 倍する
   
   \item  
      strokeWeight(a) 太さaで描く
     
   \item
      background(a,b,c);  背後の画面をRGB [a,b,c] で埋める．
      ただし$a, b, c$は$0$から$255=2^8-1$ の間の整数．
   
   \item directionalLight(a, b, c, d,e,f) RGB[a, b, c]の光で向き(d,e,f) $-1\le d,e,f\le 1$ で光を照らす
   
   \item
       ambientLight(a, b, c) RGB[a, b, c]の光で背景となる光を照らす
       
    \item 
       ambientMaterial(a, b, c) 次の図形の壁の色をRGB [a,b,c] で設定
   
   \item 
      fill(a,b,c) 次の図形をRGB [a,b,c] で埋める．
   
   \item 
      plane(a,b), box(c) などは3次元の組込立体図形
   
   \item
      camera(a, b, c, a', b', c') カメラの位置$(a, b, c)$ 対象の中心$(a', b', c')$ に向けてウィンドウを描画
   
   \item
      push(), pop() はその間に変えられた設定を元に戻すためのスタックに書き込むための関数
   
   \item
      model(modelName) modelNameを描く

\end{itemize}



\subsection{一般のJavaScriptの動かし方}
JavaScriptの処理の仕方は幾つかある．
JavaScriptはブラウザーに組み込まれているので，まず
\begin{verbatim}

   <!DOCTYPE html>
   <html>
      <head>
         <meta charset="UTF-8" />
         <title>Hello JavaScript</title>
      </head>
      <body>
         <script type="text/javascript" src="test.js"></script>
      </body>
   </html>
   
\end{verbatim}
のようなhtml のファイルを hoge.html として保存し，次に
\begin{verbatim}
   window.alert('hello world!');
\end{verbatim}
という行のみからなるファイルを test.js として hoge.htmlと同じディレクトリに保存する．
ブラウザー Chrome の
メニューバーのFileのメニューOpen Files... からhoge.htmlを開くと

\begin{figure}[htb]
   \centering
   \includegraphics[width=12cm, trim=0 1000  0 0, clip]{~/Documents/GitHub/bEuclideanGeo/appa/images/ScreenCapture1.pdf}
   \caption{}\label{alart}
\end{figure}

\begin{figure}[htb]
   \centering
   \includegraphics[width=13cm, trim=0 800  0 0, clip]{~/Documents/GitHub/bEuclideanGeo/appa/images/ScreenCapture2.pdf}
   \caption{}\label{console}
\end{figure}

\noindent
図 \ref{alart}のような結果が得られる．あるいは
Chrome の
メニューバーのFileのメニュー View から Developer の Developer Tools を開き
\begin{verbatim}
   console.log('hello world!');
\end{verbatim}
をtest.jsとして保存し，またhoge.htmlを開くと図 \ref{console}
のような結果が得られる．
%\newpage

他にも Terminal でtest.jsのあるディレクトリに移動し
\begin{verbatim}
   > osascript test.js
\end{verbatim}
でTerminalに結果が表示される．
ただし，これはMacのみで，逆に言うとMacのスクリプト言語としてJavaScriptを使うことが出来る．
一般的には Node.js を\href{https://nodejs.org/en}{ダウンロード}
しインストールすると
\begin{verbatim}
   > node test.js
\end{verbatim}
でTerminalに結果が表示される．したがってMac を直接コントロールする必要が無いときは osascriptより node.js を使うことをお奨めしたい．
%Node.js + Electron + Chromium で macOS, Windows, スマートホンに対応したデスクトップアプリケーション(One page application)を一度に作ることが出来る (例 VScode 図 \ref{vscode})．
%
%\begin{figure}[htb]
%\centering
%\includegraphics[width=7cm, trim=0 40  0 30, clip]{ScreenCapture4.pdf}
%\caption{}\label{vscode}
%\end{figure}
%
%\begin{thebibliography}{9}
%   \bibitem{P5}
%      松田晃一, p5.jsプログラミングガイド  改訂版, 2021, カットシステム
%   
%   \bibitem{T}
%      山田祥寛, 
%      改訂3版JavaScript本格入門, 2023, 技術評論社
%\end{thebibliography}
%
%
%\end{document}
%
%
%\section{GPUとCPU}
%
%\section{剛体的な図形の作図}





\backmatter

\begin{thebibliography}{99}
\bibitem{3DT}
         3DTotal.com, 
         デジタルアーティストが知っておくべきアートの原則 改訂版 -色、光、構図、解剖学、遠近法、奥行き 改訂版, 
         2021.
      
\bibitem{A} Ahlfors, Complex analysis 3rd edition, McGraw-Hill, 1979

   \bibitem{B}
Buerger,
\newblock {\em Intorduction to crystal geometry}.
\newblock  McGraw-Hill, 1974.

   \bibitem{C}
H.S.M. コクセター, 幾何学入門〈上, 下〉 (ちくま学芸文庫) 2009 

\bibitem{Ch}
L.~Charlap, 
\newblock {\em Bieberbach groups and flat manifolds}.
\newblock  Springer,
\newblock 1986, 

      \bibitem{CG}
         画像情報教育振興協会, コンピューターグラフィックス[改訂新版], 2015.
      
      \bibitem{D}
         下村 耕史訳編, アルブレヒト・デューラー「測定法教則」注解, 中央公論美術出版, 2008.


   \bibitem{F}
   W.~Fulton, 
   ALGEBRAIC CURVES---An Introduction to Algebraic Geometry (3rd edition), 2008, 
   \href{https://dept.math.lsa.umich.edu/~wfulton/CurveBook.pdf}{ダウンロード}

   \bibitem{GB}
   L.C.~Grove, C.T~Benson, \textit{Finite reflection groups 2nd Ed.}, GTM99, 1996, Springer
   
   \bibitem{H} 
AMS,
\newblock  {\em Mathematical developments arising from Hilbert Problems, part 2}, 
\newblock 1974

   \bibitem{HC}
D・ヒルベルト, S・コーン=フォッセン, 直観幾何学【新装版】, 2019 , 
みすず書房 %(D. Hilbert, S. Cohn-Vossen, Geometry and the Imagination, 2020, Ams Chelsea Publishing) 

%%　2025/07/21
\bibitem{JJS} J.J,~Sakurai, Modern quantum mechanics, 1985.

\bibitem{I} 
日本数学会監修, 
\newblock 岩波数学辞典 第四版, 
\newblock 岩波書店, 
\newblock 2007

\bibitem{IT}
\newblock {\em International Tables for Crystallography: vol. A, Space-Group Symmetry, new edition},
\newblock Springer, 
\newblock 2017.

\bibitem{N}西山 享, 射影幾何学の考え方 (数学のかんどころ 19), 2013 


\bibitem{E}
C.H.マックギラフィ, 有馬 朗人 , エッシャー—シンメトリーの世界


\bibitem{NS}
Viacheslav V. Nikulin, Igor R. Shafarevich, \textit{Geometries and Groups}, UTX, 1996, Springer, 

            \bibitem{RS}
         スコット・ロバートソン, 
         スコット・ロバートソンのHow to Draw -オブジェクトに構造を与え、実現可能なモデルとして描く方法, 
         2014.
         
      \bibitem{RSII}
         スコット・ロバートソン, 
         スコット・ロバートソンのHow to Render：アイデアを明確に伝える 光と影、反射の描き方, 2015.

\bibitem{K}
スピーロ ; 青木 薫 訳,
\newblock {ケプラー予想 〜四百年の難問が解けるまで〜}
\newblock  新潮社.

   \bibitem{KKS}
   加藤和也, 黒川信重, 斎藤毅, 数論 I Fermatの夢と類体論, 岩波書店, 2016, 
   黒川 信重, 栗原 将人, 斎藤 毅, 数論II 岩澤理論と保型形式, 岩波書店,  2017
   

\bibitem{Kl}F.~Kline, クライン:19世紀の数学, 彌永 昌吉監修, 1995

   \bibitem{KO} 小林昭七, ユークリッド幾何から現代幾何へ, 日本評論社, 1990.

   \bibitem{Ko} Koblitz, A course in Number theory and Cryptography, GTM.
   
\bibitem{L}
Lohne, J. A., Thomas Harriot als Mathematiker, Centaurus 11(1), 19--45, 1965.

   \bibitem{P5}
      松田晃一, p5.jsプログラミングガイド  改訂版, 2021, カットシステム
  
  \bibitem{N} Nydam, Visual complex analysis, 1996

   \bibitem{Se}
   セール, 数論講義, 岩波書店, 1979
 
   \bibitem{Si} Siegel, Topics in complex function theory I, II, III, 
   Wiley-Intersicence, 1969

\bibitem{S}J.  Stillwell, Mathematics and Its History  3rd ed, UTM, Springer, 2010.

\bibitem{S2}J. Stillwell, 
      Sources of Hyperbolic Geometry, History of math.10, AMS. 1996.

   \bibitem{ST} Silverman, Tate,
   Rational Points on Elliptic Curves, Springer, 2015 
  

   \bibitem{T}高木貞治, 復刻版 近世数学史談・数学雑談, 共立出版

\bibitem{Y}
P.~Yale,
\newblock \textit{Geometry and symmetry}.
\newblock  Dover.

  \bibitem{Ya}
      山田祥寛, 
      改訂3版JavaScript本格入門, 2023, 技術評論社

\bibitem{U} 上野健爾, 代数幾何, 岩波書店, 2005

\bibitem{W}
ヘルマン・ワイル, シンメトリー (ちくま学芸文庫) 2022
(Hermann Weyl, Symmetry, Princeton Science Library, 2016)
\end{thebibliography}

\printindex

\end{document}

